\documentclass[11pt]{article}
\usepackage{ifthen}
\usepackage[utf8]{inputenc}
\usepackage[T1]{fontenc}
\usepackage[brazil]{babel}
\usepackage{geometry}
\geometry{a4paper,margin=2cm}
\usepackage{longtable,graphicx,multirow,enumitem,tabularx,setspace,ragged2e}
\setlist{nosep,leftmargin=*}
\renewcommand\arraystretch{1.15}
% Caixa de seleção
\newcommand{\chk}[1]{\ifx#1X\setlength\fboxsep{1pt}\fbox{\rule{1.3ex}{0pt}X}%
\else\setlength\fboxsep{1pt}\fbox{\rule{1.3ex}{0pt}\rule{1.3ex}{0pt}}\fi}
\pagestyle{empty}
\begin{document}
\begin{center}
\begin{tabular}{lccr}
 \multirow{3}{*}{\includegraphics[height=2.7cm]{brasao.png}} &
 \multicolumn{2}{c}{\bfseries UNIVERSIDADE FEDERAL DE OURO PRETO} &
 \ \ \ \ \multirow{3}{*}{\includegraphics[height=2.7cm]{ufop.png}} \\
 & \multicolumn{2}{c}{\bfseries PRÓ-REITORIA DE GRADUAÇÃO} & \\
 & \multicolumn{2}{c}{} & \\
 & \multicolumn{2}{c}{\Large\bfseries PROGRAMA DE DISCIPLINA} & \\
\end{tabular}
\end{center}

\begin{center}
\begin{longtable}{|p{4cm}|p{4cm}|p{4cm}|p{4cm}|}
\hline
\multicolumn{3}{|p{12cm}|}{Nome do Componente Curricular em Português:} &
\multicolumn{1}{p{4cm}|}{Código:} \\ 
\multicolumn{3}{|p{12cm}|}{\textbf{Divulgação Científica}} &
\textbf{BCC132}\\ 
\multicolumn{3}{|p{12cm}|}{Nome do Componente Curricular em Inglês:} & \\ 
\multicolumn{3}{|p{12cm}|}{\textbf{Dissemination of Scientific Knowledge}} & \\ 
\hline
\multicolumn{3}{|p{12cm}|}{Nome e Sigla do Departamento} & Unidade Acadêmica: \\ 
\multicolumn{3}{|p{12cm}|}{Departamento de Computação (DECOM)} & {ICEB} \\ 
\hline
\multicolumn{4}{|p{16cm}|}{Modalidade de Oferta:
[X] presencial \hspace{0.6cm}
[ ] semipresencial \hspace{0.6cm}
[ ] à distância}\\
\hline
\multicolumn{2}{|p{8cm}|}{Carga horária semestral} &
\multicolumn{2}{p{8cm}|}{Carga horária semanal}\\
\hline
\multicolumn{1}{|p{4cm}|}{Total} &
\multicolumn{1}{p{4cm}|}{Extensionista} &
\multicolumn{1}{p{4cm}|}{Teórica} &
\multicolumn{1}{p{4cm}|}{Prática} \\ 
\multicolumn{1}{|p{4cm}|}{120\,horas} &
\multicolumn{1}{p{4cm}|}{60\;horas} &
\multicolumn{1}{p{4cm}|}{0\;horas/aula} &
\multicolumn{1}{p{4cm}|}{8\;horas/aula} \\ 
\hline
\multicolumn{4}{|p{16cm}|}{Ementa:}\\
\multicolumn{4}{|p{16cm}|}{}\\
\multicolumn{4}{|p{\dimexpr 16cm + 6\tabcolsep\relax}|}{Divulgação Científica, Conceituação e aspectos práticos das Redes Sociais; Leitura Científica; Conceitos de produção de textos jornalísticos.}\\
\multicolumn{4}{|p{16cm}|}{}\\
\hline
\multicolumn{4}{|p{16cm}|}{Conteúdo programático:}\\*
\multicolumn{4}{|p{\dimexpr 16cm + 6\tabcolsep\relax}|}{\begin{enumerate}[start=1]\item Introdução à Divulgação Científica e sua Relevância Social: o papel da ciência na sociedade e a importância da comunicação científica; a divulgação científica como ferramenta de inclusão e acesso ao conhecimento; o impacto das redes sociais e mídias digitais na disseminação da ciência\end{enumerate}}\\
\multicolumn{4}{|p{\dimexpr 16cm + 6\tabcolsep\relax}|}{\begin{enumerate}[start=2]\item Estratégias de Comunicação e Dialogicidade na Divulgação Científica: técnicas para tornar conteúdos científicos acessíveis a diferentes públicos; interação com a comunidade para identificação de temas relevantes; linguagem e adaptação de conteúdo para diversos meios (vídeo, texto, áudio)\end{enumerate}}\\
\multicolumn{4}{|p{\dimexpr 16cm + 6\tabcolsep\relax}|}{\begin{enumerate}[start=3]\item Co-criação e Produção de Conteúdo Científico para a Comunidade: oficinas de divulgação científica em escolas e espaços comunitários; produção de materiais didáticos sobre tecnologia e computação para públicos diversos; planejamento e execução de campanhas educativas sobre temas como privacidade digital, inteligência artificial e impacto social da computação\end{enumerate}}\\
\multicolumn{4}{|p{\dimexpr 16cm + 6\tabcolsep\relax}|}{\begin{enumerate}[start=4]\item Uso de Tecnologias Digitais para Divulgação Científica: ferramentas para produção de vídeos, podcasts e infográficos; estratégias para impulsionamento de conteúdo nas redes sociais; monitoramento e análise do impacto das iniciativas de divulgação\end{enumerate}}\\
\multicolumn{4}{|p{\dimexpr 16cm + 6\tabcolsep\relax}|}{\begin{enumerate}[start=5]\item Avaliação do Impacto e Reflexão Crítica sobre a Divulgação Científica: coleta de feedbacks da comunidade atendida; reflexão sobre a transformação do conhecimento acadêmico em informação acessível; registros de boas práticas e sistematização de aprendizados\end{enumerate}}\\
\multicolumn{4}{|p{\dimexpr 16cm + 6\tabcolsep\relax}|}{\begin{enumerate}[start=6]\item \textbf{Perfil da Comunidade:} A disciplina BCC132 - Divulgação Científica busca aproximar a universidade da sociedade por meio da comunicação acessível do conhecimento científico e tecnológico. Seu público-alvo inclui escolas, organizações não-governamentais (ONGs), coletivos culturais, empreendedores, professores e comunidades com pouco acesso à informação científica. Os discentes atuarão como mediadores do conhecimento, transformando conceitos da ciência da computação em conteúdos compreensíveis e aplicáveis ao cotidiano da população. A interação com a comunidade será participativa e dialogada, garantindo que os temas tratados na disciplina sejam definidos com base nos interesses e necessidades da sociedade. Dessa forma, a disciplina contribui para a democratização do conhecimento e o fortalecimento da inclusão digital e científica.\end{enumerate}}\\
\multicolumn{4}{|p{\dimexpr 16cm + 6\tabcolsep\relax}|}{\begin{enumerate}[start=7]\item \textbf{Objetivos Extensionistas:} A disciplina visa (i) promover a democratização do conhecimento científico e tecnológico, tornando a computação acessível para públicos diversos por meio de estratégias de divulgação científica, (ii) fomentar a interação dialógica entre universidade e comunidade, permitindo que os temas abordados na disciplina sejam construídos coletivamente com base nas demandas sociais, (iii) capacitar os discentes para comunicar conceitos científicos de forma acessível, desenvolvendo habilidades de escrita, oralidade e uso de tecnologias para popularizar a ciência, (iv) incentivar a produção de conteúdos educativos, promovendo oficinas, materiais digitais e campanhas informativas sobre temas relevantes da computação, e (v) avaliar o impacto da divulgação científica na comunidade, garantindo que as ações realizadas tenham um efeito concreto na disseminação do conhecimento e no engajamento social.\end{enumerate}}\\
\multicolumn{4}{|p{\dimexpr 16cm + 6\tabcolsep\relax}|}{}\\
\hline
\multicolumn{4}{|p{16cm}|}{Bibliografia Básica:}\\*
\multicolumn{4}{|p{\dimexpr 16cm + 6\tabcolsep\relax}|}{\begin{itemize}\item BUENO, W.C. Estratégias de Comunicação nas Mídias Sociais. (disponível na biblioteca virtual “Minha Biblioteca”)\end{itemize}}\\
\multicolumn{4}{|p{\dimexpr 16cm + 6\tabcolsep\relax}|}{\begin{itemize}\item MARQUES, V. Redes Sociais 360: Como comunicar online. (disponível na biblioteca virtual “Minha Biblioteca”)\end{itemize}}\\
\multicolumn{4}{|p{\dimexpr 16cm + 6\tabcolsep\relax}|}{\begin{itemize}\item MORAIS, F. Planejamento Estratégico Digital. (disponível na biblioteca virtual “Minha Biblioteca”)\end{itemize}}\\
\multicolumn{4}{|p{\dimexpr 16cm + 6\tabcolsep\relax}|}{\begin{itemize}\item OLIVEIRA, J.P.M.; MOTTA, C.A.P. Como Escrever Textos Técnicos. (disponível na biblioteca virtual “Minha Biblioteca”)\end{itemize}}\\
\multicolumn{4}{|p{\dimexpr 16cm + 6\tabcolsep\relax}|}{\begin{itemize}\item TERRA, E. Práticas de Leitura e Escrita. (disponível na biblioteca virtual “Minha Biblioteca”)\end{itemize}}\\
\multicolumn{4}{|p{16cm}|}{}\\
\hline
\multicolumn{4}{|p{16cm}|}{Bibliografia Complementar:}\\*
\multicolumn{4}{|p{\dimexpr 16cm + 6\tabcolsep\relax}|}{\begin{itemize}\item BRASILEIRO, A.M.M. UniA: Leitura e Produção Textual. (disponível na biblioteca virtual “Minha Biblioteca”)\end{itemize}}\\
\hline
\end{longtable}
\end{center}

\clearpage
\begin{center}
\begin{tabular}{lccr}
 \multirow{3}{*}{\includegraphics[height=2.7cm]{brasao.png}} &
 \multicolumn{2}{c}{\bfseries UNIVERSIDADE FEDERAL DE OURO PRETO} &
 \ \ \ \ \multirow{3}{*}{\includegraphics[height=2.7cm]{ufop.png}} \\
 & \multicolumn{2}{c}{\bfseries PRÓ-REITORIA DE GRADUAÇÃO} & \\
 & \multicolumn{2}{c}{} & \\
 & \multicolumn{2}{c}{\Large\bfseries PROGRAMA DE DISCIPLINA} & \\
\end{tabular}
\end{center}

\begin{center}
\begin{longtable}{|p{4cm}|p{4cm}|p{4cm}|p{4cm}|}
\hline
\multicolumn{3}{|p{12cm}|}{Nome do Componente Curricular em Português:} &
\multicolumn{1}{p{4cm}|}{Código:} \\ 
\multicolumn{3}{|p{12cm}|}{\textbf{Programação Funcional}} &
\textbf{BCC111}\\ 
\multicolumn{3}{|p{12cm}|}{Nome do Componente Curricular em Inglês:} & \\ 
\multicolumn{3}{|p{12cm}|}{\textbf{Functional Programming}} & \\ 
\hline
\multicolumn{3}{|p{12cm}|}{Nome e Sigla do Departamento} & Unidade Acadêmica: \\ 
\multicolumn{3}{|p{12cm}|}{Departamento de Computação (DECOM)} & {ICEB} \\ 
\hline
\multicolumn{4}{|p{16cm}|}{Modalidade de Oferta:
[X] presencial \hspace{0.6cm}
[ ] semipresencial \hspace{0.6cm}
[ ] à distância}\\
\hline
\multicolumn{2}{|p{8cm}|}{Carga horária semestral} &
\multicolumn{2}{p{8cm}|}{Carga horária semanal}\\
\hline
\multicolumn{1}{|p{4cm}|}{Total} &
\multicolumn{1}{p{4cm}|}{Extensionista} &
\multicolumn{1}{p{4cm}|}{Teórica} &
\multicolumn{1}{p{4cm}|}{Prática} \\ 
\multicolumn{1}{|p{4cm}|}{60\,horas} &
\multicolumn{1}{p{4cm}|}{0\;horas} &
\multicolumn{1}{p{4cm}|}{2\;horas/aula} &
\multicolumn{1}{p{4cm}|}{2\;horas/aula} \\ 
\hline
\multicolumn{4}{|p{16cm}|}{Ementa:}\\
\multicolumn{4}{|p{16cm}|}{}\\
\multicolumn{4}{|p{\dimexpr 16cm + 6\tabcolsep\relax}|}{Características dos principais paradigmas de programação; princípios do paradigma de programação funcional; principais características de linguagens de programação funcional: recursão, abstração funcional, funções de ordem superior, tipos de dados algébricos, polimorfismo, inferência de tipos, avaliação estrita e avaliação lazy, sobrecarga; demonstração de correção de programas funcionais.}\\
\multicolumn{4}{|p{16cm}|}{}\\
\hline
\multicolumn{4}{|p{16cm}|}{Conteúdo programático:}\\*
\multicolumn{4}{|p{\dimexpr 16cm + 6\tabcolsep\relax}|}{\begin{enumerate}[start=1]\item Introdução\end{enumerate}}\\
\multicolumn{4}{|p{\dimexpr 16cm + 6\tabcolsep\relax}|}{\begin{enumerate}[start=2]\item Paradigmas de programação\end{enumerate}}\\
\multicolumn{4}{|p{\dimexpr 16cm + 6\tabcolsep\relax}|}{\begin{enumerate}[start=3]\item Definindo funções\end{enumerate}}\\
\multicolumn{4}{|p{\dimexpr 16cm + 6\tabcolsep\relax}|}{\begin{enumerate}[start=4]\item Funções recursivas\end{enumerate}}\\
\multicolumn{4}{|p{\dimexpr 16cm + 6\tabcolsep\relax}|}{\begin{enumerate}[start=5]\item Tuplas, listas e polimorfismo paramétrico e de sobrecarga\end{enumerate}}\\
\multicolumn{4}{|p{\dimexpr 16cm + 6\tabcolsep\relax}|}{\begin{enumerate}[start=6]\item Casamento de padrão\end{enumerate}}\\
\multicolumn{4}{|p{\dimexpr 16cm + 6\tabcolsep\relax}|}{\begin{enumerate}[start=7]\item Entrada e saída\end{enumerate}}\\
\multicolumn{4}{|p{\dimexpr 16cm + 6\tabcolsep\relax}|}{\begin{enumerate}[start=8]\item Funções anônimas e de ordem superior\end{enumerate}}\\
\multicolumn{4}{|p{\dimexpr 16cm + 6\tabcolsep\relax}|}{\begin{enumerate}[start=9]\item Tipos de dados algébricos\end{enumerate}}\\
\multicolumn{4}{|p{\dimexpr 16cm + 6\tabcolsep\relax}|}{\begin{enumerate}[start=10]\item Functores e Mônadas\end{enumerate}}\\
\multicolumn{4}{|p{\dimexpr 16cm + 6\tabcolsep\relax}|}{\begin{enumerate}[start=11]\item Teste baseado em propriedades\end{enumerate}}\\
\multicolumn{4}{|p{\dimexpr 16cm + 6\tabcolsep\relax}|}{\begin{enumerate}[start=12]\item Raciocínio sobre programas funcionais\end{enumerate}}\\
\multicolumn{4}{|p{\dimexpr 16cm + 6\tabcolsep\relax}|}{\begin{enumerate}[start=13]\item Noções sobre tipos dependentes em linguagens de programação\end{enumerate}}\\
\multicolumn{4}{|p{\dimexpr 16cm + 6\tabcolsep\relax}|}{}\\
\hline
\multicolumn{4}{|p{16cm}|}{Bibliografia Básica:}\\*
\multicolumn{4}{|p{\dimexpr 16cm + 6\tabcolsep\relax}|}{\begin{itemize}\item Lipovaca, Miran. Learn you a Haskell for the great good! No Starch Press. Disponível gratuitamente on-line: http://learnyouahaskell.com\end{itemize}}\\
\multicolumn{4}{|p{\dimexpr 16cm + 6\tabcolsep\relax}|}{\begin{itemize}\item O'Sullivan, Bryan; Stewart, Don; Goerzen, John. Real World Haskell. Disponível gratuitamente on-line: http://book.realworldhaskell.org\end{itemize}}\\
\multicolumn{4}{|p{\dimexpr 16cm + 6\tabcolsep\relax}|}{\begin{itemize}\item WADLER, Philip; KOKKE, Wen; SIEK, Jeremy. Programming Languages Foundations in Agda. Disponível gratuitamente on-line: https://plfa.github.io\end{itemize}}\\
\multicolumn{4}{|p{\dimexpr 16cm + 6\tabcolsep\relax}|}{\begin{itemize}\item MARLOW, Simon. Parallel and Concurrent Programming in Haskell. Disponível gratuitamente on-line: https://simonmar.github.io/pages/pcph.html\end{itemize}}\\
\multicolumn{4}{|p{16cm}|}{}\\
\hline
\multicolumn{4}{|p{16cm}|}{Bibliografia Complementar:}\\*
\multicolumn{4}{|p{\dimexpr 16cm + 6\tabcolsep\relax}|}{\begin{itemize}\item Jhala, Ranjit; Seidel, Eric; Vazou, Niki. Programming with Refinement Types: An introduction to Liquid Haskell. Disponível gratuitamente on-line: http://ucsd-progsys.github.io/liquidhaskell-tutorial/\end{itemize}}\\
\multicolumn{4}{|p{\dimexpr 16cm + 6\tabcolsep\relax}|}{\begin{itemize}\item BIRD, Richard. Introduction to Functional Programming using Haskell. 2. ed. London: Prentice Hall Press, 1998.\end{itemize}}\\
\multicolumn{4}{|p{\dimexpr 16cm + 6\tabcolsep\relax}|}{\begin{itemize}\item MILNER, R.; TOFTE, Mads; HARPER, Robert. The definition of standard ML. Cambridge: MIT, 1997.\end{itemize}}\\
\multicolumn{4}{|p{\dimexpr 16cm + 6\tabcolsep\relax}|}{\begin{itemize}\item ULLMAN, Jeffrey D. Elements of ML programming. New Jersey: Prentice-Hall, 1998.\end{itemize}}\\
\multicolumn{4}{|p{\dimexpr 16cm + 6\tabcolsep\relax}|}{\begin{itemize}\item PAULSON, Laurence C. ML for the working programming. 2. ed. Cambridge: Cambridge University, 1996.\end{itemize}}\\
\multicolumn{4}{|p{\dimexpr 16cm + 6\tabcolsep\relax}|}{\begin{itemize}\item PEYTON JONES, Simon L. The implementation of functional programming languages. New York: Prentice-Hall, 1987.\end{itemize}}\\
\hline
\end{longtable}
\end{center}

\clearpage
\begin{center}
\begin{tabular}{lccr}
 \multirow{3}{*}{\includegraphics[height=2.7cm]{brasao.png}} &
 \multicolumn{2}{c}{\bfseries UNIVERSIDADE FEDERAL DE OURO PRETO} &
 \ \ \ \ \multirow{3}{*}{\includegraphics[height=2.7cm]{ufop.png}} \\
 & \multicolumn{2}{c}{\bfseries PRÓ-REITORIA DE GRADUAÇÃO} & \\
 & \multicolumn{2}{c}{} & \\
 & \multicolumn{2}{c}{\Large\bfseries PROGRAMA DE DISCIPLINA} & \\
\end{tabular}
\end{center}

\begin{center}
\begin{longtable}{|p{4cm}|p{4cm}|p{4cm}|p{4cm}|}
\hline
\multicolumn{3}{|p{12cm}|}{Nome do Componente Curricular em Português:} &
\multicolumn{1}{p{4cm}|}{Código:} \\ 
\multicolumn{3}{|p{12cm}|}{\textbf{Engenharia de Software II}} &
\textbf{BCC114}\\ 
\multicolumn{3}{|p{12cm}|}{Nome do Componente Curricular em Inglês:} & \\ 
\multicolumn{3}{|p{12cm}|}{\textbf{Software Engineering II}} & \\ 
\hline
\multicolumn{3}{|p{12cm}|}{Nome e Sigla do Departamento} & Unidade Acadêmica: \\ 
\multicolumn{3}{|p{12cm}|}{Departamento de Computação (DECOM)} & {ICEB} \\ 
\hline
\multicolumn{4}{|p{16cm}|}{Modalidade de Oferta:
[X] presencial \hspace{0.6cm}
[ ] semipresencial \hspace{0.6cm}
[ ] à distância}\\
\hline
\multicolumn{2}{|p{8cm}|}{Carga horária semestral} &
\multicolumn{2}{p{8cm}|}{Carga horária semanal}\\
\hline
\multicolumn{1}{|p{4cm}|}{Total} &
\multicolumn{1}{p{4cm}|}{Extensionista} &
\multicolumn{1}{p{4cm}|}{Teórica} &
\multicolumn{1}{p{4cm}|}{Prática} \\ 
\multicolumn{1}{|p{4cm}|}{60\,horas} &
\multicolumn{1}{p{4cm}|}{0\;horas} &
\multicolumn{1}{p{4cm}|}{4\;horas/aula} &
\multicolumn{1}{p{4cm}|}{0\;horas/aula} \\ 
\hline
\multicolumn{4}{|p{16cm}|}{Ementa:}\\
\multicolumn{4}{|p{16cm}|}{}\\
\multicolumn{4}{|p{\dimexpr 16cm + 6\tabcolsep\relax}|}{Processos de Desenvolvimento de Software. Engenharia de Requisitos de Software. Gestão de Projetos de Software. Software como Serviço (SaaS). Devops.}\\
\multicolumn{4}{|p{16cm}|}{}\\
\hline
\multicolumn{4}{|p{16cm}|}{Conteúdo programático:}\\*
\multicolumn{4}{|p{\dimexpr 16cm + 6\tabcolsep\relax}|}{\begin{enumerate}[start=1]\item Processos de Desenvolvimento: importância dos processos para produtividade da equipe e garantia de qualidade do produto; tipos de processos (sequenciais ou cíclicos, incrementais e/ou evolutivos); processos dirigidos por planos (cascata, espiral, Rational Unified Process); processos ágeis (Manifesto Ágil, principais características, principais diferenças entre XP, Scrum e Kanban, críticas aos métodos ágeis); Extreme Programming (contexto histórico, princípios e valores, representante do cliente e histórias de usuário, iterações e releases, estimativa de esforço, práticas de programação, gerenciamento de projetos); Scrum (contexto histórico, ciclos de desenvolvimento, papéis, artefatos, rituais); Kanban (contexto histórico, backlog do produto, quadro Kanban, limites WIP)\end{enumerate}}\\
\multicolumn{4}{|p{\dimexpr 16cm + 6\tabcolsep\relax}|}{\begin{enumerate}[start=2]\item Engenharia de Requisitos: tipos de requisitos (funcionais, não funcionais, de interface); técnicas de levantamento e análise de requisitos orientada a objetos; especificação de requisitos (desenvolvimento dirigido por comportamento - BDD, histórias de usuário e cenários de teste, casos de uso, diagramas de caso de uso - UML); produto mínimo viável (MVP); testes A/B\end{enumerate}}\\
\multicolumn{4}{|p{\dimexpr 16cm + 6\tabcolsep\relax}|}{\begin{enumerate}[start=3]\item Gestão de Projetos: métodos e técnicas para gestão de projetos baseados nas 10 áreas do PMI - Project Management Institute (escopo, tempo, custos, qualidade, recursos humanos, comunicações, risco, aquisições, partes interessadas e integração); técnicas para definição de backlog do produto e estrutura de decomposição do trabalho (WBS); mobilização de equipes (processos, organização, papéis e responsabilidades, rastreabilidade, solução de problemas); estimativa de esforço e precificação; ferramentas para monitoramento e controle de processos de desenvolvimento de software (Jira e Git)\end{enumerate}}\\
\multicolumn{4}{|p{\dimexpr 16cm + 6\tabcolsep\relax}|}{\begin{enumerate}[start=4]\item Software como Serviço (SaaS): introdução e contexto histórico; conceitos básicos (Arquitetura Orientada a Serviços - SOA, Software como Serviço, Computação em Nuvem); arquitetura de aplicação SaaS (arquitetura cliente-servidor versus peer-to-peer, cliente web e linguagens de script, comunicação via HTTP e URIs, representação via HTML e CSS, arquitetura em 3 camadas e escalabilidade horizontal, arquitetura Model-View-Controller - MVC); desempenho de aplicações SaaS (responsividade e disponibilidade, monitoramento e identificação de gargalos, uso de caching, consultas ao banco de dados); segurança: proteção dos dados do cliente e Lei Geral de Proteção de Dados (LGPD)\end{enumerate}}\\
\multicolumn{4}{|p{\dimexpr 16cm + 6\tabcolsep\relax}|}{\begin{enumerate}[start=5]\item Devops: controle de versões (conceitos, ferramentas, multirrepositórios versus monorrepositórios); Integração Contínua - CI (motivação, conceitos como master/trunk, branches, merges, conflitos de integração, gitflow); Entrega Contínua - CD (implantação/deployment, entrega/delivery, feature flags, release flags e business flags); ferramentas para controle de versões de código e automação dos pipelines CI/CD (Git)\end{enumerate}}\\
\multicolumn{4}{|p{\dimexpr 16cm + 6\tabcolsep\relax}|}{}\\
\hline
\multicolumn{4}{|p{16cm}|}{Bibliografia Básica:}\\*
\multicolumn{4}{|p{\dimexpr 16cm + 6\tabcolsep\relax}|}{\begin{itemize}\item Valente, M. T. Engenharia de Software Moderna: Princípios e Práticas para Desenvolvimento de Software com Produtividade. 1ª edição. 2020. ISBN 978-6500019506. Disponível on-line: https://engsoftmoderna.info/\end{itemize}}\\
\multicolumn{4}{|p{\dimexpr 16cm + 6\tabcolsep\relax}|}{\begin{itemize}\item Patterson, D. A.; Fox, A. Construindo Software como Serviço (SaaS): Uma abordagem Ágil Usando Computação em Nuvem. 1st ed. Strawberry Canyon LLC, 2015. ISBN 9780984881277. Disponível on-line: http://www.saasbook.info/translations/br ou http://www.saasbook.info/first-edition\end{itemize}}\\
\multicolumn{4}{|p{\dimexpr 16cm + 6\tabcolsep\relax}|}{\begin{itemize}\item HULL, Elizabeth; JACKSON, Ken; DICK, Jeremy. Requirements Engineering. 3. ed. Springer. ISBN 978184996404. https://pergamum.ufop.br/pergamum/biblioteca/fotos.php?cod\_acervo=1415964\end{itemize}}\\
\multicolumn{4}{|p{16cm}|}{}\\
\hline
\multicolumn{4}{|p{16cm}|}{Bibliografia Complementar:}\\*
\multicolumn{4}{|p{\dimexpr 16cm + 6\tabcolsep\relax}|}{\begin{itemize}\item Sommerville, I. Engenharia de Software. 10º edição. São Paulo: Pearson, 2019. ISBN 9788543024974. https://plataforma.bvirtual.com.br/Leitor/Loader/168127/pdf/6?keep=True\end{itemize}}\\
\multicolumn{4}{|p{\dimexpr 16cm + 6\tabcolsep\relax}|}{\begin{itemize}\item Fowler, M. UML distilled: a brief guide to the standard object modeling language. 3nd ed. Boston: Addison Wesley, c2004. xxx, 175 p. (Object technology series). ISBN 9780321193681. https://integrada.minhabiblioteca.com.br/\#/books/9788560031382/cfi/6/2!/4/18/52@0:76.8\end{itemize}}\\
\multicolumn{4}{|p{\dimexpr 16cm + 6\tabcolsep\relax}|}{\begin{itemize}\item Gamma, E. Padrões de Projetos: Soluções Reutilizáveis de Software Orientados a Objetos / Erich Gamma... [et al.]. Reading: Addison Wesley, 1995. 395 p. (Addison-Wesley Professional Computing Series). ISBN 9788573076103. https://doku.pub/download/padroes-de-projeto-soluoes-reutilizaveis-de-software-orientado-a-objetos-erich-gammapdf-6lkvge8em804\end{itemize}}\\
\multicolumn{4}{|p{\dimexpr 16cm + 6\tabcolsep\relax}|}{\begin{itemize}\item PMBOK® Guide (2021). https://www.pmi.org/pmbok-guide-standards/\end{itemize}}\\
\multicolumn{4}{|p{\dimexpr 16cm + 6\tabcolsep\relax}|}{\begin{itemize}\item CRISPIN, Lisa; GREGORY, Janet. Agile testing: a practical guide for testers and agile teams. Upper Saddle River, NJ: Addison Wesley, c2009. xli, 533 p. (The Addison-Wesley signature series). ISBN 9780321534460.\end{itemize}}\\
\hline
\end{longtable}
\end{center}

\clearpage
\begin{center}
\begin{tabular}{lccr}
 \multirow{3}{*}{\includegraphics[height=2.7cm]{brasao.png}} &
 \multicolumn{2}{c}{\bfseries UNIVERSIDADE FEDERAL DE OURO PRETO} &
 \ \ \ \ \multirow{3}{*}{\includegraphics[height=2.7cm]{ufop.png}} \\
 & \multicolumn{2}{c}{\bfseries PRÓ-REITORIA DE GRADUAÇÃO} & \\
 & \multicolumn{2}{c}{} & \\
 & \multicolumn{2}{c}{\Large\bfseries PROGRAMA DE DISCIPLINA} & \\
\end{tabular}
\end{center}

\begin{center}
\begin{longtable}{|p{4cm}|p{4cm}|p{4cm}|p{4cm}|}
\hline
\multicolumn{3}{|p{12cm}|}{Nome do Componente Curricular em Português:} &
\multicolumn{1}{p{4cm}|}{Código:} \\ 
\multicolumn{3}{|p{12cm}|}{\textbf{Sistemas Distribuídos}} &
\textbf{BCC115}\\ 
\multicolumn{3}{|p{12cm}|}{Nome do Componente Curricular em Inglês:} & \\ 
\multicolumn{3}{|p{12cm}|}{\textbf{Distributed Systems}} & \\ 
\hline
\multicolumn{3}{|p{12cm}|}{Nome e Sigla do Departamento} & Unidade Acadêmica: \\ 
\multicolumn{3}{|p{12cm}|}{Departamento de Computação (DECOM)} & {ICEB} \\ 
\hline
\multicolumn{4}{|p{16cm}|}{Modalidade de Oferta:
[X] presencial \hspace{0.6cm}
[ ] semipresencial \hspace{0.6cm}
[ ] à distância}\\
\hline
\multicolumn{2}{|p{8cm}|}{Carga horária semestral} &
\multicolumn{2}{p{8cm}|}{Carga horária semanal}\\
\hline
\multicolumn{1}{|p{4cm}|}{Total} &
\multicolumn{1}{p{4cm}|}{Extensionista} &
\multicolumn{1}{p{4cm}|}{Teórica} &
\multicolumn{1}{p{4cm}|}{Prática} \\ 
\multicolumn{1}{|p{4cm}|}{60\,horas} &
\multicolumn{1}{p{4cm}|}{0\;horas} &
\multicolumn{1}{p{4cm}|}{4\;horas/aula} &
\multicolumn{1}{p{4cm}|}{0\;horas/aula} \\ 
\hline
\multicolumn{4}{|p{16cm}|}{Ementa:}\\
\multicolumn{4}{|p{16cm}|}{}\\
\multicolumn{4}{|p{\dimexpr 16cm + 6\tabcolsep\relax}|}{Conceito de sistemas distribuídos. Arquitetura de sistemas distribuídos. Processos e comunicação entre processos. Serviço de nomeação em sistemas distribuídos. Serviço de sincronização em sistemas distribuídos. Serviço de consistência e replicação em sistemas distribuídos. Serviço de tolerância a falhas em sistemas distribuídos. Plataformas e aplicações distribuídas.}\\
\multicolumn{4}{|p{16cm}|}{}\\
\hline
\multicolumn{4}{|p{16cm}|}{Conteúdo programático:}\\*
\multicolumn{4}{|p{\dimexpr 16cm + 6\tabcolsep\relax}|}{\begin{enumerate}[start=1]\item Introdução aos sistemas distribuídos: arquiteturas, exemplos, evolução e os desafios ainda em aberto\end{enumerate}}\\
\multicolumn{4}{|p{\dimexpr 16cm + 6\tabcolsep\relax}|}{\begin{enumerate}[start=2]\item Revisional: processos, comunicação entre processos, programação concorrente e programação em redes\end{enumerate}}\\
\multicolumn{4}{|p{\dimexpr 16cm + 6\tabcolsep\relax}|}{\begin{enumerate}[start=3]\item Conceito de virtualização e introdução aos ambientes virtuais\end{enumerate}}\\
\multicolumn{4}{|p{\dimexpr 16cm + 6\tabcolsep\relax}|}{\begin{enumerate}[start=4]\item Exemplos de plataformas distribuídas que operam em nuvens computacionais\end{enumerate}}\\
\multicolumn{4}{|p{\dimexpr 16cm + 6\tabcolsep\relax}|}{\begin{enumerate}[start=5]\item Como nomear e localizar serviços em ambientes distribuídos\end{enumerate}}\\
\multicolumn{4}{|p{\dimexpr 16cm + 6\tabcolsep\relax}|}{\begin{enumerate}[start=6]\item Exclusão mútua e protocolos de sincronização distribuída\end{enumerate}}\\
\multicolumn{4}{|p{\dimexpr 16cm + 6\tabcolsep\relax}|}{\begin{enumerate}[start=7]\item Problema de consistência e os protocolos de replicação\end{enumerate}}\\
\multicolumn{4}{|p{\dimexpr 16cm + 6\tabcolsep\relax}|}{\begin{enumerate}[start=8]\item Tipos de falha e protocolos para tolerância a falhas e recuperação de falhas\end{enumerate}}\\
\multicolumn{4}{|p{\dimexpr 16cm + 6\tabcolsep\relax}|}{\begin{enumerate}[start=9]\item Exemplos de aplicações distribuídas (roteamento, ordenação, problemas em grafos, recuperação de informação, mineração de dados, banco de dados, entre outros)\end{enumerate}}\\
\multicolumn{4}{|p{\dimexpr 16cm + 6\tabcolsep\relax}|}{}\\
\hline
\multicolumn{4}{|p{16cm}|}{Bibliografia Básica:}\\*
\multicolumn{4}{|p{\dimexpr 16cm + 6\tabcolsep\relax}|}{\begin{itemize}\item TANENBAUM, Andrew S.; STEEN, Maarten Van. Sistemas distribuídos: princípios e paradigmas. 2. ed. São Paulo: Pearson Prentice Hall, 2007.\end{itemize}}\\
\multicolumn{4}{|p{\dimexpr 16cm + 6\tabcolsep\relax}|}{\begin{itemize}\item DOLLIMORE, Jean; KINDBERG, Tim; TORTELLO, João; CARISSIMI, Alexandre; COULOURIS, George. Sistemas distribuídos: conceitos e projeto. 5. ed. Porto Alegre: Bookman, 2013.\end{itemize}}\\
\multicolumn{4}{|p{16cm}|}{}\\
\hline
\multicolumn{4}{|p{16cm}|}{Bibliografia Complementar:}\\*
\multicolumn{4}{|p{\dimexpr 16cm + 6\tabcolsep\relax}|}{\begin{itemize}\item KLEPPMANN, Martin. Designing Data-Intensive Applications: The Big Ideas Behind Reliable, Scalable, and Maintainable Systems. 1. ed. São Paulo: Oreilly \& Assoc, 2015.\end{itemize}}\\
\multicolumn{4}{|p{\dimexpr 16cm + 6\tabcolsep\relax}|}{\begin{itemize}\item VERAS, Manoel. Computação em Nuvem. 1. ed. Brasport, 2015. Disponível online em MinhaUFOP/BibliotecaDigital/BVirtualPearson: https://plataforma.bvirtual.com.br/Acervo/Publicacao/160695\end{itemize}}\\
\multicolumn{4}{|p{\dimexpr 16cm + 6\tabcolsep\relax}|}{\begin{itemize}\item SANTOS, Ricardo Ribeiro dos. FACOM – UFMS, 2020. Material de Sistemas Distribuídos. Disponível em: https://www.facom.ufms.br/~ricardo/Courses/DisSys/Material/ . Acesso em: 07/12/2020.\end{itemize}}\\
\multicolumn{4}{|p{\dimexpr 16cm + 6\tabcolsep\relax}|}{\begin{itemize}\item RODRIGUEZ, Noemi. PUC-RJ, 2020. Material de Sistemas Distribuídos. Disponível em: http://www.inf.puc-rio.br/~noemi/sd-19/ . Acesso em: 07/12/2020.\end{itemize}}\\
\multicolumn{4}{|p{\dimexpr 16cm + 6\tabcolsep\relax}|}{\begin{itemize}\item FERNANDEZ, Marcial Porto. UECE, 2020. Material de Sistemas Distribuídos. Disponível em: http://marcial.larces.uece.br/cursos/sistemas-distribuidos-2018-1 . Acesso em: 07/12/2020.\end{itemize}}\\
\multicolumn{4}{|p{\dimexpr 16cm + 6\tabcolsep\relax}|}{\begin{itemize}\item FIGUEIREDO, Daniel Ratton. UFRJ, 2020. Material de sistemas distribuídos. Disponível em: https://www.cos.ufrj.br/~daniel/sd/ . Acesso em: 07/12/2020.\end{itemize}}\\
\hline
\end{longtable}
\end{center}

\clearpage
\begin{center}
\begin{tabular}{lccr}
 \multirow{3}{*}{\includegraphics[height=2.7cm]{brasao.png}} &
 \multicolumn{2}{c}{\bfseries UNIVERSIDADE FEDERAL DE OURO PRETO} &
 \ \ \ \ \multirow{3}{*}{\includegraphics[height=2.7cm]{ufop.png}} \\
 & \multicolumn{2}{c}{\bfseries PRÓ-REITORIA DE GRADUAÇÃO} & \\
 & \multicolumn{2}{c}{} & \\
 & \multicolumn{2}{c}{\Large\bfseries PROGRAMA DE DISCIPLINA} & \\
\end{tabular}
\end{center}

\begin{center}
\begin{longtable}{|p{4cm}|p{4cm}|p{4cm}|p{4cm}|}
\hline
\multicolumn{3}{|p{12cm}|}{Nome do Componente Curricular em Português:} &
\multicolumn{1}{p{4cm}|}{Código:} \\ 
\multicolumn{3}{|p{12cm}|}{\textbf{Banco de Dados II}} &
\textbf{BCC121}\\ 
\multicolumn{3}{|p{12cm}|}{Nome do Componente Curricular em Inglês:} & \\ 
\multicolumn{3}{|p{12cm}|}{\textbf{Database II}} & \\ 
\hline
\multicolumn{3}{|p{12cm}|}{Nome e Sigla do Departamento} & Unidade Acadêmica: \\ 
\multicolumn{3}{|p{12cm}|}{Departamento de Computação (DECOM)} & {ICEB} \\ 
\hline
\multicolumn{4}{|p{16cm}|}{Modalidade de Oferta:
[X] presencial \hspace{0.6cm}
[ ] semipresencial \hspace{0.6cm}
[ ] à distância}\\
\hline
\multicolumn{2}{|p{8cm}|}{Carga horária semestral} &
\multicolumn{2}{p{8cm}|}{Carga horária semanal}\\
\hline
\multicolumn{1}{|p{4cm}|}{Total} &
\multicolumn{1}{p{4cm}|}{Extensionista} &
\multicolumn{1}{p{4cm}|}{Teórica} &
\multicolumn{1}{p{4cm}|}{Prática} \\ 
\multicolumn{1}{|p{4cm}|}{60\,horas} &
\multicolumn{1}{p{4cm}|}{0\;horas} &
\multicolumn{1}{p{4cm}|}{4\;horas/aula} &
\multicolumn{1}{p{4cm}|}{0\;horas/aula} \\ 
\hline
\multicolumn{4}{|p{16cm}|}{Ementa:}\\
\multicolumn{4}{|p{16cm}|}{}\\
\multicolumn{4}{|p{\dimexpr 16cm + 6\tabcolsep\relax}|}{Processamento de transações. Controle de concorrência. Recuperação de falhas. Funções armazenadas e gatilhos. Otimização de consultas. Bancos de dados orientados a objetos. Introdução de XML. Tópicos avançados de banco de dados.}\\
\multicolumn{4}{|p{16cm}|}{}\\
\hline
\multicolumn{4}{|p{16cm}|}{Conteúdo programático:}\\*
\multicolumn{4}{|p{\dimexpr 16cm + 6\tabcolsep\relax}|}{\begin{enumerate}[start=1]\item Processamento de transações: teoria relativa ao processamento de transações, planos de execução de transações\end{enumerate}}\\
\multicolumn{4}{|p{\dimexpr 16cm + 6\tabcolsep\relax}|}{\begin{enumerate}[start=2]\item Controle de concorrência: teoria relativa ao controle de concorrência, técnicas de controle de concorrência\end{enumerate}}\\
\multicolumn{4}{|p{\dimexpr 16cm + 6\tabcolsep\relax}|}{\begin{enumerate}[start=3]\item Recuperação de falhas: teoria relativa à recuperação de falhas, técnicas de recuperação de falhas\end{enumerate}}\\
\multicolumn{4}{|p{\dimexpr 16cm + 6\tabcolsep\relax}|}{\begin{enumerate}[start=4]\item Funções armazenadas e gatilhos: criação e utilização de funções armazenadas, criação e utilização de gatilhos\end{enumerate}}\\
\multicolumn{4}{|p{\dimexpr 16cm + 6\tabcolsep\relax}|}{\begin{enumerate}[start=5]\item Otimização de consultas: algoritmos para processamento e otimização de consultas, estimativa de custo de processamento\end{enumerate}}\\
\multicolumn{4}{|p{\dimexpr 16cm + 6\tabcolsep\relax}|}{\begin{enumerate}[start=6]\item Bancos de dados orientados a objetos: teoria relativa a bancos de dados orientados a objetos, linguagem de definição de objetos, linguagem de consulta a objetos\end{enumerate}}\\
\multicolumn{4}{|p{\dimexpr 16cm + 6\tabcolsep\relax}|}{\begin{enumerate}[start=7]\item Introdução de XML: sintaxe de XML, XML e dados semi-estruturados, XML e bancos de dados relacionais\end{enumerate}}\\
\multicolumn{4}{|p{\dimexpr 16cm + 6\tabcolsep\relax}|}{\begin{enumerate}[start=8]\item Tópicos avançados de banco de dados: armazenamento de dados e indexação, bancos de dados espaciais, bancos de dados distribuídos, bancos de dados textuais, bancos de dados não relacionais, biblioteca digital, armazém de informações\end{enumerate}}\\
\multicolumn{4}{|p{\dimexpr 16cm + 6\tabcolsep\relax}|}{}\\
\hline
\multicolumn{4}{|p{16cm}|}{Bibliografia Básica:}\\*
\multicolumn{4}{|p{\dimexpr 16cm + 6\tabcolsep\relax}|}{\begin{itemize}\item ELMASRI, Ramez; NAVATHE, Sham. Sistemas de banco de dados. 7. ed. São Paulo: Pearson, 2018. 1152 p.\end{itemize}}\\
\multicolumn{4}{|p{\dimexpr 16cm + 6\tabcolsep\relax}|}{\begin{itemize}\item SILBERSCHATZ, Abraham; KORTH, Henry F.; SUDARSHAN, S. Sistema de banco de dados. 5. ed. Rio de Janeiro: Elsevier, 2006. 805 p.\end{itemize}}\\
\multicolumn{4}{|p{\dimexpr 16cm + 6\tabcolsep\relax}|}{\begin{itemize}\item DATE, Christopher J. Introdução a sistemas de bancos de dados. 8. ed. Rio de Janeiro: Elsevier, 2004. 1623 p.\end{itemize}}\\
\multicolumn{4}{|p{16cm}|}{}\\
\hline
\multicolumn{4}{|p{16cm}|}{Bibliografia Complementar:}\\*
\multicolumn{4}{|p{\dimexpr 16cm + 6\tabcolsep\relax}|}{\begin{itemize}\item RAMAKRISHNAN, Raghu; GEHRKE, Johannes. Sistemas de gerenciamento de banco de dados. 3. ed. São Paulo: McGraw-Hill, 2008. 905 p.\end{itemize}}\\
\multicolumn{4}{|p{\dimexpr 16cm + 6\tabcolsep\relax}|}{\begin{itemize}\item ELMASRI, Ramez; NAVATHE, Sham. Sistemas de banco de dados. 6. ed. São Paulo: Pearson, Addison Wesley, 2011. 788 p.\end{itemize}}\\
\multicolumn{4}{|p{\dimexpr 16cm + 6\tabcolsep\relax}|}{\begin{itemize}\item AMADEU, Claudia V. Banco de Dados. 1. ed. São Paulo: Pearson, 2015. 208 p.\end{itemize}}\\
\multicolumn{4}{|p{\dimexpr 16cm + 6\tabcolsep\relax}|}{\begin{itemize}\item MEDEIROS, Luciano F. Banco de Dados: princípios e prática. 1. ed. Editora Intersaberes, 2013. 188 p.\end{itemize}}\\
\multicolumn{4}{|p{\dimexpr 16cm + 6\tabcolsep\relax}|}{\begin{itemize}\item PUGA, Sandra; FRANÇA, Edson; GOYA, Milton. Banco de dados: implementação em SQL, PL/SQL e Oracle 11g. 1. ed. São Paulo: Pearson, 2013. 356 p.\end{itemize}}\\
\multicolumn{4}{|p{\dimexpr 16cm + 6\tabcolsep\relax}|}{\begin{itemize}\item SETZER, Valdemar W.; SILVA, Flávio S. C. Bancos de dados. 1. ed. Editora Blucher, 2005. 393 p.\end{itemize}}\\
\multicolumn{4}{|p{\dimexpr 16cm + 6\tabcolsep\relax}|}{\begin{itemize}\item CASANOVA, Marco; CÂMARA, Gilberto; DAVIS, Clodoveu; VINHAS, Lúbia; QUEIROZ, Gilberto R. Bancos de Dados Geográficos. 1. ed. Curitiba: MundoGEO, 2005.\end{itemize}}\\
\hline
\end{longtable}
\end{center}

\clearpage
\begin{center}
\begin{tabular}{lccr}
 \multirow{3}{*}{\includegraphics[height=2.7cm]{brasao.png}} &
 \multicolumn{2}{c}{\bfseries UNIVERSIDADE FEDERAL DE OURO PRETO} &
 \ \ \ \ \multirow{3}{*}{\includegraphics[height=2.7cm]{ufop.png}} \\
 & \multicolumn{2}{c}{\bfseries PRÓ-REITORIA DE GRADUAÇÃO} & \\
 & \multicolumn{2}{c}{} & \\
 & \multicolumn{2}{c}{\Large\bfseries PROGRAMA DE DISCIPLINA} & \\
\end{tabular}
\end{center}

\begin{center}
\begin{longtable}{|p{4cm}|p{4cm}|p{4cm}|p{4cm}|}
\hline
\multicolumn{3}{|p{12cm}|}{Nome do Componente Curricular em Português:} &
\multicolumn{1}{p{4cm}|}{Código:} \\ 
\multicolumn{3}{|p{12cm}|}{\textbf{Programação Web}} &
\textbf{BCC123}\\ 
\multicolumn{3}{|p{12cm}|}{Nome do Componente Curricular em Inglês:} & \\ 
\multicolumn{3}{|p{12cm}|}{\textbf{Web Programming}} & \\ 
\hline
\multicolumn{3}{|p{12cm}|}{Nome e Sigla do Departamento} & Unidade Acadêmica: \\ 
\multicolumn{3}{|p{12cm}|}{Departamento de Computação (DECOM)} & {ICEB} \\ 
\hline
\multicolumn{4}{|p{16cm}|}{Modalidade de Oferta:
[X] presencial \hspace{0.6cm}
[ ] semipresencial \hspace{0.6cm}
[ ] à distância}\\
\hline
\multicolumn{2}{|p{8cm}|}{Carga horária semestral} &
\multicolumn{2}{p{8cm}|}{Carga horária semanal}\\
\hline
\multicolumn{1}{|p{4cm}|}{Total} &
\multicolumn{1}{p{4cm}|}{Extensionista} &
\multicolumn{1}{p{4cm}|}{Teórica} &
\multicolumn{1}{p{4cm}|}{Prática} \\ 
\multicolumn{1}{|p{4cm}|}{60\,horas} &
\multicolumn{1}{p{4cm}|}{0\;horas} &
\multicolumn{1}{p{4cm}|}{4\;horas/aula} &
\multicolumn{1}{p{4cm}|}{0\;horas/aula} \\ 
\hline
\multicolumn{4}{|p{16cm}|}{Ementa:}\\
\multicolumn{4}{|p{16cm}|}{}\\
\multicolumn{4}{|p{\dimexpr 16cm + 6\tabcolsep\relax}|}{Linguagens de marcação e folhas de estilos; Tópicos de desenvolvimento Front-end; Programação do lado do servidor; Funcionamento dos Browsers; Bibliotecas e frameworks para desenvolvimento de aplicações Web.}\\
\multicolumn{4}{|p{16cm}|}{}\\
\hline
\multicolumn{4}{|p{16cm}|}{Conteúdo programático:}\\*
\multicolumn{4}{|p{\dimexpr 16cm + 6\tabcolsep\relax}|}{\begin{enumerate}[start=1]\item Linguagens de marcação e folhas de estilos: HTML (conceitos básicos e boas práticas, web semântica, responsividade); CSS (conceitos básicos e boas práticas, block vs inline, seletores, box model, flexbox, position, responsividade, animações)\end{enumerate}}\\
\multicolumn{4}{|p{\dimexpr 16cm + 6\tabcolsep\relax}|}{\begin{enumerate}[start=2]\item Tópicos de desenvolvimento Front-end: conceitos básicos de linguagem de script, Document Object Model (DOM), eventos, extensões de browser, eventos assíncronos (API fetch, Promise), JSON\end{enumerate}}\\
\multicolumn{4}{|p{\dimexpr 16cm + 6\tabcolsep\relax}|}{\begin{enumerate}[start=3]\item Programação do lado do servidor: requisições HTTP, CORS, rotas, bancos de dados, autenticação\end{enumerate}}\\
\multicolumn{4}{|p{\dimexpr 16cm + 6\tabcolsep\relax}|}{\begin{enumerate}[start=4]\item Funcionamento dos Browsers: história, princípios de design\end{enumerate}}\\
\multicolumn{4}{|p{\dimexpr 16cm + 6\tabcolsep\relax}|}{\begin{enumerate}[start=5]\item Bibliotecas e frameworks para desenvolvimento de aplicações Web\end{enumerate}}\\
\multicolumn{4}{|p{\dimexpr 16cm + 6\tabcolsep\relax}|}{}\\
\hline
\multicolumn{4}{|p{16cm}|}{Bibliografia Básica:}\\*
\multicolumn{4}{|p{\dimexpr 16cm + 6\tabcolsep\relax}|}{\begin{itemize}\item FLANAGAN, David. JavaScript: o guia definitivo. Porto Alegre: Bookman, 2013. Disponível em: https://integrada.minhabiblioteca.com.br/\#/books/9788565837484/. Acesso em: 31 jul. 2020.\end{itemize}}\\
\multicolumn{4}{|p{\dimexpr 16cm + 6\tabcolsep\relax}|}{\begin{itemize}\item FLATSCHART, Fábio. HTML5 Embarque imediato. Rio de Janeiro: Brasport, 2011. Disponível em: https://plataforma.bvirtual.com.br/Leitor/Publicacao/160673/epub/0. Acesso em: 31 jul. 2020.\end{itemize}}\\
\multicolumn{4}{|p{\dimexpr 16cm + 6\tabcolsep\relax}|}{\begin{itemize}\item SEGURADO, Valquiria Santos. Projeto de interface com o usuário. 1. ed. São Paulo: Pearson Education do Brasil, 2016. Disponível em: https://plataforma.bvirtual.com.br/Leitor/Publicacao/124143/pdf/0. Acesso em: 31 jul. 2020.\end{itemize}}\\
\multicolumn{4}{|p{16cm}|}{}\\
\hline
\multicolumn{4}{|p{16cm}|}{Bibliografia Complementar:}\\*
\multicolumn{4}{|p{\dimexpr 16cm + 6\tabcolsep\relax}|}{\begin{itemize}\item BANIN, Sérgio Luiz. Python 3: conceitos e aplicações: uma abordagem didática. São Paulo: Érica, 2018. Disponível em: https://integrada.minhabiblioteca.com.br/\#/books/9788536530253. Acesso em: 02 ago. 2020.\end{itemize}}\\
\multicolumn{4}{|p{\dimexpr 16cm + 6\tabcolsep\relax}|}{\begin{itemize}\item CAELUM. Apostila do curso: Desenvolvimento Web com HTML, CSS e JavaScript. 2020. Disponível em: https://www.caelum.com.br/apostila-html-css-javascript/. Acesso em: 31 jul. 2020.\end{itemize}}\\
\multicolumn{4}{|p{\dimexpr 16cm + 6\tabcolsep\relax}|}{\begin{itemize}\item HAVERBEKE, Marijn. Eloquent JavaScript: A Modern Introduction to Programming. 3. ed., 2018. Disponível em: https://eloquentjavascript.net/. Acesso em: 02 ago. 2020.\end{itemize}}\\
\multicolumn{4}{|p{\dimexpr 16cm + 6\tabcolsep\relax}|}{\begin{itemize}\item MILETTO, Evandro Manara; BERTAGNOLLI, Silvia de Castro. Desenvolvimento de Software II: Introdução ao Desenvolvimento Web com HTML, CSS, JavaScript e PHP. Porto Alegre: Bookman, 2014. Disponível em: https://integrada.minhabiblioteca.com.br/\#/books/9788582601969. Acesso em: 02 ago. 2020.\end{itemize}}\\
\multicolumn{4}{|p{\dimexpr 16cm + 6\tabcolsep\relax}|}{\begin{itemize}\item SOUSA, Roque Fernando Marcos. Canvas HTML5. Rio de Janeiro: Brasport, 2013. Disponível em: https://plataforma.bvirtual.com.br/Leitor/Publicacao/160686/epub/0. Acesso em: 02 ago. 2020.\end{itemize}}\\
\hline
\end{longtable}
\end{center}

\clearpage
\begin{center}
\begin{tabular}{lccr}
 \multirow{3}{*}{\includegraphics[height=2.7cm]{brasao.png}} &
 \multicolumn{2}{c}{\bfseries UNIVERSIDADE FEDERAL DE OURO PRETO} &
 \ \ \ \ \multirow{3}{*}{\includegraphics[height=2.7cm]{ufop.png}} \\
 & \multicolumn{2}{c}{\bfseries PRÓ-REITORIA DE GRADUAÇÃO} & \\
 & \multicolumn{2}{c}{} & \\
 & \multicolumn{2}{c}{\Large\bfseries PROGRAMA DE DISCIPLINA} & \\
\end{tabular}
\end{center}

\begin{center}
\begin{longtable}{|p{4cm}|p{4cm}|p{4cm}|p{4cm}|}
\hline
\multicolumn{3}{|p{12cm}|}{Nome do Componente Curricular em Português:} &
\multicolumn{1}{p{4cm}|}{Código:} \\ 
\multicolumn{3}{|p{12cm}|}{\textbf{Redes Complexas}} &
\textbf{BCC124}\\ 
\multicolumn{3}{|p{12cm}|}{Nome do Componente Curricular em Inglês:} & \\ 
\multicolumn{3}{|p{12cm}|}{\textbf{Complex Networks}} & \\ 
\hline
\multicolumn{3}{|p{12cm}|}{Nome e Sigla do Departamento} & Unidade Acadêmica: \\ 
\multicolumn{3}{|p{12cm}|}{Departamento de Computação (DECOM)} & {ICEB} \\ 
\hline
\multicolumn{4}{|p{16cm}|}{Modalidade de Oferta:
[X] presencial \hspace{0.6cm}
[ ] semipresencial \hspace{0.6cm}
[ ] à distância}\\
\hline
\multicolumn{2}{|p{8cm}|}{Carga horária semestral} &
\multicolumn{2}{p{8cm}|}{Carga horária semanal}\\
\hline
\multicolumn{1}{|p{4cm}|}{Total} &
\multicolumn{1}{p{4cm}|}{Extensionista} &
\multicolumn{1}{p{4cm}|}{Teórica} &
\multicolumn{1}{p{4cm}|}{Prática} \\ 
\multicolumn{1}{|p{4cm}|}{60\,horas} &
\multicolumn{1}{p{4cm}|}{0\;horas} &
\multicolumn{1}{p{4cm}|}{4\;horas/aula} &
\multicolumn{1}{p{4cm}|}{0\;horas/aula} \\ 
\hline
\multicolumn{4}{|p{16cm}|}{Ementa:}\\
\multicolumn{4}{|p{16cm}|}{}\\
\multicolumn{4}{|p{\dimexpr 16cm + 6\tabcolsep\relax}|}{Introdução e conceitos básicos; Redes direcionadas e com pesos; Modelos e algoritmos de geração de redes complexas; Correlações de grau; Robustez em redes; Detecção de Comunidades; Fenômenos dinâmicos em rede; Tópicos avançados.}\\
\multicolumn{4}{|p{16cm}|}{}\\
\hline
\multicolumn{4}{|p{16cm}|}{Conteúdo programático:}\\*
\multicolumn{4}{|p{\dimexpr 16cm + 6\tabcolsep\relax}|}{\begin{enumerate}[start=1]\item Introdução e conceitos básicos: tipos de redes e aplicações; matriz de adjacência, lista de adjacência, redes com pesos; caminhos e distâncias, medidas de centralidade, conectividade e transitividade\end{enumerate}}\\
\multicolumn{4}{|p{\dimexpr 16cm + 6\tabcolsep\relax}|}{\begin{enumerate}[start=2]\item Redes direcionadas e com pesos: redes direcionadas, a Web, PageRank\end{enumerate}}\\
\multicolumn{4}{|p{\dimexpr 16cm + 6\tabcolsep\relax}|}{\begin{enumerate}[start=3]\item Modelos e algoritmos de geração de redes complexas: redes aleatórias, redes de pequeno mundo, propriedade livre de escala, modelo Barabási-Albert, redes dinâmicas\end{enumerate}}\\
\multicolumn{4}{|p{\dimexpr 16cm + 6\tabcolsep\relax}|}{\begin{enumerate}[start=4]\item Correlações de grau: assortatividade\end{enumerate}}\\
\multicolumn{4}{|p{\dimexpr 16cm + 6\tabcolsep\relax}|}{\begin{enumerate}[start=5]\item Robustez em redes: percolação, falhas e ataques\end{enumerate}}\\
\multicolumn{4}{|p{\dimexpr 16cm + 6\tabcolsep\relax}|}{\begin{enumerate}[start=6]\item Detecção de Comunidades: conceitos e algoritmos\end{enumerate}}\\
\multicolumn{4}{|p{\dimexpr 16cm + 6\tabcolsep\relax}|}{\begin{enumerate}[start=7]\item Fenômenos dinâmicos em rede: espalhamento viral e propagação de informação, sincronização de osciladores acoplados\end{enumerate}}\\
\multicolumn{4}{|p{\dimexpr 16cm + 6\tabcolsep\relax}|}{\begin{enumerate}[start=8]\item Tópicos avançados\end{enumerate}}\\
\multicolumn{4}{|p{\dimexpr 16cm + 6\tabcolsep\relax}|}{}\\
\hline
\multicolumn{4}{|p{16cm}|}{Bibliografia Básica:}\\*
\multicolumn{4}{|p{\dimexpr 16cm + 6\tabcolsep\relax}|}{\begin{itemize}\item BARABÁSI, A.-L. Network science. Cambridge: Cambridge University Press, 2016. 475 p.\end{itemize}}\\
\multicolumn{4}{|p{\dimexpr 16cm + 6\tabcolsep\relax}|}{\begin{itemize}\item MENCZER, F.; FORTUNATO, S.; DAVIS, C. A First Course in Network Science. Cambridge: Cambridge University Press, 2020. 300 p.\end{itemize}}\\
\multicolumn{4}{|p{\dimexpr 16cm + 6\tabcolsep\relax}|}{\begin{itemize}\item NEWMAN, M. Networks. 2. ed. Oxford: Oxford University Press, 2018. 800 p.\end{itemize}}\\
\multicolumn{4}{|p{16cm}|}{}\\
\hline
\multicolumn{4}{|p{16cm}|}{Bibliografia Complementar:}\\*
\multicolumn{4}{|p{\dimexpr 16cm + 6\tabcolsep\relax}|}{\begin{itemize}\item BARABÁSI, A.-L. Linked: how everything is connected to everything else and what it means for business, science and everyday life. New York: Plume, 2003. 294 p.\end{itemize}}\\
\multicolumn{4}{|p{\dimexpr 16cm + 6\tabcolsep\relax}|}{\begin{itemize}\item BARRAT, A.; MELY, M. B.; VESPIGNANI, A. Dynamical Processes on Complex Networks. Cambridge: Cambridge University Press, 2012. 361 p.\end{itemize}}\\
\multicolumn{4}{|p{\dimexpr 16cm + 6\tabcolsep\relax}|}{\begin{itemize}\item EASLEY, D.; KLEINBERG, J. Networks, Crowds, and Markets: Reasoning about a Highly Connected World. Cambridge: Cambridge University Press, 2010. 727 p.\end{itemize}}\\
\multicolumn{4}{|p{\dimexpr 16cm + 6\tabcolsep\relax}|}{\begin{itemize}\item ESTRADA, E. The structure of complex networks: theory and applications. Oxford: Oxford University Press, 2012. 478 p.\end{itemize}}\\
\multicolumn{4}{|p{\dimexpr 16cm + 6\tabcolsep\relax}|}{\begin{itemize}\item JACKSON, M. O. Social and Economic Networks. Princeton: Princeton University Press, 2010. 504 p.\end{itemize}}\\
\hline
\end{longtable}
\end{center}

\clearpage
\begin{center}
\begin{tabular}{lccr}
 \multirow{3}{*}{\includegraphics[height=2.7cm]{brasao.png}} &
 \multicolumn{2}{c}{\bfseries UNIVERSIDADE FEDERAL DE OURO PRETO} &
 \ \ \ \ \multirow{3}{*}{\includegraphics[height=2.7cm]{ufop.png}} \\
 & \multicolumn{2}{c}{\bfseries PRÓ-REITORIA DE GRADUAÇÃO} & \\
 & \multicolumn{2}{c}{} & \\
 & \multicolumn{2}{c}{\Large\bfseries PROGRAMA DE DISCIPLINA} & \\
\end{tabular}
\end{center}

\begin{center}
\begin{longtable}{|p{4cm}|p{4cm}|p{4cm}|p{4cm}|}
\hline
\multicolumn{3}{|p{12cm}|}{Nome do Componente Curricular em Português:} &
\multicolumn{1}{p{4cm}|}{Código:} \\ 
\multicolumn{3}{|p{12cm}|}{\textbf{Engenharia de Dados}} &
\textbf{BCC125}\\ 
\multicolumn{3}{|p{12cm}|}{Nome do Componente Curricular em Inglês:} & \\ 
\multicolumn{3}{|p{12cm}|}{\textbf{Data Engineering}} & \\ 
\hline
\multicolumn{3}{|p{12cm}|}{Nome e Sigla do Departamento} & Unidade Acadêmica: \\ 
\multicolumn{3}{|p{12cm}|}{Departamento de Computação (DECOM)} & {ICEB} \\ 
\hline
\multicolumn{4}{|p{16cm}|}{Modalidade de Oferta:
[X] presencial \hspace{0.6cm}
[ ] semipresencial \hspace{0.6cm}
[ ] à distância}\\
\hline
\multicolumn{2}{|p{8cm}|}{Carga horária semestral} &
\multicolumn{2}{p{8cm}|}{Carga horária semanal}\\
\hline
\multicolumn{1}{|p{4cm}|}{Total} &
\multicolumn{1}{p{4cm}|}{Extensionista} &
\multicolumn{1}{p{4cm}|}{Teórica} &
\multicolumn{1}{p{4cm}|}{Prática} \\ 
\multicolumn{1}{|p{4cm}|}{60\,horas} &
\multicolumn{1}{p{4cm}|}{0\;horas} &
\multicolumn{1}{p{4cm}|}{2\;horas/aula} &
\multicolumn{1}{p{4cm}|}{2\;horas/aula} \\ 
\hline
\multicolumn{4}{|p{16cm}|}{Ementa:}\\
\multicolumn{4}{|p{16cm}|}{}\\
\multicolumn{4}{|p{\dimexpr 16cm + 6\tabcolsep\relax}|}{Virtualização. Nuvens computacionais. Clusters, máquinas virtuais, contêineres e orquestradores de contêineres. Ingestão de dados. Modelagem, armazenamento e consulta de dados. Pipeline de dados: desenho, implementação e testes.}\\
\multicolumn{4}{|p{16cm}|}{}\\
\hline
\multicolumn{4}{|p{16cm}|}{Conteúdo programático:}\\*
\multicolumn{4}{|p{\dimexpr 16cm + 6\tabcolsep\relax}|}{\begin{enumerate}[start=1]\item Apresentação do curso: programa, objetivos, bibliografia\end{enumerate}}\\
\multicolumn{4}{|p{\dimexpr 16cm + 6\tabcolsep\relax}|}{\begin{enumerate}[start=2]\item Conceitos de virtualização e apresentação dos principais ambientes de nuvem\end{enumerate}}\\
\multicolumn{4}{|p{\dimexpr 16cm + 6\tabcolsep\relax}|}{\begin{enumerate}[start=3]\item Operações básicas para instalação e gerenciamento de um cluster de máquinas virtuais\end{enumerate}}\\
\multicolumn{4}{|p{\dimexpr 16cm + 6\tabcolsep\relax}|}{\begin{enumerate}[start=4]\item Operações básicas para instalação e gerenciamento de contêineres e seus orquestradores\end{enumerate}}\\
\multicolumn{4}{|p{\dimexpr 16cm + 6\tabcolsep\relax}|}{\begin{enumerate}[start=5]\item Revisional: conceitos básicos de Python, o padrão pub/sub e mensageria\end{enumerate}}\\
\multicolumn{4}{|p{\dimexpr 16cm + 6\tabcolsep\relax}|}{\begin{enumerate}[start=6]\item Notebooks: conceito, ferramentas e exemplos de boas práticas\end{enumerate}}\\
\multicolumn{4}{|p{\dimexpr 16cm + 6\tabcolsep\relax}|}{\begin{enumerate}[start=7]\item Ingestão de dados: tipos de dados, arquiteturas existentes e ferramental de apoio\end{enumerate}}\\
\multicolumn{4}{|p{\dimexpr 16cm + 6\tabcolsep\relax}|}{\begin{enumerate}[start=8]\item Modelagem de dados: relacional, não-relacional e multidimensional\end{enumerate}}\\
\multicolumn{4}{|p{\dimexpr 16cm + 6\tabcolsep\relax}|}{\begin{enumerate}[start=9]\item Armazenamento de dados: ingestão de dados em data lakes ou data lakehouses\end{enumerate}}\\
\multicolumn{4}{|p{\dimexpr 16cm + 6\tabcolsep\relax}|}{\begin{enumerate}[start=10]\item Consulta de dados: ferramentas, exemplos e operacionalização\end{enumerate}}\\
\multicolumn{4}{|p{\dimexpr 16cm + 6\tabcolsep\relax}|}{\begin{enumerate}[start=11]\item Pipeline de dados: conceitos e etapas que compõem um pipeline de dados; principais ferramentas para construção e gerenciamento de um pipeline de dados em clusters virtuais; gargalos no pipeline de dados e como atenuá-los; exemplos de pipeline de dados; desenho, implementação e teste de alguns pipelines de dados\end{enumerate}}\\
\multicolumn{4}{|p{\dimexpr 16cm + 6\tabcolsep\relax}|}{}\\
\hline
\multicolumn{4}{|p{16cm}|}{Bibliografia Básica:}\\*
\multicolumn{4}{|p{\dimexpr 16cm + 6\tabcolsep\relax}|}{\begin{itemize}\item KLEPPMANN, Martin. Designing Data-Intensive Applications: The Big Ideas Behind Reliable, Scalable, and Maintainable Systems. 1. ed. São Paulo: Oreilly \& Assoc, 2015.\end{itemize}}\\
\multicolumn{4}{|p{\dimexpr 16cm + 6\tabcolsep\relax}|}{\begin{itemize}\item KUKREJA, Manoj. Data Engineering with Apache Spark, Delta Lake, and Lakehouse: Create scalable pipelines that ingest, curate, and aggregate complex data in a timely and secure way. 1. ed. Packt Publishing, 2021.\end{itemize}}\\
\multicolumn{4}{|p{\dimexpr 16cm + 6\tabcolsep\relax}|}{\begin{itemize}\item EAGAR, Gareth. Data Engineering with AWS: Learn how to design and build cloud-based data transformation pipelines using AWS. 1. ed. Packt Publishing, 2021.\end{itemize}}\\
\multicolumn{4}{|p{\dimexpr 16cm + 6\tabcolsep\relax}|}{\begin{itemize}\item JÚNIOR, Armando Kolbe. Computação em Nuvem. 1. ed. Contentus, 2020. Disponível online em MinhaUFOP/BibliotecaDigital/BVirtualPearson: https://plataforma.bvirtual.com.br/Acervo/Publicacao/184851\end{itemize}}\\
\multicolumn{4}{|p{\dimexpr 16cm + 6\tabcolsep\relax}|}{\begin{itemize}\item VITALINO, Jeferson Fernando Noronha; CASTRO, Marcus André Nunes. Descomplicando o Docker. 2. ed. Brasport, 2018. Disponível online em MinhaUFOP/BibliotecaDigital/BVirtualPearson: https://plataforma.bvirtual.com.br/Acervo/Publicacao/164690\end{itemize}}\\
\multicolumn{4}{|p{16cm}|}{}\\
\hline
\multicolumn{4}{|p{16cm}|}{Bibliografia Complementar:}\\*
\multicolumn{4}{|p{\dimexpr 16cm + 6\tabcolsep\relax}|}{\begin{itemize}\item HARENSLAK, Bas P.; RUITER, Julian Rutger. Data Pipelines with Apache Airflow. 1. ed. Manning Publications, 2021.\end{itemize}}\\
\multicolumn{4}{|p{\dimexpr 16cm + 6\tabcolsep\relax}|}{\begin{itemize}\item Vários tutoriais sobre Jupyter notebooks. Disponível em: https://docs.jupyter.org/en/latest/ . Acesso em: 07/12/2021.\end{itemize}}\\
\multicolumn{4}{|p{\dimexpr 16cm + 6\tabcolsep\relax}|}{\begin{itemize}\item Vários tutoriais sobre Kubernetes. Disponível em: https://kubernetes.io/pt-br/docs/tutorials/ . Acesso em: 07/12/2021.\end{itemize}}\\
\multicolumn{4}{|p{\dimexpr 16cm + 6\tabcolsep\relax}|}{\begin{itemize}\item KOEN, Semi. Architecting a Machine Learning Pipeline. TowardsDataScience, 2019. Disponível em: https://towardsdatascience.com/architecting-a-machine-learning-pipeline-a847f094d1c7 . Acesso em: 07/12/2021.\end{itemize}}\\
\multicolumn{4}{|p{\dimexpr 16cm + 6\tabcolsep\relax}|}{\begin{itemize}\item YOSKOVICH, Tomer. Build your Data-Lake using AWS S3 \& Athena. TowardsDataScience, 2019. Disponível em: https://medium.com/swlh/tutorial-build-your-data-lake-using-aws-s3-athena-150c1aaa44cf . Acesso em: 07/12/2021.\end{itemize}}\\
\multicolumn{4}{|p{\dimexpr 16cm + 6\tabcolsep\relax}|}{\begin{itemize}\item CARAMELLO, Ivan. Keeping Data Under (Your) Control with Apache Airflow. TowardsDataScience, 2022. Disponível em: https://medium.com/daitan-tech/keeping-data-under-your-control-with-apache-airflowa46f884f8c5d . Acesso em: 07/12/2021.\end{itemize}}\\
\multicolumn{4}{|p{\dimexpr 16cm + 6\tabcolsep\relax}|}{\begin{itemize}\item BLANCAS, Eduardo. On writing clean Jupyter notebooks. TowardsDataScience, 2021. Disponível em: https://medium.com/towards-data-science/on-writing-clean-jupyter-notebooks-abdf6c708c75 . Acesso em: 07/12/2021.\end{itemize}}\\
\hline
\end{longtable}
\end{center}

\clearpage
\begin{center}
\begin{tabular}{lccr}
 \multirow{3}{*}{\includegraphics[height=2.7cm]{brasao.png}} &
 \multicolumn{2}{c}{\bfseries UNIVERSIDADE FEDERAL DE OURO PRETO} &
 \ \ \ \ \multirow{3}{*}{\includegraphics[height=2.7cm]{ufop.png}} \\
 & \multicolumn{2}{c}{\bfseries PRÓ-REITORIA DE GRADUAÇÃO} & \\
 & \multicolumn{2}{c}{} & \\
 & \multicolumn{2}{c}{\Large\bfseries PROGRAMA DE DISCIPLINA} & \\
\end{tabular}
\end{center}

\begin{center}
\begin{longtable}{|p{4cm}|p{4cm}|p{4cm}|p{4cm}|}
\hline
\multicolumn{3}{|p{12cm}|}{Nome do Componente Curricular em Português:} &
\multicolumn{1}{p{4cm}|}{Código:} \\ 
\multicolumn{3}{|p{12cm}|}{\textbf{Visualização de Dados}} &
\textbf{BCC127}\\ 
\multicolumn{3}{|p{12cm}|}{Nome do Componente Curricular em Inglês:} & \\ 
\multicolumn{3}{|p{12cm}|}{\textbf{Data Visualization}} & \\ 
\hline
\multicolumn{3}{|p{12cm}|}{Nome e Sigla do Departamento} & Unidade Acadêmica: \\ 
\multicolumn{3}{|p{12cm}|}{Departamento de Computação (DECOM)} & {ICEB} \\ 
\hline
\multicolumn{4}{|p{16cm}|}{Modalidade de Oferta:
[X] presencial \hspace{0.6cm}
[ ] semipresencial \hspace{0.6cm}
[ ] à distância}\\
\hline
\multicolumn{2}{|p{8cm}|}{Carga horária semestral} &
\multicolumn{2}{p{8cm}|}{Carga horária semanal}\\
\hline
\multicolumn{1}{|p{4cm}|}{Total} &
\multicolumn{1}{p{4cm}|}{Extensionista} &
\multicolumn{1}{p{4cm}|}{Teórica} &
\multicolumn{1}{p{4cm}|}{Prática} \\ 
\multicolumn{1}{|p{4cm}|}{60\,horas} &
\multicolumn{1}{p{4cm}|}{0\;horas} &
\multicolumn{1}{p{4cm}|}{4\;horas/aula} &
\multicolumn{1}{p{4cm}|}{0\;horas/aula} \\ 
\hline
\multicolumn{4}{|p{16cm}|}{Ementa:}\\
\multicolumn{4}{|p{16cm}|}{}\\
\multicolumn{4}{|p{\dimexpr 16cm + 6\tabcolsep\relax}|}{Introdução à visualização de dados. Representação de dados. Visualização de dados escalares. Visualização de dados vetoriais. Visualização de dados tensoriais. Visualização de informação.}\\
\multicolumn{4}{|p{16cm}|}{}\\
\hline
\multicolumn{4}{|p{16cm}|}{Conteúdo programático:}\\*
\multicolumn{4}{|p{\dimexpr 16cm + 6\tabcolsep\relax}|}{\begin{enumerate}[start=1]\item Introdução à visualização de dados: definições e tipos de visualização; pipeline de visualização; ferramentas computacionais\end{enumerate}}\\
\multicolumn{4}{|p{\dimexpr 16cm + 6\tabcolsep\relax}|}{\begin{enumerate}[start=2]\item Representação de dados: dados contínuos; amostragem e reconstrução de dados; conjunto de dados discretos (célula, malha e atributo)\end{enumerate}}\\
\multicolumn{4}{|p{\dimexpr 16cm + 6\tabcolsep\relax}|}{\begin{enumerate}[start=3]\item Visualização de dados escalares: mapeamento de cores; contorno; gráfico de altura\end{enumerate}}\\
\multicolumn{4}{|p{\dimexpr 16cm + 6\tabcolsep\relax}|}{\begin{enumerate}[start=4]\item Visualização de dados vetoriais: glifos vetoriais; codificação de cores; objetos de corrente (linha e tubo)\end{enumerate}}\\
\multicolumn{4}{|p{\dimexpr 16cm + 6\tabcolsep\relax}|}{\begin{enumerate}[start=5]\item Visualização de dados tensoriais: técnicas de redução de dimensionalidade; glifos tensoriais; rastreamento de fibras\end{enumerate}}\\
\multicolumn{4}{|p{\dimexpr 16cm + 6\tabcolsep\relax}|}{\begin{enumerate}[start=6]\item Visualização de informação: visualização de dados relacionais; visualização de dados multidimensionais\end{enumerate}}\\
\multicolumn{4}{|p{\dimexpr 16cm + 6\tabcolsep\relax}|}{}\\
\hline
\multicolumn{4}{|p{16cm}|}{Bibliografia Básica:}\\*
\multicolumn{4}{|p{\dimexpr 16cm + 6\tabcolsep\relax}|}{\begin{itemize}\item FOLEY, James D. Computer graphics: principles and practice in C. 2nd ed. Reading, Mass.: Addison Wesley, 1996. 1175 p.\end{itemize}}\\
\multicolumn{4}{|p{\dimexpr 16cm + 6\tabcolsep\relax}|}{\begin{itemize}\item GONZALEZ, Rafael C.; WOODS, Richard E. Digital image processing. 3. ed. Upper Saddle River, NJ: Pearson Prentice Hall, 2008. 954 p.\end{itemize}}\\
\multicolumn{4}{|p{\dimexpr 16cm + 6\tabcolsep\relax}|}{\begin{itemize}\item HAN, Jiawei; KAMBER, Micheline; PEI, Jian. Data mining: concepts and techniques. 3. ed. Boston: Elsevier, 2012. 703 p.\end{itemize}}\\
\multicolumn{4}{|p{\dimexpr 16cm + 6\tabcolsep\relax}|}{\begin{itemize}\item PEDRINI, Hélio; SCHWARTZ, William Robson. Análise de imagens digitais: princípios, algoritmos e aplicações. São Paulo: Thomson Learning, 2008. 508 p.\end{itemize}}\\
\multicolumn{4}{|p{\dimexpr 16cm + 6\tabcolsep\relax}|}{\begin{itemize}\item SCHROEDER, Will; MARTIN, Ken; LORENSE, Bill. The visualization toolkit. New Jersey: Prentice-Hall, 1996. 826 p.\end{itemize}}\\
\multicolumn{4}{|p{16cm}|}{}\\
\hline
\multicolumn{4}{|p{16cm}|}{Bibliografia Complementar:}\\*
\multicolumn{4}{|p{\dimexpr 16cm + 6\tabcolsep\relax}|}{\begin{itemize}\item HAYKIN, Simon S. Redes neurais: princípios e prática. 2. ed. Porto Alegre: Bookman, 2001. 900 p.\end{itemize}}\\
\multicolumn{4}{|p{\dimexpr 16cm + 6\tabcolsep\relax}|}{\begin{itemize}\item KIRK, Andy. Data visualisation: a handbook for data driven design. 2. ed. London: Sage Publications Ltd, 2019. 328 p.\end{itemize}}\\
\multicolumn{4}{|p{\dimexpr 16cm + 6\tabcolsep\relax}|}{\begin{itemize}\item LIM, Tee T.; SMITS, Alexander J. Flow visualization: techniques and examples. 2. ed. London: Imperial College Press, 2012. 442 p.\end{itemize}}\\
\multicolumn{4}{|p{\dimexpr 16cm + 6\tabcolsep\relax}|}{\begin{itemize}\item MUNZNER, Tamara. Visualization analysis and design. 1. ed. New York: A K Peters/CRC Press, 2014. 428 p.\end{itemize}}\\
\multicolumn{4}{|p{\dimexpr 16cm + 6\tabcolsep\relax}|}{\begin{itemize}\item SPENCE, Robert. Information visualization: an introduction. 3. ed. London: Springer International Publishing, 2014. 346 p.\end{itemize}}\\
\multicolumn{4}{|p{\dimexpr 16cm + 6\tabcolsep\relax}|}{\begin{itemize}\item TELEA, Alexandru C. Data visualization: principles and practice. 2. ed. New York: A K Peters/CRC Press, 2014. 617 p.\end{itemize}}\\
\multicolumn{4}{|p{\dimexpr 16cm + 6\tabcolsep\relax}|}{\begin{itemize}\item TIERNY, Julien. Topological data analysis for scientific visualization. 1. ed. Switzerland: Springer International Publishing AG, 2017. 165 p.\end{itemize}}\\
\hline
\end{longtable}
\end{center}

\clearpage
\begin{center}
\begin{tabular}{lccr}
 \multirow{3}{*}{\includegraphics[height=2.7cm]{brasao.png}} &
 \multicolumn{2}{c}{\bfseries UNIVERSIDADE FEDERAL DE OURO PRETO} &
 \ \ \ \ \multirow{3}{*}{\includegraphics[height=2.7cm]{ufop.png}} \\
 & \multicolumn{2}{c}{\bfseries PRÓ-REITORIA DE GRADUAÇÃO} & \\
 & \multicolumn{2}{c}{} & \\
 & \multicolumn{2}{c}{\Large\bfseries PROGRAMA DE DISCIPLINA} & \\
\end{tabular}
\end{center}

\begin{center}
\begin{longtable}{|p{4cm}|p{4cm}|p{4cm}|p{4cm}|}
\hline
\multicolumn{3}{|p{12cm}|}{Nome do Componente Curricular em Português:} &
\multicolumn{1}{p{4cm}|}{Código:} \\ 
\multicolumn{3}{|p{12cm}|}{\textbf{Métodos Numéricos para Equações Diferenciais}} &
\textbf{BCC134}\\ 
\multicolumn{3}{|p{12cm}|}{Nome do Componente Curricular em Inglês:} & \\ 
\multicolumn{3}{|p{12cm}|}{\textbf{Numerical Methods for Differential Equations}} & \\ 
\hline
\multicolumn{3}{|p{12cm}|}{Nome e Sigla do Departamento} & Unidade Acadêmica: \\ 
\multicolumn{3}{|p{12cm}|}{Departamento de Computação (DECOM)} & {ICEB} \\ 
\hline
\multicolumn{4}{|p{16cm}|}{Modalidade de Oferta:
[X] presencial \hspace{0.6cm}
[ ] semipresencial \hspace{0.6cm}
[ ] à distância}\\
\hline
\multicolumn{2}{|p{8cm}|}{Carga horária semestral} &
\multicolumn{2}{p{8cm}|}{Carga horária semanal}\\
\hline
\multicolumn{1}{|p{4cm}|}{Total} &
\multicolumn{1}{p{4cm}|}{Extensionista} &
\multicolumn{1}{p{4cm}|}{Teórica} &
\multicolumn{1}{p{4cm}|}{Prática} \\ 
\multicolumn{1}{|p{4cm}|}{60\,horas} &
\multicolumn{1}{p{4cm}|}{0\;horas} &
\multicolumn{1}{p{4cm}|}{4\;horas/aula} &
\multicolumn{1}{p{4cm}|}{0\;horas/aula} \\ 
\hline
\multicolumn{4}{|p{16cm}|}{Ementa:}\\
\multicolumn{4}{|p{16cm}|}{}\\
\multicolumn{4}{|p{\dimexpr 16cm + 6\tabcolsep\relax}|}{Problemas de valor inicial para equações diferenciais ordinárias. Problemas de valor de contorno para equações diferenciais ordinárias. Soluções numéricas para equações diferenciais parciais.}\\
\multicolumn{4}{|p{16cm}|}{}\\
\hline
\multicolumn{4}{|p{16cm}|}{Conteúdo programático:}\\*
\multicolumn{4}{|p{\dimexpr 16cm + 6\tabcolsep\relax}|}{\begin{enumerate}[start=1]\item Problemas de valor inicial para equações diferenciais ordinárias: introdução aos problemas de valor inicial; método de Euler; métodos de Taylor de ordem superior; métodos de Runge-Kutta; métodos de passo múltiplo; métodos de passo múltiplo com tamanho de passo variável; equações de ordem superior e sistemas de equações diferenciais; exemplos de aplicação\end{enumerate}}\\
\multicolumn{4}{|p{\dimexpr 16cm + 6\tabcolsep\relax}|}{\begin{enumerate}[start=2]\item Problemas de valor de contorno para equações diferenciais ordinárias: introdução aos problemas de valor de contorno; método shooting linear; método shooting para problemas não lineares; método de diferenças finitas para problemas lineares; método de diferenças finitas para problemas não lineares; exemplos de aplicação\end{enumerate}}\\
\multicolumn{4}{|p{\dimexpr 16cm + 6\tabcolsep\relax}|}{\begin{enumerate}[start=3]\item Soluções numéricas para equações diferenciais parciais: introdução às equações diferenciais parciais; equações diferenciais parciais elípticas; equações diferenciais parciais parabólicas; equações diferenciais parciais hiperbólicas; exemplos de aplicação\end{enumerate}}\\
\multicolumn{4}{|p{\dimexpr 16cm + 6\tabcolsep\relax}|}{}\\
\hline
\multicolumn{4}{|p{16cm}|}{Bibliografia Básica:}\\*
\multicolumn{4}{|p{\dimexpr 16cm + 6\tabcolsep\relax}|}{\begin{itemize}\item BURDEN, R. L.; FAIRES, J. D.; BURDEN, A. M. Análise numérica. 10. ed. São Paulo: Cengage Learning, c2016. 879 p. ISBN 9788522123407.\end{itemize}}\\
\multicolumn{4}{|p{\dimexpr 16cm + 6\tabcolsep\relax}|}{\begin{itemize}\item CAMPOS FILHO, F. F. Algoritmos numéricos. Rio de Janeiro: LTC, 2001. xv, 383 p. ISBN 8521612656.\end{itemize}}\\
\multicolumn{4}{|p{\dimexpr 16cm + 6\tabcolsep\relax}|}{\begin{itemize}\item RUGGIERO, M. A. G.; LOPES, V. L. R. Cálculo numérico: aspectos teóricos e computacionais. 2. ed. São Paulo: Makron Books, c1997. xvi, 406 p. ISBN 8534602042.\end{itemize}}\\
\multicolumn{4}{|p{\dimexpr 16cm + 6\tabcolsep\relax}|}{\begin{itemize}\item FRANCO, N. B. Cálculo numérico. São Paulo: Pearson Prentice Hall, 2006. 505 p. ISBN 8576050870.\end{itemize}}\\
\multicolumn{4}{|p{\dimexpr 16cm + 6\tabcolsep\relax}|}{\begin{itemize}\item DOERING, C. I.; LOPES, A. O. Equações diferenciais ordinárias. 2. ed. Rio de Janeiro: IMPA, 2007. 423 p. ISBN 9788524402395.\end{itemize}}\\
\multicolumn{4}{|p{\dimexpr 16cm + 6\tabcolsep\relax}|}{\begin{itemize}\item SMITH, G. D. Numerical solution of partial differential equations: finite difference methods. 3. ed. Oxford: Clarendon, 1992. 337 p. ISBN 0198596502.\end{itemize}}\\
\multicolumn{4}{|p{16cm}|}{}\\
\hline
\multicolumn{4}{|p{16cm}|}{Bibliografia Complementar:}\\*
\multicolumn{4}{|p{\dimexpr 16cm + 6\tabcolsep\relax}|}{\begin{itemize}\item QUARTERONI, A.; SACCO, R.; SALERI, F. Numerical Mathematics. 2. ed. New York: Springer, 2010. 657 p. ISBN 3642071015.\end{itemize}}\\
\multicolumn{4}{|p{\dimexpr 16cm + 6\tabcolsep\relax}|}{\begin{itemize}\item QUARTERONI, A.; VALLI, A. Numerical Approximation of Partial Differential Equations. 1. ed. New York: Springer, 2008. 562 p. ISBN 3540852670.\end{itemize}}\\
\multicolumn{4}{|p{\dimexpr 16cm + 6\tabcolsep\relax}|}{\begin{itemize}\item GOCKENBACH, M. S. Partial differential equations: analytical and numerical methods. Philadelphia: Society for Industrial and Applied Mathematics, c2002. xxii, 614 p. ISBN 0898715180.\end{itemize}}\\
\hline
\end{longtable}
\end{center}

\clearpage
\begin{center}
\begin{tabular}{lccr}
 \multirow{3}{*}{\includegraphics[height=2.7cm]{brasao.png}} &
 \multicolumn{2}{c}{\bfseries UNIVERSIDADE FEDERAL DE OURO PRETO} &
 \ \ \ \ \multirow{3}{*}{\includegraphics[height=2.7cm]{ufop.png}} \\
 & \multicolumn{2}{c}{\bfseries PRÓ-REITORIA DE GRADUAÇÃO} & \\
 & \multicolumn{2}{c}{} & \\
 & \multicolumn{2}{c}{\Large\bfseries PROGRAMA DE DISCIPLINA} & \\
\end{tabular}
\end{center}

\begin{center}
\begin{longtable}{|p{4cm}|p{4cm}|p{4cm}|p{4cm}|}
\hline
\multicolumn{3}{|p{12cm}|}{Nome do Componente Curricular em Português:} &
\multicolumn{1}{p{4cm}|}{Código:} \\ 
\multicolumn{3}{|p{12cm}|}{\textbf{Lógica Aplicada à Computação}} &
\textbf{BCC135}\\ 
\multicolumn{3}{|p{12cm}|}{Nome do Componente Curricular em Inglês:} & \\ 
\multicolumn{3}{|p{12cm}|}{\textbf{Logics in Computer Science}} & \\ 
\hline
\multicolumn{3}{|p{12cm}|}{Nome e Sigla do Departamento} & Unidade Acadêmica: \\ 
\multicolumn{3}{|p{12cm}|}{Departamento de Computação (DECOM)} & {ICEB} \\ 
\hline
\multicolumn{4}{|p{16cm}|}{Modalidade de Oferta:
[X] presencial \hspace{0.6cm}
[ ] semipresencial \hspace{0.6cm}
[ ] à distância}\\
\hline
\multicolumn{2}{|p{8cm}|}{Carga horária semestral} &
\multicolumn{2}{p{8cm}|}{Carga horária semanal}\\
\hline
\multicolumn{1}{|p{4cm}|}{Total} &
\multicolumn{1}{p{4cm}|}{Extensionista} &
\multicolumn{1}{p{4cm}|}{Teórica} &
\multicolumn{1}{p{4cm}|}{Prática} \\ 
\multicolumn{1}{|p{4cm}|}{60\,horas} &
\multicolumn{1}{p{4cm}|}{0\;horas} &
\multicolumn{1}{p{4cm}|}{4\;horas/aula} &
\multicolumn{1}{p{4cm}|}{0\;horas/aula} \\ 
\hline
\multicolumn{4}{|p{16cm}|}{Ementa:}\\
\multicolumn{4}{|p{16cm}|}{}\\
\multicolumn{4}{|p{\dimexpr 16cm + 6\tabcolsep\relax}|}{Lógica de primeira ordem, o algoritmo DPLL e uso de SAT e SMT solvers. Lógica intuicionista e teoria de tipos. Uso de assistentes de prova. Indução e co-indução e seu uso para verificação de programas. Lógica Temporal e verificação de modelos.}\\
\multicolumn{4}{|p{16cm}|}{}\\
\hline
\multicolumn{4}{|p{16cm}|}{Conteúdo programático:}\\*
\multicolumn{4}{|p{\dimexpr 16cm + 6\tabcolsep\relax}|}{\begin{enumerate}[start=1]\item Apresentação da disciplina: motivação e critérios de avaliação\end{enumerate}}\\
\multicolumn{4}{|p{\dimexpr 16cm + 6\tabcolsep\relax}|}{\begin{enumerate}[start=2]\item Revisão de lógica proposicional: sintaxe, semântica e dedução natural\end{enumerate}}\\
\multicolumn{4}{|p{\dimexpr 16cm + 6\tabcolsep\relax}|}{\begin{enumerate}[start=3]\item Revisão de lógica de predicados: sintaxe, semântica e dedução natural\end{enumerate}}\\
\multicolumn{4}{|p{\dimexpr 16cm + 6\tabcolsep\relax}|}{\begin{enumerate}[start=4]\item O problema SAT e o algoritmo DPLL\end{enumerate}}\\
\multicolumn{4}{|p{\dimexpr 16cm + 6\tabcolsep\relax}|}{\begin{enumerate}[start=5]\item Modelagem de problemas usando SAT e uso de resolvedores SAT\end{enumerate}}\\
\multicolumn{4}{|p{\dimexpr 16cm + 6\tabcolsep\relax}|}{\begin{enumerate}[start=6]\item Satisfiability Modulo Theories (SMT) e o algoritmo DPLL-T\end{enumerate}}\\
\multicolumn{4}{|p{\dimexpr 16cm + 6\tabcolsep\relax}|}{\begin{enumerate}[start=7]\item Modelagem de problemas usando SMT e resolvedores SMT\end{enumerate}}\\
\multicolumn{4}{|p{\dimexpr 16cm + 6\tabcolsep\relax}|}{\begin{enumerate}[start=8]\item Lógica intuicionista e lambda-cálculo tipado: o isomorfismo de Curry-Howard\end{enumerate}}\\
\multicolumn{4}{|p{\dimexpr 16cm + 6\tabcolsep\relax}|}{\begin{enumerate}[start=9]\item Introdução ao uso de assistentes de provas baseado em teoria de tipos\end{enumerate}}\\
\multicolumn{4}{|p{\dimexpr 16cm + 6\tabcolsep\relax}|}{\begin{enumerate}[start=10]\item Revisão: indução matemática, provas por indução estrutural\end{enumerate}}\\
\multicolumn{4}{|p{\dimexpr 16cm + 6\tabcolsep\relax}|}{\begin{enumerate}[start=11]\item Indução sobre relações bem formadas\end{enumerate}}\\
\multicolumn{4}{|p{\dimexpr 16cm + 6\tabcolsep\relax}|}{\begin{enumerate}[start=12]\item Provas por co-indução; uso de co-indução em assistentes de provas\end{enumerate}}\\
\multicolumn{4}{|p{\dimexpr 16cm + 6\tabcolsep\relax}|}{\begin{enumerate}[start=13]\item Lógica temporal e verificação de modelos\end{enumerate}}\\
\multicolumn{4}{|p{\dimexpr 16cm + 6\tabcolsep\relax}|}{}\\
\hline
\multicolumn{4}{|p{16cm}|}{Bibliografia Básica:}\\*
\multicolumn{4}{|p{\dimexpr 16cm + 6\tabcolsep\relax}|}{\begin{itemize}\item HUTH, Michael; RYAN, Mark. Lógica para ciência da computação: modelagem e argumentação sobre sistemas. 2. ed. LTC, 2008.\end{itemize}}\\
\multicolumn{4}{|p{\dimexpr 16cm + 6\tabcolsep\relax}|}{\begin{itemize}\item BERTOT, Yves; CASTÉRAN, Pierre. Interactive Theorem Proving and Program Development: The Coq'Art – The calculus of inductive constructions. 1. ed. Springer-Verlag, 2002.\end{itemize}}\\
\multicolumn{4}{|p{\dimexpr 16cm + 6\tabcolsep\relax}|}{\begin{itemize}\item MITCHELL, John. Foundations for Programming Languages. MIT Press, 1996.\end{itemize}}\\
\multicolumn{4}{|p{16cm}|}{}\\
\hline
\multicolumn{4}{|p{16cm}|}{Bibliografia Complementar:}\\*
\multicolumn{4}{|p{\dimexpr 16cm + 6\tabcolsep\relax}|}{\begin{itemize}\item BIERE, Armin; HEULE, Martin. Handbook of Satisfiability: Volume 185, Frontiers in Artificial Intelligence and Applications. 1. ed. IOS Press.\end{itemize}}\\
\multicolumn{4}{|p{\dimexpr 16cm + 6\tabcolsep\relax}|}{\begin{itemize}\item CLARKE, Edmund; GRUMBERG, Orna. Model Checking. 1. ed. MIT Press, 1999.\end{itemize}}\\
\multicolumn{4}{|p{\dimexpr 16cm + 6\tabcolsep\relax}|}{\begin{itemize}\item BAIER, Christel; KATOEN, Joost-Pieter. Principles of Model Checking. 1. ed. MIT Press, 2008.\end{itemize}}\\
\multicolumn{4}{|p{\dimexpr 16cm + 6\tabcolsep\relax}|}{\begin{itemize}\item SANGIORGI, Davide. Introduction to Bisimulation and Coinduction. 1. ed. Cambridge University Press, 2011.\end{itemize}}\\
\hline
\end{longtable}
\end{center}

\clearpage
\begin{center}
\begin{tabular}{lccr}
 \multirow{3}{*}{\includegraphics[height=2.7cm]{brasao.png}} &
 \multicolumn{2}{c}{\bfseries UNIVERSIDADE FEDERAL DE OURO PRETO} &
 \ \ \ \ \multirow{3}{*}{\includegraphics[height=2.7cm]{ufop.png}} \\
 & \multicolumn{2}{c}{\bfseries PRÓ-REITORIA DE GRADUAÇÃO} & \\
 & \multicolumn{2}{c}{} & \\
 & \multicolumn{2}{c}{\Large\bfseries PROGRAMA DE DISCIPLINA} & \\
\end{tabular}
\end{center}

\begin{center}
\begin{longtable}{|p{4cm}|p{4cm}|p{4cm}|p{4cm}|}
\hline
\multicolumn{3}{|p{12cm}|}{Nome do Componente Curricular em Português:} &
\multicolumn{1}{p{4cm}|}{Código:} \\ 
\multicolumn{3}{|p{12cm}|}{\textbf{Projeto de Software de Tempo Real}} &
\textbf{BCC142}\\ 
\multicolumn{3}{|p{12cm}|}{Nome do Componente Curricular em Inglês:} & \\ 
\multicolumn{3}{|p{12cm}|}{\textbf{Real-Time Software Design}} & \\ 
\hline
\multicolumn{3}{|p{12cm}|}{Nome e Sigla do Departamento} & Unidade Acadêmica: \\ 
\multicolumn{3}{|p{12cm}|}{Departamento de Computação (DECOM)} & {ICEB} \\ 
\hline
\multicolumn{4}{|p{16cm}|}{Modalidade de Oferta:
[X] presencial \hspace{0.6cm}
[ ] semipresencial \hspace{0.6cm}
[ ] à distância}\\
\hline
\multicolumn{2}{|p{8cm}|}{Carga horária semestral} &
\multicolumn{2}{p{8cm}|}{Carga horária semanal}\\
\hline
\multicolumn{1}{|p{4cm}|}{Total} &
\multicolumn{1}{p{4cm}|}{Extensionista} &
\multicolumn{1}{p{4cm}|}{Teórica} &
\multicolumn{1}{p{4cm}|}{Prática} \\ 
\multicolumn{1}{|p{4cm}|}{60\,horas} &
\multicolumn{1}{p{4cm}|}{0\;horas} &
\multicolumn{1}{p{4cm}|}{2\;horas/aula} &
\multicolumn{1}{p{4cm}|}{2\;horas/aula} \\ 
\hline
\multicolumn{4}{|p{16cm}|}{Ementa:}\\
\multicolumn{4}{|p{16cm}|}{}\\
\multicolumn{4}{|p{\dimexpr 16cm + 6\tabcolsep\relax}|}{Introdução aos Sistemas Operacionais. Conceitos básicos de Sistemas Operacionais. Introdução aos Sistemas Operacionais de Tempo Real. Classificação dos Sistemas Operacionais. Concorrência. Escalonamento para Sistemas Operacionais. Escalonamento para Sistemas Operacionais de Tempo Real. Objetos de Sincronização e Seção Crítica. Comunicação entre Tarefas. Gerenciamento de Tempo. Métodos formais para Tempo Real: modelagem, linguagens, verificação. Desenvolvimento de sistemas de tempo real.}\\
\multicolumn{4}{|p{16cm}|}{}\\
\hline
\multicolumn{4}{|p{16cm}|}{Conteúdo programático:}\\*
\multicolumn{4}{|p{\dimexpr 16cm + 6\tabcolsep\relax}|}{\begin{enumerate}[start=1]\item Conceitos de Sistemas Operacionais: características dos sistemas operacionais, processos, tarefas, escalonadores, políticas de escalonamento, sincronização entre tarefas/processos, estados das tarefas e dos processos\end{enumerate}}\\
\multicolumn{4}{|p{\dimexpr 16cm + 6\tabcolsep\relax}|}{\begin{enumerate}[start=2]\item Introdução aos Sistemas Operacionais de Tempo Real: definições e características; motivação para sistemas operacionais de tempo real; tarefas periódicas e aperiódicas (tempo de chegada, tempo de pronto, tempo de execução, período, deadline); prioridades; classificação (Hard e Soft RTOS); variáveis globais voláteis e estáticas; preemptivo e não preemptivo; interrupções e exceções em tempo real; estados das tarefas de tempo real; bloco de controle das tarefas (BCT); medidas de desempenho\end{enumerate}}\\
\multicolumn{4}{|p{\dimexpr 16cm + 6\tabcolsep\relax}|}{\begin{enumerate}[start=3]\item Sistemas Concorrentes: conceitos e problemas\end{enumerate}}\\
\multicolumn{4}{|p{\dimexpr 16cm + 6\tabcolsep\relax}|}{\begin{enumerate}[start=4]\item Políticas de Escalonamento de Sistemas de Tempo Real\end{enumerate}}\\
\multicolumn{4}{|p{\dimexpr 16cm + 6\tabcolsep\relax}|}{\begin{enumerate}[start=5]\item Objetos de Sincronização e Seção Crítica\end{enumerate}}\\
\multicolumn{4}{|p{\dimexpr 16cm + 6\tabcolsep\relax}|}{\begin{enumerate}[start=6]\item Sincronismo entre Tarefas: espera ocupada, semáforos, mutex, deadlock (bloqueio perpétuo)\end{enumerate}}\\
\multicolumn{4}{|p{\dimexpr 16cm + 6\tabcolsep\relax}|}{\begin{enumerate}[start=7]\item Comunicação entre Tarefas: memória compartilhada, troca de mensagens\end{enumerate}}\\
\multicolumn{4}{|p{\dimexpr 16cm + 6\tabcolsep\relax}|}{\begin{enumerate}[start=8]\item Comunicação de Tempo Real: arquitetura, métodos de acesso, protocolos de tempo real\end{enumerate}}\\
\multicolumn{4}{|p{\dimexpr 16cm + 6\tabcolsep\relax}|}{\begin{enumerate}[start=9]\item Métodos Formais para Tempo Real: modelagem, linguagem, verificação\end{enumerate}}\\
\multicolumn{4}{|p{\dimexpr 16cm + 6\tabcolsep\relax}|}{\begin{enumerate}[start=10]\item Desenvolvimento de sistemas de tempo real\end{enumerate}}\\
\multicolumn{4}{|p{\dimexpr 16cm + 6\tabcolsep\relax}|}{}\\
\hline
\multicolumn{4}{|p{16cm}|}{Bibliografia Básica:}\\*
\multicolumn{4}{|p{\dimexpr 16cm + 6\tabcolsep\relax}|}{\begin{itemize}\item BERNARDIN, Gustavo Weber; BARRIQUELLO, Carlos Henrique. Sistemas operacionais de tempo real e sua aplicação em sistemas embarcados. Editora Blucher, 2019. ISBN 9788521213970. Disponível na Minha Biblioteca: https://plataforma.bvirtual.com.br/Leitor/Publicacao/169968/pdf/ (acessível via Minha UFOP – Biblioteca Digital).\end{itemize}}\\
\multicolumn{4}{|p{\dimexpr 16cm + 6\tabcolsep\relax}|}{\begin{itemize}\item LABROSSE, J. J. MicroC/OS-III, The Real-Time Kernel, or a High Performance, Scalable, ROMable, Preemptive, Multitasking Kernel for Microprocessors, Microcontrollers \& DSPs. Micrium Press, 2009. ISBN 978-0982337530 / 9788587918574. Disponível em: https://www.micrium.com/books/ucosiii/\end{itemize}}\\
\multicolumn{4}{|p{\dimexpr 16cm + 6\tabcolsep\relax}|}{\begin{itemize}\item TANENBAUM, Andrew S. Sistemas Operacionais Modernos. Editora Pearson, 2003. Disponível na Minha Biblioteca: https://plataforma.bvirtual.com.br/Leitor/Publicacao/36876/pdf/ (acessível via Minha UFOP – Biblioteca Digital).\end{itemize}}\\
\multicolumn{4}{|p{16cm}|}{}\\
\hline
\multicolumn{4}{|p{16cm}|}{Bibliografia Complementar:}\\*
\multicolumn{4}{|p{\dimexpr 16cm + 6\tabcolsep\relax}|}{\begin{itemize}\item SILBERSCHATZ, Abraham; GAGNE, Greg; GALVIN, Peter B. Fundamentos de Sistemas Operacionais. John Wiley \& Sons, 2008. ISBN 9780470128725. Disponível na Minha Biblioteca: https://integrada.minhabiblioteca.com.br/\#/books/978-85-216-3001-2 (acessível via Minha UFOP – Biblioteca Digital).\end{itemize}}\\
\multicolumn{4}{|p{\dimexpr 16cm + 6\tabcolsep\relax}|}{\begin{itemize}\item DEITEL, H. M.; DEITEL, P. J.; CHOFFNES, D. R. Sistemas Operacionais. Editora Pearson, 2005. ISBN 9788576050117. Disponível na Minha Biblioteca: https://plataforma.bvirtual.com.br/Leitor/Loader/315/pdf (acessível via Minha UFOP – Biblioteca Digital Virtual Pearson).\end{itemize}}\\
\multicolumn{4}{|p{\dimexpr 16cm + 6\tabcolsep\relax}|}{\begin{itemize}\item TANENBAUM, Andrew S. Organização estruturada de computadores. 6. ed. Disponível na Biblioteca Pearson: https://plataforma.bvirtual.com.br/Acervo/Publicacao/3825 (acessível via Minha UFOP – Biblioteca Digital).\end{itemize}}\\
\multicolumn{4}{|p{\dimexpr 16cm + 6\tabcolsep\relax}|}{\begin{itemize}\item HALLINAN, Christopher. Embedded Linux Primer: A Practical Real-World Approach. Prentice-Hall, 2010. ISBN 9780137017836.\end{itemize}}\\
\multicolumn{4}{|p{\dimexpr 16cm + 6\tabcolsep\relax}|}{\begin{itemize}\item ABBOTT, Doug. Linux for Embedded and Real-time Applications. Newnes, 2006. ISBN 9780750679329.\end{itemize}}\\
\multicolumn{4}{|p{\dimexpr 16cm + 6\tabcolsep\relax}|}{\begin{itemize}\item BURNS, Alan; WELLINGS, Andy. Real-Time Systems \& Programming Languages. Addison-Wesley, 1. ed. ISBN 9780201729887.\end{itemize}}\\
\hline
\end{longtable}
\end{center}

\clearpage
\begin{center}
\begin{tabular}{lccr}
 \multirow{3}{*}{\includegraphics[height=2.7cm]{brasao.png}} &
 \multicolumn{2}{c}{\bfseries UNIVERSIDADE FEDERAL DE OURO PRETO} &
 \ \ \ \ \multirow{3}{*}{\includegraphics[height=2.7cm]{ufop.png}} \\
 & \multicolumn{2}{c}{\bfseries PRÓ-REITORIA DE GRADUAÇÃO} & \\
 & \multicolumn{2}{c}{} & \\
 & \multicolumn{2}{c}{\Large\bfseries PROGRAMA DE DISCIPLINA} & \\
\end{tabular}
\end{center}

\begin{center}
\begin{longtable}{|p{4cm}|p{4cm}|p{4cm}|p{4cm}|}
\hline
\multicolumn{3}{|p{12cm}|}{Nome do Componente Curricular em Português:} &
\multicolumn{1}{p{4cm}|}{Código:} \\ 
\multicolumn{3}{|p{12cm}|}{\textbf{Estrutura de Dados II}} &
\textbf{BCC203}\\ 
\multicolumn{3}{|p{12cm}|}{Nome do Componente Curricular em Inglês:} & \\ 
\multicolumn{3}{|p{12cm}|}{\textbf{Data Structures II}} & \\ 
\hline
\multicolumn{3}{|p{12cm}|}{Nome e Sigla do Departamento} & Unidade Acadêmica: \\ 
\multicolumn{3}{|p{12cm}|}{Departamento de Computação (DECOM)} & {ICEB} \\ 
\hline
\multicolumn{4}{|p{16cm}|}{Modalidade de Oferta:
[X] presencial \hspace{0.6cm}
[ ] semipresencial \hspace{0.6cm}
[ ] à distância}\\
\hline
\multicolumn{2}{|p{8cm}|}{Carga horária semestral} &
\multicolumn{2}{p{8cm}|}{Carga horária semanal}\\
\hline
\multicolumn{1}{|p{4cm}|}{Total} &
\multicolumn{1}{p{4cm}|}{Extensionista} &
\multicolumn{1}{p{4cm}|}{Teórica} &
\multicolumn{1}{p{4cm}|}{Prática} \\ 
\multicolumn{1}{|p{4cm}|}{60\,horas} &
\multicolumn{1}{p{4cm}|}{0\;horas} &
\multicolumn{1}{p{4cm}|}{4\;horas/aula} &
\multicolumn{1}{p{4cm}|}{0\;horas/aula} \\ 
\hline
\multicolumn{4}{|p{16cm}|}{Ementa:}\\
\multicolumn{4}{|p{16cm}|}{}\\
\multicolumn{4}{|p{\dimexpr 16cm + 6\tabcolsep\relax}|}{Armazenamento em memória secundária. Ordenação em memória secundária. Pesquisa em memória secundária. Casamento de cadeias. Compressão de textos. Estruturas de dados espaciais.}\\
\multicolumn{4}{|p{16cm}|}{}\\
\hline
\multicolumn{4}{|p{16cm}|}{Conteúdo programático:}\\*
\multicolumn{4}{|p{\dimexpr 16cm + 6\tabcolsep\relax}|}{\begin{enumerate}[start=1]\item Armazenamento em memória secundária: dispositivos de armazenamento, acesso sequencial e aleatório\end{enumerate}}\\
\multicolumn{4}{|p{\dimexpr 16cm + 6\tabcolsep\relax}|}{\begin{enumerate}[start=2]\item Ordenação em memória secundária: intercalação balanceada de vários caminhos, seleção por substituição, intercalação polifásica, quicksort externo\end{enumerate}}\\
\multicolumn{4}{|p{\dimexpr 16cm + 6\tabcolsep\relax}|}{\begin{enumerate}[start=3]\item Pesquisa em memória secundária: acesso sequencial indexado, árvore B, árvore B*\end{enumerate}}\\
\multicolumn{4}{|p{\dimexpr 16cm + 6\tabcolsep\relax}|}{\begin{enumerate}[start=4]\item Casamento de cadeias: casamento exato de cadeias, casamento aproximado de cadeias\end{enumerate}}\\
\multicolumn{4}{|p{\dimexpr 16cm + 6\tabcolsep\relax}|}{\begin{enumerate}[start=5]\item Compressão de textos: compressão de textos em linguagem natural, codificação de Huffman\end{enumerate}}\\
\multicolumn{4}{|p{\dimexpr 16cm + 6\tabcolsep\relax}|}{\begin{enumerate}[start=6]\item Estruturas de dados espaciais: árvore de quadrante, árvore kD, árvore R\end{enumerate}}\\
\multicolumn{4}{|p{\dimexpr 16cm + 6\tabcolsep\relax}|}{}\\
\hline
\multicolumn{4}{|p{16cm}|}{Bibliografia Básica:}\\*
\multicolumn{4}{|p{\dimexpr 16cm + 6\tabcolsep\relax}|}{\begin{itemize}\item SEDGEWICK, Robert. Algorithms. 4. ed. Upper Saddle River: Addison Wesley, 2011.\end{itemize}}\\
\multicolumn{4}{|p{\dimexpr 16cm + 6\tabcolsep\relax}|}{\begin{itemize}\item ZIVIANI, Nivio. Projeto de algoritmos: com implementações em Pascal e C. 5. ed. São Paulo: Pioneira, 2000.\end{itemize}}\\
\multicolumn{4}{|p{\dimexpr 16cm + 6\tabcolsep\relax}|}{\begin{itemize}\item CORMEN, Thomas H. ... [et al.]; tradução Vandenberg D. de Souza. Algoritmos: teoria e prática. Rio de Janeiro: Campus, 2002.\end{itemize}}\\
\multicolumn{4}{|p{16cm}|}{}\\
\hline
\multicolumn{4}{|p{16cm}|}{Bibliografia Complementar:}\\*
\multicolumn{4}{|p{\dimexpr 16cm + 6\tabcolsep\relax}|}{\begin{itemize}\item ZIVIANI, Nivio; BOTELHO, Fabiano C. Projeto de algoritmos: com implementações em Java e C++. São Paulo: Cengage Learning, 2015. (item inexistente na UFOP)\end{itemize}}\\
\multicolumn{4}{|p{\dimexpr 16cm + 6\tabcolsep\relax}|}{\begin{itemize}\item Bancos de dados geográficos. Curitiba: MundoGeo, 2005.\end{itemize}}\\
\multicolumn{4}{|p{\dimexpr 16cm + 6\tabcolsep\relax}|}{\begin{itemize}\item DEITEL, Harvey M. C++ how to program. Harvey M. Deitel, Paul J. Deitel. 6. ed. Boston: Prentice Hall, c2008.\end{itemize}}\\
\multicolumn{4}{|p{\dimexpr 16cm + 6\tabcolsep\relax}|}{\begin{itemize}\item MIZRAHI, Victorine Viviane. Treinamento em linguagem C. 2. ed. São Paulo: Pearson Education, c2010.\end{itemize}}\\
\multicolumn{4}{|p{\dimexpr 16cm + 6\tabcolsep\relax}|}{\begin{itemize}\item TENENBAUM, Aaron M. Estruturas de dados usando C. São Paulo: Makron Books, 1995.\end{itemize}}\\
\hline
\end{longtable}
\end{center}

\clearpage
\begin{center}
\begin{tabular}{lccr}
 \multirow{3}{*}{\includegraphics[height=2.7cm]{brasao.png}} &
 \multicolumn{2}{c}{\bfseries UNIVERSIDADE FEDERAL DE OURO PRETO} &
 \ \ \ \ \multirow{3}{*}{\includegraphics[height=2.7cm]{ufop.png}} \\
 & \multicolumn{2}{c}{\bfseries PRÓ-REITORIA DE GRADUAÇÃO} & \\
 & \multicolumn{2}{c}{} & \\
 & \multicolumn{2}{c}{\Large\bfseries PROGRAMA DE DISCIPLINA} & \\
\end{tabular}
\end{center}

\begin{center}
\begin{longtable}{|p{4cm}|p{4cm}|p{4cm}|p{4cm}|}
\hline
\multicolumn{3}{|p{12cm}|}{Nome do Componente Curricular em Português:} &
\multicolumn{1}{p{4cm}|}{Código:} \\ 
\multicolumn{3}{|p{12cm}|}{\textbf{Teoria da Computação}} &
\textbf{BCC244}\\ 
\multicolumn{3}{|p{12cm}|}{Nome do Componente Curricular em Inglês:} & \\ 
\multicolumn{3}{|p{12cm}|}{\textbf{Theory of Computation}} & \\ 
\hline
\multicolumn{3}{|p{12cm}|}{Nome e Sigla do Departamento} & Unidade Acadêmica: \\ 
\multicolumn{3}{|p{12cm}|}{Departamento de Computação (DECOM)} & {ICEB} \\ 
\hline
\multicolumn{4}{|p{16cm}|}{Modalidade de Oferta:
[X] presencial \hspace{0.6cm}
[ ] semipresencial \hspace{0.6cm}
[ ] à distância}\\
\hline
\multicolumn{2}{|p{8cm}|}{Carga horária semestral} &
\multicolumn{2}{p{8cm}|}{Carga horária semanal}\\
\hline
\multicolumn{1}{|p{4cm}|}{Total} &
\multicolumn{1}{p{4cm}|}{Extensionista} &
\multicolumn{1}{p{4cm}|}{Teórica} &
\multicolumn{1}{p{4cm}|}{Prática} \\ 
\multicolumn{1}{|p{4cm}|}{60\,horas} &
\multicolumn{1}{p{4cm}|}{0\;horas} &
\multicolumn{1}{p{4cm}|}{4\;horas/aula} &
\multicolumn{1}{p{4cm}|}{0\;horas/aula} \\ 
\hline
\multicolumn{4}{|p{16cm}|}{Ementa:}\\
\multicolumn{4}{|p{16cm}|}{}\\
\multicolumn{4}{|p{\dimexpr 16cm + 6\tabcolsep\relax}|}{Linguagens regulares, expressões regulares, autômatos de estados finitos. Linguagens e gramáticas livres de contexto, autômatos de pilha. Linguagens e gramáticas sensíveis ao contexto. Máquinas de Turing. Tese de Church-Turing. Computabilidade e decidibilidade. Hierarquia de Chomsky.}\\
\multicolumn{4}{|p{16cm}|}{}\\
\hline
\multicolumn{4}{|p{16cm}|}{Conteúdo programático:}\\*
\multicolumn{4}{|p{\dimexpr 16cm + 6\tabcolsep\relax}|}{\begin{enumerate}[start=1]\item Introdução: alfabetos, strings e linguagens\end{enumerate}}\\
\multicolumn{4}{|p{\dimexpr 16cm + 6\tabcolsep\relax}|}{\begin{enumerate}[start=2]\item Autômatos de Estados Finitos Deterministas e não Deterministas\end{enumerate}}\\
\multicolumn{4}{|p{\dimexpr 16cm + 6\tabcolsep\relax}|}{\begin{enumerate}[start=3]\item Expressões Regulares\end{enumerate}}\\
\multicolumn{4}{|p{\dimexpr 16cm + 6\tabcolsep\relax}|}{\begin{enumerate}[start=4]\item Minimização de Autômatos Finitos\end{enumerate}}\\
\multicolumn{4}{|p{\dimexpr 16cm + 6\tabcolsep\relax}|}{\begin{enumerate}[start=5]\item Propriedades de Linguagens Regulares\end{enumerate}}\\
\multicolumn{4}{|p{\dimexpr 16cm + 6\tabcolsep\relax}|}{\begin{enumerate}[start=6]\item Lema do Bombeamento para Linguagens Regulares (LRs)\end{enumerate}}\\
\multicolumn{4}{|p{\dimexpr 16cm + 6\tabcolsep\relax}|}{\begin{enumerate}[start=7]\item Gramáticas e Linguagens Livres de Contexto (LLC)\end{enumerate}}\\
\multicolumn{4}{|p{\dimexpr 16cm + 6\tabcolsep\relax}|}{\begin{enumerate}[start=8]\item Ambiguidade\end{enumerate}}\\
\multicolumn{4}{|p{\dimexpr 16cm + 6\tabcolsep\relax}|}{\begin{enumerate}[start=9]\item Propriedades de LLCs\end{enumerate}}\\
\multicolumn{4}{|p{\dimexpr 16cm + 6\tabcolsep\relax}|}{\begin{enumerate}[start=10]\item Autômatos de Pilha\end{enumerate}}\\
\multicolumn{4}{|p{\dimexpr 16cm + 6\tabcolsep\relax}|}{\begin{enumerate}[start=11]\item Forma normal de Chomsky\end{enumerate}}\\
\multicolumn{4}{|p{\dimexpr 16cm + 6\tabcolsep\relax}|}{\begin{enumerate}[start=12]\item Gramáticas Regulares e Gramáticas Sensíveis ao Contexto\end{enumerate}}\\
\multicolumn{4}{|p{\dimexpr 16cm + 6\tabcolsep\relax}|}{\begin{enumerate}[start=13]\item Lema do Bombeamento para LLCs\end{enumerate}}\\
\multicolumn{4}{|p{\dimexpr 16cm + 6\tabcolsep\relax}|}{\begin{enumerate}[start=14]\item Máquinas de Turing\end{enumerate}}\\
\multicolumn{4}{|p{\dimexpr 16cm + 6\tabcolsep\relax}|}{\begin{enumerate}[start=15]\item Tese de Church-Turing\end{enumerate}}\\
\multicolumn{4}{|p{\dimexpr 16cm + 6\tabcolsep\relax}|}{\begin{enumerate}[start=16]\item Problemas de Decisão\end{enumerate}}\\
\multicolumn{4}{|p{\dimexpr 16cm + 6\tabcolsep\relax}|}{\begin{enumerate}[start=17]\item Indecidibilidade do Problema da Parada\end{enumerate}}\\
\multicolumn{4}{|p{\dimexpr 16cm + 6\tabcolsep\relax}|}{\begin{enumerate}[start=18]\item Problemas decidíveis e não decidíveis sobre linguagens\end{enumerate}}\\
\multicolumn{4}{|p{\dimexpr 16cm + 6\tabcolsep\relax}|}{}\\
\hline
\multicolumn{4}{|p{16cm}|}{Bibliografia Básica:}\\*
\multicolumn{4}{|p{\dimexpr 16cm + 6\tabcolsep\relax}|}{\begin{itemize}\item VIEIRA, Newton José. Introdução aos fundamentos da computação: linguagens e máquinas. São Paulo: Thomson Pioneira, 2006.\end{itemize}}\\
\multicolumn{4}{|p{\dimexpr 16cm + 6\tabcolsep\relax}|}{\begin{itemize}\item SIPSER, Michael. Introdução à teoria da computação. São Paulo: Thomson Learning, 2007.\end{itemize}}\\
\multicolumn{4}{|p{\dimexpr 16cm + 6\tabcolsep\relax}|}{\begin{itemize}\item SUDKAMP, Thomas A. Languages and machines: an introduction to the theory of computer science. 2. ed. Reading: Addison Wesley, c1997.\end{itemize}}\\
\multicolumn{4}{|p{16cm}|}{}\\
\hline
\multicolumn{4}{|p{16cm}|}{Bibliografia Complementar:}\\*
\multicolumn{4}{|p{\dimexpr 16cm + 6\tabcolsep\relax}|}{\begin{itemize}\item MENDELSON, Elliott. Introduction to mathematical logic. 3rd ed. Pacific Grove, CA: Wadsworth \& Brooks/Cole Advanced Books \& Software, 1987.\end{itemize}}\\
\multicolumn{4}{|p{\dimexpr 16cm + 6\tabcolsep\relax}|}{\begin{itemize}\item LEWIS, Harry R. Elementos de teoria da computação. 2. ed. Porto Alegre: Bookman, 2000.\end{itemize}}\\
\multicolumn{4}{|p{\dimexpr 16cm + 6\tabcolsep\relax}|}{\begin{itemize}\item EPSTEIN, Richard L. Computability: computable functions, logic and the foundations of mathematics. Pacific Grove: Wadsworth, c1989.\end{itemize}}\\
\multicolumn{4}{|p{\dimexpr 16cm + 6\tabcolsep\relax}|}{\begin{itemize}\item DAVIS, Martin. Computability and unsolvability. New York: Dover, c1982.\end{itemize}}\\
\multicolumn{4}{|p{\dimexpr 16cm + 6\tabcolsep\relax}|}{\begin{itemize}\item Handbook of theoretical computer science. Amsterdã: Elsevier.\end{itemize}}\\
\hline
\end{longtable}
\end{center}

\clearpage
\begin{center}
\begin{tabular}{lccr}
 \multirow{3}{*}{\includegraphics[height=2.7cm]{brasao.png}} &
 \multicolumn{2}{c}{\bfseries UNIVERSIDADE FEDERAL DE OURO PRETO} &
 \ \ \ \ \multirow{3}{*}{\includegraphics[height=2.7cm]{ufop.png}} \\
 & \multicolumn{2}{c}{\bfseries PRÓ-REITORIA DE GRADUAÇÃO} & \\
 & \multicolumn{2}{c}{} & \\
 & \multicolumn{2}{c}{\Large\bfseries PROGRAMA DE DISCIPLINA} & \\
\end{tabular}
\end{center}

\begin{center}
\begin{longtable}{|p{4cm}|p{4cm}|p{4cm}|p{4cm}|}
\hline
\multicolumn{3}{|p{12cm}|}{Nome do Componente Curricular em Português:} &
\multicolumn{1}{p{4cm}|}{Código:} \\ 
\multicolumn{3}{|p{12cm}|}{\textbf{Sistemas Operacionais}} &
\textbf{BCC264}\\ 
\multicolumn{3}{|p{12cm}|}{Nome do Componente Curricular em Inglês:} & \\ 
\multicolumn{3}{|p{12cm}|}{\textbf{Operational Systems}} & \\ 
\hline
\multicolumn{3}{|p{12cm}|}{Nome e Sigla do Departamento} & Unidade Acadêmica: \\ 
\multicolumn{3}{|p{12cm}|}{Departamento de Computação (DECOM)} & {ICEB} \\ 
\hline
\multicolumn{4}{|p{16cm}|}{Modalidade de Oferta:
[X] presencial \hspace{0.6cm}
[ ] semipresencial \hspace{0.6cm}
[ ] à distância}\\
\hline
\multicolumn{2}{|p{8cm}|}{Carga horária semestral} &
\multicolumn{2}{p{8cm}|}{Carga horária semanal}\\
\hline
\multicolumn{1}{|p{4cm}|}{Total} &
\multicolumn{1}{p{4cm}|}{Extensionista} &
\multicolumn{1}{p{4cm}|}{Teórica} &
\multicolumn{1}{p{4cm}|}{Prática} \\ 
\multicolumn{1}{|p{4cm}|}{60\,horas} &
\multicolumn{1}{p{4cm}|}{0\;horas} &
\multicolumn{1}{p{4cm}|}{4\;horas/aula} &
\multicolumn{1}{p{4cm}|}{0\;horas/aula} \\ 
\hline
\multicolumn{4}{|p{16cm}|}{Ementa:}\\
\multicolumn{4}{|p{16cm}|}{}\\
\multicolumn{4}{|p{\dimexpr 16cm + 6\tabcolsep\relax}|}{Visão geral e histórico de sistemas operacionais; princípios de sistemas operacionais; gerenciamento de processos; uso de API (Application Programming Interface) de threads; gerenciamento de memória; gerenciamento de dispositivos; segurança e proteção; sistemas de arquivos.}\\
\multicolumn{4}{|p{16cm}|}{}\\
\hline
\multicolumn{4}{|p{16cm}|}{Conteúdo programático:}\\*
\multicolumn{4}{|p{\dimexpr 16cm + 6\tabcolsep\relax}|}{\begin{enumerate}[start=1]\item Introdução: definição, história, conceitos básicos, as chamadas de sistema e estrutura\end{enumerate}}\\
\multicolumn{4}{|p{\dimexpr 16cm + 6\tabcolsep\relax}|}{\begin{enumerate}[start=2]\item Processos e Threads: comunicação inter processos (IPC), scheduling e impasses (deadlocks)\end{enumerate}}\\
\multicolumn{4}{|p{\dimexpr 16cm + 6\tabcolsep\relax}|}{\begin{enumerate}[start=3]\item Gerenciamento de memória: swap, memória virtual, algoritmos de substituição e segmentação\end{enumerate}}\\
\multicolumn{4}{|p{\dimexpr 16cm + 6\tabcolsep\relax}|}{\begin{enumerate}[start=4]\item Entrada/Saída: discos, clocks, redes e terminais\end{enumerate}}\\
\multicolumn{4}{|p{\dimexpr 16cm + 6\tabcolsep\relax}|}{\begin{enumerate}[start=5]\item Sistemas de Arquivos: diretórios, implementação do sistema de arquivos e exemplos\end{enumerate}}\\
\multicolumn{4}{|p{\dimexpr 16cm + 6\tabcolsep\relax}|}{\begin{enumerate}[start=6]\item Segurança de redes: conceitos básicos de criptografia, autenticação, ataques e mecanismos de proteção\end{enumerate}}\\
\multicolumn{4}{|p{\dimexpr 16cm + 6\tabcolsep\relax}|}{\begin{enumerate}[start=7]\item Estudos de Casos: visão geral, processos, gerenciamento de memória, E/S, sistema de arquivos e segurança\end{enumerate}}\\
\multicolumn{4}{|p{\dimexpr 16cm + 6\tabcolsep\relax}|}{}\\
\hline
\multicolumn{4}{|p{16cm}|}{Bibliografia Básica:}\\*
\multicolumn{4}{|p{\dimexpr 16cm + 6\tabcolsep\relax}|}{\begin{itemize}\item SILBERSCHATZ, Abraham. Fundamentos de sistemas operacionais. 8. ed. Rio de Janeiro: LTC, c2010.\end{itemize}}\\
\multicolumn{4}{|p{\dimexpr 16cm + 6\tabcolsep\relax}|}{\begin{itemize}\item SILBERSCHATZ, Abraham; GALVIN, Peter Baer; tradução e revisão técnica Carlos Camarão de Figueiredo. Sistemas operacionais: conceitos. São Paulo: Prentice-Hall, 2000.\end{itemize}}\\
\multicolumn{4}{|p{\dimexpr 16cm + 6\tabcolsep\relax}|}{\begin{itemize}\item TANENBAUM, Andrew S. Sistemas operacionais modernos. 2. ed. São Paulo: Prentice-Hall, 2003.\end{itemize}}\\
\multicolumn{4}{|p{16cm}|}{}\\
\hline
\multicolumn{4}{|p{16cm}|}{Bibliografia Complementar:}\\*
\multicolumn{4}{|p{\dimexpr 16cm + 6\tabcolsep\relax}|}{\begin{itemize}\item BACH, Maurice J. The design of the UNIX operating system. Englewood Cliffs, N.J.: Prentice-Hall, c1986.\end{itemize}}\\
\multicolumn{4}{|p{\dimexpr 16cm + 6\tabcolsep\relax}|}{\begin{itemize}\item STEVENS, W. Richard. Advanced programming in the UNIX environment. Reading, MA: Addison Wesley, c1992.\end{itemize}}\\
\multicolumn{4}{|p{\dimexpr 16cm + 6\tabcolsep\relax}|}{\begin{itemize}\item MACHADO, Francis B. Arquitetura de sistemas operacionais. 3. ed. Rio de Janeiro: LTC, 2002.\end{itemize}}\\
\multicolumn{4}{|p{\dimexpr 16cm + 6\tabcolsep\relax}|}{\begin{itemize}\item CÔRTES, Pedro Luiz. Sistemas operacionais: fundamentos. São Paulo: Érica, 2003.\end{itemize}}\\
\multicolumn{4}{|p{\dimexpr 16cm + 6\tabcolsep\relax}|}{\begin{itemize}\item DEITEL, Harvey M. An Introduction to operating systems. 2nd ed. Reading: Addison Wesley, 1990.\end{itemize}}\\
\hline
\end{longtable}
\end{center}

\clearpage
\begin{center}
\begin{tabular}{lccr}
 \multirow{3}{*}{\includegraphics[height=2.7cm]{brasao.png}} &
 \multicolumn{2}{c}{\bfseries UNIVERSIDADE FEDERAL DE OURO PRETO} &
 \ \ \ \ \multirow{3}{*}{\includegraphics[height=2.7cm]{ufop.png}} \\
 & \multicolumn{2}{c}{\bfseries PRÓ-REITORIA DE GRADUAÇÃO} & \\
 & \multicolumn{2}{c}{} & \\
 & \multicolumn{2}{c}{\Large\bfseries PROGRAMA DE DISCIPLINA} & \\
\end{tabular}
\end{center}

\begin{center}
\begin{longtable}{|p{4cm}|p{4cm}|p{4cm}|p{4cm}|}
\hline
\multicolumn{3}{|p{12cm}|}{Nome do Componente Curricular em Português:} &
\multicolumn{1}{p{4cm}|}{Código:} \\ 
\multicolumn{3}{|p{12cm}|}{\textbf{Interação Humano-Computador}} &
\textbf{BCC324}\\ 
\multicolumn{3}{|p{12cm}|}{Nome do Componente Curricular em Inglês:} & \\ 
\multicolumn{3}{|p{12cm}|}{\textbf{Human-Computer Interaction}} & \\ 
\hline
\multicolumn{3}{|p{12cm}|}{Nome e Sigla do Departamento} & Unidade Acadêmica: \\ 
\multicolumn{3}{|p{12cm}|}{Departamento de Computação (DECOM)} & {ICEB} \\ 
\hline
\multicolumn{4}{|p{16cm}|}{Modalidade de Oferta:
[X] presencial \hspace{0.6cm}
[ ] semipresencial \hspace{0.6cm}
[ ] à distância}\\
\hline
\multicolumn{2}{|p{8cm}|}{Carga horária semestral} &
\multicolumn{2}{p{8cm}|}{Carga horária semanal}\\
\hline
\multicolumn{1}{|p{4cm}|}{Total} &
\multicolumn{1}{p{4cm}|}{Extensionista} &
\multicolumn{1}{p{4cm}|}{Teórica} &
\multicolumn{1}{p{4cm}|}{Prática} \\ 
\multicolumn{1}{|p{4cm}|}{60\,horas} &
\multicolumn{1}{p{4cm}|}{0\;horas} &
\multicolumn{1}{p{4cm}|}{4\;horas/aula} &
\multicolumn{1}{p{4cm}|}{0\;horas/aula} \\ 
\hline
\multicolumn{4}{|p{16cm}|}{Ementa:}\\
\multicolumn{4}{|p{16cm}|}{}\\
\multicolumn{4}{|p{\dimexpr 16cm + 6\tabcolsep\relax}|}{Interface e interação com o usuário. Qualidade em IHC. Abordagens teóricas em IHC. Introdução ao processo de design de sistemas interativos. Introdução ao processo de avaliação de sistemas interativos.}\\
\multicolumn{4}{|p{16cm}|}{}\\
\hline
\multicolumn{4}{|p{16cm}|}{Conteúdo programático:}\\*
\multicolumn{4}{|p{\dimexpr 16cm + 6\tabcolsep\relax}|}{\begin{enumerate}[start=1]\item Interface e Interação com o Usuário: conceito de interface e interação, importância da área de IHC no desenvolvimento de software, histórico, IHC como área multidisciplinar, benefícios de IHC, problemas de interação\end{enumerate}}\\
\multicolumn{4}{|p{\dimexpr 16cm + 6\tabcolsep\relax}|}{\begin{enumerate}[start=2]\item Qualidade em IHC: usabilidade, comunicabilidade, acessibilidade, user experience\end{enumerate}}\\
\multicolumn{4}{|p{\dimexpr 16cm + 6\tabcolsep\relax}|}{\begin{enumerate}[start=3]\item Abordagens Teóricas em IHC: psicologia experimental (Lei de Hick-Hyman e Lei de Fitts), engenharia cognitiva, engenharia semiótica\end{enumerate}}\\
\multicolumn{4}{|p{\dimexpr 16cm + 6\tabcolsep\relax}|}{\begin{enumerate}[start=4]\item Introdução ao Processo de Design de Sistemas Interativos: modelagem de tarefas, modelagem de interação, prototipagem, princípios e diretrizes para o design de IHC, estilos de interação\end{enumerate}}\\
\multicolumn{4}{|p{\dimexpr 16cm + 6\tabcolsep\relax}|}{\begin{enumerate}[start=5]\item Introdução ao Processo de Avaliação de Sistemas Interativos: planejamento da avaliação de IHC, avaliação por inspeção, testes com usuários\end{enumerate}}\\
\multicolumn{4}{|p{\dimexpr 16cm + 6\tabcolsep\relax}|}{}\\
\hline
\multicolumn{4}{|p{16cm}|}{Bibliografia Básica:}\\*
\multicolumn{4}{|p{\dimexpr 16cm + 6\tabcolsep\relax}|}{\begin{itemize}\item PREECE, Jennifer; ROGERS, Yvonne; SHARP, Helen; tradução de Viviane Possamai. Design de interação: além da interação homem-computador. Porto Alegre: Bookman, 2005.\end{itemize}}\\
\multicolumn{4}{|p{\dimexpr 16cm + 6\tabcolsep\relax}|}{\begin{itemize}\item BARBOSA, Simone Diniz Junqueira. Interação humano-computador. Rio de Janeiro: Campus, 2010.\end{itemize}}\\
\multicolumn{4}{|p{\dimexpr 16cm + 6\tabcolsep\relax}|}{\begin{itemize}\item SOUZA, Clarisse da Silva Vieira Camelo de. The semiotic engineering of human-computer interaction. Cambridge, Mass.: MIT Press, 2005.\end{itemize}}\\
\multicolumn{4}{|p{16cm}|}{}\\
\hline
\multicolumn{4}{|p{16cm}|}{Bibliografia Complementar:}\\*
\multicolumn{4}{|p{\dimexpr 16cm + 6\tabcolsep\relax}|}{\begin{itemize}\item PREECE, Jenny. Human-computer interaction. Harlow: Addison Wesley, c1994.\end{itemize}}\\
\multicolumn{4}{|p{\dimexpr 16cm + 6\tabcolsep\relax}|}{\begin{itemize}\item CYBIS, Walter; BETIOL, Adriana Holtz; FAUST, Richard. Ergonomia e usabilidade: conhecimentos, métodos e aplicações. 2. ed. São Paulo: Novatec, 2007.\end{itemize}}\\
\multicolumn{4}{|p{\dimexpr 16cm + 6\tabcolsep\relax}|}{\begin{itemize}\item BENYON, David. Interação humano-computador. 2. ed. São Paulo: Pearson, 2011.\end{itemize}}\\
\multicolumn{4}{|p{\dimexpr 16cm + 6\tabcolsep\relax}|}{\begin{itemize}\item SHNEIDERMAN, Ben; PLAISANT, Catherine. Designing the user interface: strategies for effective human-computer interaction. 5th ed. Boston: Addison Wesley, 2010.\end{itemize}}\\
\multicolumn{4}{|p{\dimexpr 16cm + 6\tabcolsep\relax}|}{\begin{itemize}\item JACKO, Julie A. (ed.). The human-computer interaction handbook: fundamentals, evolving technologies, and emerging applications. 3rd ed. Boca Raton, FL: CRC Press, c2012.\end{itemize}}\\
\hline
\end{longtable}
\end{center}

\clearpage
\begin{center}
\begin{tabular}{lccr}
 \multirow{3}{*}{\includegraphics[height=2.7cm]{brasao.png}} &
 \multicolumn{2}{c}{\bfseries UNIVERSIDADE FEDERAL DE OURO PRETO} &
 \ \ \ \ \multirow{3}{*}{\includegraphics[height=2.7cm]{ufop.png}} \\
 & \multicolumn{2}{c}{\bfseries PRÓ-REITORIA DE GRADUAÇÃO} & \\
 & \multicolumn{2}{c}{} & \\
 & \multicolumn{2}{c}{\Large\bfseries PROGRAMA DE DISCIPLINA} & \\
\end{tabular}
\end{center}

\begin{center}
\begin{longtable}{|p{4cm}|p{4cm}|p{4cm}|p{4cm}|}
\hline
\multicolumn{3}{|p{12cm}|}{Nome do Componente Curricular em Português:} &
\multicolumn{1}{p{4cm}|}{Código:} \\ 
\multicolumn{3}{|p{12cm}|}{\textbf{Computação Gráfica}} &
\textbf{BCC327}\\ 
\multicolumn{3}{|p{12cm}|}{Nome do Componente Curricular em Inglês:} & \\ 
\multicolumn{3}{|p{12cm}|}{\textbf{Computer Graphics}} & \\ 
\hline
\multicolumn{3}{|p{12cm}|}{Nome e Sigla do Departamento} & Unidade Acadêmica: \\ 
\multicolumn{3}{|p{12cm}|}{Departamento de Computação (DECOM)} & {ICEB} \\ 
\hline
\multicolumn{4}{|p{16cm}|}{Modalidade de Oferta:
[X] presencial \hspace{0.6cm}
[ ] semipresencial \hspace{0.6cm}
[ ] à distância}\\
\hline
\multicolumn{2}{|p{8cm}|}{Carga horária semestral} &
\multicolumn{2}{p{8cm}|}{Carga horária semanal}\\
\hline
\multicolumn{1}{|p{4cm}|}{Total} &
\multicolumn{1}{p{4cm}|}{Extensionista} &
\multicolumn{1}{p{4cm}|}{Teórica} &
\multicolumn{1}{p{4cm}|}{Prática} \\ 
\multicolumn{1}{|p{4cm}|}{60\,horas} &
\multicolumn{1}{p{4cm}|}{0\;horas} &
\multicolumn{1}{p{4cm}|}{4\;horas/aula} &
\multicolumn{1}{p{4cm}|}{0\;horas/aula} \\ 
\hline
\multicolumn{4}{|p{16cm}|}{Ementa:}\\
\multicolumn{4}{|p{16cm}|}{}\\
\multicolumn{4}{|p{\dimexpr 16cm + 6\tabcolsep\relax}|}{Sistemas gráficos; programação gráfica interativa; modelagem de objetos e cenas tridimensionais; transformações geométricas; visualização; iluminação e colorização; transparência e opacidade; mapeamento de textura; recorte; remoção de superfícies escondidas; rasterização.}\\
\multicolumn{4}{|p{16cm}|}{}\\
\hline
\multicolumn{4}{|p{16cm}|}{Conteúdo programático:}\\*
\multicolumn{4}{|p{\dimexpr 16cm + 6\tabcolsep\relax}|}{\begin{enumerate}[start=1]\item Sistemas gráficos\end{enumerate}}\\
\multicolumn{4}{|p{\dimexpr 16cm + 6\tabcolsep\relax}|}{\begin{enumerate}[start=2]\item Dispositivos de entrada e saída\end{enumerate}}\\
\multicolumn{4}{|p{\dimexpr 16cm + 6\tabcolsep\relax}|}{\begin{enumerate}[start=3]\item Arquitetura de sistemas gráficos\end{enumerate}}\\
\multicolumn{4}{|p{\dimexpr 16cm + 6\tabcolsep\relax}|}{\begin{enumerate}[start=4]\item Cor\end{enumerate}}\\
\multicolumn{4}{|p{\dimexpr 16cm + 6\tabcolsep\relax}|}{\begin{enumerate}[start=5]\item Programação gráfica: interface para o programador de aplicações, primitivas e atributos, técnicas interativas\end{enumerate}}\\
\multicolumn{4}{|p{\dimexpr 16cm + 6\tabcolsep\relax}|}{\begin{enumerate}[start=6]\item Modelagem de objetos e cenas tridimensionais: objetos geométricos elementares, operações com dados geométricos, primitivas tridimensionais, sistemas de coordenadas, modelagem de objetos\end{enumerate}}\\
\multicolumn{4}{|p{\dimexpr 16cm + 6\tabcolsep\relax}|}{\begin{enumerate}[start=7]\item Estruturas de dados\end{enumerate}}\\
\multicolumn{4}{|p{\dimexpr 16cm + 6\tabcolsep\relax}|}{\begin{enumerate}[start=8]\item Transformações afins\end{enumerate}}\\
\multicolumn{4}{|p{\dimexpr 16cm + 6\tabcolsep\relax}|}{\begin{enumerate}[start=9]\item Movimentação de objetos\end{enumerate}}\\
\multicolumn{4}{|p{\dimexpr 16cm + 6\tabcolsep\relax}|}{\begin{enumerate}[start=10]\item Visualização: posicionamento do observador em relação à cena, transformação de visualização, tipos de projeção, transformação de projeção\end{enumerate}}\\
\multicolumn{4}{|p{\dimexpr 16cm + 6\tabcolsep\relax}|}{\begin{enumerate}[start=11]\item Iluminação e colorização: interação entre luz e superfícies, fontes de luz, modelo de reflexão de Phong, colorização de superfícies poligonais, transparência, mapeamento de textura\end{enumerate}}\\
\multicolumn{4}{|p{\dimexpr 16cm + 6\tabcolsep\relax}|}{\begin{enumerate}[start=12]\item Recorte: de segmentos de retas, de polígonos, de outras primitivas\end{enumerate}}\\
\multicolumn{4}{|p{\dimexpr 16cm + 6\tabcolsep\relax}|}{\begin{enumerate}[start=13]\item Remoção de superfícies escondidas: eliminação de faces posteriores\end{enumerate}}\\
\multicolumn{4}{|p{\dimexpr 16cm + 6\tabcolsep\relax}|}{\begin{enumerate}[start=14]\item Algoritmo de ordenação z-buffer: ordenação por profundidade\end{enumerate}}\\
\multicolumn{4}{|p{\dimexpr 16cm + 6\tabcolsep\relax}|}{\begin{enumerate}[start=15]\item Rasterização: de segmentos de retas, de polígonos\end{enumerate}}\\
\multicolumn{4}{|p{\dimexpr 16cm + 6\tabcolsep\relax}|}{\begin{enumerate}[start=16]\item Aliasing\end{enumerate}}\\
\multicolumn{4}{|p{\dimexpr 16cm + 6\tabcolsep\relax}|}{}\\
\hline
\multicolumn{4}{|p{16cm}|}{Bibliografia Básica:}\\*
\multicolumn{4}{|p{\dimexpr 16cm + 6\tabcolsep\relax}|}{\begin{itemize}\item CONCI, Aura; AZEVEDO, Eduardo. Computação gráfica: geração de imagens. 8. ed. Rio de Janeiro: Campus, 2003.\end{itemize}}\\
\multicolumn{4}{|p{\dimexpr 16cm + 6\tabcolsep\relax}|}{\begin{itemize}\item ANGEL, Edward. Interactive computer graphics: a top-down approach with OpenGL. 6th ed. Boston: Addison Wesley, 2012.\end{itemize}}\\
\multicolumn{4}{|p{\dimexpr 16cm + 6\tabcolsep\relax}|}{\begin{itemize}\item HILL, F. S. Jr. Computer graphics: using OpenGL. 2nd ed. Upper Saddle River: Prentice-Hall, 2001.\end{itemize}}\\
\multicolumn{4}{|p{16cm}|}{}\\
\hline
\multicolumn{4}{|p{16cm}|}{Bibliografia Complementar:}\\*
\multicolumn{4}{|p{\dimexpr 16cm + 6\tabcolsep\relax}|}{\begin{itemize}\item GOMES, Jonas de Miranda. Fundamentos da computação gráfica. Rio de Janeiro: Instituto Nacional de Matemática Pura e Aplicada, 2008.\end{itemize}}\\
\multicolumn{4}{|p{\dimexpr 16cm + 6\tabcolsep\relax}|}{\begin{itemize}\item HEARN, Donald. Computer graphics with OpenGL. 4th ed. Boston: Prentice-Hall, 2011.\end{itemize}}\\
\multicolumn{4}{|p{\dimexpr 16cm + 6\tabcolsep\relax}|}{\begin{itemize}\item FOLEY, James D. Fundamentals of interactive computer graphics. Reading, Mass.: Addison Wesley, 1984.\end{itemize}}\\
\multicolumn{4}{|p{\dimexpr 16cm + 6\tabcolsep\relax}|}{\begin{itemize}\item TORI, Romero et al. Fundamentos da computação gráfica: Compugrafia. São Paulo: Livros Técnicos e Científicos.\end{itemize}}\\
\multicolumn{4}{|p{\dimexpr 16cm + 6\tabcolsep\relax}|}{\begin{itemize}\item AZEVEDO, Eduardo; CONCI, Aura. Computação gráfica: teoria e prática. Rio de Janeiro: Elsevier.\end{itemize}}\\
\hline
\end{longtable}
\end{center}

\clearpage
\begin{center}
\begin{tabular}{lccr}
 \multirow{3}{*}{\includegraphics[height=2.7cm]{brasao.png}} &
 \multicolumn{2}{c}{\bfseries UNIVERSIDADE FEDERAL DE OURO PRETO} &
 \ \ \ \ \multirow{3}{*}{\includegraphics[height=2.7cm]{ufop.png}} \\
 & \multicolumn{2}{c}{\bfseries PRÓ-REITORIA DE GRADUAÇÃO} & \\
 & \multicolumn{2}{c}{} & \\
 & \multicolumn{2}{c}{\Large\bfseries PROGRAMA DE DISCIPLINA} & \\
\end{tabular}
\end{center}

\begin{center}
\begin{longtable}{|p{4cm}|p{4cm}|p{4cm}|p{4cm}|}
\hline
\multicolumn{3}{|p{12cm}|}{Nome do Componente Curricular em Português:} &
\multicolumn{1}{p{4cm}|}{Código:} \\ 
\multicolumn{3}{|p{12cm}|}{\textbf{Construção de Compiladores I}} &
\textbf{BCC328}\\ 
\multicolumn{3}{|p{12cm}|}{Nome do Componente Curricular em Inglês:} & \\ 
\multicolumn{3}{|p{12cm}|}{\textbf{Compiler Construction I}} & \\ 
\hline
\multicolumn{3}{|p{12cm}|}{Nome e Sigla do Departamento} & Unidade Acadêmica: \\ 
\multicolumn{3}{|p{12cm}|}{Departamento de Computação (DECOM)} & {ICEB} \\ 
\hline
\multicolumn{4}{|p{16cm}|}{Modalidade de Oferta:
[X] presencial \hspace{0.6cm}
[ ] semipresencial \hspace{0.6cm}
[ ] à distância}\\
\hline
\multicolumn{2}{|p{8cm}|}{Carga horária semestral} &
\multicolumn{2}{p{8cm}|}{Carga horária semanal}\\
\hline
\multicolumn{1}{|p{4cm}|}{Total} &
\multicolumn{1}{p{4cm}|}{Extensionista} &
\multicolumn{1}{p{4cm}|}{Teórica} &
\multicolumn{1}{p{4cm}|}{Prática} \\ 
\multicolumn{1}{|p{4cm}|}{60\,horas} &
\multicolumn{1}{p{4cm}|}{0\;horas} &
\multicolumn{1}{p{4cm}|}{4\;horas/aula} &
\multicolumn{1}{p{4cm}|}{0\;horas/aula} \\ 
\hline
\multicolumn{4}{|p{16cm}|}{Ementa:}\\
\multicolumn{4}{|p{16cm}|}{}\\
\multicolumn{4}{|p{\dimexpr 16cm + 6\tabcolsep\relax}|}{Implementação de linguagens de programação: compilação e interpretação; análise léxica; análise sintática; análise semântica.}\\
\multicolumn{4}{|p{16cm}|}{}\\
\hline
\multicolumn{4}{|p{16cm}|}{Conteúdo programático:}\\*
\multicolumn{4}{|p{\dimexpr 16cm + 6\tabcolsep\relax}|}{\begin{enumerate}[start=1]\item Implementação de linguagens de programação: compilação e interpretação\end{enumerate}}\\
\multicolumn{4}{|p{\dimexpr 16cm + 6\tabcolsep\relax}|}{\begin{enumerate}[start=2]\item Análise léxica: análise léxica, analisador léxico ad hoc, expressões regulares, autômatos finitos e geradores de analisadores léxicos\end{enumerate}}\\
\multicolumn{4}{|p{\dimexpr 16cm + 6\tabcolsep\relax}|}{\begin{enumerate}[start=3]\item Análise sintática: análise sintática, gramáticas livres de contexto, análise sintática descendente, análise sintática ascendente, geradores de analisadores sintáticos e ações semânticas\end{enumerate}}\\
\multicolumn{4}{|p{\dimexpr 16cm + 6\tabcolsep\relax}|}{\begin{enumerate}[start=4]\item Análise semântica: árvore abstrata, tabelas de símbolos, checagem de tipos e regras semânticas\end{enumerate}}\\
\multicolumn{4}{|p{\dimexpr 16cm + 6\tabcolsep\relax}|}{}\\
\hline
\multicolumn{4}{|p{16cm}|}{Bibliografia Básica:}\\*
\multicolumn{4}{|p{\dimexpr 16cm + 6\tabcolsep\relax}|}{\begin{itemize}\item APPEL, Andrew W. Modern compiler implementation in Java. 2nd ed. Cambridge: Cambridge University Press, c2002.\end{itemize}}\\
\multicolumn{4}{|p{\dimexpr 16cm + 6\tabcolsep\relax}|}{\begin{itemize}\item AHO, Alfred V. Compiladores, princípios, técnicas e ferramentas. São Paulo: Pearson Addison-Wesley, 2007.\end{itemize}}\\
\multicolumn{4}{|p{\dimexpr 16cm + 6\tabcolsep\relax}|}{\begin{itemize}\item APPEL, Andrew W. Modern compiler implementation in ML. New York: Cambridge University Press, c1998.\end{itemize}}\\
\multicolumn{4}{|p{16cm}|}{}\\
\hline
\multicolumn{4}{|p{16cm}|}{Bibliografia Complementar:}\\*
\multicolumn{4}{|p{\dimexpr 16cm + 6\tabcolsep\relax}|}{\begin{itemize}\item MOGENSEN, Torben Ægidius. Introduction to Compiler Design. 2. ed. Springer, 2011. (item inexistente na UFOP)\end{itemize}}\\
\multicolumn{4}{|p{\dimexpr 16cm + 6\tabcolsep\relax}|}{\begin{itemize}\item GRÜNE, Dick; REEUWIJK, Kees van; BAL, Henri E.; JACOBS, Ceriel J. H.; LANGENDOEN, Koen. Modern Compiler Design. 2. ed. Springer, 2012. (item inexistente na UFOP)\end{itemize}}\\
\multicolumn{4}{|p{\dimexpr 16cm + 6\tabcolsep\relax}|}{\begin{itemize}\item WILHELM, Reinhard; SEIDL, Helmut; HACK, Sebastian. Compiler Design: Syntactic and Semantic Analysis. Springer, 2013. (item inexistente na UFOP)\end{itemize}}\\
\multicolumn{4}{|p{\dimexpr 16cm + 6\tabcolsep\relax}|}{\begin{itemize}\item FISCHER, Charles N.; CYTRON, Ron K.; LEBLANC, Richard J. Jr. Crafting a Compiler. Pearson, 2010. (item inexistente na UFOP)\end{itemize}}\\
\multicolumn{4}{|p{\dimexpr 16cm + 6\tabcolsep\relax}|}{\begin{itemize}\item COOPER, Keith D.; TORCZON, Linda. Engineering a Compiler. 2. ed. Elsevier, 2011. (item inexistente na UFOP)\end{itemize}}\\
\hline
\end{longtable}
\end{center}

\clearpage
\begin{center}
\begin{tabular}{lccr}
 \multirow{3}{*}{\includegraphics[height=2.7cm]{brasao.png}} &
 \multicolumn{2}{c}{\bfseries UNIVERSIDADE FEDERAL DE OURO PRETO} &
 \ \ \ \ \multirow{3}{*}{\includegraphics[height=2.7cm]{ufop.png}} \\
 & \multicolumn{2}{c}{\bfseries PRÓ-REITORIA DE GRADUAÇÃO} & \\
 & \multicolumn{2}{c}{} & \\
 & \multicolumn{2}{c}{\Large\bfseries PROGRAMA DE DISCIPLINA} & \\
\end{tabular}
\end{center}

\begin{center}
\begin{longtable}{|p{4cm}|p{4cm}|p{4cm}|p{4cm}|}
\hline
\multicolumn{3}{|p{12cm}|}{Nome do Componente Curricular em Português:} &
\multicolumn{1}{p{4cm}|}{Código:} \\ 
\multicolumn{3}{|p{12cm}|}{\textbf{Projeto de Circuitos Lógicos Integrados usando HDL (Hardware Description Language)}} &
\textbf{BCC408}\\ 
\multicolumn{3}{|p{12cm}|}{Nome do Componente Curricular em Inglês:} & \\ 
\multicolumn{3}{|p{12cm}|}{\textbf{Design of Integrated Logic Circuits Using HDL (Hardware Description Language)}} & \\ 
\hline
\multicolumn{3}{|p{12cm}|}{Nome e Sigla do Departamento} & Unidade Acadêmica: \\ 
\multicolumn{3}{|p{12cm}|}{Departamento de Computação (DECOM)} & {ICEB} \\ 
\hline
\multicolumn{4}{|p{16cm}|}{Modalidade de Oferta:
[X] presencial \hspace{0.6cm}
[ ] semipresencial \hspace{0.6cm}
[ ] à distância}\\
\hline
\multicolumn{2}{|p{8cm}|}{Carga horária semestral} &
\multicolumn{2}{p{8cm}|}{Carga horária semanal}\\
\hline
\multicolumn{1}{|p{4cm}|}{Total} &
\multicolumn{1}{p{4cm}|}{Extensionista} &
\multicolumn{1}{p{4cm}|}{Teórica} &
\multicolumn{1}{p{4cm}|}{Prática} \\ 
\multicolumn{1}{|p{4cm}|}{60\,horas} &
\multicolumn{1}{p{4cm}|}{0\;horas} &
\multicolumn{1}{p{4cm}|}{1\;horas/aula} &
\multicolumn{1}{p{4cm}|}{3\;horas/aula} \\ 
\hline
\multicolumn{4}{|p{16cm}|}{Ementa:}\\
\multicolumn{4}{|p{16cm}|}{}\\
\multicolumn{4}{|p{\dimexpr 16cm + 6\tabcolsep\relax}|}{Visão geral do fluxo de Projeto de Circuitos Integrados Digitais e sua importância para o desenvolvimento econômico do país. Conceituação, aplicabilidade e aplicabilidade de circuitos integrados de propósito específico (ASIC – Application Specific Integrated Circuits). Principais linguagens de descrição de hardware (HDL – Hardware Description Language). Técnicas de Projeto de Circuitos Integrados. Uso de HDL para o projeto de circuitos integrados.}\\
\multicolumn{4}{|p{16cm}|}{}\\
\hline
\multicolumn{4}{|p{16cm}|}{Conteúdo programático:}\\*
\multicolumn{4}{|p{\dimexpr 16cm + 6\tabcolsep\relax}|}{\begin{enumerate}[start=1]\item Introdução: fluxo de projeto de circuito integrado digital (arquitetura, design, verificação, síntese lógica e síntese física)\end{enumerate}}\\
\multicolumn{4}{|p{\dimexpr 16cm + 6\tabcolsep\relax}|}{\begin{enumerate}[start=2]\item Introdução à FPGA (arranjo de portas, capacidade, LookUp Table - LUT's)\end{enumerate}}\\
\multicolumn{4}{|p{\dimexpr 16cm + 6\tabcolsep\relax}|}{\begin{enumerate}[start=3]\item Arquitetura de Circuitos Integrados Digitais (processadores, barramentos de comunicação destinados a System-on-Chip - SoC, periféricos e acesso a memórias)\end{enumerate}}\\
\multicolumn{4}{|p{\dimexpr 16cm + 6\tabcolsep\relax}|}{\begin{enumerate}[start=4]\item Linguagem de Descrição de Hardware: Verilog 2001\end{enumerate}}\\
\multicolumn{4}{|p{\dimexpr 16cm + 6\tabcolsep\relax}|}{\begin{enumerate}[start=5]\item Descrição de Circuitos Integrados Digitais combinacionais e sequenciais básicos (portas lógicas, meio-somador, somador completo, multiplexadores e demultiplexadores, codificadores e decodificadores, contadores, registradores e flip-flops) utilizando Verilog 2001\end{enumerate}}\\
\multicolumn{4}{|p{\dimexpr 16cm + 6\tabcolsep\relax}|}{\begin{enumerate}[start=6]\item Arquitetura CISC e Arquitetura RISC\end{enumerate}}\\
\multicolumn{4}{|p{\dimexpr 16cm + 6\tabcolsep\relax}|}{\begin{enumerate}[start=7]\item Projetos de circuitos integrados usando HDL\end{enumerate}}\\
\multicolumn{4}{|p{\dimexpr 16cm + 6\tabcolsep\relax}|}{}\\
\hline
\multicolumn{4}{|p{16cm}|}{Bibliografia Básica:}\\*
\multicolumn{4}{|p{\dimexpr 16cm + 6\tabcolsep\relax}|}{\begin{itemize}\item PEDRONI, Volnei A. Eletrônica digital moderna e VHDL. Rio de Janeiro: Elsevier, 2010. ISBN 9788535234657.\end{itemize}}\\
\multicolumn{4}{|p{\dimexpr 16cm + 6\tabcolsep\relax}|}{\begin{itemize}\item COFFMAN, Ken. Real world FPGA design with Verilog. Pearson Education, 1999. ISBN 0130998516.\end{itemize}}\\
\multicolumn{4}{|p{\dimexpr 16cm + 6\tabcolsep\relax}|}{\begin{itemize}\item CHU, Pong P. FPGA prototyping by VHDL examples: Xilinx Spartan-3 version. John Wiley \& Sons, 2011. ISBN 9780470185315.\end{itemize}}\\
\multicolumn{4}{|p{16cm}|}{}\\
\hline
\multicolumn{4}{|p{16cm}|}{Bibliografia Complementar:}\\*
\multicolumn{4}{|p{\dimexpr 16cm + 6\tabcolsep\relax}|}{\begin{itemize}\item ERCEGOVAC, Miloš Dragutin; LANG, Tomás; MORENO, Jaime H. Introdução aos sistemas digitais. Bookman, 2000. ISBN 0471527998.\end{itemize}}\\
\multicolumn{4}{|p{\dimexpr 16cm + 6\tabcolsep\relax}|}{\begin{itemize}\item TOCCI, Ronald J.; WIDMER, Neal S.; MOSS, Gregory L. Sistemas Digitais: Princípios e Aplicações. 11. ed. Rio de Janeiro: MZ Editora, 2007. ISBN 9788576059226.\end{itemize}}\\
\multicolumn{4}{|p{\dimexpr 16cm + 6\tabcolsep\relax}|}{\begin{itemize}\item VAHID, Frank. Sistemas Digitais. Bookman Editora, 2009. ISBN 9788577801909.\end{itemize}}\\
\multicolumn{4}{|p{\dimexpr 16cm + 6\tabcolsep\relax}|}{\begin{itemize}\item TANENBAUM, Andrew S. Organização estruturada de computadores. 3. ed. Rio de Janeiro: Prentice-Hall do Brasil, 1992. 460 p. ISBN 857054040X.\end{itemize}}\\
\multicolumn{4}{|p{\dimexpr 16cm + 6\tabcolsep\relax}|}{\begin{itemize}\item STALLINGS, William. Arquitetura e organização de computadores. 8. ed. São Paulo: Pearson, c2010. 624 p. ISBN 9788576055648.\end{itemize}}\\
\hline
\end{longtable}
\end{center}

\clearpage
\begin{center}
\begin{tabular}{lccr}
 \multirow{3}{*}{\includegraphics[height=2.7cm]{brasao.png}} &
 \multicolumn{2}{c}{\bfseries UNIVERSIDADE FEDERAL DE OURO PRETO} &
 \ \ \ \ \multirow{3}{*}{\includegraphics[height=2.7cm]{ufop.png}} \\
 & \multicolumn{2}{c}{\bfseries PRÓ-REITORIA DE GRADUAÇÃO} & \\
 & \multicolumn{2}{c}{} & \\
 & \multicolumn{2}{c}{\Large\bfseries PROGRAMA DE DISCIPLINA} & \\
\end{tabular}
\end{center}

\begin{center}
\begin{longtable}{|p{4cm}|p{4cm}|p{4cm}|p{4cm}|}
\hline
\multicolumn{3}{|p{12cm}|}{Nome do Componente Curricular em Português:} &
\multicolumn{1}{p{4cm}|}{Código:} \\ 
\multicolumn{3}{|p{12cm}|}{\textbf{Computação Móvel}} &
\textbf{BCC421}\\ 
\multicolumn{3}{|p{12cm}|}{Nome do Componente Curricular em Inglês:} & \\ 
\multicolumn{3}{|p{12cm}|}{\textbf{Mobile Computing}} & \\ 
\hline
\multicolumn{3}{|p{12cm}|}{Nome e Sigla do Departamento} & Unidade Acadêmica: \\ 
\multicolumn{3}{|p{12cm}|}{Departamento de Computação (DECOM)} & {ICEB} \\ 
\hline
\multicolumn{4}{|p{16cm}|}{Modalidade de Oferta:
[X] presencial \hspace{0.6cm}
[ ] semipresencial \hspace{0.6cm}
[ ] à distância}\\
\hline
\multicolumn{2}{|p{8cm}|}{Carga horária semestral} &
\multicolumn{2}{p{8cm}|}{Carga horária semanal}\\
\hline
\multicolumn{1}{|p{4cm}|}{Total} &
\multicolumn{1}{p{4cm}|}{Extensionista} &
\multicolumn{1}{p{4cm}|}{Teórica} &
\multicolumn{1}{p{4cm}|}{Prática} \\ 
\multicolumn{1}{|p{4cm}|}{60\,horas} &
\multicolumn{1}{p{4cm}|}{0\;horas} &
\multicolumn{1}{p{4cm}|}{4\;horas/aula} &
\multicolumn{1}{p{4cm}|}{0\;horas/aula} \\ 
\hline
\multicolumn{4}{|p{16cm}|}{Ementa:}\\
\multicolumn{4}{|p{16cm}|}{}\\
\multicolumn{4}{|p{\dimexpr 16cm + 6\tabcolsep\relax}|}{Propagação de sinais, redes de comunicação sem fio, serviços baseados na localização, protocolos de comunicação, gerência de informação, redes estruturadas e não estruturadas. Desenvolvimento de aplicações móveis nas plataformas correntes.}\\
\multicolumn{4}{|p{16cm}|}{}\\
\hline
\multicolumn{4}{|p{16cm}|}{Conteúdo programático:}\\*
\multicolumn{4}{|p{\dimexpr 16cm + 6\tabcolsep\relax}|}{\begin{enumerate}[start=1]\item Propagação de sinais\end{enumerate}}\\
\multicolumn{4}{|p{\dimexpr 16cm + 6\tabcolsep\relax}|}{\begin{enumerate}[start=2]\item Redes de comunicação sem fio\end{enumerate}}\\
\multicolumn{4}{|p{\dimexpr 16cm + 6\tabcolsep\relax}|}{\begin{enumerate}[start=3]\item Serviços baseados na localização\end{enumerate}}\\
\multicolumn{4}{|p{\dimexpr 16cm + 6\tabcolsep\relax}|}{\begin{enumerate}[start=4]\item Protocolos de comunicação\end{enumerate}}\\
\multicolumn{4}{|p{\dimexpr 16cm + 6\tabcolsep\relax}|}{\begin{enumerate}[start=5]\item Gerência de informação\end{enumerate}}\\
\multicolumn{4}{|p{\dimexpr 16cm + 6\tabcolsep\relax}|}{\begin{enumerate}[start=6]\item Redes estruturadas e não estruturadas\end{enumerate}}\\
\multicolumn{4}{|p{\dimexpr 16cm + 6\tabcolsep\relax}|}{\begin{enumerate}[start=7]\item Desenvolvimento de aplicações móveis: plataforma Android, plataforma Windows, plataforma iOS\end{enumerate}}\\
\multicolumn{4}{|p{\dimexpr 16cm + 6\tabcolsep\relax}|}{}\\
\hline
\multicolumn{4}{|p{16cm}|}{Bibliografia Básica:}\\*
\multicolumn{4}{|p{\dimexpr 16cm + 6\tabcolsep\relax}|}{\begin{itemize}\item RAPPAPORT, Theodore. Wireless Personal Communications. 2. ed. Kluwer Academic, 2007. (item inexistente na UFOP)\end{itemize}}\\
\multicolumn{4}{|p{\dimexpr 16cm + 6\tabcolsep\relax}|}{\begin{itemize}\item MATEUS, G. R. Introdução à Computação Móvel. 11. ed. Rio de Janeiro: Escola de Computação, 1998. (item inexistente na UFOP)\end{itemize}}\\
\multicolumn{4}{|p{\dimexpr 16cm + 6\tabcolsep\relax}|}{\begin{itemize}\item TERUEL, E. Web Mobile: Desenvolva Sites para Dispositivos Móveis com Tecnologias de Uso Livre. Ciência Moderna, 2010. (item inexistente na UFOP)\end{itemize}}\\
\multicolumn{4}{|p{\dimexpr 16cm + 6\tabcolsep\relax}|}{\begin{itemize}\item SCHILLER, Jochen H. Mobile communications. 2nd ed. London: Addison Wesley, 2003.\end{itemize}}\\
\multicolumn{4}{|p{16cm}|}{}\\
\hline
\multicolumn{4}{|p{16cm}|}{Bibliografia Complementar:}\\*
\multicolumn{4}{|p{\dimexpr 16cm + 6\tabcolsep\relax}|}{\begin{itemize}\item LECHETA, Ricardo R. Google Android: aprenda a criar aplicações para dispositivos móveis com o Android SDK. 2. ed. São Paulo: Novatec Editora, 2010. (item inexistente na UFOP)\end{itemize}}\\
\multicolumn{4}{|p{\dimexpr 16cm + 6\tabcolsep\relax}|}{\begin{itemize}\item BORGES JÚNIOR, M. Aplicativos Móveis: Aplicativos para Dispositivos Móveis usando C\#.Net. Rio de Janeiro: Ciência Moderna, 2005. (item inexistente na UFOP)\end{itemize}}\\
\multicolumn{4}{|p{\dimexpr 16cm + 6\tabcolsep\relax}|}{\begin{itemize}\item PETERSON, Larry L.; DAVIE, Bruce S.; tradução Daniel Vieira; revisão técnica Marcos Simplício. Redes de computadores: uma abordagem de sistemas. Rio de Janeiro: Campus.\end{itemize}}\\
\multicolumn{4}{|p{\dimexpr 16cm + 6\tabcolsep\relax}|}{\begin{itemize}\item OPPENHEIM, Alan V. Discrete-time signal processing. 3rd ed. New Delhi: Pearson, c2010.\end{itemize}}\\
\hline
\end{longtable}
\end{center}

\clearpage
\begin{center}
\begin{tabular}{lccr}
 \multirow{3}{*}{\includegraphics[height=2.7cm]{brasao.png}} &
 \multicolumn{2}{c}{\bfseries UNIVERSIDADE FEDERAL DE OURO PRETO} &
 \ \ \ \ \multirow{3}{*}{\includegraphics[height=2.7cm]{ufop.png}} \\
 & \multicolumn{2}{c}{\bfseries PRÓ-REITORIA DE GRADUAÇÃO} & \\
 & \multicolumn{2}{c}{} & \\
 & \multicolumn{2}{c}{\Large\bfseries PROGRAMA DE DISCIPLINA} & \\
\end{tabular}
\end{center}

\begin{center}
\begin{longtable}{|p{4cm}|p{4cm}|p{4cm}|p{4cm}|}
\hline
\multicolumn{3}{|p{12cm}|}{Nome do Componente Curricular em Português:} &
\multicolumn{1}{p{4cm}|}{Código:} \\ 
\multicolumn{3}{|p{12cm}|}{\textbf{Computação nas Nuvens}} &
\textbf{BCC422}\\ 
\multicolumn{3}{|p{12cm}|}{Nome do Componente Curricular em Inglês:} & \\ 
\multicolumn{3}{|p{12cm}|}{\textbf{Cloud Computing}} & \\ 
\hline
\multicolumn{3}{|p{12cm}|}{Nome e Sigla do Departamento} & Unidade Acadêmica: \\ 
\multicolumn{3}{|p{12cm}|}{Departamento de Computação (DECOM)} & {ICEB} \\ 
\hline
\multicolumn{4}{|p{16cm}|}{Modalidade de Oferta:
[X] presencial \hspace{0.6cm}
[ ] semipresencial \hspace{0.6cm}
[ ] à distância}\\
\hline
\multicolumn{2}{|p{8cm}|}{Carga horária semestral} &
\multicolumn{2}{p{8cm}|}{Carga horária semanal}\\
\hline
\multicolumn{1}{|p{4cm}|}{Total} &
\multicolumn{1}{p{4cm}|}{Extensionista} &
\multicolumn{1}{p{4cm}|}{Teórica} &
\multicolumn{1}{p{4cm}|}{Prática} \\ 
\multicolumn{1}{|p{4cm}|}{60\,horas} &
\multicolumn{1}{p{4cm}|}{0\;horas} &
\multicolumn{1}{p{4cm}|}{4\;horas/aula} &
\multicolumn{1}{p{4cm}|}{0\;horas/aula} \\ 
\hline
\multicolumn{4}{|p{16cm}|}{Ementa:}\\
\multicolumn{4}{|p{16cm}|}{}\\
\multicolumn{4}{|p{\dimexpr 16cm + 6\tabcolsep\relax}|}{Fundamentos de Cloud Computing: terminologias e conceitos; virtualização; elasticidade, resiliência, on-demand e uso medido; benefícios, desafios e riscos das plataformas e serviços; administração e regras; Software as a Service (SaaS), Platform as a Service (PaaS) e Infrastructure as a Service (IaaS) Cloud Delivery Models; Service Level Agreements (SLAs) for Cloud-based IT Resources.}\\
\multicolumn{4}{|p{16cm}|}{}\\
\hline
\multicolumn{4}{|p{16cm}|}{Conteúdo programático:}\\*
\multicolumn{4}{|p{\dimexpr 16cm + 6\tabcolsep\relax}|}{\begin{enumerate}[start=1]\item Fundamentos de Cloud Computing: terminologias e conceitos\end{enumerate}}\\
\multicolumn{4}{|p{\dimexpr 16cm + 6\tabcolsep\relax}|}{\begin{enumerate}[start=2]\item Virtualização: XEN, QEMU, Eucalyptus, Cloud Privada\end{enumerate}}\\
\multicolumn{4}{|p{\dimexpr 16cm + 6\tabcolsep\relax}|}{\begin{enumerate}[start=3]\item Elasticidade, Resiliência, On-Demand e Uso Medido\end{enumerate}}\\
\multicolumn{4}{|p{\dimexpr 16cm + 6\tabcolsep\relax}|}{\begin{enumerate}[start=4]\item Software as a Service (SaaS), Platform as a Service (PaaS) e Infrastructure as a Service (IaaS) Cloud Delivery\end{enumerate}}\\
\multicolumn{4}{|p{\dimexpr 16cm + 6\tabcolsep\relax}|}{\begin{enumerate}[start=5]\item Benefícios, desafios e riscos das plataformas e serviços\end{enumerate}}\\
\multicolumn{4}{|p{\dimexpr 16cm + 6\tabcolsep\relax}|}{\begin{enumerate}[start=6]\item Criando uma solução na Cloud: Ruby on Rails como PaaS; Roteamento definido por Software como IaaS; SOA como SaaS\end{enumerate}}\\
\multicolumn{4}{|p{\dimexpr 16cm + 6\tabcolsep\relax}|}{\begin{enumerate}[start=7]\item Administração e Regras\end{enumerate}}\\
\multicolumn{4}{|p{\dimexpr 16cm + 6\tabcolsep\relax}|}{\begin{enumerate}[start=8]\item Service Level Agreements (SLAs) for Cloud-based IT Resources\end{enumerate}}\\
\multicolumn{4}{|p{\dimexpr 16cm + 6\tabcolsep\relax}|}{\begin{enumerate}[start=9]\item Análise de desempenho\end{enumerate}}\\
\multicolumn{4}{|p{\dimexpr 16cm + 6\tabcolsep\relax}|}{}\\
\hline
\multicolumn{4}{|p{16cm}|}{Bibliografia Básica:}\\*
\multicolumn{4}{|p{\dimexpr 16cm + 6\tabcolsep\relax}|}{\begin{itemize}\item BIRMAN, Kenneth. Guide to Reliable Distributed Systems: Building High-Assurance Applications and Cloud-Hosted Services. Springer, 2012. (item inexistente na UFOP)\end{itemize}}\\
\multicolumn{4}{|p{\dimexpr 16cm + 6\tabcolsep\relax}|}{\begin{itemize}\item DEITEL, Paul J. Internet Applications e desenvolvimento Web para programadores. São Paulo: Prentice-Hall, 2009. (item inexistente na UFOP)\end{itemize}}\\
\multicolumn{4}{|p{\dimexpr 16cm + 6\tabcolsep\relax}|}{\begin{itemize}\item KUROSE, James F. Redes de computadores e a Internet: uma abordagem top-down. 5. ed. São Paulo: Pearson Addison Wesley, 2010.\end{itemize}}\\
\multicolumn{4}{|p{16cm}|}{}\\
\hline
\multicolumn{4}{|p{16cm}|}{Bibliografia Complementar:}\\*
\multicolumn{4}{|p{\dimexpr 16cm + 6\tabcolsep\relax}|}{\begin{itemize}\item COMER, Douglas E.; tradução ARX Publicações. Interligação em redes com TCP/IP, volume 1: princípios, protocolos e arquitetura. 6. ed. Rio de Janeiro: Campus, c2015.\end{itemize}}\\
\multicolumn{4}{|p{\dimexpr 16cm + 6\tabcolsep\relax}|}{\begin{itemize}\item STALLINGS, William. Arquitetura e organização de computadores. 8. ed. São Paulo: Pearson, c2010.\end{itemize}}\\
\multicolumn{4}{|p{\dimexpr 16cm + 6\tabcolsep\relax}|}{\begin{itemize}\item ERL, Thomas; tradução Carlos Schafranski e Edson Furmankiewicz; revisão técnica Francisco Calmon N. da Gama, Ricardo A. de Castro Barbosa. SOA: princípios de design de serviços. São Paulo: Prentice Hall, c2009.\end{itemize}}\\
\multicolumn{4}{|p{\dimexpr 16cm + 6\tabcolsep\relax}|}{\begin{itemize}\item PETERSON, Larry L.; DAVIE, Bruce S.; tradução Daniel Vieira; revisão técnica Marcos Simplício. Redes de computadores: uma abordagem de sistemas. Rio de Janeiro: Campus.\end{itemize}}\\
\hline
\end{longtable}
\end{center}

\clearpage
\begin{center}
\begin{tabular}{lccr}
 \multirow{3}{*}{\includegraphics[height=2.7cm]{brasao.png}} &
 \multicolumn{2}{c}{\bfseries UNIVERSIDADE FEDERAL DE OURO PRETO} &
 \ \ \ \ \multirow{3}{*}{\includegraphics[height=2.7cm]{ufop.png}} \\
 & \multicolumn{2}{c}{\bfseries PRÓ-REITORIA DE GRADUAÇÃO} & \\
 & \multicolumn{2}{c}{} & \\
 & \multicolumn{2}{c}{\Large\bfseries PROGRAMA DE DISCIPLINA} & \\
\end{tabular}
\end{center}

\begin{center}
\begin{longtable}{|p{4cm}|p{4cm}|p{4cm}|p{4cm}|}
\hline
\multicolumn{3}{|p{12cm}|}{Nome do Componente Curricular em Português:} &
\multicolumn{1}{p{4cm}|}{Código:} \\ 
\multicolumn{3}{|p{12cm}|}{\textbf{Criptografia e Segurança de Sistemas}} &
\textbf{BCC423}\\ 
\multicolumn{3}{|p{12cm}|}{Nome do Componente Curricular em Inglês:} & \\ 
\multicolumn{3}{|p{12cm}|}{\textbf{Cryptography and Security Systems}} & \\ 
\hline
\multicolumn{3}{|p{12cm}|}{Nome e Sigla do Departamento} & Unidade Acadêmica: \\ 
\multicolumn{3}{|p{12cm}|}{Departamento de Computação (DECOM)} & {ICEB} \\ 
\hline
\multicolumn{4}{|p{16cm}|}{Modalidade de Oferta:
[X] presencial \hspace{0.6cm}
[ ] semipresencial \hspace{0.6cm}
[ ] à distância}\\
\hline
\multicolumn{2}{|p{8cm}|}{Carga horária semestral} &
\multicolumn{2}{p{8cm}|}{Carga horária semanal}\\
\hline
\multicolumn{1}{|p{4cm}|}{Total} &
\multicolumn{1}{p{4cm}|}{Extensionista} &
\multicolumn{1}{p{4cm}|}{Teórica} &
\multicolumn{1}{p{4cm}|}{Prática} \\ 
\multicolumn{1}{|p{4cm}|}{60\,horas} &
\multicolumn{1}{p{4cm}|}{0\;horas} &
\multicolumn{1}{p{4cm}|}{4\;horas/aula} &
\multicolumn{1}{p{4cm}|}{0\;horas/aula} \\ 
\hline
\multicolumn{4}{|p{16cm}|}{Ementa:}\\
\multicolumn{4}{|p{16cm}|}{}\\
\multicolumn{4}{|p{\dimexpr 16cm + 6\tabcolsep\relax}|}{Segurança de redes, de sistemas, da informação, segurança cibernética. Criptografia. Assinaturas digitais, certificados digitais, certificados de atributos. Segurança de redes. Mídias criptográficas: identificadores biométricos. Cybersegurança e impactos na sociedade contemporânea.}\\
\multicolumn{4}{|p{16cm}|}{}\\
\hline
\multicolumn{4}{|p{16cm}|}{Conteúdo programático:}\\*
\multicolumn{4}{|p{\dimexpr 16cm + 6\tabcolsep\relax}|}{\begin{enumerate}[start=1]\item Introdução à segurança: segurança da informação, segurança de redes, de sistemas, da informação e segurança cibernética\end{enumerate}}\\
\multicolumn{4}{|p{\dimexpr 16cm + 6\tabcolsep\relax}|}{\begin{enumerate}[start=2]\item Criptografia simétrica e assimétrica\end{enumerate}}\\
\multicolumn{4}{|p{\dimexpr 16cm + 6\tabcolsep\relax}|}{\begin{enumerate}[start=3]\item Algoritmos criptográficos\end{enumerate}}\\
\multicolumn{4}{|p{\dimexpr 16cm + 6\tabcolsep\relax}|}{\begin{enumerate}[start=4]\item Assinaturas digitais: certificados digitais e o padrão PKI (ICP)\end{enumerate}}\\
\multicolumn{4}{|p{\dimexpr 16cm + 6\tabcolsep\relax}|}{\begin{enumerate}[start=5]\item ICP-Brasil e ICP-Edu\end{enumerate}}\\
\multicolumn{4}{|p{\dimexpr 16cm + 6\tabcolsep\relax}|}{\begin{enumerate}[start=6]\item Segurança de redes: monitoramento e invasão de redes\end{enumerate}}\\
\multicolumn{4}{|p{\dimexpr 16cm + 6\tabcolsep\relax}|}{\begin{enumerate}[start=7]\item Princípios de hacking\end{enumerate}}\\
\multicolumn{4}{|p{\dimexpr 16cm + 6\tabcolsep\relax}|}{\begin{enumerate}[start=8]\item Mídias criptográficas e identificadores biométricos\end{enumerate}}\\
\multicolumn{4}{|p{\dimexpr 16cm + 6\tabcolsep\relax}|}{\begin{enumerate}[start=9]\item ISO27000\end{enumerate}}\\
\multicolumn{4}{|p{\dimexpr 16cm + 6\tabcolsep\relax}|}{\begin{enumerate}[start=10]\item ITIL e COBIT, dentro do contexto da segurança da informação\end{enumerate}}\\
\multicolumn{4}{|p{\dimexpr 16cm + 6\tabcolsep\relax}|}{\begin{enumerate}[start=11]\item Certificações IISP e outros\end{enumerate}}\\
\multicolumn{4}{|p{\dimexpr 16cm + 6\tabcolsep\relax}|}{\begin{enumerate}[start=12]\item Cibersegurança e impactos na sociedade contemporânea\end{enumerate}}\\
\multicolumn{4}{|p{\dimexpr 16cm + 6\tabcolsep\relax}|}{}\\
\hline
\multicolumn{4}{|p{16cm}|}{Bibliografia Básica:}\\*
\multicolumn{4}{|p{\dimexpr 16cm + 6\tabcolsep\relax}|}{\begin{itemize}\item STALLINGS, William. Cryptography and Network Security: Principles and Practice. 7. ed. N.Y.: Pearson, 2016. (item inexistente na UFOP)\end{itemize}}\\
\multicolumn{4}{|p{\dimexpr 16cm + 6\tabcolsep\relax}|}{\begin{itemize}\item GOLDREICH, Oded. Foundations of cryptography. Cambridge: Cambridge University Press, 2009. (item inexistente na UFOP)\end{itemize}}\\
\multicolumn{4}{|p{\dimexpr 16cm + 6\tabcolsep\relax}|}{\begin{itemize}\item MENEZES, A. J. Handbook of applied cryptography. Boca Raton: CRC Press, c1997.\end{itemize}}\\
\multicolumn{4}{|p{16cm}|}{}\\
\hline
\multicolumn{4}{|p{16cm}|}{Bibliografia Complementar:}\\*
\multicolumn{4}{|p{\dimexpr 16cm + 6\tabcolsep\relax}|}{\begin{itemize}\item TRAPPE, Wade; WASHINGTON, Lawrence C. Introduction to cryptography: with coding theory. 2. ed. Upper Saddle River, N.J.: Pearson Prentice Hall, c2006.\end{itemize}}\\
\multicolumn{4}{|p{\dimexpr 16cm + 6\tabcolsep\relax}|}{\begin{itemize}\item COUTINHO, S. C. Números inteiros e criptografia RSA. 2. ed. Rio de Janeiro: IMPA, 2009.\end{itemize}}\\
\multicolumn{4}{|p{\dimexpr 16cm + 6\tabcolsep\relax}|}{\begin{itemize}\item KATZ, Jonathan. Introduction to modern cryptography. New York: Chapman \& Hall/CRC, 2008.\end{itemize}}\\
\multicolumn{4}{|p{\dimexpr 16cm + 6\tabcolsep\relax}|}{\begin{itemize}\item GOLDREICH, Oded. Foundations of cryptography: basic tools. New York: Cambridge University Press, 2001.\end{itemize}}\\
\hline
\end{longtable}
\end{center}

\clearpage
\begin{center}
\begin{tabular}{lccr}
 \multirow{3}{*}{\includegraphics[height=2.7cm]{brasao.png}} &
 \multicolumn{2}{c}{\bfseries UNIVERSIDADE FEDERAL DE OURO PRETO} &
 \ \ \ \ \multirow{3}{*}{\includegraphics[height=2.7cm]{ufop.png}} \\
 & \multicolumn{2}{c}{\bfseries PRÓ-REITORIA DE GRADUAÇÃO} & \\
 & \multicolumn{2}{c}{} & \\
 & \multicolumn{2}{c}{\Large\bfseries PROGRAMA DE DISCIPLINA} & \\
\end{tabular}
\end{center}

\begin{center}
\begin{longtable}{|p{4cm}|p{4cm}|p{4cm}|p{4cm}|}
\hline
\multicolumn{3}{|p{12cm}|}{Nome do Componente Curricular em Português:} &
\multicolumn{1}{p{4cm}|}{Código:} \\ 
\multicolumn{3}{|p{12cm}|}{\textbf{Sistemas Embutidos}} &
\textbf{BCC425}\\ 
\multicolumn{3}{|p{12cm}|}{Nome do Componente Curricular em Inglês:} & \\ 
\multicolumn{3}{|p{12cm}|}{\textbf{Embedded Systems}} & \\ 
\hline
\multicolumn{3}{|p{12cm}|}{Nome e Sigla do Departamento} & Unidade Acadêmica: \\ 
\multicolumn{3}{|p{12cm}|}{Departamento de Computação (DECOM)} & {ICEB} \\ 
\hline
\multicolumn{4}{|p{16cm}|}{Modalidade de Oferta:
[X] presencial \hspace{0.6cm}
[ ] semipresencial \hspace{0.6cm}
[ ] à distância}\\
\hline
\multicolumn{2}{|p{8cm}|}{Carga horária semestral} &
\multicolumn{2}{p{8cm}|}{Carga horária semanal}\\
\hline
\multicolumn{1}{|p{4cm}|}{Total} &
\multicolumn{1}{p{4cm}|}{Extensionista} &
\multicolumn{1}{p{4cm}|}{Teórica} &
\multicolumn{1}{p{4cm}|}{Prática} \\ 
\multicolumn{1}{|p{4cm}|}{60\,horas} &
\multicolumn{1}{p{4cm}|}{0\;horas} &
\multicolumn{1}{p{4cm}|}{2\;horas/aula} &
\multicolumn{1}{p{4cm}|}{2\;horas/aula} \\ 
\hline
\multicolumn{4}{|p{16cm}|}{Ementa:}\\
\multicolumn{4}{|p{16cm}|}{}\\
\multicolumn{4}{|p{\dimexpr 16cm + 6\tabcolsep\relax}|}{Aritmética inteira de ponto fixo e ponto flutuante. Arquitetura de computadores e micro-controladores. Programação de microcontroladores em C e Assembly. Conceitos de entrada e saída para micro-controladores. Uso de máquinas de estado finito na solução de problemas de programação. Programação de aplicações com teclados, displays de cristal líquido e sete segmentos, interface com atuadores (motores). Conversor analógico para digital (ADC) e digital para analógico (DAC). Conceito de interrupção e Timers. Transmissão de dados. Projeto de Sistemas Embutidos.}\\
\multicolumn{4}{|p{16cm}|}{}\\
\hline
\multicolumn{4}{|p{16cm}|}{Conteúdo programático:}\\*
\multicolumn{4}{|p{\dimexpr 16cm + 6\tabcolsep\relax}|}{\begin{enumerate}[start=1]\item Apresentação do plano de curso e Introdução: o que é um sistema embarcado, tipos de sistemas embarcados\end{enumerate}}\\
\multicolumn{4}{|p{\dimexpr 16cm + 6\tabcolsep\relax}|}{\begin{enumerate}[start=2]\item Representação numérica: inteira, ponto fixo, ponto flutuante\end{enumerate}}\\
\multicolumn{4}{|p{\dimexpr 16cm + 6\tabcolsep\relax}|}{\begin{enumerate}[start=3]\item Linguagem C e ferramentas: compilador, linker, loader\end{enumerate}}\\
\multicolumn{4}{|p{\dimexpr 16cm + 6\tabcolsep\relax}|}{\begin{enumerate}[start=4]\item Revisão de arquitetura de computadores e introdução à arquitetura dos microcontroladores: arquitetura, conjunto de instruções, programação, Assembly para microcontroladores\end{enumerate}}\\
\multicolumn{4}{|p{\dimexpr 16cm + 6\tabcolsep\relax}|}{\begin{enumerate}[start=5]\item Interface básica de Entrada/Saída: botões, LEDs, displays, segmentos, LCD; interface básica de entrada/saída com atuadores (motor de C.C., motor de passo, servo motor); interface analógica (ADC, DAC, PWM)\end{enumerate}}\\
\multicolumn{4}{|p{\dimexpr 16cm + 6\tabcolsep\relax}|}{\begin{enumerate}[start=6]\item Interrupções e Timers: conceito, programação, problema com compartilhamento de recursos\end{enumerate}}\\
\multicolumn{4}{|p{\dimexpr 16cm + 6\tabcolsep\relax}|}{\begin{enumerate}[start=7]\item Comunicação Serial: RS232, I2C, SPI, CAN\end{enumerate}}\\
\multicolumn{4}{|p{\dimexpr 16cm + 6\tabcolsep\relax}|}{\begin{enumerate}[start=8]\item Projeto de sistemas embarcados: metodologia para modelagem, máquinas de estados finitos, ciclo de desenvolvimento de software, diagramas de fluxo de dados, statecharts\end{enumerate}}\\
\multicolumn{4}{|p{\dimexpr 16cm + 6\tabcolsep\relax}|}{}\\
\hline
\multicolumn{4}{|p{16cm}|}{Bibliografia Básica:}\\*
\multicolumn{4}{|p{\dimexpr 16cm + 6\tabcolsep\relax}|}{\begin{itemize}\item LEE, Edward Ashford. Introduction to embedded systems: A cyber-physical systems approach. MIT Press, 2016. (item inexistente na UFOP)\end{itemize}}\\
\multicolumn{4}{|p{\dimexpr 16cm + 6\tabcolsep\relax}|}{\begin{itemize}\item NOERGAARD, Tammy. Embedded systems architecture: a comprehensive guide for engineers and programmers. Amsterdam.\end{itemize}}\\
\multicolumn{4}{|p{\dimexpr 16cm + 6\tabcolsep\relax}|}{\begin{itemize}\item BALL, Stuart R. Embedded microprocessor systems: real world design. 2. ed. Boston: Newnes, 2002.\end{itemize}}\\
\multicolumn{4}{|p{16cm}|}{}\\
\hline
\multicolumn{4}{|p{16cm}|}{Bibliografia Complementar:}\\*
\multicolumn{4}{|p{\dimexpr 16cm + 6\tabcolsep\relax}|}{\begin{itemize}\item WILMSHURST, Tim. Designing embedded systems with PIC microcontrollers: principles and applications. Amsterdã.\end{itemize}}\\
\multicolumn{4}{|p{\dimexpr 16cm + 6\tabcolsep\relax}|}{\begin{itemize}\item PONT, Michael J. Embedded C. Boston, MA: Addison Wesley, 2002.\end{itemize}}\\
\multicolumn{4}{|p{\dimexpr 16cm + 6\tabcolsep\relax}|}{\begin{itemize}\item PEREIRA, Fábio. Microcontroladores PIC: programação em C. 7. ed. São Paulo: Érica, 2007.\end{itemize}}\\
\multicolumn{4}{|p{\dimexpr 16cm + 6\tabcolsep\relax}|}{\begin{itemize}\item HALLINAN, Christopher. Embedded Linux Primer: a practical, real-world approach. 2. ed. New York: Prentice Hall, 2011.\end{itemize}}\\
\multicolumn{4}{|p{\dimexpr 16cm + 6\tabcolsep\relax}|}{\begin{itemize}\item SOUZA, David José de. Desbravando o PIC: ampliado e atualizado para PIC16F628A. 9. ed. São Paulo: Érica, 2005.\end{itemize}}\\
\hline
\end{longtable}
\end{center}

\clearpage
\begin{center}
\begin{tabular}{lccr}
 \multirow{3}{*}{\includegraphics[height=2.7cm]{brasao.png}} &
 \multicolumn{2}{c}{\bfseries UNIVERSIDADE FEDERAL DE OURO PRETO} &
 \ \ \ \ \multirow{3}{*}{\includegraphics[height=2.7cm]{ufop.png}} \\
 & \multicolumn{2}{c}{\bfseries PRÓ-REITORIA DE GRADUAÇÃO} & \\
 & \multicolumn{2}{c}{} & \\
 & \multicolumn{2}{c}{\Large\bfseries PROGRAMA DE DISCIPLINA} & \\
\end{tabular}
\end{center}

\begin{center}
\begin{longtable}{|p{4cm}|p{4cm}|p{4cm}|p{4cm}|}
\hline
\multicolumn{3}{|p{12cm}|}{Nome do Componente Curricular em Português:} &
\multicolumn{1}{p{4cm}|}{Código:} \\ 
\multicolumn{3}{|p{12cm}|}{\textbf{Geoprocessamento e Sistemas de Informação Geográfica}} &
\textbf{BCC443}\\ 
\multicolumn{3}{|p{12cm}|}{Nome do Componente Curricular em Inglês:} & \\ 
\multicolumn{3}{|p{12cm}|}{\textbf{Geoprocessing and Geographic Information Systems}} & \\ 
\hline
\multicolumn{3}{|p{12cm}|}{Nome e Sigla do Departamento} & Unidade Acadêmica: \\ 
\multicolumn{3}{|p{12cm}|}{Departamento de Computação (DECOM)} & {ICEB} \\ 
\hline
\multicolumn{4}{|p{16cm}|}{Modalidade de Oferta:
[X] presencial \hspace{0.6cm}
[ ] semipresencial \hspace{0.6cm}
[ ] à distância}\\
\hline
\multicolumn{2}{|p{8cm}|}{Carga horária semestral} &
\multicolumn{2}{p{8cm}|}{Carga horária semanal}\\
\hline
\multicolumn{1}{|p{4cm}|}{Total} &
\multicolumn{1}{p{4cm}|}{Extensionista} &
\multicolumn{1}{p{4cm}|}{Teórica} &
\multicolumn{1}{p{4cm}|}{Prática} \\ 
\multicolumn{1}{|p{4cm}|}{60\,horas} &
\multicolumn{1}{p{4cm}|}{0\;horas} &
\multicolumn{1}{p{4cm}|}{4\;horas/aula} &
\multicolumn{1}{p{4cm}|}{0\;horas/aula} \\ 
\hline
\multicolumn{4}{|p{16cm}|}{Ementa:}\\
\multicolumn{4}{|p{16cm}|}{}\\
\multicolumn{4}{|p{\dimexpr 16cm + 6\tabcolsep\relax}|}{Representações computacionais do espaço geográfico; introdução à cartografia para geoinformática; introdução aos Sistemas de Informação Geográfica; introdução à modelagem numérica de terreno; introdução ao sistema de posicionamento global; introdução ao sensoriamento remoto; introdução à modelagem de Bancos de Dados Geográficos; introdução à Análise de Dados Geográficos; aplicações em Geoinformática.}\\
\multicolumn{4}{|p{16cm}|}{}\\
\hline
\multicolumn{4}{|p{16cm}|}{Conteúdo programático:}\\*
\multicolumn{4}{|p{\dimexpr 16cm + 6\tabcolsep\relax}|}{\begin{enumerate}[start=1]\item Representações computacionais do espaço geográfico: conceituação de modelos matemático-computacionais, tempo, espaço, escala (extensão e resolução) e processos de mudança; modelos conceituais de campos e objetos; estruturas de dados matriciais (grades numéricas e imagens) e estruturas de dados vetoriais (pontos, linhas e polígonos)\end{enumerate}}\\
\multicolumn{4}{|p{\dimexpr 16cm + 6\tabcolsep\relax}|}{\begin{enumerate}[start=2]\item Introdução à Cartografia para Geoinformática: datum, projeções cartográficas, sistemas de coordenadas geográficas, sistemas de referência geográfica padronizados, escala cartográfica\end{enumerate}}\\
\multicolumn{4}{|p{\dimexpr 16cm + 6\tabcolsep\relax}|}{\begin{enumerate}[start=3]\item Introdução aos Sistemas de Informação Geográfica: arquitetura (dual e integrada), tipologia (desktop, cliente-servidor, web), uso de ferramentas livres (TerraView e QGIS); camada de dados (layers), visões (views), temas (themes) e legendas; aquisição de dados vetoriais e matriciais; interoperabilidade (DXF, SHP, GML); consultas alfanuméricas e espaciais\end{enumerate}}\\
\multicolumn{4}{|p{\dimexpr 16cm + 6\tabcolsep\relax}|}{\begin{enumerate}[start=4]\item Introdução ao Sensoriamento Remoto: conceituação de satélites e sensores remotos; tipos de imagens (óticas, radar e termais); bandas e composição falsa cor (RGB); registro e georreferenciamento de imagens; técnicas de processamento digital de imagens (contraste, reamostragem, filtragem, segmentação e classificação)\end{enumerate}}\\
\multicolumn{4}{|p{\dimexpr 16cm + 6\tabcolsep\relax}|}{\begin{enumerate}[start=5]\item Introdução à Modelagem de Bancos de Dados Geográficos: notação OMT-G; implementação de bancos de dados geográficos em ferramentas livres; consultas espaciais (relatórios) em ferramentas livres; indexação de dados espaciais (Quadtree e R-tree)\end{enumerate}}\\
\multicolumn{4}{|p{\dimexpr 16cm + 6\tabcolsep\relax}|}{\begin{enumerate}[start=6]\item Introdução à Modelagem Numérica de Terreno: modelo digital de elevação (DEM) e redes irregulares de triângulos (TIN); construção de modelos digitais de terrenos a partir de curvas de nível, linhas de drenagem e linhas de crista\end{enumerate}}\\
\multicolumn{4}{|p{\dimexpr 16cm + 6\tabcolsep\relax}|}{\begin{enumerate}[start=7]\item Introdução ao Sistema de Posicionamento Global: sistemas GPS e Galileo; calibração e uso de equipamentos GPS; fontes de erro e falta de precisão; tipologia de equipamentos GPS (navegação, mapeamento, diferencial)\end{enumerate}}\\
\multicolumn{4}{|p{\dimexpr 16cm + 6\tabcolsep\relax}|}{\begin{enumerate}[start=8]\item Introdução à Análise de Dados Geográficos: análise espacial (dados pontuais e de área), geoestatística (krigagem ordinária) e álgebra de mapas (operadores de Tomlin e Egenhofer)\end{enumerate}}\\
\multicolumn{4}{|p{\dimexpr 16cm + 6\tabcolsep\relax}|}{\begin{enumerate}[start=9]\item Aplicações em Geoinformática: mudança de uso e cobertura do solo, dispersão de doenças e cadastro técnico multifinalitário\end{enumerate}}\\
\multicolumn{4}{|p{\dimexpr 16cm + 6\tabcolsep\relax}|}{}\\
\hline
\multicolumn{4}{|p{16cm}|}{Bibliografia Básica:}\\*
\multicolumn{4}{|p{\dimexpr 16cm + 6\tabcolsep\relax}|}{\begin{itemize}\item BLASCHKE, Thomas; KUX, Hermann (org.); tradução Hermann Kux. Sensoriamento remoto e SIG avançados: novos sistemas sensores, métodos inovadores. São Paulo: Oficina de Textos, c2005.\end{itemize}}\\
\multicolumn{4}{|p{\dimexpr 16cm + 6\tabcolsep\relax}|}{\begin{itemize}\item Anatomia de sistemas de informação geográfica. Campinas: UNICAMP, 1996.\end{itemize}}\\
\multicolumn{4}{|p{\dimexpr 16cm + 6\tabcolsep\relax}|}{\begin{itemize}\item GONZALEZ, Rafael C. Digital image processing. 3rd ed. Upper Saddle River, NJ: Pearson Prentice Hall, c2008.\end{itemize}}\\
\multicolumn{4}{|p{\dimexpr 16cm + 6\tabcolsep\relax}|}{\begin{itemize}\item JENSEN, John R.; tradução José Carlos Neves Epiphanio ... [et al.]. Sensoriamento remoto do ambiente: uma perspectiva em recursos terrestres. São José dos Campos: Parêntese, 2011.\end{itemize}}\\
\multicolumn{4}{|p{\dimexpr 16cm + 6\tabcolsep\relax}|}{\begin{itemize}\item CÂMARA, Gilberto. Geoprocessamento para projetos ambientais. São José dos Campos: INPE, 1996.\end{itemize}}\\
\multicolumn{4}{|p{\dimexpr 16cm + 6\tabcolsep\relax}|}{\begin{itemize}\item Conceitos básicos de sistemas de informação geográfica e cartografia aplicados à saúde. Brasília: Organização Pan-Americana de Saúde, 2000.\end{itemize}}\\
\multicolumn{4}{|p{16cm}|}{}\\
\hline
\multicolumn{4}{|p{16cm}|}{Bibliografia Complementar:}\\*
\multicolumn{4}{|p{\dimexpr 16cm + 6\tabcolsep\relax}|}{\begin{itemize}\item Bancos de dados geográficos. Curitiba: MundoGeo, 2005.\end{itemize}}\\
\multicolumn{4}{|p{\dimexpr 16cm + 6\tabcolsep\relax}|}{\begin{itemize}\item Análise espacial de dados geográficos. Planaltina, DF: EMBRAPA Cerrados, 2004.\end{itemize}}\\
\multicolumn{4}{|p{\dimexpr 16cm + 6\tabcolsep\relax}|}{\begin{itemize}\item CANTY, Morton J. Image analysis, classification and change detection in remote sensing: with algorithms for ENVI/IDL. Boca Raton, FL: CRC/Taylor \& Francis, 2007.\end{itemize}}\\
\multicolumn{4}{|p{\dimexpr 16cm + 6\tabcolsep\relax}|}{\begin{itemize}\item NOVO, Evlyn M. L. de Moraes. Sensoriamento remoto: princípios e aplicações. 2. ed. São Paulo: E. Blucher, c1992.\end{itemize}}\\
\multicolumn{4}{|p{\dimexpr 16cm + 6\tabcolsep\relax}|}{\begin{itemize}\item MOURA, Ana Clara Mourão. Geoprocessamento na gestão e planejamento urbano. Belo Horizonte: Difusora, 2003.\end{itemize}}\\
\hline
\end{longtable}
\end{center}

\clearpage
\begin{center}
\begin{tabular}{lccr}
 \multirow{3}{*}{\includegraphics[height=2.7cm]{brasao.png}} &
 \multicolumn{2}{c}{\bfseries UNIVERSIDADE FEDERAL DE OURO PRETO} &
 \ \ \ \ \multirow{3}{*}{\includegraphics[height=2.7cm]{ufop.png}} \\
 & \multicolumn{2}{c}{\bfseries PRÓ-REITORIA DE GRADUAÇÃO} & \\
 & \multicolumn{2}{c}{} & \\
 & \multicolumn{2}{c}{\Large\bfseries PROGRAMA DE DISCIPLINA} & \\
\end{tabular}
\end{center}

\begin{center}
\begin{longtable}{|p{4cm}|p{4cm}|p{4cm}|p{4cm}|}
\hline
\multicolumn{3}{|p{12cm}|}{Nome do Componente Curricular em Português:} &
\multicolumn{1}{p{4cm}|}{Código:} \\ 
\multicolumn{3}{|p{12cm}|}{\textbf{Modelagem e Simulação de Sistemas Terrestres}} &
\textbf{BCC445}\\ 
\multicolumn{3}{|p{12cm}|}{Nome do Componente Curricular em Inglês:} & \\ 
\multicolumn{3}{|p{12cm}|}{\textbf{Modeling and Simulation of Land Systems}} & \\ 
\hline
\multicolumn{3}{|p{12cm}|}{Nome e Sigla do Departamento} & Unidade Acadêmica: \\ 
\multicolumn{3}{|p{12cm}|}{Departamento de Computação (DECOM)} & {ICEB} \\ 
\hline
\multicolumn{4}{|p{16cm}|}{Modalidade de Oferta:
[X] presencial \hspace{0.6cm}
[ ] semipresencial \hspace{0.6cm}
[ ] à distância}\\
\hline
\multicolumn{2}{|p{8cm}|}{Carga horária semestral} &
\multicolumn{2}{p{8cm}|}{Carga horária semanal}\\
\hline
\multicolumn{1}{|p{4cm}|}{Total} &
\multicolumn{1}{p{4cm}|}{Extensionista} &
\multicolumn{1}{p{4cm}|}{Teórica} &
\multicolumn{1}{p{4cm}|}{Prática} \\ 
\multicolumn{1}{|p{4cm}|}{60\,horas} &
\multicolumn{1}{p{4cm}|}{0\;horas} &
\multicolumn{1}{p{4cm}|}{4\;horas/aula} &
\multicolumn{1}{p{4cm}|}{0\;horas/aula} \\ 
\hline
\multicolumn{4}{|p{16cm}|}{Ementa:}\\
\multicolumn{4}{|p{16cm}|}{}\\
\multicolumn{4}{|p{\dimexpr 16cm + 6\tabcolsep\relax}|}{Introdução à modelagem computacional de fenômenos geográficos. Fundamentação teórica da ciência de sistemas terrestres. Conceituação e representações computacionais para escala, espaço, tempo e comportamento. Introdução aos modelos dinâmicos espacialmente-explícitos. Introdução à simulação computacional. Introdução aos paradigmas de modelagem dos sistemas terrestres. Ferramentas livres para modelagem dos sistemas terrestres. Modelagem de sistemas terrestres em múltiplas escalas. Aplicações em modelagem e simulação dos sistemas terrestres. Métodos matemáticos aplicados à modelagem dos sistemas terrestres.}\\
\multicolumn{4}{|p{16cm}|}{}\\
\hline
\multicolumn{4}{|p{16cm}|}{Conteúdo programático:}\\*
\multicolumn{4}{|p{\dimexpr 16cm + 6\tabcolsep\relax}|}{\begin{enumerate}[start=1]\item Introdução à Modelagem Computacional de Fenômenos Geográficos\end{enumerate}}\\
\multicolumn{4}{|p{\dimexpr 16cm + 6\tabcolsep\relax}|}{\begin{enumerate}[start=2]\item Fundamentação teórica da ciência dos sistemas terrestres: modelos matemático-computacionais, o processo de modelagem, simulação de processos, estoques e fluxos de energia, feedbacks, a ciência de sistemas terrestres\end{enumerate}}\\
\multicolumn{4}{|p{\dimexpr 16cm + 6\tabcolsep\relax}|}{\begin{enumerate}[start=3]\item Modelos dinâmicos espacialmente-explícitos: definição, taxonomia, metodologia de desenvolvimento, verificação, calibração, validação\end{enumerate}}\\
\multicolumn{4}{|p{\dimexpr 16cm + 6\tabcolsep\relax}|}{\begin{enumerate}[start=4]\item Introdução à simulação computacional: simulação discreta e simulação por eventos discretos, vantagens e desvantagens, implementação de simuladores\end{enumerate}}\\
\multicolumn{4}{|p{\dimexpr 16cm + 6\tabcolsep\relax}|}{\begin{enumerate}[start=5]\item Conceituação e representação computacional de escala, espaço, tempo e comportamento\end{enumerate}}\\
\multicolumn{4}{|p{\dimexpr 16cm + 6\tabcolsep\relax}|}{\begin{enumerate}[start=6]\item Paradigmas de modelagem dos sistemas terrestres: dinâmica de sistemas, autômatos celulares, múltiplos agentes\end{enumerate}}\\
\multicolumn{4}{|p{\dimexpr 16cm + 6\tabcolsep\relax}|}{\begin{enumerate}[start=7]\item Ferramentas livres para modelagem ambiental: Vensim, NetLogo, Repast, TerraME\end{enumerate}}\\
\multicolumn{4}{|p{\dimexpr 16cm + 6\tabcolsep\relax}|}{\begin{enumerate}[start=8]\item Modelagem de sistemas terrestres em múltiplas escalas: requisitos de software, integração a sistemas de informação geográfica, plataformas de desenvolvimento\end{enumerate}}\\
\multicolumn{4}{|p{\dimexpr 16cm + 6\tabcolsep\relax}|}{\begin{enumerate}[start=9]\item Aplicações: modelos hidrológicos, modelos climáticos, modelos de dinâmica populacional, modelos de mudança de uso e cobertura da Terra, modelos de dispersão de espécies, modelos atmosféricos, modelos epidemiológicos, modelos econométricos\end{enumerate}}\\
\multicolumn{4}{|p{\dimexpr 16cm + 6\tabcolsep\relax}|}{\begin{enumerate}[start=10]\item Métodos matemáticos aplicados à modelagem de sistemas terrestres: inferência estatística, probabilidade e análise espacial\end{enumerate}}\\
\multicolumn{4}{|p{\dimexpr 16cm + 6\tabcolsep\relax}|}{}\\
\hline
\multicolumn{4}{|p{16cm}|}{Bibliografia Básica:}\\*
\multicolumn{4}{|p{\dimexpr 16cm + 6\tabcolsep\relax}|}{\begin{itemize}\item AXELROD, Robert M. The complexity of cooperation: agent-based models of competition and collaboration. Princeton, N.J.: Princeton University Press, 1997. (item inexistente na UFOP)\end{itemize}}\\
\multicolumn{4}{|p{\dimexpr 16cm + 6\tabcolsep\relax}|}{\begin{itemize}\item GOODCHILD, Michael F.; STEYAERT, Louis T.; PARKS, Bradley O.; JOHNSTON, Carol ... [et al.]. GIS and environmental modeling: progress and research issues. Fort Collins: GIS World Books, c1996.\end{itemize}}\\
\multicolumn{4}{|p{\dimexpr 16cm + 6\tabcolsep\relax}|}{\begin{itemize}\item WAINWRIGHT, John; MULLIGAN, Mark (eds.). Environmental modelling: finding simplicity in complexity. Chichester, West Sussex, England.\end{itemize}}\\
\multicolumn{4}{|p{\dimexpr 16cm + 6\tabcolsep\relax}|}{\begin{itemize}\item ZEIGLER, Bernard P.; PRAEHOFER, Herbert; KIM, Tag Gon. Theory of modeling and simulation: integrating discrete event and continuous complex dynamic systems. 2nd ed. San Diego: Academic, c2000.\end{itemize}}\\
\multicolumn{4}{|p{\dimexpr 16cm + 6\tabcolsep\relax}|}{\begin{itemize}\item FORD, Andrew. Modeling the environment: an introduction to system dynamics models of environmental systems. 2. ed. Washington, D.C.: Island Press, c1999.\end{itemize}}\\
\multicolumn{4}{|p{\dimexpr 16cm + 6\tabcolsep\relax}|}{\begin{itemize}\item EPSTEIN, Joshua M. Generative social science: studies in agent-based computational modeling. Princeton, NJ.\end{itemize}}\\
\multicolumn{4}{|p{\dimexpr 16cm + 6\tabcolsep\relax}|}{\begin{itemize}\item Integrating geographic information systems and agent-based modeling techniques for simulating social and ecological processes. Oxford.\end{itemize}}\\
\multicolumn{4}{|p{\dimexpr 16cm + 6\tabcolsep\relax}|}{\begin{itemize}\item MEADOWS, Donella H.; ed. Diana Wright. Thinking in systems: a primer. White River Junction, Vt.: Chelsea Green Pub., c2008.\end{itemize}}\\
\multicolumn{4}{|p{\dimexpr 16cm + 6\tabcolsep\relax}|}{\begin{itemize}\item LANTUÉJOUL, Christian. Geostatistical simulation: models and algorithms. Berlin: Springer-Verlag, 2002.\end{itemize}}\\
\multicolumn{4}{|p{16cm}|}{}\\
\hline
\multicolumn{4}{|p{16cm}|}{Bibliografia Complementar:}\\*
\multicolumn{4}{|p{\dimexpr 16cm + 6\tabcolsep\relax}|}{\begin{itemize}\item BATTY, Michael. Cities and complexity: understanding cities with cellular automata, agent-based models, and fractals. Cambridge, Mass.: MIT Press, 2005. (item inexistente na UFOP)\end{itemize}}\\
\multicolumn{4}{|p{\dimexpr 16cm + 6\tabcolsep\relax}|}{\begin{itemize}\item SKIDMORE, Andrew. Environmental modelling with GIS and remote sensing. London: Taylor \& Francis, 2002. (item inexistente na UFOP)\end{itemize}}\\
\multicolumn{4}{|p{\dimexpr 16cm + 6\tabcolsep\relax}|}{\begin{itemize}\item SMITH, Jo U. Introduction to environmental modelling. Oxford: Oxford University Press, 2007.\end{itemize}}\\
\hline
\end{longtable}
\end{center}

\clearpage
\begin{center}
\begin{tabular}{lccr}
 \multirow{3}{*}{\includegraphics[height=2.7cm]{brasao.png}} &
 \multicolumn{2}{c}{\bfseries UNIVERSIDADE FEDERAL DE OURO PRETO} &
 \ \ \ \ \multirow{3}{*}{\includegraphics[height=2.7cm]{ufop.png}} \\
 & \multicolumn{2}{c}{\bfseries PRÓ-REITORIA DE GRADUAÇÃO} & \\
 & \multicolumn{2}{c}{} & \\
 & \multicolumn{2}{c}{\Large\bfseries PROGRAMA DE DISCIPLINA} & \\
\end{tabular}
\end{center}

\begin{center}
\begin{longtable}{|p{4cm}|p{4cm}|p{4cm}|p{4cm}|}
\hline
\multicolumn{3}{|p{12cm}|}{Nome do Componente Curricular em Português:} &
\multicolumn{1}{p{4cm}|}{Código:} \\ 
\multicolumn{3}{|p{12cm}|}{\textbf{Programação em Lógica}} &
\textbf{BCC446}\\ 
\multicolumn{3}{|p{12cm}|}{Nome do Componente Curricular em Inglês:} & \\ 
\multicolumn{3}{|p{12cm}|}{\textbf{Logic Programming}} & \\ 
\hline
\multicolumn{3}{|p{12cm}|}{Nome e Sigla do Departamento} & Unidade Acadêmica: \\ 
\multicolumn{3}{|p{12cm}|}{Departamento de Computação (DECOM)} & {ICEB} \\ 
\hline
\multicolumn{4}{|p{16cm}|}{Modalidade de Oferta:
[X] presencial \hspace{0.6cm}
[ ] semipresencial \hspace{0.6cm}
[ ] à distância}\\
\hline
\multicolumn{2}{|p{8cm}|}{Carga horária semestral} &
\multicolumn{2}{p{8cm}|}{Carga horária semanal}\\
\hline
\multicolumn{1}{|p{4cm}|}{Total} &
\multicolumn{1}{p{4cm}|}{Extensionista} &
\multicolumn{1}{p{4cm}|}{Teórica} &
\multicolumn{1}{p{4cm}|}{Prática} \\ 
\multicolumn{1}{|p{4cm}|}{60\,horas} &
\multicolumn{1}{p{4cm}|}{0\;horas} &
\multicolumn{1}{p{4cm}|}{4\;horas/aula} &
\multicolumn{1}{p{4cm}|}{0\;horas/aula} \\ 
\hline
\multicolumn{4}{|p{16cm}|}{Ementa:}\\
\multicolumn{4}{|p{16cm}|}{}\\
\multicolumn{4}{|p{\dimexpr 16cm + 6\tabcolsep\relax}|}{Fundamentos teóricos da programação lógica: unificação e resolução; conceitos de linguagens de programação lógica; estudo de uma linguagem de programação lógica; aplicações.}\\
\multicolumn{4}{|p{16cm}|}{}\\
\hline
\multicolumn{4}{|p{16cm}|}{Conteúdo programático:}\\*
\multicolumn{4}{|p{\dimexpr 16cm + 6\tabcolsep\relax}|}{\begin{enumerate}[start=1]\item Programação declarativa\end{enumerate}}\\
\multicolumn{4}{|p{\dimexpr 16cm + 6\tabcolsep\relax}|}{\begin{enumerate}[start=2]\item Introdução a Prolog\end{enumerate}}\\
\multicolumn{4}{|p{\dimexpr 16cm + 6\tabcolsep\relax}|}{\begin{enumerate}[start=3]\item Sintaxe e Semântica de programas Prolog\end{enumerate}}\\
\multicolumn{4}{|p{\dimexpr 16cm + 6\tabcolsep\relax}|}{\begin{enumerate}[start=4]\item Listas, operadores, aritmética\end{enumerate}}\\
\multicolumn{4}{|p{\dimexpr 16cm + 6\tabcolsep\relax}|}{\begin{enumerate}[start=5]\item Estruturas\end{enumerate}}\\
\multicolumn{4}{|p{\dimexpr 16cm + 6\tabcolsep\relax}|}{\begin{enumerate}[start=6]\item Backtracking\end{enumerate}}\\
\multicolumn{4}{|p{\dimexpr 16cm + 6\tabcolsep\relax}|}{\begin{enumerate}[start=7]\item Entrada e saída\end{enumerate}}\\
\multicolumn{4}{|p{\dimexpr 16cm + 6\tabcolsep\relax}|}{\begin{enumerate}[start=8]\item Predicados pré-definidos\end{enumerate}}\\
\multicolumn{4}{|p{\dimexpr 16cm + 6\tabcolsep\relax}|}{\begin{enumerate}[start=9]\item Estilo e técnicas de programação\end{enumerate}}\\
\multicolumn{4}{|p{\dimexpr 16cm + 6\tabcolsep\relax}|}{\begin{enumerate}[start=10]\item Operações em estruturas de dados\end{enumerate}}\\
\multicolumn{4}{|p{\dimexpr 16cm + 6\tabcolsep\relax}|}{\begin{enumerate}[start=11]\item Relação de Prolog com Lógica: cálculo de predicados, forma clausal, resolução, prova de teoremas e cláusulas de Horn\end{enumerate}}\\
\multicolumn{4}{|p{\dimexpr 16cm + 6\tabcolsep\relax}|}{}\\
\hline
\multicolumn{4}{|p{16cm}|}{Bibliografia Básica:}\\*
\multicolumn{4}{|p{\dimexpr 16cm + 6\tabcolsep\relax}|}{\begin{itemize}\item BRATKO, Ivan. Prolog programming for artificial intelligence. 3rd ed. Harlow: Pearson Education, 2001.\end{itemize}}\\
\multicolumn{4}{|p{\dimexpr 16cm + 6\tabcolsep\relax}|}{\begin{itemize}\item CASANOVA, Marco Antonio. Programação em lógica e a linguagem PROLOG. São Paulo: E. Blucher, c1987.\end{itemize}}\\
\multicolumn{4}{|p{\dimexpr 16cm + 6\tabcolsep\relax}|}{\begin{itemize}\item CLOCKSIN, W. F. Programming in Prolog. 4th ed. Berlin: Springer-Verlag, c1994.\end{itemize}}\\
\multicolumn{4}{|p{16cm}|}{}\\
\hline
\multicolumn{4}{|p{16cm}|}{Bibliografia Complementar:}\\*
\multicolumn{4}{|p{\dimexpr 16cm + 6\tabcolsep\relax}|}{\begin{itemize}\item O'KEEFE, Richard. The Craft of Prolog. MIT Press, 2009. (item inexistente na UFOP)\end{itemize}}\\
\multicolumn{4}{|p{\dimexpr 16cm + 6\tabcolsep\relax}|}{\begin{itemize}\item BRAMER, Max. Logic Programming with Prolog. 2. ed. Springer, 2014. (item inexistente na UFOP)\end{itemize}}\\
\multicolumn{4}{|p{\dimexpr 16cm + 6\tabcolsep\relax}|}{\begin{itemize}\item MAIER, David; WARREN, David S. Computing with logic: logic programming with prolog. Menlo Park, Calif.: Benjamin/Cummings, c1988.\end{itemize}}\\
\multicolumn{4}{|p{\dimexpr 16cm + 6\tabcolsep\relax}|}{\begin{itemize}\item DERANSART, Pierre; ED-DBALI, AbdelAli; CERVONI, Laurent. Prolog: the standard: reference manual. Berlin: Springer-Verlag, 1996.\end{itemize}}\\
\multicolumn{4}{|p{\dimexpr 16cm + 6\tabcolsep\relax}|}{\begin{itemize}\item STERLING, Leon; SHAPIRO, Ehud; foreword by David H. D. Warren. The art of prolog: advanced programming techniques. 2nd ed. Cambridge: The MIT Press, 1994.\end{itemize}}\\
\hline
\end{longtable}
\end{center}

\clearpage
\begin{center}
\begin{tabular}{lccr}
 \multirow{3}{*}{\includegraphics[height=2.7cm]{brasao.png}} &
 \multicolumn{2}{c}{\bfseries UNIVERSIDADE FEDERAL DE OURO PRETO} &
 \ \ \ \ \multirow{3}{*}{\includegraphics[height=2.7cm]{ufop.png}} \\
 & \multicolumn{2}{c}{\bfseries PRÓ-REITORIA DE GRADUAÇÃO} & \\
 & \multicolumn{2}{c}{} & \\
 & \multicolumn{2}{c}{\Large\bfseries PROGRAMA DE DISCIPLINA} & \\
\end{tabular}
\end{center}

\begin{center}
\begin{longtable}{|p{4cm}|p{4cm}|p{4cm}|p{4cm}|}
\hline
\multicolumn{3}{|p{12cm}|}{Nome do Componente Curricular em Português:} &
\multicolumn{1}{p{4cm}|}{Código:} \\ 
\multicolumn{3}{|p{12cm}|}{\textbf{Recuperação de Informação na Web}} &
\textbf{BCC449}\\ 
\multicolumn{3}{|p{12cm}|}{Nome do Componente Curricular em Inglês:} & \\ 
\multicolumn{3}{|p{12cm}|}{\textbf{Web Information Retrieval}} & \\ 
\hline
\multicolumn{3}{|p{12cm}|}{Nome e Sigla do Departamento} & Unidade Acadêmica: \\ 
\multicolumn{3}{|p{12cm}|}{Departamento de Computação (DECOM)} & {ICEB} \\ 
\hline
\multicolumn{4}{|p{16cm}|}{Modalidade de Oferta:
[X] presencial \hspace{0.6cm}
[ ] semipresencial \hspace{0.6cm}
[ ] à distância}\\
\hline
\multicolumn{2}{|p{8cm}|}{Carga horária semestral} &
\multicolumn{2}{p{8cm}|}{Carga horária semanal}\\
\hline
\multicolumn{1}{|p{4cm}|}{Total} &
\multicolumn{1}{p{4cm}|}{Extensionista} &
\multicolumn{1}{p{4cm}|}{Teórica} &
\multicolumn{1}{p{4cm}|}{Prática} \\ 
\multicolumn{1}{|p{4cm}|}{60\,horas} &
\multicolumn{1}{p{4cm}|}{0\;horas} &
\multicolumn{1}{p{4cm}|}{4\;horas/aula} &
\multicolumn{1}{p{4cm}|}{0\;horas/aula} \\ 
\hline
\multicolumn{4}{|p{16cm}|}{Ementa:}\\
\multicolumn{4}{|p{16cm}|}{}\\
\multicolumn{4}{|p{\dimexpr 16cm + 6\tabcolsep\relax}|}{Modelagem de sistemas de recuperação de informação. Caracterização da web. Funcionamento de sistemas de coleta. Indexação distribuída. Linguagens de consulta. Operações sobre consultas. Geração de funções para ordenação de resultados. Avaliação de resultados.}\\
\multicolumn{4}{|p{16cm}|}{}\\
\hline
\multicolumn{4}{|p{16cm}|}{Conteúdo programático:}\\*
\multicolumn{4}{|p{\dimexpr 16cm + 6\tabcolsep\relax}|}{\begin{enumerate}[start=1]\item Introdução\end{enumerate}}\\
\multicolumn{4}{|p{\dimexpr 16cm + 6\tabcolsep\relax}|}{\begin{enumerate}[start=2]\item Recuperação usando modelo booleano\end{enumerate}}\\
\multicolumn{4}{|p{\dimexpr 16cm + 6\tabcolsep\relax}|}{\begin{enumerate}[start=3]\item Vocabulário de termos e listas invertidas\end{enumerate}}\\
\multicolumn{4}{|p{\dimexpr 16cm + 6\tabcolsep\relax}|}{\begin{enumerate}[start=4]\item Dicionários e recuperação com tolerância\end{enumerate}}\\
\multicolumn{4}{|p{\dimexpr 16cm + 6\tabcolsep\relax}|}{\begin{enumerate}[start=5]\item Construção de índice\end{enumerate}}\\
\multicolumn{4}{|p{\dimexpr 16cm + 6\tabcolsep\relax}|}{\begin{enumerate}[start=6]\item Compressão de índice\end{enumerate}}\\
\multicolumn{4}{|p{\dimexpr 16cm + 6\tabcolsep\relax}|}{\begin{enumerate}[start=7]\item Pontuação, peso de termo e modelo de espaço vetorial\end{enumerate}}\\
\multicolumn{4}{|p{\dimexpr 16cm + 6\tabcolsep\relax}|}{\begin{enumerate}[start=8]\item Computação de pontos em um sistema de busca completo\end{enumerate}}\\
\multicolumn{4}{|p{\dimexpr 16cm + 6\tabcolsep\relax}|}{\begin{enumerate}[start=9]\item Avaliação em recuperação de informação\end{enumerate}}\\
\multicolumn{4}{|p{\dimexpr 16cm + 6\tabcolsep\relax}|}{\begin{enumerate}[start=10]\item Feedback de relevância e expansão de consultas\end{enumerate}}\\
\multicolumn{4}{|p{\dimexpr 16cm + 6\tabcolsep\relax}|}{\begin{enumerate}[start=11]\item Introdução à busca Web\end{enumerate}}\\
\multicolumn{4}{|p{\dimexpr 16cm + 6\tabcolsep\relax}|}{\begin{enumerate}[start=12]\item Coleta na Web e indexação\end{enumerate}}\\
\multicolumn{4}{|p{\dimexpr 16cm + 6\tabcolsep\relax}|}{\begin{enumerate}[start=13]\item Análise de links\end{enumerate}}\\
\multicolumn{4}{|p{\dimexpr 16cm + 6\tabcolsep\relax}|}{}\\
\hline
\multicolumn{4}{|p{16cm}|}{Bibliografia Básica:}\\*
\multicolumn{4}{|p{\dimexpr 16cm + 6\tabcolsep\relax}|}{\begin{itemize}\item MANNING, Christopher D. Introduction to information retrieval. New York: Cambridge University Press, 2008.\end{itemize}}\\
\multicolumn{4}{|p{\dimexpr 16cm + 6\tabcolsep\relax}|}{\begin{itemize}\item LIU, Bing. Web data mining: exploring hyperlinks, contents, and usage data. Berlin.\end{itemize}}\\
\multicolumn{4}{|p{\dimexpr 16cm + 6\tabcolsep\relax}|}{\begin{itemize}\item BAEZA-YATES, Ricardo; RIBEIRO-NETO, Berthier; tradução Leandro Krug Wives, Viviane Pereira Moreira. Recuperação de informação: conceitos e tecnologia das máquinas de busca. 2. ed. Porto Alegre: Bookman, 2013.\end{itemize}}\\
\multicolumn{4}{|p{16cm}|}{}\\
\hline
\multicolumn{4}{|p{16cm}|}{Bibliografia Complementar:}\\*
\multicolumn{4}{|p{\dimexpr 16cm + 6\tabcolsep\relax}|}{\begin{itemize}\item GOLDSCHMIDT, Ronaldo; PASSOS, Emmanuel. Data Mining: um guia prático: conceitos, técnicas, ferramentas, orientações e aplicações. Rio de Janeiro: Elsevier.\end{itemize}}\\
\multicolumn{4}{|p{\dimexpr 16cm + 6\tabcolsep\relax}|}{\begin{itemize}\item CHAKRABARTI, Soumen. Mining the Web: discovering knowledge from hypertext data. San Francisco, CA: Morgan Kaufmann Publishers, c2003.\end{itemize}}\\
\multicolumn{4}{|p{\dimexpr 16cm + 6\tabcolsep\relax}|}{\begin{itemize}\item HAN, Jiawei; KAMBER, Micheline; PEI, Jian. Data mining: concepts and techniques. 3. ed. Boston: Elsevier, 2012.\end{itemize}}\\
\multicolumn{4}{|p{\dimexpr 16cm + 6\tabcolsep\relax}|}{\begin{itemize}\item WITTEN, Ian H.; MOFFAT, Alistair; BELL, Timothy C. Managing gigabytes: compressing and indexing documents and images. 2. ed. San Francisco, Calif.: Morgan Kaufmann Publishers, 1999.\end{itemize}}\\
\multicolumn{4}{|p{\dimexpr 16cm + 6\tabcolsep\relax}|}{\begin{itemize}\item HEAPS, H. S. Information retrieval, computational and theoretical aspects. New York.\end{itemize}}\\
\hline
\end{longtable}
\end{center}

\clearpage
\begin{center}
\begin{tabular}{lccr}
 \multirow{3}{*}{\includegraphics[height=2.7cm]{brasao.png}} &
 \multicolumn{2}{c}{\bfseries UNIVERSIDADE FEDERAL DE OURO PRETO} &
 \ \ \ \ \multirow{3}{*}{\includegraphics[height=2.7cm]{ufop.png}} \\
 & \multicolumn{2}{c}{\bfseries PRÓ-REITORIA DE GRADUAÇÃO} & \\
 & \multicolumn{2}{c}{} & \\
 & \multicolumn{2}{c}{\Large\bfseries PROGRAMA DE DISCIPLINA} & \\
\end{tabular}
\end{center}

\begin{center}
\begin{longtable}{|p{4cm}|p{4cm}|p{4cm}|p{4cm}|}
\hline
\multicolumn{3}{|p{12cm}|}{Nome do Componente Curricular em Português:} &
\multicolumn{1}{p{4cm}|}{Código:} \\ 
\multicolumn{3}{|p{12cm}|}{\textbf{Gerência de Dados na Web}} &
\textbf{BCC450}\\ 
\multicolumn{3}{|p{12cm}|}{Nome do Componente Curricular em Inglês:} & \\ 
\multicolumn{3}{|p{12cm}|}{\textbf{Web Data Management}} & \\ 
\hline
\multicolumn{3}{|p{12cm}|}{Nome e Sigla do Departamento} & Unidade Acadêmica: \\ 
\multicolumn{3}{|p{12cm}|}{Departamento de Computação (DECOM)} & {ICEB} \\ 
\hline
\multicolumn{4}{|p{16cm}|}{Modalidade de Oferta:
[X] presencial \hspace{0.6cm}
[ ] semipresencial \hspace{0.6cm}
[ ] à distância}\\
\hline
\multicolumn{2}{|p{8cm}|}{Carga horária semestral} &
\multicolumn{2}{p{8cm}|}{Carga horária semanal}\\
\hline
\multicolumn{1}{|p{4cm}|}{Total} &
\multicolumn{1}{p{4cm}|}{Extensionista} &
\multicolumn{1}{p{4cm}|}{Teórica} &
\multicolumn{1}{p{4cm}|}{Prática} \\ 
\multicolumn{1}{|p{4cm}|}{60\,horas} &
\multicolumn{1}{p{4cm}|}{0\;horas} &
\multicolumn{1}{p{4cm}|}{4\;horas/aula} &
\multicolumn{1}{p{4cm}|}{0\;horas/aula} \\ 
\hline
\multicolumn{4}{|p{16cm}|}{Ementa:}\\
\multicolumn{4}{|p{16cm}|}{}\\
\multicolumn{4}{|p{\dimexpr 16cm + 6\tabcolsep\relax}|}{Caracterização dos dados da Web; modelagem de dados semiestruturados; linguagens de consulta; extração de dados da Web; integração de dados da web; aplicações.}\\
\multicolumn{4}{|p{16cm}|}{}\\
\hline
\multicolumn{4}{|p{16cm}|}{Conteúdo programático:}\\*
\multicolumn{4}{|p{\dimexpr 16cm + 6\tabcolsep\relax}|}{\begin{enumerate}[start=1]\item Dados da Web: características dos dados da Web, Web versus bancos de dados, arquitetura das aplicações Web\end{enumerate}}\\
\multicolumn{4}{|p{\dimexpr 16cm + 6\tabcolsep\relax}|}{\begin{enumerate}[start=2]\item Modelagem de dados semi-estruturados e XML: características dos dados semi-estruturados, modelos para representação de dados semi-estruturados, XML, representação relacional de dados XML, publicação de dados relacionais através de XML, descoberta de estrutura de dados XML, acessórios para gerência de dados XML, compressão de dados XML\end{enumerate}}\\
\multicolumn{4}{|p{\dimexpr 16cm + 6\tabcolsep\relax}|}{\begin{enumerate}[start=3]\item Linguagens de consulta: linguagens de consulta para a Web, linguagens de consulta para dados semi-estruturados, linguagens de consulta e manipulação de dados XML\end{enumerate}}\\
\multicolumn{4}{|p{\dimexpr 16cm + 6\tabcolsep\relax}|}{\begin{enumerate}[start=4]\item Extração de dados de fontes da Web: caracterização do problema, principais abordagens, ferramentas, avaliação da qualidade das extrações\end{enumerate}}\\
\multicolumn{4}{|p{\dimexpr 16cm + 6\tabcolsep\relax}|}{\begin{enumerate}[start=5]\item Integração de dados da Web: caracterização do problema de integração, abordagens e arquiteturas, fusão de dados, medidas de similaridade\end{enumerate}}\\
\multicolumn{4}{|p{\dimexpr 16cm + 6\tabcolsep\relax}|}{\begin{enumerate}[start=6]\item Tópicos avançados: processamento de consultas em documentos XML, acesso a dados da Web oculta, reconhecimento e extração de entidades, Web semântica, Web de dados, coleta de páginas da Web, bibliotecas digitais\end{enumerate}}\\
\multicolumn{4}{|p{\dimexpr 16cm + 6\tabcolsep\relax}|}{}\\
\hline
\multicolumn{4}{|p{16cm}|}{Bibliografia Básica:}\\*
\multicolumn{4}{|p{\dimexpr 16cm + 6\tabcolsep\relax}|}{\begin{itemize}\item GRAVES, Mark. Projeto de banco de dados com XML. São Paulo: Pearson Education, 2003.\end{itemize}}\\
\multicolumn{4}{|p{\dimexpr 16cm + 6\tabcolsep\relax}|}{\begin{itemize}\item BREITMAN, Karin Koogan. Web semântica: a internet do futuro. Rio de Janeiro: LTC, 2006.\end{itemize}}\\
\multicolumn{4}{|p{\dimexpr 16cm + 6\tabcolsep\relax}|}{\begin{itemize}\item BAEZA-YATES, R. Modern information retrieval. New York: Addison Wesley, 2011.\end{itemize}}\\
\multicolumn{4}{|p{\dimexpr 16cm + 6\tabcolsep\relax}|}{\begin{itemize}\item ABITEBOUL, S. Gerenciando dados na WEB. Rio de Janeiro: Campus, 2000.\end{itemize}}\\
\multicolumn{4}{|p{16cm}|}{}\\
\hline
\multicolumn{4}{|p{16cm}|}{Bibliografia Complementar:}\\*
\multicolumn{4}{|p{\dimexpr 16cm + 6\tabcolsep\relax}|}{\begin{itemize}\item ANTONIOU, G. A semantic web primer. Cambridge: MIT Press, 2008.\end{itemize}}\\
\multicolumn{4}{|p{\dimexpr 16cm + 6\tabcolsep\relax}|}{\begin{itemize}\item ELMASRI, Ramez. Sistemas de banco de dados. 6. ed. São Paulo: Pearson.\end{itemize}}\\
\multicolumn{4}{|p{\dimexpr 16cm + 6\tabcolsep\relax}|}{\begin{itemize}\item KANTARDZIC, Mehmed. Data mining: concepts, models, methods, and algorithms. Hoboken, NJ: Wiley-Interscience.\end{itemize}}\\
\multicolumn{4}{|p{\dimexpr 16cm + 6\tabcolsep\relax}|}{\begin{itemize}\item LIU, Bing. Web data mining: exploring hyperlinks, contents, and usage data. Berlin.\end{itemize}}\\
\multicolumn{4}{|p{\dimexpr 16cm + 6\tabcolsep\relax}|}{\begin{itemize}\item GARCIA-MOLINA, Hector; ULLMAN, Jeffrey D.; WIDOM, Jennifer. Database systems: the complete book. 2nd ed. Upper Saddle River: Pearson Prentice Hall, c2009.\end{itemize}}\\
\hline
\end{longtable}
\end{center}

\clearpage
\begin{center}
\begin{tabular}{lccr}
 \multirow{3}{*}{\includegraphics[height=2.7cm]{brasao.png}} &
 \multicolumn{2}{c}{\bfseries UNIVERSIDADE FEDERAL DE OURO PRETO} &
 \ \ \ \ \multirow{3}{*}{\includegraphics[height=2.7cm]{ufop.png}} \\
 & \multicolumn{2}{c}{\bfseries PRÓ-REITORIA DE GRADUAÇÃO} & \\
 & \multicolumn{2}{c}{} & \\
 & \multicolumn{2}{c}{\Large\bfseries PROGRAMA DE DISCIPLINA} & \\
\end{tabular}
\end{center}

\begin{center}
\begin{longtable}{|p{4cm}|p{4cm}|p{4cm}|p{4cm}|}
\hline
\multicolumn{3}{|p{12cm}|}{Nome do Componente Curricular em Português:} &
\multicolumn{1}{p{4cm}|}{Código:} \\ 
\multicolumn{3}{|p{12cm}|}{\textbf{Computação Evolutiva}} &
\textbf{BCC461}\\ 
\multicolumn{3}{|p{12cm}|}{Nome do Componente Curricular em Inglês:} & \\ 
\multicolumn{3}{|p{12cm}|}{\textbf{Evolutionary Computing}} & \\ 
\hline
\multicolumn{3}{|p{12cm}|}{Nome e Sigla do Departamento} & Unidade Acadêmica: \\ 
\multicolumn{3}{|p{12cm}|}{Departamento de Computação (DECOM)} & {ICEB} \\ 
\hline
\multicolumn{4}{|p{16cm}|}{Modalidade de Oferta:
[X] presencial \hspace{0.6cm}
[ ] semipresencial \hspace{0.6cm}
[ ] à distância}\\
\hline
\multicolumn{2}{|p{8cm}|}{Carga horária semestral} &
\multicolumn{2}{p{8cm}|}{Carga horária semanal}\\
\hline
\multicolumn{1}{|p{4cm}|}{Total} &
\multicolumn{1}{p{4cm}|}{Extensionista} &
\multicolumn{1}{p{4cm}|}{Teórica} &
\multicolumn{1}{p{4cm}|}{Prática} \\ 
\multicolumn{1}{|p{4cm}|}{60\,horas} &
\multicolumn{1}{p{4cm}|}{0\;horas} &
\multicolumn{1}{p{4cm}|}{4\;horas/aula} &
\multicolumn{1}{p{4cm}|}{0\;horas/aula} \\ 
\hline
\multicolumn{4}{|p{16cm}|}{Ementa:}\\
\multicolumn{4}{|p{16cm}|}{}\\
\multicolumn{4}{|p{\dimexpr 16cm + 6\tabcolsep\relax}|}{Introdução à computação evolutiva; evolução por seleção natural; conceitos básicos e componentes de algoritmos evolutivos; algoritmos genéticos; estratégias evolutivas; programação evolucionária; programação genética; sistemas classificadores; algoritmos híbridos; tratamento de restrições; otimização multiobjetivo com algoritmos evolutivos.}\\
\multicolumn{4}{|p{16cm}|}{}\\
\hline
\multicolumn{4}{|p{16cm}|}{Conteúdo programático:}\\*
\multicolumn{4}{|p{\dimexpr 16cm + 6\tabcolsep\relax}|}{\begin{enumerate}[start=1]\item Introdução à computação evolutiva e relação com modelos biológicos\end{enumerate}}\\
\multicolumn{4}{|p{\dimexpr 16cm + 6\tabcolsep\relax}|}{\begin{enumerate}[start=2]\item Componentes básicos comuns a algoritmos evolutivos\end{enumerate}}\\
\multicolumn{4}{|p{\dimexpr 16cm + 6\tabcolsep\relax}|}{\begin{enumerate}[start=3]\item Categorias de algoritmos evolutivos\end{enumerate}}\\
\multicolumn{4}{|p{\dimexpr 16cm + 6\tabcolsep\relax}|}{\begin{enumerate}[start=4]\item Sistemas evolutivos aplicados a sistemas de classificação\end{enumerate}}\\
\multicolumn{4}{|p{\dimexpr 16cm + 6\tabcolsep\relax}|}{\begin{enumerate}[start=5]\item Algoritmos híbridos\end{enumerate}}\\
\multicolumn{4}{|p{\dimexpr 16cm + 6\tabcolsep\relax}|}{\begin{enumerate}[start=6]\item Introdução à otimização multiobjetivo\end{enumerate}}\\
\multicolumn{4}{|p{\dimexpr 16cm + 6\tabcolsep\relax}|}{}\\
\hline
\multicolumn{4}{|p{16cm}|}{Bibliografia Básica:}\\*
\multicolumn{4}{|p{\dimexpr 16cm + 6\tabcolsep\relax}|}{\begin{itemize}\item BÄCK, Thomas. Handbook of evolutionary computation. CRC Press, 1997.\end{itemize}}\\
\multicolumn{4}{|p{\dimexpr 16cm + 6\tabcolsep\relax}|}{\begin{itemize}\item David B. Fogel. Evolutionary computation: toward a new philosophy of machine intelligence. 3rd ed. New Jersey: Wiley Interscience, 2006.\end{itemize}}\\
\multicolumn{4}{|p{16cm}|}{}\\
\hline
\multicolumn{4}{|p{16cm}|}{Bibliografia Complementar:}\\*
\multicolumn{4}{|p{\dimexpr 16cm + 6\tabcolsep\relax}|}{\begin{itemize}\item BACK, Thomas. Evolutionary algorithms in theory and practice. Oxford Univ. Press, 1996.\end{itemize}}\\
\multicolumn{4}{|p{\dimexpr 16cm + 6\tabcolsep\relax}|}{\begin{itemize}\item DEB, Kalyanmoy. Multi-objective optimization using evolutionary algorithms. John Wiley \& Sons, 2001.\end{itemize}}\\
\multicolumn{4}{|p{\dimexpr 16cm + 6\tabcolsep\relax}|}{\begin{itemize}\item YAO, Xin. Evolutionary computation: Theory and applications. World Scientific, 1999.\end{itemize}}\\
\multicolumn{4}{|p{\dimexpr 16cm + 6\tabcolsep\relax}|}{\begin{itemize}\item JOHNSON, R. P. Automatic generation of neural network architecture using evolutionary computation. World Scientific, 1997.\end{itemize}}\\
\multicolumn{4}{|p{\dimexpr 16cm + 6\tabcolsep\relax}|}{\begin{itemize}\item Kenneth A. De Jong. Evolutionary computation: a unified approach. Cambridge, Mass.: MIT Press, 2006.\end{itemize}}\\
\hline
\end{longtable}
\end{center}

\clearpage
\begin{center}
\begin{tabular}{lccr}
 \multirow{3}{*}{\includegraphics[height=2.7cm]{brasao.png}} &
 \multicolumn{2}{c}{\bfseries UNIVERSIDADE FEDERAL DE OURO PRETO} &
 \ \ \ \ \multirow{3}{*}{\includegraphics[height=2.7cm]{ufop.png}} \\
 & \multicolumn{2}{c}{\bfseries PRÓ-REITORIA DE GRADUAÇÃO} & \\
 & \multicolumn{2}{c}{} & \\
 & \multicolumn{2}{c}{\Large\bfseries PROGRAMA DE DISCIPLINA} & \\
\end{tabular}
\end{center}

\begin{center}
\begin{longtable}{|p{4cm}|p{4cm}|p{4cm}|p{4cm}|}
\hline
\multicolumn{3}{|p{12cm}|}{Nome do Componente Curricular em Português:} &
\multicolumn{1}{p{4cm}|}{Código:} \\ 
\multicolumn{3}{|p{12cm}|}{\textbf{Otimização em Redes}} &
\textbf{BCC463}\\ 
\multicolumn{3}{|p{12cm}|}{Nome do Componente Curricular em Inglês:} & \\ 
\multicolumn{3}{|p{12cm}|}{\textbf{Network Optimization}} & \\ 
\hline
\multicolumn{3}{|p{12cm}|}{Nome e Sigla do Departamento} & Unidade Acadêmica: \\ 
\multicolumn{3}{|p{12cm}|}{Departamento de Computação (DECOM)} & {ICEB} \\ 
\hline
\multicolumn{4}{|p{16cm}|}{Modalidade de Oferta:
[X] presencial \hspace{0.6cm}
[ ] semipresencial \hspace{0.6cm}
[ ] à distância}\\
\hline
\multicolumn{2}{|p{8cm}|}{Carga horária semestral} &
\multicolumn{2}{p{8cm}|}{Carga horária semanal}\\
\hline
\multicolumn{1}{|p{4cm}|}{Total} &
\multicolumn{1}{p{4cm}|}{Extensionista} &
\multicolumn{1}{p{4cm}|}{Teórica} &
\multicolumn{1}{p{4cm}|}{Prática} \\ 
\multicolumn{1}{|p{4cm}|}{60\,horas} &
\multicolumn{1}{p{4cm}|}{0\;horas} &
\multicolumn{1}{p{4cm}|}{4\;horas/aula} &
\multicolumn{1}{p{4cm}|}{0\;horas/aula} \\ 
\hline
\multicolumn{4}{|p{16cm}|}{Ementa:}\\
\multicolumn{4}{|p{16cm}|}{}\\
\multicolumn{4}{|p{\dimexpr 16cm + 6\tabcolsep\relax}|}{Conceitos básicos sobre grafos; modelos de fluxos em redes; algoritmos do caminho mínimo, do fluxo máximo e do fluxo com custo mínimo; aplicações e implementação de algoritmos especializados.}\\
\multicolumn{4}{|p{16cm}|}{}\\
\hline
\multicolumn{4}{|p{16cm}|}{Conteúdo programático:}\\*
\multicolumn{4}{|p{\dimexpr 16cm + 6\tabcolsep\relax}|}{\begin{enumerate}[start=1]\item Conceitos básicos sobre grafos: Notação e definições; Representação em redes; Exemplos de modelos em redes\end{enumerate}}\\
\multicolumn{4}{|p{\dimexpr 16cm + 6\tabcolsep\relax}|}{\begin{enumerate}[start=2]\item Problemas de caminho mínimo: Introdução; Aplicações; Algoritmos de solução e implementações\end{enumerate}}\\
\multicolumn{4}{|p{\dimexpr 16cm + 6\tabcolsep\relax}|}{\begin{enumerate}[start=3]\item Problema de fluxo máximo: Introdução; Aplicações; Fluxos e cortes; Teorema do fluxo máximo – corte mínimo; Algoritmos de solução e implementações\end{enumerate}}\\
\multicolumn{4}{|p{\dimexpr 16cm + 6\tabcolsep\relax}|}{\begin{enumerate}[start=4]\item Problemas de fluxo com custo mínimo: Introdução; Aplicações; Algoritmos de solução e implementações\end{enumerate}}\\
\multicolumn{4}{|p{\dimexpr 16cm + 6\tabcolsep\relax}|}{\begin{enumerate}[start=5]\item Aplicações: Problemas de designação; Problemas de emparelhamento; Problemas de roteamento; Problemas de sequenciamento\end{enumerate}}\\
\multicolumn{4}{|p{\dimexpr 16cm + 6\tabcolsep\relax}|}{}\\
\hline
\multicolumn{4}{|p{16cm}|}{Bibliografia Básica:}\\*
\multicolumn{4}{|p{\dimexpr 16cm + 6\tabcolsep\relax}|}{\begin{itemize}\item Marco Goldbarg; Elizabeth Goldbarg. Grafos: conceitos, algoritmos e aplicações. Rio de Janeiro: Elsevier, 2012.\end{itemize}}\\
\multicolumn{4}{|p{\dimexpr 16cm + 6\tabcolsep\relax}|}{\begin{itemize}\item Chvatal, Vasek. Linear programming. New York: W. H. Freeman, 1983.\end{itemize}}\\
\multicolumn{4}{|p{\dimexpr 16cm + 6\tabcolsep\relax}|}{\begin{itemize}\item Frazer, J. Ronald. Applied linear programming. Englewood Cliffs, N.J.: Prentice-Hall, 1968.\end{itemize}}\\
\multicolumn{4}{|p{\dimexpr 16cm + 6\tabcolsep\relax}|}{\begin{itemize}\item Ravindra K. Ahuja, Thomas L. Magnanti, James B. Orlin. Network flows: theory, algorithms and applications. Englewood Cliffs: Prentice-Hall, 1993.\end{itemize}}\\
\multicolumn{4}{|p{\dimexpr 16cm + 6\tabcolsep\relax}|}{\begin{itemize}\item Bazaraa, M. S. Linear programming and network flows. 2. ed. New York: J. Wiley, 1990.\end{itemize}}\\
\multicolumn{4}{|p{16cm}|}{}\\
\hline
\multicolumn{4}{|p{16cm}|}{Bibliografia Complementar:}\\*
\multicolumn{4}{|p{\dimexpr 16cm + 6\tabcolsep\relax}|}{\begin{itemize}\item Hu, T. C. Integer programming and network flows. Reading: Addison Wesley, 1970.\end{itemize}}\\
\multicolumn{4}{|p{\dimexpr 16cm + 6\tabcolsep\relax}|}{\begin{itemize}\item Marcos Arenales, Vinicius Armentano, Reinaldo Morabito, Horacio Yanasse. Pesquisa Operacional: para cursos de engenharia. Rio de Janeiro: Elsevier, 2007.\end{itemize}}\\
\multicolumn{4}{|p{\dimexpr 16cm + 6\tabcolsep\relax}|}{\begin{itemize}\item Wolsey, Laurence A. Integer programming. New York: John Wiley \& Sons, 1998.\end{itemize}}\\
\multicolumn{4}{|p{\dimexpr 16cm + 6\tabcolsep\relax}|}{\begin{itemize}\item Marco Cesar Goldbarg, Henrique Pacca L. Luna. Otimização combinatória e programação linear: modelos e algoritmos. 2. ed. rev. e atual. Rio de Janeiro: Campus, 2005.\end{itemize}}\\
\hline
\end{longtable}
\end{center}

\clearpage
\begin{center}
\begin{tabular}{lccr}
 \multirow{3}{*}{\includegraphics[height=2.7cm]{brasao.png}} &
 \multicolumn{2}{c}{\bfseries UNIVERSIDADE FEDERAL DE OURO PRETO} &
 \ \ \ \ \multirow{3}{*}{\includegraphics[height=2.7cm]{ufop.png}} \\
 & \multicolumn{2}{c}{\bfseries PRÓ-REITORIA DE GRADUAÇÃO} & \\
 & \multicolumn{2}{c}{} & \\
 & \multicolumn{2}{c}{\Large\bfseries PROGRAMA DE DISCIPLINA} & \\
\end{tabular}
\end{center}

\begin{center}
\begin{longtable}{|p{4cm}|p{4cm}|p{4cm}|p{4cm}|}
\hline
\multicolumn{3}{|p{12cm}|}{Nome do Componente Curricular em Português:} &
\multicolumn{1}{p{4cm}|}{Código:} \\ 
\multicolumn{3}{|p{12cm}|}{\textbf{Otimização Linear e Inteira}} &
\textbf{BCC464}\\ 
\multicolumn{3}{|p{12cm}|}{Nome do Componente Curricular em Inglês:} & \\ 
\multicolumn{3}{|p{12cm}|}{\textbf{Linear and Integer Optimization}} & \\ 
\hline
\multicolumn{3}{|p{12cm}|}{Nome e Sigla do Departamento} & Unidade Acadêmica: \\ 
\multicolumn{3}{|p{12cm}|}{Departamento de Computação (DECOM)} & {ICEB} \\ 
\hline
\multicolumn{4}{|p{16cm}|}{Modalidade de Oferta:
[X] presencial \hspace{0.6cm}
[ ] semipresencial \hspace{0.6cm}
[ ] à distância}\\
\hline
\multicolumn{2}{|p{8cm}|}{Carga horária semestral} &
\multicolumn{2}{p{8cm}|}{Carga horária semanal}\\
\hline
\multicolumn{1}{|p{4cm}|}{Total} &
\multicolumn{1}{p{4cm}|}{Extensionista} &
\multicolumn{1}{p{4cm}|}{Teórica} &
\multicolumn{1}{p{4cm}|}{Prática} \\ 
\multicolumn{1}{|p{4cm}|}{60\,horas} &
\multicolumn{1}{p{4cm}|}{0\;horas} &
\multicolumn{1}{p{4cm}|}{4\;horas/aula} &
\multicolumn{1}{p{4cm}|}{0\;horas/aula} \\ 
\hline
\multicolumn{4}{|p{16cm}|}{Ementa:}\\
\multicolumn{4}{|p{16cm}|}{}\\
\multicolumn{4}{|p{\dimexpr 16cm + 6\tabcolsep\relax}|}{Modelagem em programação linear. Algoritmo simplex. Dualidade. Análise de sensibilidade. Geração de colunas. Métodos de decomposição de Dantzig-Wolfe. Modelagem em programação inteira. Enumeração implícita. Planos de corte. Limites e relaxações.}\\
\multicolumn{4}{|p{16cm}|}{}\\
\hline
\multicolumn{4}{|p{16cm}|}{Conteúdo programático:}\\*
\multicolumn{4}{|p{\dimexpr 16cm + 6\tabcolsep\relax}|}{\begin{enumerate}[start=1]\item Introdução à Programação Linear (Histórico)\end{enumerate}}\\
\multicolumn{4}{|p{\dimexpr 16cm + 6\tabcolsep\relax}|}{\begin{enumerate}[start=2]\item Modelagem de problemas simples\end{enumerate}}\\
\multicolumn{4}{|p{\dimexpr 16cm + 6\tabcolsep\relax}|}{\begin{enumerate}[start=3]\item Método gráfico\end{enumerate}}\\
\multicolumn{4}{|p{\dimexpr 16cm + 6\tabcolsep\relax}|}{\begin{enumerate}[start=4]\item Modelagem de problemas e aplicações de exemplo\end{enumerate}}\\
\multicolumn{4}{|p{\dimexpr 16cm + 6\tabcolsep\relax}|}{\begin{enumerate}[start=5]\item Linguagem de modelagem MathProg\end{enumerate}}\\
\multicolumn{4}{|p{\dimexpr 16cm + 6\tabcolsep\relax}|}{\begin{enumerate}[start=6]\item Algoritmo Simplex\end{enumerate}}\\
\multicolumn{4}{|p{\dimexpr 16cm + 6\tabcolsep\relax}|}{\begin{enumerate}[start=7]\item Dualidade\end{enumerate}}\\
\multicolumn{4}{|p{\dimexpr 16cm + 6\tabcolsep\relax}|}{\begin{enumerate}[start=8]\item Análise de sensibilidade\end{enumerate}}\\
\multicolumn{4}{|p{\dimexpr 16cm + 6\tabcolsep\relax}|}{\begin{enumerate}[start=9]\item Geração de colunas\end{enumerate}}\\
\multicolumn{4}{|p{\dimexpr 16cm + 6\tabcolsep\relax}|}{\begin{enumerate}[start=10]\item Decomposição Dantzig-Wolfe\end{enumerate}}\\
\multicolumn{4}{|p{\dimexpr 16cm + 6\tabcolsep\relax}|}{\begin{enumerate}[start=11]\item Modelagem em Programação Inteira\end{enumerate}}\\
\multicolumn{4}{|p{\dimexpr 16cm + 6\tabcolsep\relax}|}{\begin{enumerate}[start=12]\item Avaliando a Qualidade das Formulações\end{enumerate}}\\
\multicolumn{4}{|p{\dimexpr 16cm + 6\tabcolsep\relax}|}{\begin{enumerate}[start=13]\item Geração de Desigualdades válidas e reformulação automática\end{enumerate}}\\
\multicolumn{4}{|p{\dimexpr 16cm + 6\tabcolsep\relax}|}{\begin{enumerate}[start=14]\item Técnicas para geração de limites duais e primais fortes\end{enumerate}}\\
\multicolumn{4}{|p{\dimexpr 16cm + 6\tabcolsep\relax}|}{\begin{enumerate}[start=15]\item Geração de colunas em Programação Inteira\end{enumerate}}\\
\multicolumn{4}{|p{\dimexpr 16cm + 6\tabcolsep\relax}|}{}\\
\hline
\multicolumn{4}{|p{16cm}|}{Bibliografia Básica:}\\*
\multicolumn{4}{|p{\dimexpr 16cm + 6\tabcolsep\relax}|}{\begin{itemize}\item Maculan, Nelson. Otimização linear. Brasília: Ed. UnB, 2006.\end{itemize}}\\
\multicolumn{4}{|p{\dimexpr 16cm + 6\tabcolsep\relax}|}{\begin{itemize}\item Chvatal, Vasek. Linear programming. New York: W. H. Freeman, 1983.\end{itemize}}\\
\multicolumn{4}{|p{\dimexpr 16cm + 6\tabcolsep\relax}|}{\begin{itemize}\item Wolsey, Laurence A. Integer programming. New York: John Wiley \& Sons, 1998.\end{itemize}}\\
\multicolumn{4}{|p{16cm}|}{}\\
\hline
\multicolumn{4}{|p{16cm}|}{Bibliografia Complementar:}\\*
\multicolumn{4}{|p{\dimexpr 16cm + 6\tabcolsep\relax}|}{\begin{itemize}\item Edited by John K. Karlof. Integer programming: theory and practice. Boca Raton, Fla.\end{itemize}}\\
\multicolumn{4}{|p{\dimexpr 16cm + 6\tabcolsep\relax}|}{\begin{itemize}\item Dasgupta, Sanjoy. Algoritmos. São Paulo: McGraw-Hill, 2009.\end{itemize}}\\
\multicolumn{4}{|p{\dimexpr 16cm + 6\tabcolsep\relax}|}{\begin{itemize}\item Michael Jünger... [et al.], editors. 50 years of integer programming, 1958-2008: the early years and state-of-the-art surveys. Heidelberg: Springer, 2010.\end{itemize}}\\
\multicolumn{4}{|p{\dimexpr 16cm + 6\tabcolsep\relax}|}{\begin{itemize}\item Lee, Jon. A first course in combinatorial optimization. Cambridge, UK.\end{itemize}}\\
\multicolumn{4}{|p{\dimexpr 16cm + 6\tabcolsep\relax}|}{\begin{itemize}\item Marco Cesar Goldbarg, Henrique Pacca L. Luna. Otimização combinatória e programação linear: modelos e algoritmos. Rio de Janeiro: Campus, 2000.\end{itemize}}\\
\multicolumn{4}{|p{\dimexpr 16cm + 6\tabcolsep\relax}|}{\begin{itemize}\item Dantzig, George B. Linear programming and extensions. Princeton, N.J.: Princeton University Press, 1963.\end{itemize}}\\
\hline
\end{longtable}
\end{center}

\clearpage
\begin{center}
\begin{tabular}{lccr}
 \multirow{3}{*}{\includegraphics[height=2.7cm]{brasao.png}} &
 \multicolumn{2}{c}{\bfseries UNIVERSIDADE FEDERAL DE OURO PRETO} &
 \ \ \ \ \multirow{3}{*}{\includegraphics[height=2.7cm]{ufop.png}} \\
 & \multicolumn{2}{c}{\bfseries PRÓ-REITORIA DE GRADUAÇÃO} & \\
 & \multicolumn{2}{c}{} & \\
 & \multicolumn{2}{c}{\Large\bfseries PROGRAMA DE DISCIPLINA} & \\
\end{tabular}
\end{center}

\begin{center}
\begin{longtable}{|p{4cm}|p{4cm}|p{4cm}|p{4cm}|}
\hline
\multicolumn{3}{|p{12cm}|}{Nome do Componente Curricular em Português:} &
\multicolumn{1}{p{4cm}|}{Código:} \\ 
\multicolumn{3}{|p{12cm}|}{\textbf{Técnicas de Otimização Multiobjetivo}} &
\textbf{BCC465}\\ 
\multicolumn{3}{|p{12cm}|}{Nome do Componente Curricular em Inglês:} & \\ 
\multicolumn{3}{|p{12cm}|}{\textbf{Multi-Objective Optimization Techniques}} & \\ 
\hline
\multicolumn{3}{|p{12cm}|}{Nome e Sigla do Departamento} & Unidade Acadêmica: \\ 
\multicolumn{3}{|p{12cm}|}{Departamento de Computação (DECOM)} & {ICEB} \\ 
\hline
\multicolumn{4}{|p{16cm}|}{Modalidade de Oferta:
[X] presencial \hspace{0.6cm}
[ ] semipresencial \hspace{0.6cm}
[ ] à distância}\\
\hline
\multicolumn{2}{|p{8cm}|}{Carga horária semestral} &
\multicolumn{2}{p{8cm}|}{Carga horária semanal}\\
\hline
\multicolumn{1}{|p{4cm}|}{Total} &
\multicolumn{1}{p{4cm}|}{Extensionista} &
\multicolumn{1}{p{4cm}|}{Teórica} &
\multicolumn{1}{p{4cm}|}{Prática} \\ 
\multicolumn{1}{|p{4cm}|}{60\,horas} &
\multicolumn{1}{p{4cm}|}{0\;horas} &
\multicolumn{1}{p{4cm}|}{4\;horas/aula} &
\multicolumn{1}{p{4cm}|}{0\;horas/aula} \\ 
\hline
\multicolumn{4}{|p{16cm}|}{Ementa:}\\
\multicolumn{4}{|p{16cm}|}{}\\
\multicolumn{4}{|p{\dimexpr 16cm + 6\tabcolsep\relax}|}{Fundamentos matemáticos; conceitos básicos de otimização não-linear; condições de otimalidade; métodos de minimização irrestrita e restrita; funcionais-objetivo e vetores de objetivos; conjuntos ordenados; soluções de Pareto: caracterização analítica; geração de soluções de Pareto; indicação de preferências; algoritmos para otimização multi-objetivo; otimização multi-objetivo com algoritmos evolutivos; decisão.}\\
\multicolumn{4}{|p{16cm}|}{}\\
\hline
\multicolumn{4}{|p{16cm}|}{Conteúdo programático:}\\*
\multicolumn{4}{|p{\dimexpr 16cm + 6\tabcolsep\relax}|}{\begin{enumerate}[start=1]\item Apresentação do plano de curso e introdução: O Problema de otimização vetorial; Notação\end{enumerate}}\\
\multicolumn{4}{|p{\dimexpr 16cm + 6\tabcolsep\relax}|}{\begin{enumerate}[start=2]\item Ordenamento de soluções: Conjunto Pareto-Ótimo; Conjunto localmente Pareto-ótimo; Solução utópica\end{enumerate}}\\
\multicolumn{4}{|p{\dimexpr 16cm + 6\tabcolsep\relax}|}{\begin{enumerate}[start=3]\item O Problema de determinação das soluções eficientes: Condições de Kuhn-Tucker para eficiência\end{enumerate}}\\
\multicolumn{4}{|p{\dimexpr 16cm + 6\tabcolsep\relax}|}{\begin{enumerate}[start=4]\item Geração de soluções eficientes: Abordagem via problema ponderado; Abordagem via problema epsilon-Restrito; Abordagem da programação-alvo; Teste de eficiência\end{enumerate}}\\
\multicolumn{4}{|p{\dimexpr 16cm + 6\tabcolsep\relax}|}{\begin{enumerate}[start=5]\item Estrutura do conjunto Pareto-Ótimo\end{enumerate}}\\
\multicolumn{4}{|p{\dimexpr 16cm + 6\tabcolsep\relax}|}{\begin{enumerate}[start=6]\item Análise Multi-objetivo: Consistência; Ordenamento e dominância; Extensão\end{enumerate}}\\
\multicolumn{4}{|p{\dimexpr 16cm + 6\tabcolsep\relax}|}{\begin{enumerate}[start=7]\item Decisão e Síntese Multi-objetivo\end{enumerate}}\\
\multicolumn{4}{|p{\dimexpr 16cm + 6\tabcolsep\relax}|}{\begin{enumerate}[start=8]\item Algoritmos Genéticos Multi-objetivo: Construção dos Algoritmos Genéticos Multi-objetivo\end{enumerate}}\\
\multicolumn{4}{|p{\dimexpr 16cm + 6\tabcolsep\relax}|}{}\\
\hline
\multicolumn{4}{|p{16cm}|}{Bibliografia Básica:}\\*
\multicolumn{4}{|p{\dimexpr 16cm + 6\tabcolsep\relax}|}{\begin{itemize}\item DEB, Kalyanmoy. Multi-Objective Optimization using Evolutionary Algorithms. John Wiley \& Sons, 2001.\end{itemize}}\\
\multicolumn{4}{|p{\dimexpr 16cm + 6\tabcolsep\relax}|}{\begin{itemize}\item COELLO COELLO, Carlos. Evolutionary Algorithms for Solving Multi-Objective Problems. 2. ed. Springer, 2007.\end{itemize}}\\
\multicolumn{4}{|p{\dimexpr 16cm + 6\tabcolsep\relax}|}{\begin{itemize}\item Vira Chankong and Yacov Y. Haimes. Multiobjective decision making: theory and methodology. New York: North-Holland, 1983.\end{itemize}}\\
\multicolumn{4}{|p{\dimexpr 16cm + 6\tabcolsep\relax}|}{\begin{itemize}\item Yann Collette, Patrick Siarry. Multiobjective optimization: principles and case studies. Berlin.\end{itemize}}\\
\multicolumn{4}{|p{\dimexpr 16cm + 6\tabcolsep\relax}|}{\begin{itemize}\item Mokhtar S. Bazaraa, Hanif D. Sherali, C. M. Shetty. Nonlinear programming: theory and algorithms. 3rd ed. Hoboken, N.J.: Wiley-Interscience, 2006.\end{itemize}}\\
\multicolumn{4}{|p{16cm}|}{}\\
\hline
\multicolumn{4}{|p{16cm}|}{Bibliografia Complementar:}\\*
\multicolumn{4}{|p{\dimexpr 16cm + 6\tabcolsep\relax}|}{\begin{itemize}\item SAWARAGI, Yoshikazu. Theory of Multiobjective Optimization. Elsevier, 1985.\end{itemize}}\\
\multicolumn{4}{|p{\dimexpr 16cm + 6\tabcolsep\relax}|}{\begin{itemize}\item ABRAHAM, Ajith. Evolutionary Multiobjective Optimization: Theoretical Advances and Applications. Springer, 2005.\end{itemize}}\\
\multicolumn{4}{|p{\dimexpr 16cm + 6\tabcolsep\relax}|}{\begin{itemize}\item EICHFELDER, Gabriele. Adaptive Scalarization Methods in Multiobjective Optimization. Springer, 2018.\end{itemize}}\\
\multicolumn{4}{|p{\dimexpr 16cm + 6\tabcolsep\relax}|}{\begin{itemize}\item MIETTINEN, Kaisa. Nonlinear Multiobjective Optimization. Springer, 1999.\end{itemize}}\\
\multicolumn{4}{|p{\dimexpr 16cm + 6\tabcolsep\relax}|}{\begin{itemize}\item TANG, Kit. Multiobjective Optimization Methodology: A Jumping Gene Approach. 1. ed. CRC Press, 2018.\end{itemize}}\\
\hline
\end{longtable}
\end{center}

\clearpage
\begin{center}
\begin{tabular}{lccr}
 \multirow{3}{*}{\includegraphics[height=2.7cm]{brasao.png}} &
 \multicolumn{2}{c}{\bfseries UNIVERSIDADE FEDERAL DE OURO PRETO} &
 \ \ \ \ \multirow{3}{*}{\includegraphics[height=2.7cm]{ufop.png}} \\
 & \multicolumn{2}{c}{\bfseries PRÓ-REITORIA DE GRADUAÇÃO} & \\
 & \multicolumn{2}{c}{} & \\
 & \multicolumn{2}{c}{\Large\bfseries PROGRAMA DE DISCIPLINA} & \\
\end{tabular}
\end{center}

\begin{center}
\begin{longtable}{|p{4cm}|p{4cm}|p{4cm}|p{4cm}|}
\hline
\multicolumn{3}{|p{12cm}|}{Nome do Componente Curricular em Português:} &
\multicolumn{1}{p{4cm}|}{Código:} \\ 
\multicolumn{3}{|p{12cm}|}{\textbf{Técnicas Metaheurísticas para Otimização Combinatória}} &
\textbf{BCC466}\\ 
\multicolumn{3}{|p{12cm}|}{Nome do Componente Curricular em Inglês:} & \\ 
\multicolumn{3}{|p{12cm}|}{\textbf{Metaheuristics Techniques in Combinatorial Optimization}} & \\ 
\hline
\multicolumn{3}{|p{12cm}|}{Nome e Sigla do Departamento} & Unidade Acadêmica: \\ 
\multicolumn{3}{|p{12cm}|}{Departamento de Computação (DECOM)} & {ICEB} \\ 
\hline
\multicolumn{4}{|p{16cm}|}{Modalidade de Oferta:
[X] presencial \hspace{0.6cm}
[ ] semipresencial \hspace{0.6cm}
[ ] à distância}\\
\hline
\multicolumn{2}{|p{8cm}|}{Carga horária semestral} &
\multicolumn{2}{p{8cm}|}{Carga horária semanal}\\
\hline
\multicolumn{1}{|p{4cm}|}{Total} &
\multicolumn{1}{p{4cm}|}{Extensionista} &
\multicolumn{1}{p{4cm}|}{Teórica} &
\multicolumn{1}{p{4cm}|}{Prática} \\ 
\multicolumn{1}{|p{4cm}|}{60\,horas} &
\multicolumn{1}{p{4cm}|}{0\;horas} &
\multicolumn{1}{p{4cm}|}{4\;horas/aula} &
\multicolumn{1}{p{4cm}|}{0\;horas/aula} \\ 
\hline
\multicolumn{4}{|p{16cm}|}{Ementa:}\\
\multicolumn{4}{|p{16cm}|}{}\\
\multicolumn{4}{|p{\dimexpr 16cm + 6\tabcolsep\relax}|}{Problemas de otimização combinatória. Heurísticas clássicas. Busca local. Métodos de múltiplos reinícios. Metaheurísticas baseadas em trajetória. Metaheurísticas populacionais. Métodos inspirados na natureza. Aplicações.}\\
\multicolumn{4}{|p{16cm}|}{}\\
\hline
\multicolumn{4}{|p{16cm}|}{Conteúdo programático:}\\*
\multicolumn{4}{|p{\dimexpr 16cm + 6\tabcolsep\relax}|}{\begin{enumerate}[start=1]\item Introdução: problemas de otimização combinatória\end{enumerate}}\\
\multicolumn{4}{|p{\dimexpr 16cm + 6\tabcolsep\relax}|}{\begin{enumerate}[start=2]\item Algoritmos gulosos\end{enumerate}}\\
\multicolumn{4}{|p{\dimexpr 16cm + 6\tabcolsep\relax}|}{\begin{enumerate}[start=3]\item Busca local\end{enumerate}}\\
\multicolumn{4}{|p{\dimexpr 16cm + 6\tabcolsep\relax}|}{\begin{enumerate}[start=4]\item Busca em grandes vizinhanças\end{enumerate}}\\
\multicolumn{4}{|p{\dimexpr 16cm + 6\tabcolsep\relax}|}{\begin{enumerate}[start=5]\item Reconexão de caminhos\end{enumerate}}\\
\multicolumn{4}{|p{\dimexpr 16cm + 6\tabcolsep\relax}|}{\begin{enumerate}[start=6]\item Simulated Annealing\end{enumerate}}\\
\multicolumn{4}{|p{\dimexpr 16cm + 6\tabcolsep\relax}|}{\begin{enumerate}[start=7]\item Busca Tabu\end{enumerate}}\\
\multicolumn{4}{|p{\dimexpr 16cm + 6\tabcolsep\relax}|}{\begin{enumerate}[start=8]\item Busca local iterada\end{enumerate}}\\
\multicolumn{4}{|p{\dimexpr 16cm + 6\tabcolsep\relax}|}{\begin{enumerate}[start=9]\item Busca em vizinhança variável\end{enumerate}}\\
\multicolumn{4}{|p{\dimexpr 16cm + 6\tabcolsep\relax}|}{\begin{enumerate}[start=10]\item Procedimentos com múltiplos reinícios\end{enumerate}}\\
\multicolumn{4}{|p{\dimexpr 16cm + 6\tabcolsep\relax}|}{\begin{enumerate}[start=11]\item Algoritmos inspirados na natureza\end{enumerate}}\\
\multicolumn{4}{|p{\dimexpr 16cm + 6\tabcolsep\relax}|}{\begin{enumerate}[start=12]\item Algoritmos evolutivos\end{enumerate}}\\
\multicolumn{4}{|p{\dimexpr 16cm + 6\tabcolsep\relax}|}{\begin{enumerate}[start=13]\item Scatter Search\end{enumerate}}\\
\multicolumn{4}{|p{\dimexpr 16cm + 6\tabcolsep\relax}|}{\begin{enumerate}[start=14]\item Aplicações em roteamento\end{enumerate}}\\
\multicolumn{4}{|p{\dimexpr 16cm + 6\tabcolsep\relax}|}{\begin{enumerate}[start=15]\item Aplicações em sequenciamento\end{enumerate}}\\
\multicolumn{4}{|p{\dimexpr 16cm + 6\tabcolsep\relax}|}{\begin{enumerate}[start=16]\item Aplicações em programação de horários\end{enumerate}}\\
\multicolumn{4}{|p{\dimexpr 16cm + 6\tabcolsep\relax}|}{\begin{enumerate}[start=17]\item Aplicações em projetos de redes\end{enumerate}}\\
\multicolumn{4}{|p{\dimexpr 16cm + 6\tabcolsep\relax}|}{}\\
\hline
\multicolumn{4}{|p{16cm}|}{Bibliografia Básica:}\\*
\multicolumn{4}{|p{\dimexpr 16cm + 6\tabcolsep\relax}|}{\begin{itemize}\item Handbook of approximation algorithms and metaheuristics. New York: Chapman \& Hall/CRC, 2007.\end{itemize}}\\
\multicolumn{4}{|p{\dimexpr 16cm + 6\tabcolsep\relax}|}{\begin{itemize}\item Voss, Stefan. Introduction to computational optimization models for production planning in a supply chain. 2nd ed. Berlin.\end{itemize}}\\
\multicolumn{4}{|p{\dimexpr 16cm + 6\tabcolsep\relax}|}{\begin{itemize}\item El-Ghazali Talbi. Metaheuristics: from design to implementation. Hoboken, N.J.: John Wiley \& Sons, 2009.\end{itemize}}\\
\multicolumn{4}{|p{\dimexpr 16cm + 6\tabcolsep\relax}|}{\begin{itemize}\item Zbigniew Michalewicz, David B. Fogel. How to solve it: modern heuristics. 2nd rev. and extended ed. Berlin.\end{itemize}}\\
\multicolumn{4}{|p{\dimexpr 16cm + 6\tabcolsep\relax}|}{\begin{itemize}\item Mauricio G.C. Resende, Jorge Pinho de Sousa, with assistance from Ana Viana (eds.). Metaheuristics: computer decision-making. Boston: Kluwer Academic Publishers, 2004.\end{itemize}}\\
\multicolumn{4}{|p{16cm}|}{}\\
\hline
\multicolumn{4}{|p{16cm}|}{Bibliografia Complementar:}\\*
\multicolumn{4}{|p{\dimexpr 16cm + 6\tabcolsep\relax}|}{\begin{itemize}\item Local search in combinatorial optimization. New Jersey: Princeton University Press, 2003.\end{itemize}}\\
\multicolumn{4}{|p{\dimexpr 16cm + 6\tabcolsep\relax}|}{\begin{itemize}\item J. Dréo... [et al.]. Metaheuristics for hard optimization: methods and case studies. Berlin: Springer, 2006.\end{itemize}}\\
\multicolumn{4}{|p{\dimexpr 16cm + 6\tabcolsep\relax}|}{\begin{itemize}\item Kenneth V. Price, Rainer M. Storn, Jouni A. Lampinen. Differential evolution: a practical approach to global optimization. Berlin.\end{itemize}}\\
\multicolumn{4}{|p{\dimexpr 16cm + 6\tabcolsep\relax}|}{\begin{itemize}\item Dorigo, Marco. Ant colony optimization. Cambridge: The MIT Press, 2004.\end{itemize}}\\
\multicolumn{4}{|p{\dimexpr 16cm + 6\tabcolsep\relax}|}{\begin{itemize}\item Edited by Vittorio Maniezzo, Thomas Stutzle, Stefan Voss. Matheuristics: hybridizing metaheuristics and mathematical programming. New York: Springer, 2009.\end{itemize}}\\
\hline
\end{longtable}
\end{center}

\clearpage
\begin{center}
\begin{tabular}{lccr}
 \multirow{3}{*}{\includegraphics[height=2.7cm]{brasao.png}} &
 \multicolumn{2}{c}{\bfseries UNIVERSIDADE FEDERAL DE OURO PRETO} &
 \ \ \ \ \multirow{3}{*}{\includegraphics[height=2.7cm]{ufop.png}} \\
 & \multicolumn{2}{c}{\bfseries PRÓ-REITORIA DE GRADUAÇÃO} & \\
 & \multicolumn{2}{c}{} & \\
 & \multicolumn{2}{c}{\Large\bfseries PROGRAMA DE DISCIPLINA} & \\
\end{tabular}
\end{center}

\begin{center}
\begin{longtable}{|p{4cm}|p{4cm}|p{4cm}|p{4cm}|}
\hline
\multicolumn{3}{|p{12cm}|}{Nome do Componente Curricular em Português:} &
\multicolumn{1}{p{4cm}|}{Código:} \\ 
\multicolumn{3}{|p{12cm}|}{\textbf{Gerência de Projetos de Software}} &
\textbf{BCC482}\\ 
\multicolumn{3}{|p{12cm}|}{Nome do Componente Curricular em Inglês:} & \\ 
\multicolumn{3}{|p{12cm}|}{\textbf{Software Project Management}} & \\ 
\hline
\multicolumn{3}{|p{12cm}|}{Nome e Sigla do Departamento} & Unidade Acadêmica: \\ 
\multicolumn{3}{|p{12cm}|}{Departamento de Computação (DECOM)} & {ICEB} \\ 
\hline
\multicolumn{4}{|p{16cm}|}{Modalidade de Oferta:
[X] presencial \hspace{0.6cm}
[ ] semipresencial \hspace{0.6cm}
[ ] à distância}\\
\hline
\multicolumn{2}{|p{8cm}|}{Carga horária semestral} &
\multicolumn{2}{p{8cm}|}{Carga horária semanal}\\
\hline
\multicolumn{1}{|p{4cm}|}{Total} &
\multicolumn{1}{p{4cm}|}{Extensionista} &
\multicolumn{1}{p{4cm}|}{Teórica} &
\multicolumn{1}{p{4cm}|}{Prática} \\ 
\multicolumn{1}{|p{4cm}|}{60\,horas} &
\multicolumn{1}{p{4cm}|}{0\;horas} &
\multicolumn{1}{p{4cm}|}{4\;horas/aula} &
\multicolumn{1}{p{4cm}|}{0\;horas/aula} \\ 
\hline
\multicolumn{4}{|p{16cm}|}{Ementa:}\\
\multicolumn{4}{|p{16cm}|}{}\\
\multicolumn{4}{|p{\dimexpr 16cm + 6\tabcolsep\relax}|}{Gerência de projetos de software, processo de software, ciclo de vida de software, gerência de projetos segundo o PMI-PMBOK, gerência de projetos ágil (SCRUM), ferramentas de gestão.}\\
\multicolumn{4}{|p{16cm}|}{}\\
\hline
\multicolumn{4}{|p{16cm}|}{Conteúdo programático:}\\*
\multicolumn{4}{|p{\dimexpr 16cm + 6\tabcolsep\relax}|}{\begin{enumerate}[start=1]\item Gerência de Projetos de Software: Gerência de Projetos; Gerência de Projetos de Software; Engenharia de Software e Gerência de Projetos; Processo de gestão de projetos; Gestão tradicional e ágil\end{enumerate}}\\
\multicolumn{4}{|p{\dimexpr 16cm + 6\tabcolsep\relax}|}{\begin{enumerate}[start=2]\item Processo de Software: Processo; Processos de desenvolvimento de software; Ciclo de vida\end{enumerate}}\\
\multicolumn{4}{|p{\dimexpr 16cm + 6\tabcolsep\relax}|}{\begin{enumerate}[start=3]\item Software: O que é um software; O que é um software enquanto um produto\end{enumerate}}\\
\multicolumn{4}{|p{\dimexpr 16cm + 6\tabcolsep\relax}|}{\begin{enumerate}[start=4]\item PMI: O que é PMI; Grupos de processos; Áreas de conhecimento; Projetos; Gestão de integração; Gestão do escopo; Gestão de tempo; Gestão de custos; Gestão da qualidade; Gestão de recursos humanos; Gestão das comunicações; Gestão de riscos; Gestão de aquisições; Gestão de envolvidos\end{enumerate}}\\
\multicolumn{4}{|p{\dimexpr 16cm + 6\tabcolsep\relax}|}{\begin{enumerate}[start=5]\item SCRUM: Definição; Modelo SCRUM; Papéis; Artefatos; Reuniões\end{enumerate}}\\
\multicolumn{4}{|p{\dimexpr 16cm + 6\tabcolsep\relax}|}{\begin{enumerate}[start=6]\item Ferramentas: Ferramentas para apoiar as diversas atividades e práticas da gestão de projetos nas abordagens tradicional e ágil\end{enumerate}}\\
\multicolumn{4}{|p{\dimexpr 16cm + 6\tabcolsep\relax}|}{}\\
\hline
\multicolumn{4}{|p{16cm}|}{Bibliografia Básica:}\\*
\multicolumn{4}{|p{\dimexpr 16cm + 6\tabcolsep\relax}|}{\begin{itemize}\item Martins, José Carlos Cordeiro. Gerenciando projetos de desenvolvimento de software com PMI, RUP e UML. 4. ed. Rio de Janeiro: Brasport, 2007.\end{itemize}}\\
\multicolumn{4}{|p{\dimexpr 16cm + 6\tabcolsep\relax}|}{\begin{itemize}\item Schwaber, Ken. Agile project management with Scrum. Redmond, Wash.: Microsoft Press, 2004.\end{itemize}}\\
\multicolumn{4}{|p{\dimexpr 16cm + 6\tabcolsep\relax}|}{\begin{itemize}\item Kim Heldman; tradução Luciana do Amaral Teixeira; revisão técnica Martius Vicente Rodriguez Y Rodriguez. Gerência de projetos = PMP Project management professional: guia para o exame oficial do PMI. 5. ed. rev. e atual. Rio de Janeiro: Elsevier, 2009.\end{itemize}}\\
\multicolumn{4}{|p{16cm}|}{}\\
\hline
\multicolumn{4}{|p{16cm}|}{Bibliografia Complementar:}\\*
\multicolumn{4}{|p{\dimexpr 16cm + 6\tabcolsep\relax}|}{\begin{itemize}\item Sommerville, Ian. Engenharia de software. 9. ed. São Paulo: Addison Wesley, 2011.\end{itemize}}\\
\multicolumn{4}{|p{\dimexpr 16cm + 6\tabcolsep\relax}|}{\begin{itemize}\item Roger S. Pressman; tradução Ariovaldo Griesi, Mario Moro Fecchio; revisão técnica Reginaldo Arakaki, Julio Arakaki, Renato Manzan de Andrade. Engenharia de software: uma abordagem profissional. 7. ed. Porto Alegre, RS: AMGH, 2011.\end{itemize}}\\
\multicolumn{4}{|p{\dimexpr 16cm + 6\tabcolsep\relax}|}{\begin{itemize}\item Daniel Capaldo Amaral... [et al.]. Gerenciamento Ágil de projetos: aplicação em produtos inovadores. São Paulo: Saraiva, 2011.\end{itemize}}\\
\multicolumn{4}{|p{\dimexpr 16cm + 6\tabcolsep\relax}|}{\begin{itemize}\item Gattoni, Roberto Luís Capuruço. Gestão do conhecimento aplicada à prática da gerência de projetos. Belo Horizonte, MG: FUMEC-FACE.\end{itemize}}\\
\multicolumn{4}{|p{\dimexpr 16cm + 6\tabcolsep\relax}|}{\begin{itemize}\item Joseph Phillips; consultoria editorial Lorenzo Ridolfi, Sérgio Colcher; tradução Ana Beatriz Tavares, Daniela Lacerda; revisão técnica Andrea Hemerly. Gerência de projetos de tecnologia da informação: no caminho certo, do início ao fim. Rio de Janeiro: Campus.\end{itemize}}\\
\hline
\end{longtable}
\end{center}

\clearpage
\begin{center}
\begin{tabular}{lccr}
 \multirow{3}{*}{\includegraphics[height=2.7cm]{brasao.png}} &
 \multicolumn{2}{c}{\bfseries UNIVERSIDADE FEDERAL DE OURO PRETO} &
 \ \ \ \ \multirow{3}{*}{\includegraphics[height=2.7cm]{ufop.png}} \\
 & \multicolumn{2}{c}{\bfseries PRÓ-REITORIA DE GRADUAÇÃO} & \\
 & \multicolumn{2}{c}{} & \\
 & \multicolumn{2}{c}{\Large\bfseries PROGRAMA DE DISCIPLINA} & \\
\end{tabular}
\end{center}

\begin{center}
\begin{longtable}{|p{4cm}|p{4cm}|p{4cm}|p{4cm}|}
\hline
\multicolumn{3}{|p{12cm}|}{Nome do Componente Curricular em Português:} &
\multicolumn{1}{p{4cm}|}{Código:} \\ 
\multicolumn{3}{|p{12cm}|}{\textbf{Qualidade de Software}} &
\textbf{BCC483}\\ 
\multicolumn{3}{|p{12cm}|}{Nome do Componente Curricular em Inglês:} & \\ 
\multicolumn{3}{|p{12cm}|}{\textbf{Software Quality}} & \\ 
\hline
\multicolumn{3}{|p{12cm}|}{Nome e Sigla do Departamento} & Unidade Acadêmica: \\ 
\multicolumn{3}{|p{12cm}|}{Departamento de Computação (DECOM)} & {ICEB} \\ 
\hline
\multicolumn{4}{|p{16cm}|}{Modalidade de Oferta:
[X] presencial \hspace{0.6cm}
[ ] semipresencial \hspace{0.6cm}
[ ] à distância}\\
\hline
\multicolumn{2}{|p{8cm}|}{Carga horária semestral} &
\multicolumn{2}{p{8cm}|}{Carga horária semanal}\\
\hline
\multicolumn{1}{|p{4cm}|}{Total} &
\multicolumn{1}{p{4cm}|}{Extensionista} &
\multicolumn{1}{p{4cm}|}{Teórica} &
\multicolumn{1}{p{4cm}|}{Prática} \\ 
\multicolumn{1}{|p{4cm}|}{60\,horas} &
\multicolumn{1}{p{4cm}|}{0\;horas} &
\multicolumn{1}{p{4cm}|}{4\;horas/aula} &
\multicolumn{1}{p{4cm}|}{0\;horas/aula} \\ 
\hline
\multicolumn{4}{|p{16cm}|}{Ementa:}\\
\multicolumn{4}{|p{16cm}|}{}\\
\multicolumn{4}{|p{\dimexpr 16cm + 6\tabcolsep\relax}|}{O que é qualidade. Processos de software. Qualidade de processo. Qualidade de produto. MPS-BR (Desenvolvimento de software e serviço). CMMI. Métricas. Teste. Verificação. Validação. Auditoria.}\\
\multicolumn{4}{|p{16cm}|}{}\\
\hline
\multicolumn{4}{|p{16cm}|}{Conteúdo programático:}\\*
\multicolumn{4}{|p{\dimexpr 16cm + 6\tabcolsep\relax}|}{\begin{enumerate}[start=1]\item O que é Qualidade: Qualidade; Qualidade de Software; Qualidade de Processo\end{enumerate}}\\
\multicolumn{4}{|p{\dimexpr 16cm + 6\tabcolsep\relax}|}{\begin{enumerate}[start=2]\item Qualidade de Processo: O que é processo e processo de software; Maturidade de processo; Normas e padrões; Melhoria do processo\end{enumerate}}\\
\multicolumn{4}{|p{\dimexpr 16cm + 6\tabcolsep\relax}|}{\begin{enumerate}[start=3]\item Qualidade de Produto: O que é produto e produto de software; Normas e padrões\end{enumerate}}\\
\multicolumn{4}{|p{\dimexpr 16cm + 6\tabcolsep\relax}|}{\begin{enumerate}[start=4]\item MPS-BR: Desenvolvimento de Software; Serviço\end{enumerate}}\\
\multicolumn{4}{|p{\dimexpr 16cm + 6\tabcolsep\relax}|}{\begin{enumerate}[start=5]\item CMMI\end{enumerate}}\\
\multicolumn{4}{|p{\dimexpr 16cm + 6\tabcolsep\relax}|}{\begin{enumerate}[start=6]\item Monitoramento e controle da qualidade: Métricas; Teste; Verificação; Validação; Auditoria\end{enumerate}}\\
\multicolumn{4}{|p{\dimexpr 16cm + 6\tabcolsep\relax}|}{}\\
\hline
\multicolumn{4}{|p{16cm}|}{Bibliografia Básica:}\\*
\multicolumn{4}{|p{\dimexpr 16cm + 6\tabcolsep\relax}|}{\begin{itemize}\item Mary Beth Chrissis; Mike Konrad; Sandy Shrum. CMMI: guidelines for process integration and product improvement. 2. ed. Boston: Addison Wesley, 2007.\end{itemize}}\\
\multicolumn{4}{|p{\dimexpr 16cm + 6\tabcolsep\relax}|}{\begin{itemize}\item Sommerville, Ian. Engenharia de software. 9. ed. São Paulo: Addison Wesley, 2011.\end{itemize}}\\
\multicolumn{4}{|p{16cm}|}{}\\
\hline
\multicolumn{4}{|p{16cm}|}{Bibliografia Complementar:}\\*
\multicolumn{4}{|p{\dimexpr 16cm + 6\tabcolsep\relax}|}{\begin{itemize}\item Bartié, Alexandre. Garantia da qualidade de software. Rio de Janeiro: Campus, 2002.\end{itemize}}\\
\multicolumn{4}{|p{\dimexpr 16cm + 6\tabcolsep\relax}|}{\begin{itemize}\item Edson Junio Rodrigues Lobo. Curso de engenharia de software: métodos e processos para garantir a qualidade no desenvolvimento de softwares. São Paulo: Digeratis Books, 2007.\end{itemize}}\\
\multicolumn{4}{|p{\dimexpr 16cm + 6\tabcolsep\relax}|}{\begin{itemize}\item Roger S. Pressman; tradução Ariovaldo Griesi, Mario Moro Fecchio; revisão técnica Reginaldo Arakaki, Julio Arakaki, Renato Manzan de Andrade. Engenharia de software: uma abordagem profissional. 7. ed. Porto Alegre, RS: AMGH, 2011.\end{itemize}}\\
\multicolumn{4}{|p{\dimexpr 16cm + 6\tabcolsep\relax}|}{\begin{itemize}\item Handbook of software quality assurance. 4th ed. Boston: Artech House, 2008.\end{itemize}}\\
\multicolumn{4}{|p{\dimexpr 16cm + 6\tabcolsep\relax}|}{\begin{itemize}\item Dennis M. Ahern, Aaron Clouse, Richard Turner. CMMI distilled: a practical introduction to integrated process improvement. 3rd ed. Upper Saddle River, NJ: Addison Wesley, 2008.\end{itemize}}\\
\hline
\end{longtable}
\end{center}

\clearpage
\begin{center}
\begin{tabular}{lccr}
 \multirow{3}{*}{\includegraphics[height=2.7cm]{brasao.png}} &
 \multicolumn{2}{c}{\bfseries UNIVERSIDADE FEDERAL DE OURO PRETO} &
 \ \ \ \ \multirow{3}{*}{\includegraphics[height=2.7cm]{ufop.png}} \\
 & \multicolumn{2}{c}{\bfseries PRÓ-REITORIA DE GRADUAÇÃO} & \\
 & \multicolumn{2}{c}{} & \\
 & \multicolumn{2}{c}{\Large\bfseries PROGRAMA DE DISCIPLINA} & \\
\end{tabular}
\end{center}

\begin{center}
\begin{longtable}{|p{4cm}|p{4cm}|p{4cm}|p{4cm}|}
\hline
\multicolumn{3}{|p{12cm}|}{Nome do Componente Curricular em Português:} &
\multicolumn{1}{p{4cm}|}{Código:} \\ 
\multicolumn{3}{|p{12cm}|}{\textbf{Design de Interação}} &
\textbf{BCC485}\\ 
\multicolumn{3}{|p{12cm}|}{Nome do Componente Curricular em Inglês:} & \\ 
\multicolumn{3}{|p{12cm}|}{\textbf{Interaction Design}} & \\ 
\hline
\multicolumn{3}{|p{12cm}|}{Nome e Sigla do Departamento} & Unidade Acadêmica: \\ 
\multicolumn{3}{|p{12cm}|}{Departamento de Computação (DECOM)} & {ICEB} \\ 
\hline
\multicolumn{4}{|p{16cm}|}{Modalidade de Oferta:
[X] presencial \hspace{0.6cm}
[ ] semipresencial \hspace{0.6cm}
[ ] à distância}\\
\hline
\multicolumn{2}{|p{8cm}|}{Carga horária semestral} &
\multicolumn{2}{p{8cm}|}{Carga horária semanal}\\
\hline
\multicolumn{1}{|p{4cm}|}{Total} &
\multicolumn{1}{p{4cm}|}{Extensionista} &
\multicolumn{1}{p{4cm}|}{Teórica} &
\multicolumn{1}{p{4cm}|}{Prática} \\ 
\multicolumn{1}{|p{4cm}|}{60\,horas} &
\multicolumn{1}{p{4cm}|}{0\;horas} &
\multicolumn{1}{p{4cm}|}{4\;horas/aula} &
\multicolumn{1}{p{4cm}|}{0\;horas/aula} \\ 
\hline
\multicolumn{4}{|p{16cm}|}{Ementa:}\\
\multicolumn{4}{|p{16cm}|}{}\\
\multicolumn{4}{|p{\dimexpr 16cm + 6\tabcolsep\relax}|}{O processo de Design de Interação; levantamento de requisitos; modelagem de tarefas e modelagem de interação; prototipagem e construção; noções de Design Visual.}\\
\multicolumn{4}{|p{16cm}|}{}\\
\hline
\multicolumn{4}{|p{16cm}|}{Conteúdo programático:}\\*
\multicolumn{4}{|p{\dimexpr 16cm + 6\tabcolsep\relax}|}{\begin{enumerate}[start=1]\item O Processo de Design de Interação: Exemplos de bons e maus projetos; O que é Design de Interação; Metas de usabilidade e metas da experiência de usuário; Modelos Conceituais; Metáforas de Interface; Paradigmas de Interação; Modelos de ciclo de vida para o Design de Interação; Modelos de ciclo de vida em Engenharia de Software; Integração de Modelos de Engenharia de Software e IHC\end{enumerate}}\\
\multicolumn{4}{|p{\dimexpr 16cm + 6\tabcolsep\relax}|}{\begin{enumerate}[start=2]\item Levantamento de Requisitos: Requisitos Funcionais e Requisitos Não Funcionais; Questionários; Entrevistas; Grupos de Foco; Classificação de Cartões; Brainstorming; Investigação Contextual\end{enumerate}}\\
\multicolumn{4}{|p{\dimexpr 16cm + 6\tabcolsep\relax}|}{\begin{enumerate}[start=3]\item Modelagem de Tarefas e Interação: Representações abstratas da interface/interação; Perfis de usuários; Personas e objetivos; Storyboards; Cenários de análise; Cenários de interação; Casos de uso; Análise hierárquica de tarefas; Diagramas de Interação; MoLIC\end{enumerate}}\\
\multicolumn{4}{|p{\dimexpr 16cm + 6\tabcolsep\relax}|}{\begin{enumerate}[start=4]\item Prototipagem e Construção: Protótipos em papel; Protótipos de baixa fidelidade; Protótipos de alta fidelidade; Ferramentas de prototipagem; Design Conceitual: dos requisitos ao primeiro design; Design Físico\end{enumerate}}\\
\multicolumn{4}{|p{\dimexpr 16cm + 6\tabcolsep\relax}|}{\begin{enumerate}[start=5]\item Noções de Design Visual: Projeto de ícones; Projeto de layout; Projeto de cores; Projeto de tipografia\end{enumerate}}\\
\multicolumn{4}{|p{\dimexpr 16cm + 6\tabcolsep\relax}|}{}\\
\hline
\multicolumn{4}{|p{16cm}|}{Bibliografia Básica:}\\*
\multicolumn{4}{|p{\dimexpr 16cm + 6\tabcolsep\relax}|}{\begin{itemize}\item Jennifer Preece, Yvonne Rogers e Helen Sharp; tradução de Viviane Possamai. Design de interação: além da interação homem-computador. Porto Alegre, RS: Bookman, 2005.\end{itemize}}\\
\multicolumn{4}{|p{\dimexpr 16cm + 6\tabcolsep\relax}|}{\begin{itemize}\item Robin Williams; tradução: Laura Karin Gillon. Design para quem não é designer: noções básicas de planejamento visual. 3. ed. São Paulo: Callis, 2009.\end{itemize}}\\
\multicolumn{4}{|p{\dimexpr 16cm + 6\tabcolsep\relax}|}{\begin{itemize}\item Ben Shneiderman \& Catherine Plaisant. Designing the user interface: strategies for effective human-computer interaction. 5th ed. Boston: Addison Wesley, 2010.\end{itemize}}\\
\multicolumn{4}{|p{16cm}|}{}\\
\hline
\multicolumn{4}{|p{16cm}|}{Bibliografia Complementar:}\\*
\multicolumn{4}{|p{\dimexpr 16cm + 6\tabcolsep\relax}|}{\begin{itemize}\item Barbosa, Simone Diniz Junqueira. Interação humano-computador. Rio de Janeiro: Campus, 2010.\end{itemize}}\\
\multicolumn{4}{|p{\dimexpr 16cm + 6\tabcolsep\relax}|}{\begin{itemize}\item Benyon, David. Interação humano-computador. 2. ed. São Paulo: Pearson, 2011.\end{itemize}}\\
\multicolumn{4}{|p{\dimexpr 16cm + 6\tabcolsep\relax}|}{\begin{itemize}\item Research methods for human-computer interaction. Cambridge, UK.\end{itemize}}\\
\multicolumn{4}{|p{\dimexpr 16cm + 6\tabcolsep\relax}|}{\begin{itemize}\item Edited by Julie A. Jacko. The human-computer interaction handbook: fundamentals, evolving technologies, and emerging applications. 3rd ed. Boca Raton, FL: CRC Press, 2012.\end{itemize}}\\
\multicolumn{4}{|p{\dimexpr 16cm + 6\tabcolsep\relax}|}{\begin{itemize}\item Souza, Clarisse da Silva Vieira Camelo de. The semiotic engineering of human-computer interaction. Cambridge, MA: MIT Press, 2005.\end{itemize}}\\
\hline
\end{longtable}
\end{center}

\clearpage
\begin{center}
\begin{tabular}{lccr}
 \multirow{3}{*}{\includegraphics[height=2.7cm]{brasao.png}} &
 \multicolumn{2}{c}{\bfseries UNIVERSIDADE FEDERAL DE OURO PRETO} &
 \ \ \ \ \multirow{3}{*}{\includegraphics[height=2.7cm]{ufop.png}} \\
 & \multicolumn{2}{c}{\bfseries PRÓ-REITORIA DE GRADUAÇÃO} & \\
 & \multicolumn{2}{c}{} & \\
 & \multicolumn{2}{c}{\Large\bfseries PROGRAMA DE DISCIPLINA} & \\
\end{tabular}
\end{center}

\begin{center}
\begin{longtable}{|p{4cm}|p{4cm}|p{4cm}|p{4cm}|}
\hline
\multicolumn{3}{|p{12cm}|}{Nome do Componente Curricular em Português:} &
\multicolumn{1}{p{4cm}|}{Código:} \\ 
\multicolumn{3}{|p{12cm}|}{\textbf{UI e UX para Inteligência Artificial}} &
\textbf{BCC851}\\ 
\multicolumn{3}{|p{12cm}|}{Nome do Componente Curricular em Inglês:} & \\ 
\multicolumn{3}{|p{12cm}|}{\textbf{UI and UX for Artificial Intelligence}} & \\ 
\hline
\multicolumn{3}{|p{12cm}|}{Nome e Sigla do Departamento} & Unidade Acadêmica: \\ 
\multicolumn{3}{|p{12cm}|}{Departamento de Computação (DECOM)} & {ICEB} \\ 
\hline
\multicolumn{4}{|p{16cm}|}{Modalidade de Oferta:
[X] presencial \hspace{0.6cm}
[ ] semipresencial \hspace{0.6cm}
[ ] à distância}\\
\hline
\multicolumn{2}{|p{8cm}|}{Carga horária semestral} &
\multicolumn{2}{p{8cm}|}{Carga horária semanal}\\
\hline
\multicolumn{1}{|p{4cm}|}{Total} &
\multicolumn{1}{p{4cm}|}{Extensionista} &
\multicolumn{1}{p{4cm}|}{Teórica} &
\multicolumn{1}{p{4cm}|}{Prática} \\ 
\multicolumn{1}{|p{4cm}|}{60\,horas} &
\multicolumn{1}{p{4cm}|}{0\;horas} &
\multicolumn{1}{p{4cm}|}{4\;horas/aula} &
\multicolumn{1}{p{4cm}|}{0\;horas/aula} \\ 
\hline
\multicolumn{4}{|p{16cm}|}{Ementa:}\\
\multicolumn{4}{|p{16cm}|}{}\\
\multicolumn{4}{|p{\dimexpr 16cm + 6\tabcolsep\relax}|}{Estudo dos princípios de design de interfaces de usuário (UI) e experiência do usuário (UX) aplicados a sistemas com Inteligência Artificial. Foco em usabilidade, acessibilidade, interação homem-máquina e avaliação de interfaces inteligentes.}\\
\multicolumn{4}{|p{16cm}|}{}\\
\hline
\multicolumn{4}{|p{16cm}|}{Conteúdo programático:}\\*
\multicolumn{4}{|p{\dimexpr 16cm + 6\tabcolsep\relax}|}{\begin{enumerate}[start=1]\item Fundamentos de UI e UX: conceitos, história e importância\end{enumerate}}\\
\multicolumn{4}{|p{\dimexpr 16cm + 6\tabcolsep\relax}|}{\begin{enumerate}[start=2]\item Princípios de design centrado no usuário (user-centered design)\end{enumerate}}\\
\multicolumn{4}{|p{\dimexpr 16cm + 6\tabcolsep\relax}|}{\begin{enumerate}[start=3]\item Especificidades de UX para sistemas baseados em IA: transparência, confiança e controle do usuário\end{enumerate}}\\
\multicolumn{4}{|p{\dimexpr 16cm + 6\tabcolsep\relax}|}{\begin{enumerate}[start=4]\item Técnicas de prototipagem e wireframing para interfaces inteligentes\end{enumerate}}\\
\multicolumn{4}{|p{\dimexpr 16cm + 6\tabcolsep\relax}|}{\begin{enumerate}[start=5]\item Interação humano-IA: design de conversação, chatbots e assistentes virtuais\end{enumerate}}\\
\multicolumn{4}{|p{\dimexpr 16cm + 6\tabcolsep\relax}|}{\begin{enumerate}[start=6]\item Avaliação de usabilidade: métodos qualitativos e quantitativos\end{enumerate}}\\
\multicolumn{4}{|p{\dimexpr 16cm + 6\tabcolsep\relax}|}{\begin{enumerate}[start=7]\item Acessibilidade digital e inclusão em sistemas inteligentes\end{enumerate}}\\
\multicolumn{4}{|p{\dimexpr 16cm + 6\tabcolsep\relax}|}{\begin{enumerate}[start=8]\item Desafios éticos em UI/UX para IA: viés, privacidade e explicabilidade\end{enumerate}}\\
\multicolumn{4}{|p{\dimexpr 16cm + 6\tabcolsep\relax}|}{\begin{enumerate}[start=9]\item Ferramentas e frameworks para desenvolvimento de UI/UX em IA\end{enumerate}}\\
\multicolumn{4}{|p{\dimexpr 16cm + 6\tabcolsep\relax}|}{\begin{enumerate}[start=10]\item Estudos de caso: aplicações reais em IA e interfaces inteligentes\end{enumerate}}\\
\multicolumn{4}{|p{\dimexpr 16cm + 6\tabcolsep\relax}|}{}\\
\hline
\multicolumn{4}{|p{16cm}|}{Bibliografia Básica:}\\*
\multicolumn{4}{|p{\dimexpr 16cm + 6\tabcolsep\relax}|}{\begin{itemize}\item[] (nenhum item informado)\end{itemize}}\\
\multicolumn{4}{|p{16cm}|}{}\\
\hline
\multicolumn{4}{|p{16cm}|}{Bibliografia Complementar:}\\*
\multicolumn{4}{|p{\dimexpr 16cm + 6\tabcolsep\relax}|}{\begin{itemize}\item[] (nenhum item informado)\end{itemize}}\\
\hline
\end{longtable}
\end{center}

\clearpage
\begin{center}
\begin{tabular}{lccr}
 \multirow{3}{*}{\includegraphics[height=2.7cm]{brasao.png}} &
 \multicolumn{2}{c}{\bfseries UNIVERSIDADE FEDERAL DE OURO PRETO} &
 \ \ \ \ \multirow{3}{*}{\includegraphics[height=2.7cm]{ufop.png}} \\
 & \multicolumn{2}{c}{\bfseries PRÓ-REITORIA DE GRADUAÇÃO} & \\
 & \multicolumn{2}{c}{} & \\
 & \multicolumn{2}{c}{\Large\bfseries PROGRAMA DE DISCIPLINA} & \\
\end{tabular}
\end{center}

\begin{center}
\begin{longtable}{|p{4cm}|p{4cm}|p{4cm}|p{4cm}|}
\hline
\multicolumn{3}{|p{12cm}|}{Nome do Componente Curricular em Português:} &
\multicolumn{1}{p{4cm}|}{Código:} \\ 
\multicolumn{3}{|p{12cm}|}{\textbf{Processamento de Áudio}} &
\textbf{BCC852}\\ 
\multicolumn{3}{|p{12cm}|}{Nome do Componente Curricular em Inglês:} & \\ 
\multicolumn{3}{|p{12cm}|}{\textbf{Audio Processing}} & \\ 
\hline
\multicolumn{3}{|p{12cm}|}{Nome e Sigla do Departamento} & Unidade Acadêmica: \\ 
\multicolumn{3}{|p{12cm}|}{Departamento de Computação (DECOM)} & {ICEB} \\ 
\hline
\multicolumn{4}{|p{16cm}|}{Modalidade de Oferta:
[X] presencial \hspace{0.6cm}
[ ] semipresencial \hspace{0.6cm}
[ ] à distância}\\
\hline
\multicolumn{2}{|p{8cm}|}{Carga horária semestral} &
\multicolumn{2}{p{8cm}|}{Carga horária semanal}\\
\hline
\multicolumn{1}{|p{4cm}|}{Total} &
\multicolumn{1}{p{4cm}|}{Extensionista} &
\multicolumn{1}{p{4cm}|}{Teórica} &
\multicolumn{1}{p{4cm}|}{Prática} \\ 
\multicolumn{1}{|p{4cm}|}{60\,horas} &
\multicolumn{1}{p{4cm}|}{0\;horas} &
\multicolumn{1}{p{4cm}|}{4\;horas/aula} &
\multicolumn{1}{p{4cm}|}{0\;horas/aula} \\ 
\hline
\multicolumn{4}{|p{16cm}|}{Ementa:}\\
\multicolumn{4}{|p{16cm}|}{}\\
\multicolumn{4}{|p{\dimexpr 16cm + 6\tabcolsep\relax}|}{Introdução aos conceitos e técnicas fundamentais para o processamento digital de sinais de áudio. Análise, transformação, filtragem e síntese de áudio aplicados em sistemas de reconhecimento, síntese sonora e IA.}\\
\multicolumn{4}{|p{16cm}|}{}\\
\hline
\multicolumn{4}{|p{16cm}|}{Conteúdo programático:}\\*
\multicolumn{4}{|p{\dimexpr 16cm + 6\tabcolsep\relax}|}{\begin{enumerate}[start=1]\item Modelos acústicos tradicionais, misturas de gaussianas e cadeias ocultas de Markov\end{enumerate}}\\
\multicolumn{4}{|p{\dimexpr 16cm + 6\tabcolsep\relax}|}{\begin{enumerate}[start=2]\item Representações de sinais para áudio e voz\end{enumerate}}\\
\multicolumn{4}{|p{\dimexpr 16cm + 6\tabcolsep\relax}|}{\begin{enumerate}[start=3]\item Arquiteturas de redes neurais para reconhecimento de fala\end{enumerate}}\\
\multicolumn{4}{|p{\dimexpr 16cm + 6\tabcolsep\relax}|}{\begin{enumerate}[start=4]\item Modelos de encoder, vocoder e arquiteturas para sintetização de voz\end{enumerate}}\\
\multicolumn{4}{|p{\dimexpr 16cm + 6\tabcolsep\relax}|}{\begin{enumerate}[start=5]\item Arquiteturas para aplicações em música\end{enumerate}}\\
\multicolumn{4}{|p{\dimexpr 16cm + 6\tabcolsep\relax}|}{}\\
\hline
\multicolumn{4}{|p{16cm}|}{Bibliografia Básica:}\\*
\multicolumn{4}{|p{\dimexpr 16cm + 6\tabcolsep\relax}|}{\begin{itemize}\item KAMATH, U. L.; WHITAKER, J. Deep learning for NLP and speech recognition. Springer Nature, 2019.\end{itemize}}\\
\multicolumn{4}{|p{\dimexpr 16cm + 6\tabcolsep\relax}|}{\begin{itemize}\item DONG, Y.; DENG, Y. Automatic speech recognition. Springer London Limited, 2016.\end{itemize}}\\
\multicolumn{4}{|p{\dimexpr 16cm + 6\tabcolsep\relax}|}{\begin{itemize}\item GOPI, E. S. Digital speech processing using Matlab. New Delhi: Springer, 2014.\end{itemize}}\\
\multicolumn{4}{|p{16cm}|}{}\\
\hline
\multicolumn{4}{|p{16cm}|}{Bibliografia Complementar:}\\*
\multicolumn{4}{|p{\dimexpr 16cm + 6\tabcolsep\relax}|}{\begin{itemize}\item CAMASTRA, F.; VINCIARELLI, A. Machine learning for audio, image and video analysis: theory and applications. Springer, 2015.\end{itemize}}\\
\multicolumn{4}{|p{\dimexpr 16cm + 6\tabcolsep\relax}|}{\begin{itemize}\item STEVENS, E.; ANTIGA, L. Deep Learning with PyTorch. Manning Publications, 2020.\end{itemize}}\\
\multicolumn{4}{|p{\dimexpr 16cm + 6\tabcolsep\relax}|}{\begin{itemize}\item GALEONE, P. Hands-On Neural Networks with TensorFlow 2.0: Understand TensorFlow, from static graph to eager execution, and design neural networks. Packt Publishing, 2019.\end{itemize}}\\
\hline
\end{longtable}
\end{center}

\clearpage
\begin{center}
\begin{tabular}{lccr}
 \multirow{3}{*}{\includegraphics[height=2.7cm]{brasao.png}} &
 \multicolumn{2}{c}{\bfseries UNIVERSIDADE FEDERAL DE OURO PRETO} &
 \ \ \ \ \multirow{3}{*}{\includegraphics[height=2.7cm]{ufop.png}} \\
 & \multicolumn{2}{c}{\bfseries PRÓ-REITORIA DE GRADUAÇÃO} & \\
 & \multicolumn{2}{c}{} & \\
 & \multicolumn{2}{c}{\Large\bfseries PROGRAMA DE DISCIPLINA} & \\
\end{tabular}
\end{center}

\begin{center}
\begin{longtable}{|p{4cm}|p{4cm}|p{4cm}|p{4cm}|}
\hline
\multicolumn{3}{|p{12cm}|}{Nome do Componente Curricular em Português:} &
\multicolumn{1}{p{4cm}|}{Código:} \\ 
\multicolumn{3}{|p{12cm}|}{\textbf{Bioinformática: Conceitos e Aplicações}} &
\textbf{CBI076}\\ 
\multicolumn{3}{|p{12cm}|}{Nome do Componente Curricular em Inglês:} & \\ 
\multicolumn{3}{|p{12cm}|}{\textbf{Bioinformatics: Concepts and Applications}} & \\ 
\hline
\multicolumn{3}{|p{12cm}|}{Nome e Sigla do Departamento} & Unidade Acadêmica: \\ 
\multicolumn{3}{|p{12cm}|}{Departamento de Ciências Biológicas (DECBI)} & {ICEB} \\ 
\hline
\multicolumn{4}{|p{16cm}|}{Modalidade de Oferta:
[X] presencial \hspace{0.6cm}
[ ] semipresencial \hspace{0.6cm}
[ ] à distância}\\
\hline
\multicolumn{2}{|p{8cm}|}{Carga horária semestral} &
\multicolumn{2}{p{8cm}|}{Carga horária semanal}\\
\hline
\multicolumn{1}{|p{4cm}|}{Total} &
\multicolumn{1}{p{4cm}|}{Extensionista} &
\multicolumn{1}{p{4cm}|}{Teórica} &
\multicolumn{1}{p{4cm}|}{Prática} \\ 
\multicolumn{1}{|p{4cm}|}{60\,horas} &
\multicolumn{1}{p{4cm}|}{0\;horas} &
\multicolumn{1}{p{4cm}|}{2\;horas/aula} &
\multicolumn{1}{p{4cm}|}{2\;horas/aula} \\ 
\hline
\multicolumn{4}{|p{16cm}|}{Ementa:}\\
\multicolumn{4}{|p{16cm}|}{}\\
\multicolumn{4}{|p{\dimexpr 16cm + 6\tabcolsep\relax}|}{Repositório de dados; Conceitos fundamentais sobre composição de genomas; Ferramentas bioinformáticas online para o estudo de genes; Matrizes de comparação; Alinhamento e comparação de sequências. Análise de cladogramas e filogenia. Evolução molecular. Análise estrutural de proteínas.}\\
\multicolumn{4}{|p{16cm}|}{}\\
\hline
\multicolumn{4}{|p{16cm}|}{Conteúdo programático:}\\*
\multicolumn{4}{|p{\dimexpr 16cm + 6\tabcolsep\relax}|}{\begin{enumerate}[start=1]\item Repositório de dados em biologia\end{enumerate}}\\
\multicolumn{4}{|p{\dimexpr 16cm + 6\tabcolsep\relax}|}{\begin{enumerate}[start=2]\item Parálogos, ortólogos, core e pan genomas\end{enumerate}}\\
\multicolumn{4}{|p{\dimexpr 16cm + 6\tabcolsep\relax}|}{\begin{enumerate}[start=3]\item Matrizes de alinhamento e comparação de sequências\end{enumerate}}\\
\multicolumn{4}{|p{\dimexpr 16cm + 6\tabcolsep\relax}|}{\begin{enumerate}[start=4]\item Alinhamento par a par\end{enumerate}}\\
\multicolumn{4}{|p{\dimexpr 16cm + 6\tabcolsep\relax}|}{\begin{enumerate}[start=5]\item Alinhamento múltiplo de sequências\end{enumerate}}\\
\multicolumn{4}{|p{\dimexpr 16cm + 6\tabcolsep\relax}|}{\begin{enumerate}[start=6]\item Construção de cladogramas e filogramas\end{enumerate}}\\
\multicolumn{4}{|p{\dimexpr 16cm + 6\tabcolsep\relax}|}{\begin{enumerate}[start=7]\item Evolução de genes e genomas\end{enumerate}}\\
\multicolumn{4}{|p{\dimexpr 16cm + 6\tabcolsep\relax}|}{\begin{enumerate}[start=8]\item Análise estrutural de proteínas\end{enumerate}}\\
\multicolumn{4}{|p{\dimexpr 16cm + 6\tabcolsep\relax}|}{}\\
\hline
\multicolumn{4}{|p{16cm}|}{Bibliografia Básica:}\\*
\multicolumn{4}{|p{\dimexpr 16cm + 6\tabcolsep\relax}|}{\begin{itemize}\item Krawetz, S.A. \& Womble, D.D. Introduction to bioinformatics – A theoretical and practical approach. Humana Press, 2003.\end{itemize}}\\
\multicolumn{4}{|p{\dimexpr 16cm + 6\tabcolsep\relax}|}{\begin{itemize}\item Lesk, A. Introdução à bioinformática. 2. ed. Artmed, 2008.\end{itemize}}\\
\multicolumn{4}{|p{\dimexpr 16cm + 6\tabcolsep\relax}|}{\begin{itemize}\item Moreira, L.M. Ciências Genômicas: fundamentos e aplicações. Cubo, 2015.\end{itemize}}\\
\multicolumn{4}{|p{\dimexpr 16cm + 6\tabcolsep\relax}|}{\begin{itemize}\item Samuelson, T. Genomics and bioinformatics. Cambridge: University Press, 2012.\end{itemize}}\\
\multicolumn{4}{|p{16cm}|}{}\\
\hline
\multicolumn{4}{|p{16cm}|}{Bibliografia Complementar:}\\*
\multicolumn{4}{|p{\dimexpr 16cm + 6\tabcolsep\relax}|}{\begin{itemize}\item Gonçalves, E.G.; Lorenzi, H. Morfologia Vegetal: organografia e dicionário de morfologia das plantas vasculares. Nova Odessa: Instituto Plantarum, 2011.\end{itemize}}\\
\multicolumn{4}{|p{\dimexpr 16cm + 6\tabcolsep\relax}|}{\begin{itemize}\item Peixoto, A.L.; Costa, M. Manual de Procedimentos para Herbários. Recife: INCT - Herbário Virtual da Flora e dos Fungos, Editora Universitária UFPE, 2013.\end{itemize}}\\
\hline
\end{longtable}
\end{center}

\clearpage
\begin{center}
\begin{tabular}{lccr}
 \multirow{3}{*}{\includegraphics[height=2.7cm]{brasao.png}} &
 \multicolumn{2}{c}{\bfseries UNIVERSIDADE FEDERAL DE OURO PRETO} &
 \ \ \ \ \multirow{3}{*}{\includegraphics[height=2.7cm]{ufop.png}} \\
 & \multicolumn{2}{c}{\bfseries PRÓ-REITORIA DE GRADUAÇÃO} & \\
 & \multicolumn{2}{c}{} & \\
 & \multicolumn{2}{c}{\Large\bfseries PROGRAMA DE DISCIPLINA} & \\
\end{tabular}
\end{center}

\begin{center}
\begin{longtable}{|p{4cm}|p{4cm}|p{4cm}|p{4cm}|}
\hline
\multicolumn{3}{|p{12cm}|}{Nome do Componente Curricular em Português:} &
\multicolumn{1}{p{4cm}|}{Código:} \\ 
\multicolumn{3}{|p{12cm}|}{\textbf{Estatística Descritiva}} &
\textbf{EST313}\\ 
\multicolumn{3}{|p{12cm}|}{Nome do Componente Curricular em Inglês:} & \\ 
\multicolumn{3}{|p{12cm}|}{\textbf{Descriptive Statistics}} & \\ 
\hline
\multicolumn{3}{|p{12cm}|}{Nome e Sigla do Departamento} & Unidade Acadêmica: \\ 
\multicolumn{3}{|p{12cm}|}{Departamento de Estatística (DEEST)} & {ICEB} \\ 
\hline
\multicolumn{4}{|p{16cm}|}{Modalidade de Oferta:
[ ] presencial \hspace{0.6cm}
[X] semipresencial \hspace{0.6cm}
[ ] à distância}\\
\hline
\multicolumn{2}{|p{8cm}|}{Carga horária semestral} &
\multicolumn{2}{p{8cm}|}{Carga horária semanal}\\
\hline
\multicolumn{1}{|p{4cm}|}{Total} &
\multicolumn{1}{p{4cm}|}{Extensionista} &
\multicolumn{1}{p{4cm}|}{Teórica} &
\multicolumn{1}{p{4cm}|}{Prática} \\ 
\multicolumn{1}{|p{4cm}|}{60\,horas} &
\multicolumn{1}{p{4cm}|}{0\;horas} &
\multicolumn{1}{p{4cm}|}{4\;horas/aula} &
\multicolumn{1}{p{4cm}|}{0\;horas/aula} \\ 
\hline
\multicolumn{4}{|p{16cm}|}{Ementa:}\\
\multicolumn{4}{|p{16cm}|}{}\\
\multicolumn{4}{|p{\dimexpr 16cm + 6\tabcolsep\relax}|}{Introdução à Estatística. A Estatística e o Trabalho Científico. Metodologia da Pesquisa. Noções de Amostragem. Levantamento de dados. Representação Tabular e Gráfica. Distribuição de Frequências. Medidas de Posição (Tendência Central e Separatrizes). Medidas de Dispersão (Absoluta e Relativa). Medidas de Forma (Momentos, Assimetria e Curtose). Outras Estratégias de Análise (Ramo-e-Folhas, Esquema dos 5 números e Boxplot).}\\
\multicolumn{4}{|p{16cm}|}{}\\
\hline
\multicolumn{4}{|p{16cm}|}{Conteúdo programático:}\\*
\multicolumn{4}{|p{\dimexpr 16cm + 6\tabcolsep\relax}|}{\begin{enumerate}[start=1]\item Introdução à Estatística: Estatística como ciência e profissão. História da Estatística e suas aplicações. Alguns conceitos: população e amostra, variáveis e classificação de variáveis. O método científico e as etapas do método estatístico. Metodologia da pesquisa\end{enumerate}}\\
\multicolumn{4}{|p{\dimexpr 16cm + 6\tabcolsep\relax}|}{\begin{enumerate}[start=2]\item Noções de Amostragem: Conceitos básicos em amostragem. Tipos de amostragem: Não-probabilística e Probabilística. Algumas técnicas não-probabilísticas: conveniência, quotas, material contínuo. Algumas técnicas probabilísticas: aleatória simples, sistemática, estratificada, por conglomerado\end{enumerate}}\\
\multicolumn{4}{|p{\dimexpr 16cm + 6\tabcolsep\relax}|}{\begin{enumerate}[start=3]\item Levantamento de Dados: Mensuração. Consistência. Tabulação de dados\end{enumerate}}\\
\multicolumn{4}{|p{\dimexpr 16cm + 6\tabcolsep\relax}|}{\begin{enumerate}[start=4]\item Representação Tabular e Gráfica: Séries homógradas: cronológicas, especificativas, geográficas e mistas; séries heterógradas\end{enumerate}}\\
\multicolumn{4}{|p{\dimexpr 16cm + 6\tabcolsep\relax}|}{\begin{enumerate}[start=5]\item Distribuição de Frequências\end{enumerate}}\\
\multicolumn{4}{|p{\dimexpr 16cm + 6\tabcolsep\relax}|}{\begin{enumerate}[start=6]\item Medidas de Posição: Tendência Central (média, mediana, moda) e Separatrizes (quartis, decis, percentis)\end{enumerate}}\\
\multicolumn{4}{|p{\dimexpr 16cm + 6\tabcolsep\relax}|}{\begin{enumerate}[start=7]\item Medidas de Dispersão: Absoluta: amplitude, desvio médio, desvio padrão, variância. Relativa: coeficiente de variação, escore padronizado\end{enumerate}}\\
\multicolumn{4}{|p{\dimexpr 16cm + 6\tabcolsep\relax}|}{\begin{enumerate}[start=8]\item Medidas de Forma: Momentos de ordem 1, 2, 3 e 4. Assimetria, utilizando momentos. Índices de curtose ou achatamento utilizando: momentos, coeficiente percentílico de curtose\end{enumerate}}\\
\multicolumn{4}{|p{\dimexpr 16cm + 6\tabcolsep\relax}|}{\begin{enumerate}[start=9]\item Outras Estratégias de Análise: Ramo-e-folhas. O esquema dos 5-números. Desenho esquemático (Box-Plot)\end{enumerate}}\\
\multicolumn{4}{|p{\dimexpr 16cm + 6\tabcolsep\relax}|}{\begin{enumerate}[start=10]\item Os conteúdos serão trabalhados de forma semipresencial, sendo: 53 horas presenciais e 07 horas à distância\end{enumerate}}\\
\multicolumn{4}{|p{\dimexpr 16cm + 6\tabcolsep\relax}|}{}\\
\hline
\multicolumn{4}{|p{16cm}|}{Bibliografia Básica:}\\*
\multicolumn{4}{|p{\dimexpr 16cm + 6\tabcolsep\relax}|}{\begin{itemize}\item COSTA, Giovani Gláucio de Oliveira. Curso de Estatística Básica: Teoria e Prática. 2. ed. São Paulo: Atlas, 2015.\end{itemize}}\\
\multicolumn{4}{|p{\dimexpr 16cm + 6\tabcolsep\relax}|}{\begin{itemize}\item VIEIRA, Sônia. Elementos de Estatística. 5. ed. São Paulo: Atlas, 2013.\end{itemize}}\\
\multicolumn{4}{|p{\dimexpr 16cm + 6\tabcolsep\relax}|}{\begin{itemize}\item BUSSAB, Wilton de O.; MORETTIN, Pedro A. Estatística Básica. 9. ed. São Paulo: Saraiva, 2017.\end{itemize}}\\
\multicolumn{4}{|p{16cm}|}{}\\
\hline
\multicolumn{4}{|p{16cm}|}{Bibliografia Complementar:}\\*
\multicolumn{4}{|p{\dimexpr 16cm + 6\tabcolsep\relax}|}{\begin{itemize}\item TRIOLA, M. F. Introdução à Estatística. 12. ed. LTC, 2017.\end{itemize}}\\
\multicolumn{4}{|p{\dimexpr 16cm + 6\tabcolsep\relax}|}{\begin{itemize}\item MORGADO, Augusto César de Oliveira et al. Análise Combinatória e Probabilidade. 10. ed. SBM, 2016.\end{itemize}}\\
\multicolumn{4}{|p{\dimexpr 16cm + 6\tabcolsep\relax}|}{\begin{itemize}\item HAZZAN, Samuel. Fundamentos de Matemática elementar: combinatória e probabilidade. 8. ed. São Paulo: Atual, 2013. v. 5\end{itemize}}\\
\multicolumn{4}{|p{\dimexpr 16cm + 6\tabcolsep\relax}|}{\begin{itemize}\item LEVIN, Jack; FOX, James Alan. Estatística para Ciências do Comportamento. 9. ed. Pearson, 2004.\end{itemize}}\\
\multicolumn{4}{|p{\dimexpr 16cm + 6\tabcolsep\relax}|}{\begin{itemize}\item KERLINGER, Fred. Metodologia da Pesquisa em Ciências Sociais. São Paulo: EDUSP, 1980.\end{itemize}}\\
\hline
\end{longtable}
\end{center}

\clearpage
\begin{center}
\begin{tabular}{lccr}
 \multirow{3}{*}{\includegraphics[height=2.7cm]{brasao.png}} &
 \multicolumn{2}{c}{\bfseries UNIVERSIDADE FEDERAL DE OURO PRETO} &
 \ \ \ \ \multirow{3}{*}{\includegraphics[height=2.7cm]{ufop.png}} \\
 & \multicolumn{2}{c}{\bfseries PRÓ-REITORIA DE GRADUAÇÃO} & \\
 & \multicolumn{2}{c}{} & \\
 & \multicolumn{2}{c}{\Large\bfseries PROGRAMA DE DISCIPLINA} & \\
\end{tabular}
\end{center}

\begin{center}
\begin{longtable}{|p{4cm}|p{4cm}|p{4cm}|p{4cm}|}
\hline
\multicolumn{3}{|p{12cm}|}{Nome do Componente Curricular em Português:} &
\multicolumn{1}{p{4cm}|}{Código:} \\ 
\multicolumn{3}{|p{12cm}|}{\textbf{Análise de Dados Categóricos}} &
\textbf{EST333}\\ 
\multicolumn{3}{|p{12cm}|}{Nome do Componente Curricular em Inglês:} & \\ 
\multicolumn{3}{|p{12cm}|}{\textbf{Categorical Data Analysis}} & \\ 
\hline
\multicolumn{3}{|p{12cm}|}{Nome e Sigla do Departamento} & Unidade Acadêmica: \\ 
\multicolumn{3}{|p{12cm}|}{Departamento de Estatística (DEEST)} & {ICEB} \\ 
\hline
\multicolumn{4}{|p{16cm}|}{Modalidade de Oferta:
[ ] presencial \hspace{0.6cm}
[X] semipresencial \hspace{0.6cm}
[ ] à distância}\\
\hline
\multicolumn{2}{|p{8cm}|}{Carga horária semestral} &
\multicolumn{2}{p{8cm}|}{Carga horária semanal}\\
\hline
\multicolumn{1}{|p{4cm}|}{Total} &
\multicolumn{1}{p{4cm}|}{Extensionista} &
\multicolumn{1}{p{4cm}|}{Teórica} &
\multicolumn{1}{p{4cm}|}{Prática} \\ 
\multicolumn{1}{|p{4cm}|}{60\,horas} &
\multicolumn{1}{p{4cm}|}{0\;horas} &
\multicolumn{1}{p{4cm}|}{4\;horas/aula} &
\multicolumn{1}{p{4cm}|}{0\;horas/aula} \\ 
\hline
\multicolumn{4}{|p{16cm}|}{Ementa:}\\
\multicolumn{4}{|p{16cm}|}{}\\
\multicolumn{4}{|p{\dimexpr 16cm + 6\tabcolsep\relax}|}{Inferência para as distribuições binomial e multinomial. Testes de Pearson e da razão de verossimilhança generalizada. Testes de adequabilidade de distribuições. Tabelas de contingência. Diferença de proporções. Risco relativo e razão de chances. Testes de independência e homogeneidade. O teste exato de Fisher. Associação parcial em tabelas estratificadas 2 x 2. O teste de Cochran-Mantel-Haenszel. Modelos lineares generalizados. Regressão de Poisson. Regressão logística. Modelos log-lineares.}\\
\multicolumn{4}{|p{16cm}|}{}\\
\hline
\multicolumn{4}{|p{16cm}|}{Conteúdo programático:}\\*
\multicolumn{4}{|p{\dimexpr 16cm + 6\tabcolsep\relax}|}{\begin{enumerate}[start=1]\item Unidade 1: Introdução à análise de dados categóricos: dados com respostas categóricas; distribuições de probabilidade para dados categóricos; inferência estatística para uma proporção; outros pontos em inferência estatística para dados discretos\end{enumerate}}\\
\multicolumn{4}{|p{\dimexpr 16cm + 6\tabcolsep\relax}|}{\begin{enumerate}[start=2]\item Unidade 2: Tabelas de contingência: estrutura de probabilidade para tabelas de contingência; comparando proporções em tabelas 2 x 2; razão de chances; testes qui-quadrado para independência; testando independência para dados ordinais; inferência exata para pequenas amostras; associação em tabelas three way\end{enumerate}}\\
\multicolumn{4}{|p{\dimexpr 16cm + 6\tabcolsep\relax}|}{\begin{enumerate}[start=3]\item Unidade 3: Modelos lineares generalizados: componentes de um modelo linear generalizado; modelos lineares generalizados para dados binários; modelos lineares generalizados para contagens\end{enumerate}}\\
\multicolumn{4}{|p{\dimexpr 16cm + 6\tabcolsep\relax}|}{\begin{enumerate}[start=4]\item Unidade 4: Regressão logística: regressão logística com preditores categóricos; regressão logística múltipla\end{enumerate}}\\
\multicolumn{4}{|p{\dimexpr 16cm + 6\tabcolsep\relax}|}{\begin{enumerate}[start=5]\item Unidade 5: Construindo e aplicando modelos de regressão logística: estratégias em seleção de modelos; verificação de modelos; efeitos de dados esparsos; regressão logística condicional e inferência exata; tamanho de amostra e poder para a regressão logística\end{enumerate}}\\
\multicolumn{4}{|p{\dimexpr 16cm + 6\tabcolsep\relax}|}{\begin{enumerate}[start=6]\item Unidade 6: Modelos log-lineares para tabelas de contingência: modelos log-lineares para tabelas two-way e three way; inferência para modelos log-lineares\end{enumerate}}\\
\multicolumn{4}{|p{\dimexpr 16cm + 6\tabcolsep\relax}|}{\begin{enumerate}[start=7]\item Os conteúdos serão trabalhados de forma semipresencial, sendo: 53 horas presenciais e 07 horas à distância\end{enumerate}}\\
\multicolumn{4}{|p{\dimexpr 16cm + 6\tabcolsep\relax}|}{}\\
\hline
\multicolumn{4}{|p{16cm}|}{Bibliografia Básica:}\\*
\multicolumn{4}{|p{\dimexpr 16cm + 6\tabcolsep\relax}|}{\begin{itemize}\item AGRESTI, Alan. An introduction to categorical data analysis. 2. ed. New Jersey: John Wiley \& Sons, 2007.\end{itemize}}\\
\multicolumn{4}{|p{\dimexpr 16cm + 6\tabcolsep\relax}|}{\begin{itemize}\item PAULINO, Carlos D.; SINGER, Júlio da M. Análise de Dados Categorizados.\end{itemize}}\\
\multicolumn{4}{|p{16cm}|}{}\\
\hline
\multicolumn{4}{|p{16cm}|}{Bibliografia Complementar:}\\*
\multicolumn{4}{|p{\dimexpr 16cm + 6\tabcolsep\relax}|}{\begin{itemize}\item HOSMER JR, D.W.; LEMESHOW, Stanley. Applied Logistic Regression. 2. ed. New York: John Wiley \& Sons, 2000.\end{itemize}}\\
\multicolumn{4}{|p{\dimexpr 16cm + 6\tabcolsep\relax}|}{\begin{itemize}\item AGRESTI, Alan. Categorical data Analysis. New York: John Wiley \& Sons, 2013.\end{itemize}}\\
\multicolumn{4}{|p{\dimexpr 16cm + 6\tabcolsep\relax}|}{\begin{itemize}\item DOBSON, Annette J.; BARNETT, Adrian G. An introduction to generalized linear models. London. 3. ed. Chapman and Hall/CRC, 2008.\end{itemize}}\\
\multicolumn{4}{|p{\dimexpr 16cm + 6\tabcolsep\relax}|}{\begin{itemize}\item CHRISTENSEN, R. Log-Linear Models and Logistic Regression. New York: Springer-Verlag, 1997.\end{itemize}}\\
\multicolumn{4}{|p{\dimexpr 16cm + 6\tabcolsep\relax}|}{\begin{itemize}\item BILDER, Christopher R.; LOUGHIN, Thomas M. Analysis of categorical data with R. Boca Raton: CRC Press, 2015. (Texts in Statistical Science). ISBN 9781439855676.\end{itemize}}\\
\multicolumn{4}{|p{\dimexpr 16cm + 6\tabcolsep\relax}|}{\begin{itemize}\item MYERS, Raymond H. et al. Generalized linear models: with applications in engineering and the sciences. 2. ed. New Jersey: Wiley, 2010. (Wiley Series in Probability and Statistics). ISBN 9780470454633.\end{itemize}}\\
\hline
\end{longtable}
\end{center}

\clearpage
\begin{center}
\begin{tabular}{lccr}
 \multirow{3}{*}{\includegraphics[height=2.7cm]{brasao.png}} &
 \multicolumn{2}{c}{\bfseries UNIVERSIDADE FEDERAL DE OURO PRETO} &
 \ \ \ \ \multirow{3}{*}{\includegraphics[height=2.7cm]{ufop.png}} \\
 & \multicolumn{2}{c}{\bfseries PRÓ-REITORIA DE GRADUAÇÃO} & \\
 & \multicolumn{2}{c}{} & \\
 & \multicolumn{2}{c}{\Large\bfseries PROGRAMA DE DISCIPLINA} & \\
\end{tabular}
\end{center}

\begin{center}
\begin{longtable}{|p{4cm}|p{4cm}|p{4cm}|p{4cm}|}
\hline
\multicolumn{3}{|p{12cm}|}{Nome do Componente Curricular em Português:} &
\multicolumn{1}{p{4cm}|}{Código:} \\ 
\multicolumn{3}{|p{12cm}|}{\textbf{Física Conceitual}} &
\textbf{FIS118}\\ 
\multicolumn{3}{|p{12cm}|}{Nome do Componente Curricular em Inglês:} & \\ 
\multicolumn{3}{|p{12cm}|}{\textbf{Conceptual Physics}} & \\ 
\hline
\multicolumn{3}{|p{12cm}|}{Nome e Sigla do Departamento} & Unidade Acadêmica: \\ 
\multicolumn{3}{|p{12cm}|}{Departamento de Física (DEFIS)} & {ICEB} \\ 
\hline
\multicolumn{4}{|p{16cm}|}{Modalidade de Oferta:
[X] presencial \hspace{0.6cm}
[ ] semipresencial \hspace{0.6cm}
[ ] à distância}\\
\hline
\multicolumn{2}{|p{8cm}|}{Carga horária semestral} &
\multicolumn{2}{p{8cm}|}{Carga horária semanal}\\
\hline
\multicolumn{1}{|p{4cm}|}{Total} &
\multicolumn{1}{p{4cm}|}{Extensionista} &
\multicolumn{1}{p{4cm}|}{Teórica} &
\multicolumn{1}{p{4cm}|}{Prática} \\ 
\multicolumn{1}{|p{4cm}|}{60\,horas} &
\multicolumn{1}{p{4cm}|}{0\;horas} &
\multicolumn{1}{p{4cm}|}{2\;horas/aula} &
\multicolumn{1}{p{4cm}|}{2\;horas/aula} \\ 
\hline
\multicolumn{4}{|p{16cm}|}{Ementa:}\\
\multicolumn{4}{|p{16cm}|}{}\\
\multicolumn{4}{|p{\dimexpr 16cm + 6\tabcolsep\relax}|}{Abordagem qualitativa dos fenômenos Físicos: sistemas referenciais e movimento relativo, movimento de rotação, hidráulica, gravitação, transferência de calor, ondas, indução eletromagnética, fenômenos ópticos, aplicações e relações desses conceitos físicos com outras ciências.}\\
\multicolumn{4}{|p{16cm}|}{}\\
\hline
\multicolumn{4}{|p{16cm}|}{Conteúdo programático:}\\*
\multicolumn{4}{|p{\dimexpr 16cm + 6\tabcolsep\relax}|}{\begin{enumerate}[start=1]\item Referenciais inerciais e não inerciais; Movimento relativo\end{enumerate}}\\
\multicolumn{4}{|p{\dimexpr 16cm + 6\tabcolsep\relax}|}{\begin{enumerate}[start=2]\item Transmissão de movimento (polias, engrenagens, discos)\end{enumerate}}\\
\multicolumn{4}{|p{\dimexpr 16cm + 6\tabcolsep\relax}|}{\begin{enumerate}[start=3]\item Hidráulica\end{enumerate}}\\
\multicolumn{4}{|p{\dimexpr 16cm + 6\tabcolsep\relax}|}{\begin{enumerate}[start=4]\item Lei da gravitação, marés, campos gravitacionais\end{enumerate}}\\
\multicolumn{4}{|p{\dimexpr 16cm + 6\tabcolsep\relax}|}{\begin{enumerate}[start=5]\item Transferência de calor\end{enumerate}}\\
\multicolumn{4}{|p{\dimexpr 16cm + 6\tabcolsep\relax}|}{\begin{enumerate}[start=6]\item Ondas e som\end{enumerate}}\\
\multicolumn{4}{|p{\dimexpr 16cm + 6\tabcolsep\relax}|}{\begin{enumerate}[start=7]\item Cor, reflexo e refração\end{enumerate}}\\
\multicolumn{4}{|p{\dimexpr 16cm + 6\tabcolsep\relax}|}{\begin{enumerate}[start=8]\item Lei de Faraday, lei de Lenz, geradores, transformadores, auto-indução\end{enumerate}}\\
\multicolumn{4}{|p{\dimexpr 16cm + 6\tabcolsep\relax}|}{\begin{enumerate}[start=9]\item Relações da Física com a Química, Biologia, Geologia\end{enumerate}}\\
\multicolumn{4}{|p{\dimexpr 16cm + 6\tabcolsep\relax}|}{}\\
\hline
\multicolumn{4}{|p{16cm}|}{Bibliografia Básica:}\\*
\multicolumn{4}{|p{\dimexpr 16cm + 6\tabcolsep\relax}|}{\begin{itemize}\item HEWITT, Paul G. Física Conceitual. 9. ed. Bookman.\end{itemize}}\\
\multicolumn{4}{|p{\dimexpr 16cm + 6\tabcolsep\relax}|}{\begin{itemize}\item FEYNMAN, R. P. Lições de Física de Feynman. Bookman, 2008.\end{itemize}}\\
\multicolumn{4}{|p{\dimexpr 16cm + 6\tabcolsep\relax}|}{\begin{itemize}\item FREIRE, Paulo. Extensão ou Comunicação. Paz e Terra, 2010.\end{itemize}}\\
\multicolumn{4}{|p{16cm}|}{}\\
\hline
\multicolumn{4}{|p{16cm}|}{Bibliografia Complementar:}\\*
\multicolumn{4}{|p{\dimexpr 16cm + 6\tabcolsep\relax}|}{\begin{itemize}\item Introductory Physics. New Jersey: Prentice Hall, Upper Saddle River, 2002.\end{itemize}}\\
\multicolumn{4}{|p{\dimexpr 16cm + 6\tabcolsep\relax}|}{\begin{itemize}\item COHEN, Ernest. Avaliação de Projetos Sociais. São Paulo: Vozes, 2013.\end{itemize}}\\
\multicolumn{4}{|p{\dimexpr 16cm + 6\tabcolsep\relax}|}{\begin{itemize}\item SOUZA NETO, João Clemente de. Extensão Universitária: Construção de Solidariedade. São Paulo: Expressão \& Arte, 2005.\end{itemize}}\\
\hline
\end{longtable}
\end{center}

\clearpage
\begin{center}
\begin{tabular}{lccr}
 \multirow{3}{*}{\includegraphics[height=2.7cm]{brasao.png}} &
 \multicolumn{2}{c}{\bfseries UNIVERSIDADE FEDERAL DE OURO PRETO} &
 \ \ \ \ \multirow{3}{*}{\includegraphics[height=2.7cm]{ufop.png}} \\
 & \multicolumn{2}{c}{\bfseries PRÓ-REITORIA DE GRADUAÇÃO} & \\
 & \multicolumn{2}{c}{} & \\
 & \multicolumn{2}{c}{\Large\bfseries PROGRAMA DE DISCIPLINA} & \\
\end{tabular}
\end{center}

\begin{center}
\begin{longtable}{|p{4cm}|p{4cm}|p{4cm}|p{4cm}|}
\hline
\multicolumn{3}{|p{12cm}|}{Nome do Componente Curricular em Português:} &
\multicolumn{1}{p{4cm}|}{Código:} \\ 
\multicolumn{3}{|p{12cm}|}{\textbf{Introdução aos Métodos da Física}} &
\textbf{FIS123}\\ 
\multicolumn{3}{|p{12cm}|}{Nome do Componente Curricular em Inglês:} & \\ 
\multicolumn{3}{|p{12cm}|}{\textbf{Introduction to Physics Methods}} & \\ 
\hline
\multicolumn{3}{|p{12cm}|}{Nome e Sigla do Departamento} & Unidade Acadêmica: \\ 
\multicolumn{3}{|p{12cm}|}{Departamento de Física (DEFIS)} & {ICEB} \\ 
\hline
\multicolumn{4}{|p{16cm}|}{Modalidade de Oferta:
[X] presencial \hspace{0.6cm}
[ ] semipresencial \hspace{0.6cm}
[ ] à distância}\\
\hline
\multicolumn{2}{|p{8cm}|}{Carga horária semestral} &
\multicolumn{2}{p{8cm}|}{Carga horária semanal}\\
\hline
\multicolumn{1}{|p{4cm}|}{Total} &
\multicolumn{1}{p{4cm}|}{Extensionista} &
\multicolumn{1}{p{4cm}|}{Teórica} &
\multicolumn{1}{p{4cm}|}{Prática} \\ 
\multicolumn{1}{|p{4cm}|}{60\,horas} &
\multicolumn{1}{p{4cm}|}{0\;horas} &
\multicolumn{1}{p{4cm}|}{4\;horas/aula} &
\multicolumn{1}{p{4cm}|}{0\;horas/aula} \\ 
\hline
\multicolumn{4}{|p{16cm}|}{Ementa:}\\
\multicolumn{4}{|p{16cm}|}{}\\
\multicolumn{4}{|p{\dimexpr 16cm + 6\tabcolsep\relax}|}{Conceitos Básicos de Metrologia; Representação gráfica das grandezas físicas; Grandezas físicas escalares e vetoriais. Relações entre grandezas: derivadas e integrais.}\\
\multicolumn{4}{|p{16cm}|}{}\\
\hline
\multicolumn{4}{|p{16cm}|}{Conteúdo programático:}\\*
\multicolumn{4}{|p{\dimexpr 16cm + 6\tabcolsep\relax}|}{\begin{enumerate}[start=1]\item Conceitos básicos de metrologia: Grandezas e unidades; Unidades de base e derivadas; Múltiplos e submúltiplos de unidades; Conversão de unidades; Análise dimensional; Medições: precisão e exatidão; Erro e incerteza; Algarismos significativos\end{enumerate}}\\
\multicolumn{4}{|p{\dimexpr 16cm + 6\tabcolsep\relax}|}{\begin{enumerate}[start=2]\item Representação gráfica de grandezas físicas\end{enumerate}}\\
\multicolumn{4}{|p{\dimexpr 16cm + 6\tabcolsep\relax}|}{\begin{enumerate}[start=3]\item Grandezas escalares e vetoriais: Operações com grandezas vetoriais; Decomposição de vetores; Produtos escalar e vetorial; Exemplos de grandezas vetoriais: posição, deslocamento, trabalho e torque\end{enumerate}}\\
\multicolumn{4}{|p{\dimexpr 16cm + 6\tabcolsep\relax}|}{\begin{enumerate}[start=4]\item Relações entre grandezas: Derivadas e integrais em 1D; Representação gráfica da derivada: velocidade e deslocamento; Representação gráfica da integração: deslocamento; Trabalho como função do produto força e deslocamento\end{enumerate}}\\
\multicolumn{4}{|p{\dimexpr 16cm + 6\tabcolsep\relax}|}{}\\
\hline
\multicolumn{4}{|p{16cm}|}{Bibliografia Básica:}\\*
\multicolumn{4}{|p{\dimexpr 16cm + 6\tabcolsep\relax}|}{\begin{itemize}\item BATISTA, E.; MOZOLEVSKI, I. Métodos de Física Matemática - V.1. Florianópolis: UFSC/EAD/CED/CFM, 2010.\end{itemize}}\\
\multicolumn{4}{|p{\dimexpr 16cm + 6\tabcolsep\relax}|}{\begin{itemize}\item LIMA JUNIOR, P. et al. O laboratório de mecânica: Subsídios para o ensino de Física Experimental. Porto Alegre: UFRGS, Instituto de Física, 2013.\end{itemize}}\\
\multicolumn{4}{|p{\dimexpr 16cm + 6\tabcolsep\relax}|}{\begin{itemize}\item Inmetro. Vocabulário Internacional de Metrologia: Conceitos fundamentais e gerais e termos associados (VIM 2012). Duque de Caxias, RJ: INMETRO, 2012.\end{itemize}}\\
\multicolumn{4}{|p{16cm}|}{}\\
\hline
\multicolumn{4}{|p{16cm}|}{Bibliografia Complementar:}\\*
\multicolumn{4}{|p{\dimexpr 16cm + 6\tabcolsep\relax}|}{\begin{itemize}\item Inmetro. Avaliação de dados de medição: uma introdução ao "Guia para a expressão de incerteza de medição" e a documentos correlatos – INTROGUM 2009. Duque de Caxias, RJ: INMETRO/CICMA/SEPIN.\end{itemize}}\\
\multicolumn{4}{|p{\dimexpr 16cm + 6\tabcolsep\relax}|}{\begin{itemize}\item Inmetro. Sistema Internacional de Unidades: SI. Duque de Caxias, RJ: INMETRO/CICMA/SEPIN, 2012.\end{itemize}}\\
\hline
\end{longtable}
\end{center}

\clearpage
\begin{center}
\begin{tabular}{lccr}
 \multirow{3}{*}{\includegraphics[height=2.7cm]{brasao.png}} &
 \multicolumn{2}{c}{\bfseries UNIVERSIDADE FEDERAL DE OURO PRETO} &
 \ \ \ \ \multirow{3}{*}{\includegraphics[height=2.7cm]{ufop.png}} \\
 & \multicolumn{2}{c}{\bfseries PRÓ-REITORIA DE GRADUAÇÃO} & \\
 & \multicolumn{2}{c}{} & \\
 & \multicolumn{2}{c}{\Large\bfseries PROGRAMA DE DISCIPLINA} & \\
\end{tabular}
\end{center}

\begin{center}
\begin{longtable}{|p{4cm}|p{4cm}|p{4cm}|p{4cm}|}
\hline
\multicolumn{3}{|p{12cm}|}{Nome do Componente Curricular em Português:} &
\multicolumn{1}{p{4cm}|}{Código:} \\ 
\multicolumn{3}{|p{12cm}|}{\textbf{Física Computacional}} &
\textbf{FIS331}\\ 
\multicolumn{3}{|p{12cm}|}{Nome do Componente Curricular em Inglês:} & \\ 
\multicolumn{3}{|p{12cm}|}{\textbf{Computational Physics}} & \\ 
\hline
\multicolumn{3}{|p{12cm}|}{Nome e Sigla do Departamento} & Unidade Acadêmica: \\ 
\multicolumn{3}{|p{12cm}|}{Departamento de Física (DEFIS)} & {ICEB} \\ 
\hline
\multicolumn{4}{|p{16cm}|}{Modalidade de Oferta:
[X] presencial \hspace{0.6cm}
[ ] semipresencial \hspace{0.6cm}
[ ] à distância}\\
\hline
\multicolumn{2}{|p{8cm}|}{Carga horária semestral} &
\multicolumn{2}{p{8cm}|}{Carga horária semanal}\\
\hline
\multicolumn{1}{|p{4cm}|}{Total} &
\multicolumn{1}{p{4cm}|}{Extensionista} &
\multicolumn{1}{p{4cm}|}{Teórica} &
\multicolumn{1}{p{4cm}|}{Prática} \\ 
\multicolumn{1}{|p{4cm}|}{60\,horas} &
\multicolumn{1}{p{4cm}|}{0\;horas} &
\multicolumn{1}{p{4cm}|}{4\;horas/aula} &
\multicolumn{1}{p{4cm}|}{0\;horas/aula} \\ 
\hline
\multicolumn{4}{|p{16cm}|}{Ementa:}\\
\multicolumn{4}{|p{16cm}|}{}\\
\multicolumn{4}{|p{\dimexpr 16cm + 6\tabcolsep\relax}|}{Análise de Experimentos e Ajuste de Curvas com Aplicações, Integração Numérica e Aplicações, Solução Numérica de Equações Ordinárias e Aplicações em problemas de Física, Lançamento de Projéteis, Oscilações e Leis de Kepler, Problemas de Autovalor e Fronteira.}\\
\multicolumn{4}{|p{16cm}|}{}\\
\hline
\multicolumn{4}{|p{16cm}|}{Conteúdo programático:}\\*
\multicolumn{4}{|p{\dimexpr 16cm + 6\tabcolsep\relax}|}{\begin{enumerate}[start=1]\item Modelos preditivos com a resolução de equações, como queda livre ou lançamento de projéteis sem atrito: Revisar estruturas de programação, inclusive com a introdução de vetores; Construção de gráficos para apresentação dos resultados; Comparar resultados com a previsão teórica para poder validar o código/algoritmo\end{enumerate}}\\
\multicolumn{4}{|p{\dimexpr 16cm + 6\tabcolsep\relax}|}{\begin{enumerate}[start=2]\item Introdução ao LaTeX como ferramenta para elaboração de relatórios e documentos\end{enumerate}}\\
\multicolumn{4}{|p{\dimexpr 16cm + 6\tabcolsep\relax}|}{\begin{enumerate}[start=3]\item Desenvolvimento de métodos para solução de problemas físicos por meio de: Achar raízes de funções e comparar eficiência e aplicabilidade dos métodos, inclusive de equações transcendentais; Ajuste linear, linearização, extrapolação; Ajuste polinomial; Solução de sistemas lineares; Interpolação por polinômios e extrapolação; Integração numérica: Aproximação por polinômios de grau 0, 1 e 2, ou seja, ponto central, Simpson 1/3 e 3/8; Quadratura de Gauss; Diferenciação numérica e aplicações; Resolução de equações diferenciais ordinárias de 1a e 2a ordem: Euler, Verlet, RK2, RK4\end{enumerate}}\\
\multicolumn{4}{|p{\dimexpr 16cm + 6\tabcolsep\relax}|}{\begin{enumerate}[start=4]\item Introdução a Monte Carlo: Cálculo de constantes por acertos e erros (Pi); Integração (erros e acertos e amostragem por importância); Estimativa de incerteza\end{enumerate}}\\
\multicolumn{4}{|p{\dimexpr 16cm + 6\tabcolsep\relax}|}{\begin{enumerate}[start=5]\item Caminhada aleatória e a distribuição binomial: a equação da difusão\end{enumerate}}\\
\multicolumn{4}{|p{\dimexpr 16cm + 6\tabcolsep\relax}|}{}\\
\hline
\multicolumn{4}{|p{16cm}|}{Bibliografia Básica:}\\*
\multicolumn{4}{|p{\dimexpr 16cm + 6\tabcolsep\relax}|}{\begin{itemize}\item Amos Gilat, Vish Subramaniam; tradução Alberto Resende de Conti. Métodos numéricos para engenheiros e cientistas: uma introdução com aplicações usando o MATLAB. Porto Alegre, RS: Bookman, 2008.\end{itemize}}\\
\multicolumn{4}{|p{\dimexpr 16cm + 6\tabcolsep\relax}|}{\begin{itemize}\item Scherer, Claudio. Métodos computacionais da física. São Paulo: Livraria da Física, 2010.\end{itemize}}\\
\multicolumn{4}{|p{\dimexpr 16cm + 6\tabcolsep\relax}|}{\begin{itemize}\item Klein, Andi. Introductory computational physics. Cambridge, UK.\end{itemize}}\\
\multicolumn{4}{|p{\dimexpr 16cm + 6\tabcolsep\relax}|}{\begin{itemize}\item The burden and the glory. New York: Harper \& Row, 1964.\end{itemize}}\\
\multicolumn{4}{|p{\dimexpr 16cm + 6\tabcolsep\relax}|}{\begin{itemize}\item BINDER, K.; HEERMAN, Dieter W. Simulation in statistical physics. 3. ed. 1997.\end{itemize}}\\
\multicolumn{4}{|p{16cm}|}{}\\
\hline
\multicolumn{4}{|p{16cm}|}{Bibliografia Complementar:}\\*
\multicolumn{4}{|p{\dimexpr 16cm + 6\tabcolsep\relax}|}{\begin{itemize}\item Burden, Richard L. Análise numérica. 10. ed. São Paulo: Cengage Learning, 2016.\end{itemize}}\\
\multicolumn{4}{|p{\dimexpr 16cm + 6\tabcolsep\relax}|}{\begin{itemize}\item Décio Sperandio, João Teixeira Mendes, Luiz Henry Monken e Silva. Cálculo numérico: características matemáticas e computacionais dos métodos numéricos. São Paulo: Prentice-Hall, 2003.\end{itemize}}\\
\multicolumn{4}{|p{\dimexpr 16cm + 6\tabcolsep\relax}|}{\begin{itemize}\item Koyre, Alexandre. Estudos de história do pensamento científico. 2. rev. ed. Rio de Janeiro: Forense Universitária, 1991.\end{itemize}}\\
\multicolumn{4}{|p{\dimexpr 16cm + 6\tabcolsep\relax}|}{\begin{itemize}\item Barbosa, Leonardo da Conceição. Métodos numéricos para resolver equações diferenciais ordinárias. 2007.\end{itemize}}\\
\multicolumn{4}{|p{\dimexpr 16cm + 6\tabcolsep\relax}|}{\begin{itemize}\item Ross, Sheldon M. Simulation. 5. ed. Amsterdã: Elsevier, 2013.\end{itemize}}\\
\hline
\end{longtable}
\end{center}

\clearpage
\begin{center}
\begin{tabular}{lccr}
 \multirow{3}{*}{\includegraphics[height=2.7cm]{brasao.png}} &
 \multicolumn{2}{c}{\bfseries UNIVERSIDADE FEDERAL DE OURO PRETO} &
 \ \ \ \ \multirow{3}{*}{\includegraphics[height=2.7cm]{ufop.png}} \\
 & \multicolumn{2}{c}{\bfseries PRÓ-REITORIA DE GRADUAÇÃO} & \\
 & \multicolumn{2}{c}{} & \\
 & \multicolumn{2}{c}{\Large\bfseries PROGRAMA DE DISCIPLINA} & \\
\end{tabular}
\end{center}

\begin{center}
\begin{longtable}{|p{4cm}|p{4cm}|p{4cm}|p{4cm}|}
\hline
\multicolumn{3}{|p{12cm}|}{Nome do Componente Curricular em Português:} &
\multicolumn{1}{p{4cm}|}{Código:} \\ 
\multicolumn{3}{|p{12cm}|}{\textbf{Física e Ciência, Sociedade e Meio Ambiente}} &
\textbf{FIS600}\\ 
\multicolumn{3}{|p{12cm}|}{Nome do Componente Curricular em Inglês:} & \\ 
\multicolumn{3}{|p{12cm}|}{\textbf{Physics and Science, Society and the Environment}} & \\ 
\hline
\multicolumn{3}{|p{12cm}|}{Nome e Sigla do Departamento} & Unidade Acadêmica: \\ 
\multicolumn{3}{|p{12cm}|}{Departamento de Física (DEFIS)} & {ICEB} \\ 
\hline
\multicolumn{4}{|p{16cm}|}{Modalidade de Oferta:
[X] presencial \hspace{0.6cm}
[ ] semipresencial \hspace{0.6cm}
[ ] à distância}\\
\hline
\multicolumn{2}{|p{8cm}|}{Carga horária semestral} &
\multicolumn{2}{p{8cm}|}{Carga horária semanal}\\
\hline
\multicolumn{1}{|p{4cm}|}{Total} &
\multicolumn{1}{p{4cm}|}{Extensionista} &
\multicolumn{1}{p{4cm}|}{Teórica} &
\multicolumn{1}{p{4cm}|}{Prática} \\ 
\multicolumn{1}{|p{4cm}|}{30\,horas} &
\multicolumn{1}{p{4cm}|}{0\;horas} &
\multicolumn{1}{p{4cm}|}{2\;horas/aula} &
\multicolumn{1}{p{4cm}|}{0\;horas/aula} \\ 
\hline
\multicolumn{4}{|p{16cm}|}{Ementa:}\\
\multicolumn{4}{|p{16cm}|}{}\\
\multicolumn{4}{|p{\dimexpr 16cm + 6\tabcolsep\relax}|}{Atividades que abordem a relação do físico com o meio ambiente e a sociedade: temas sobre relações étnico-raciais; direitos humanos, questões de gênero, meio ambiente e sustentabilidade, trabalho e consumo, pluralidade cultural, inclusão da pessoa com deficiência.}\\
\multicolumn{4}{|p{16cm}|}{}\\
\hline
\multicolumn{4}{|p{16cm}|}{Conteúdo programático:}\\*
\multicolumn{4}{|p{\dimexpr 16cm + 6\tabcolsep\relax}|}{\begin{enumerate}[start=1]\item Atividades ministradas por professores da UFOP ou convidados externos abordando os seguintes tópicos: temas sobre relações étnico-raciais; direitos humanos; questões de gênero; meio ambiente e sustentabilidade; trabalho e consumo; pluralidade cultural; inclusão da pessoa com deficiência\end{enumerate}}\\
\multicolumn{4}{|p{\dimexpr 16cm + 6\tabcolsep\relax}|}{}\\
\hline
\multicolumn{4}{|p{16cm}|}{Bibliografia Básica:}\\*
\multicolumn{4}{|p{\dimexpr 16cm + 6\tabcolsep\relax}|}{\begin{itemize}\item Nilma Lino Gomes. Sem perder a raiz: corpo e cabelo como símbolos da identidade negra. Belo Horizonte, MG: Autêntica, 2006.\end{itemize}}\\
\multicolumn{4}{|p{\dimexpr 16cm + 6\tabcolsep\relax}|}{\begin{itemize}\item Mulheres na Física: casos históricos, panorama e perspectivas. São Paulo: Livraria da Física, 2015.\end{itemize}}\\
\multicolumn{4}{|p{\dimexpr 16cm + 6\tabcolsep\relax}|}{\begin{itemize}\item Robert E. Ricklefs; tradução Cecília Bueno, Pedro P. de Lima e Silva com a colaboração de Rogério Ribeiro de Oliveira. A economia da natureza: um livro-texto em ecologia básica. 3. ed. Rio de Janeiro: Guanabara Koogan, 1996.\end{itemize}}\\
\multicolumn{4}{|p{\dimexpr 16cm + 6\tabcolsep\relax}|}{\begin{itemize}\item Vernant, Jean Pierre. As origens do pensamento grego. 10. ed. São Paulo: DIFEL, 1998.\end{itemize}}\\
\multicolumn{4}{|p{16cm}|}{}\\
\hline
\multicolumn{4}{|p{16cm}|}{Bibliografia Complementar:}\\*
\multicolumn{4}{|p{\dimexpr 16cm + 6\tabcolsep\relax}|}{\begin{itemize}\item FREIRE, Paulo. Extensão ou Comunicação. São Paulo: Paz e Terra, 2010.\end{itemize}}\\
\multicolumn{4}{|p{\dimexpr 16cm + 6\tabcolsep\relax}|}{\begin{itemize}\item ARROYO, Miguel. Outras Pedagogias: Outros Sujeitos. 1. ed. Petrópolis, 2012.\end{itemize}}\\
\multicolumn{4}{|p{\dimexpr 16cm + 6\tabcolsep\relax}|}{\begin{itemize}\item ODUM. Fundamentos de ecologia. São Paulo: Thomson Learning, 2007.\end{itemize}}\\
\multicolumn{4}{|p{\dimexpr 16cm + 6\tabcolsep\relax}|}{\begin{itemize}\item SECRETARIA ESPECIAL DE POLÍTICAS PARA MULHERES; CENTRO LATINO-AMERICANO EM SEXUALIDADE E DIREITOS HUMANOS. Brasília: CepeSC/SPM, 2009.\end{itemize}}\\
\multicolumn{4}{|p{\dimexpr 16cm + 6\tabcolsep\relax}|}{\begin{itemize}\item COHEN, Ernesto. Avaliação de Projetos Sociais. São Paulo: Vozes, 2013.\end{itemize}}\\
\multicolumn{4}{|p{\dimexpr 16cm + 6\tabcolsep\relax}|}{\begin{itemize}\item SOUZA NETO. Construção de Solidariedade: Extensão Universitária. São Paulo: Expressão \& Arte, 2005.\end{itemize}}\\
\multicolumn{4}{|p{\dimexpr 16cm + 6\tabcolsep\relax}|}{\begin{itemize}\item Donna Haraway; Tomaz Tadeu da Silva (org. e trad.). Antropologia do ciborgue: as vertigens do pós-humano. 2. ed. Belo Horizonte, MG: Autêntica, 2009.\end{itemize}}\\
\hline
\end{longtable}
\end{center}

\clearpage
\begin{center}
\begin{tabular}{lccr}
 \multirow{3}{*}{\includegraphics[height=2.7cm]{brasao.png}} &
 \multicolumn{2}{c}{\bfseries UNIVERSIDADE FEDERAL DE OURO PRETO} &
 \ \ \ \ \multirow{3}{*}{\includegraphics[height=2.7cm]{ufop.png}} \\
 & \multicolumn{2}{c}{\bfseries PRÓ-REITORIA DE GRADUAÇÃO} & \\
 & \multicolumn{2}{c}{} & \\
 & \multicolumn{2}{c}{\Large\bfseries PROGRAMA DE DISCIPLINA} & \\
\end{tabular}
\end{center}

\begin{center}
\begin{longtable}{|p{4cm}|p{4cm}|p{4cm}|p{4cm}|}
\hline
\multicolumn{3}{|p{12cm}|}{Nome do Componente Curricular em Português:} &
\multicolumn{1}{p{4cm}|}{Código:} \\ 
\multicolumn{3}{|p{12cm}|}{\textbf{Física Geral}} &
\textbf{FIS701}\\ 
\multicolumn{3}{|p{12cm}|}{Nome do Componente Curricular em Inglês:} & \\ 
\multicolumn{3}{|p{12cm}|}{\textbf{General Physics}} & \\ 
\hline
\multicolumn{3}{|p{12cm}|}{Nome e Sigla do Departamento} & Unidade Acadêmica: \\ 
\multicolumn{3}{|p{12cm}|}{Departamento de Física (DEFIS)} & {ICEB} \\ 
\hline
\multicolumn{4}{|p{16cm}|}{Modalidade de Oferta:
[X] presencial \hspace{0.6cm}
[ ] semipresencial \hspace{0.6cm}
[ ] à distância}\\
\hline
\multicolumn{2}{|p{8cm}|}{Carga horária semestral} &
\multicolumn{2}{p{8cm}|}{Carga horária semanal}\\
\hline
\multicolumn{1}{|p{4cm}|}{Total} &
\multicolumn{1}{p{4cm}|}{Extensionista} &
\multicolumn{1}{p{4cm}|}{Teórica} &
\multicolumn{1}{p{4cm}|}{Prática} \\ 
\multicolumn{1}{|p{4cm}|}{75\,horas} &
\multicolumn{1}{p{4cm}|}{0\;horas} &
\multicolumn{1}{p{4cm}|}{3\;horas/aula} &
\multicolumn{1}{p{4cm}|}{2\;horas/aula} \\ 
\hline
\multicolumn{4}{|p{16cm}|}{Ementa:}\\
\multicolumn{4}{|p{16cm}|}{}\\
\multicolumn{4}{|p{\dimexpr 16cm + 6\tabcolsep\relax}|}{}\\
\multicolumn{4}{|p{16cm}|}{}\\
\hline
\multicolumn{4}{|p{16cm}|}{Conteúdo programático:}\\*
\multicolumn{4}{|p{\dimexpr 16cm + 6\tabcolsep\relax}|}{\begin{itemize}\item[] (nenhum item informado)\end{itemize}}\\
\multicolumn{4}{|p{\dimexpr 16cm + 6\tabcolsep\relax}|}{}\\
\hline
\multicolumn{4}{|p{16cm}|}{Bibliografia Básica:}\\*
\multicolumn{4}{|p{\dimexpr 16cm + 6\tabcolsep\relax}|}{\begin{itemize}\item David Halliday, Robert Resnick, Jearl Walker; tradução Flávio Menezes de Aguiar, José Wellington Rocha Tabosa. Fundamentos de física 1: mecânica. 7. ed. Rio de Janeiro: LTC, 2006.\end{itemize}}\\
\multicolumn{4}{|p{\dimexpr 16cm + 6\tabcolsep\relax}|}{\begin{itemize}\item David Halliday, Robert Resnick, Jearl Walker; tradução e revisão técnica Ronaldo Sérgio de Biasi. Fundamentos de física 3: eletromagnetismo. 8. ed. Rio de Janeiro: LTC, 2009.\end{itemize}}\\
\multicolumn{4}{|p{\dimexpr 16cm + 6\tabcolsep\relax}|}{\begin{itemize}\item David Halliday, Robert Resnick, Jearl Walker; tradução e revisão técnica Ronaldo Sérgio de Biasi. Fundamentos de física 2: gravitação, ondas e termodinâmica. 8. ed. Rio de Janeiro: LTC, 2009.\end{itemize}}\\
\multicolumn{4}{|p{\dimexpr 16cm + 6\tabcolsep\relax}|}{\begin{itemize}\item David Halliday, Robert Resnick, Jearl Walker; tradução e revisão técnica. Fundamentos de física 1: mecânica. 8. ed. Rio de Janeiro: LTC, 2009.\end{itemize}}\\
\multicolumn{4}{|p{\dimexpr 16cm + 6\tabcolsep\relax}|}{\begin{itemize}\item David Halliday, Robert Resnick, Jearl Walker; tradução e revisão técnica Ronaldo Sérgio de Biasi. Fundamentos de física 2: gravitação, ondas e termodinâmica. 9. ed. Rio de Janeiro: LTC, 2012.\end{itemize}}\\
\multicolumn{4}{|p{\dimexpr 16cm + 6\tabcolsep\relax}|}{\begin{itemize}\item David Halliday, Robert Resnick, Jearl Walker; tradução e revisão técnica Ronaldo Sérgio de Biasi. Fundamentos de física 1: mecânica. 9. ed. Rio de Janeiro: LTC, 2012.\end{itemize}}\\
\multicolumn{4}{|p{\dimexpr 16cm + 6\tabcolsep\relax}|}{\begin{itemize}\item David Halliday, Robert Resnick, Jearl Walker; tradução e revisão técnica Ronaldo Sérgio de Biasi. Fundamentos de física 3: eletromagnetismo. 9. ed. Rio de Janeiro: LTC, 2012.\end{itemize}}\\
\multicolumn{4}{|p{\dimexpr 16cm + 6\tabcolsep\relax}|}{\begin{itemize}\item David Halliday, Robert Resnick, Jearl Walker; tradução Ronaldo Sérgio de Biasi. Fundamentos de física 2: gravitação, ondas e termodinâmica. 10. ed. Rio de Janeiro: Gen.\end{itemize}}\\
\multicolumn{4}{|p{\dimexpr 16cm + 6\tabcolsep\relax}|}{\begin{itemize}\item David Halliday, Robert Resnick; com a assistência de John Merrill; tradução Adir Moysés Luiz... [et al.]. Fundamentos de física 3: eletromagnetismo. 3. ed. Livros Técnicos e Científicos, 1994.\end{itemize}}\\
\multicolumn{4}{|p{16cm}|}{}\\
\hline
\multicolumn{4}{|p{16cm}|}{Bibliografia Complementar:}\\*
\multicolumn{4}{|p{\dimexpr 16cm + 6\tabcolsep\relax}|}{\begin{itemize}\item Okuno, Emico. Física para ciências biológicas e biomédicas. São Paulo: HARBRA, 1986.\end{itemize}}\\
\multicolumn{4}{|p{\dimexpr 16cm + 6\tabcolsep\relax}|}{\begin{itemize}\item Paul A. Tipler, Gene Mosca; tradução: Fernando Ribeiro da Silva, Gisele Maria Ribeiro Vieira. Física para cientistas e engenheiros: volume 1 mecânica, oscilações e ondas, termodinâmica. 5. ed. Rio de Janeiro: LTC, 2006.\end{itemize}}\\
\hline
\end{longtable}
\end{center}

\clearpage
\begin{center}
\begin{tabular}{lccr}
 \multirow{3}{*}{\includegraphics[height=2.7cm]{brasao.png}} &
 \multicolumn{2}{c}{\bfseries UNIVERSIDADE FEDERAL DE OURO PRETO} &
 \ \ \ \ \multirow{3}{*}{\includegraphics[height=2.7cm]{ufop.png}} \\
 & \multicolumn{2}{c}{\bfseries PRÓ-REITORIA DE GRADUAÇÃO} & \\
 & \multicolumn{2}{c}{} & \\
 & \multicolumn{2}{c}{\Large\bfseries PROGRAMA DE DISCIPLINA} & \\
\end{tabular}
\end{center}

\begin{center}
\begin{longtable}{|p{4cm}|p{4cm}|p{4cm}|p{4cm}|}
\hline
\multicolumn{3}{|p{12cm}|}{Nome do Componente Curricular em Português:} &
\multicolumn{1}{p{4cm}|}{Código:} \\ 
\multicolumn{3}{|p{12cm}|}{\textbf{Introdução à Informação Quântica}} &
\textbf{FIS827}\\ 
\multicolumn{3}{|p{12cm}|}{Nome do Componente Curricular em Inglês:} & \\ 
\multicolumn{3}{|p{12cm}|}{\textbf{Introduction to Quantum Information}} & \\ 
\hline
\multicolumn{3}{|p{12cm}|}{Nome e Sigla do Departamento} & Unidade Acadêmica: \\ 
\multicolumn{3}{|p{12cm}|}{Departamento de Física (DEFIS)} & {ICEB} \\ 
\hline
\multicolumn{4}{|p{16cm}|}{Modalidade de Oferta:
[X] presencial \hspace{0.6cm}
[ ] semipresencial \hspace{0.6cm}
[ ] à distância}\\
\hline
\multicolumn{2}{|p{8cm}|}{Carga horária semestral} &
\multicolumn{2}{p{8cm}|}{Carga horária semanal}\\
\hline
\multicolumn{1}{|p{4cm}|}{Total} &
\multicolumn{1}{p{4cm}|}{Extensionista} &
\multicolumn{1}{p{4cm}|}{Teórica} &
\multicolumn{1}{p{4cm}|}{Prática} \\ 
\multicolumn{1}{|p{4cm}|}{60\,horas} &
\multicolumn{1}{p{4cm}|}{0\;horas} &
\multicolumn{1}{p{4cm}|}{4\;horas/aula} &
\multicolumn{1}{p{4cm}|}{0\;horas/aula} \\ 
\hline
\multicolumn{4}{|p{16cm}|}{Ementa:}\\
\multicolumn{4}{|p{16cm}|}{}\\
\multicolumn{4}{|p{\dimexpr 16cm + 6\tabcolsep\relax}|}{Conceitos fundamentais de Mecânica Quântica, medida e evolução e emaranhamento ou entrelaçamento quântico. Introdução e informação clássica quântica. Objetivo: Esse curso pretende ser uma introdução aos conceitos básicos da informação Quântica para alunos dos cursos de Física e Computação da UFOP, com ênfase em recentes resultados obtidos em computação quântica pra sistemas discretos e contínuos.}\\
\multicolumn{4}{|p{16cm}|}{}\\
\hline
\multicolumn{4}{|p{16cm}|}{Conteúdo programático:}\\*
\multicolumn{4}{|p{\dimexpr 16cm + 6\tabcolsep\relax}|}{\begin{enumerate}[start=1]\item Conceitos Fundamentais: O experimento de Stern-Gerlach. Axiomas e formalismo da MQ. O q-bit. Combinação de sistemas quânticos: O produto tensorial. Matriz densidade. Teorema de Gleason, entrelaçamento quântico e de composição de Schmidt. Interpretação de "Ensemble". Teorema GHJW. Medida e evolução: Medidas ortogonais e generalizadas. Superoperadores. Representação de Kraus. Equação-mestre. Emaranhamento ou entrelaçamento quântico: Localidade de Einstein e "variáveis escondidas". Desigualdades de Bell. Codificação densa, teorema de não-clonagem e teleporte quântico\end{enumerate}}\\
\multicolumn{4}{|p{\dimexpr 16cm + 6\tabcolsep\relax}|}{\begin{enumerate}[start=2]\item Computação Quântica: Circuitos clássicos. Circuitos quânticos. Exemplos de algoritmos quânticos. Algoritmo de Deutsch-Josza, algoritmo de Shor e algoritmo de busca de Grover\end{enumerate}}\\
\multicolumn{4}{|p{\dimexpr 16cm + 6\tabcolsep\relax}|}{\begin{enumerate}[start=3]\item Informação clássica e quântica: Entropia de Shannon e compressão de dados. Informação mútua e o teorema do canal com ruído. Entropia de Von-Neumann. Compressão de dados quânticos. Codificação de estados misturados. Informação de Holevo. Teoremas de acessibilidade de informação. Noções sobre correção de erros quânticos\end{enumerate}}\\
\multicolumn{4}{|p{\dimexpr 16cm + 6\tabcolsep\relax}|}{}\\
\hline
\multicolumn{4}{|p{16cm}|}{Bibliografia Básica:}\\*
\multicolumn{4}{|p{\dimexpr 16cm + 6\tabcolsep\relax}|}{\begin{itemize}\item SAKURAI, J.J. Modern Quantum Mechanics.\end{itemize}}\\
\multicolumn{4}{|p{\dimexpr 16cm + 6\tabcolsep\relax}|}{\begin{itemize}\item PRESKILL, J. Quantum Information and Computation.\end{itemize}}\\
\multicolumn{4}{|p{\dimexpr 16cm + 6\tabcolsep\relax}|}{\begin{itemize}\item NIELSEN, M.A.; CHUANG, I.L. Quantum Computation and Quantum Information.\end{itemize}}\\
\multicolumn{4}{|p{16cm}|}{}\\
\hline
\multicolumn{4}{|p{16cm}|}{Bibliografia Complementar:}\\*
\multicolumn{4}{|p{\dimexpr 16cm + 6\tabcolsep\relax}|}{\begin{itemize}\item[] (nenhum item informado)\end{itemize}}\\
\hline
\end{longtable}
\end{center}

\clearpage
\begin{center}
\begin{tabular}{lccr}
 \multirow{3}{*}{\includegraphics[height=2.7cm]{brasao.png}} &
 \multicolumn{2}{c}{\bfseries UNIVERSIDADE FEDERAL DE OURO PRETO} &
 \ \ \ \ \multirow{3}{*}{\includegraphics[height=2.7cm]{ufop.png}} \\
 & \multicolumn{2}{c}{\bfseries PRÓ-REITORIA DE GRADUAÇÃO} & \\
 & \multicolumn{2}{c}{} & \\
 & \multicolumn{2}{c}{\Large\bfseries PROGRAMA DE DISCIPLINA} & \\
\end{tabular}
\end{center}

\begin{center}
\begin{longtable}{|p{4cm}|p{4cm}|p{4cm}|p{4cm}|}
\hline
\multicolumn{3}{|p{12cm}|}{Nome do Componente Curricular em Português:} &
\multicolumn{1}{p{4cm}|}{Código:} \\ 
\multicolumn{3}{|p{12cm}|}{\textbf{Introdução a LIBRAS}} &
\textbf{LET966}\\ 
\multicolumn{3}{|p{12cm}|}{Nome do Componente Curricular em Inglês:} & \\ 
\multicolumn{3}{|p{12cm}|}{\textbf{Introduction to the Brazilian Sign Language}} & \\ 
\hline
\multicolumn{3}{|p{12cm}|}{Nome e Sigla do Departamento} & Unidade Acadêmica: \\ 
\multicolumn{3}{|p{12cm}|}{Departamento de Letras (DELET)} & {ICHS} \\ 
\hline
\multicolumn{4}{|p{16cm}|}{Modalidade de Oferta:
[X] presencial \hspace{0.6cm}
[ ] semipresencial \hspace{0.6cm}
[ ] à distância}\\
\hline
\multicolumn{2}{|p{8cm}|}{Carga horária semestral} &
\multicolumn{2}{p{8cm}|}{Carga horária semanal}\\
\hline
\multicolumn{1}{|p{4cm}|}{Total} &
\multicolumn{1}{p{4cm}|}{Extensionista} &
\multicolumn{1}{p{4cm}|}{Teórica} &
\multicolumn{1}{p{4cm}|}{Prática} \\ 
\multicolumn{1}{|p{4cm}|}{60\,horas} &
\multicolumn{1}{p{4cm}|}{0\;horas} &
\multicolumn{1}{p{4cm}|}{2\;horas/aula} &
\multicolumn{1}{p{4cm}|}{2\;horas/aula} \\ 
\hline
\multicolumn{4}{|p{16cm}|}{Ementa:}\\
\multicolumn{4}{|p{16cm}|}{}\\
\multicolumn{4}{|p{\dimexpr 16cm + 6\tabcolsep\relax}|}{Princípios básicos do funcionamento da Língua Brasileira de Sinais – LIBRAS. Estrutura linguística em contextos comunicativos. Aspectos peculiares da cultura das pessoas surdas.}\\
\multicolumn{4}{|p{16cm}|}{}\\
\hline
\multicolumn{4}{|p{16cm}|}{Conteúdo programático:}\\*
\multicolumn{4}{|p{\dimexpr 16cm + 6\tabcolsep\relax}|}{\begin{enumerate}[start=1]\item A) Conceitual: 1) Adquirir conhecimentos básicos de um conjunto lexical envolvendo a variação dialetal da LIBRAS praticada em Minas Gerais; 2) Compreender o código gestual do Alfabeto Manual ou escrita manual datilológica e como a mesma é utilizada em situações comunicativas; 3) Adquirir noções básicas da organização fonológica da LIBRAS, expressas através dos Parâmetros Fonológicos da LIBRAS; 4) Adquirir noções básicas da organização morfossintática da LIBRAS; 5) Refletir criticamente sobre a concepção da LIBRAS enquanto língua com status linguístico equivalente ao das línguas orais; 6) Adquirir noções básicas de dialeto, variação dialetal, idioleto, empréstimo linguístico e regionalismo em LIBRAS\end{enumerate}}\\
\multicolumn{4}{|p{\dimexpr 16cm + 6\tabcolsep\relax}|}{\begin{enumerate}[start=2]\item B) Procedimental: 1) Desenvolver estratégias de leitura, interação e compreensão de textos sinalizados e registrados em vídeos; 2) Desenvolver estratégias de conversação em LIBRAS; 3) Desenvolver estratégias de conversação que utilizem o Alfabeto Manual; 4) Desenvolver a habilidade de reconhecer e produzir enunciados básicos em situações comunicativas envolvendo as seguintes temáticas: saudação, apresentação, escolaridade, organização espacial e temporal; 5) Principiar o desenvolvimento da habilidade de produção do sentido em LIBRAS; 6) Desenvolver estratégias para aprimorar as habilidades gestuais/motoras e visuais\end{enumerate}}\\
\multicolumn{4}{|p{\dimexpr 16cm + 6\tabcolsep\relax}|}{\begin{enumerate}[start=3]\item C) Atitudinal: 1) Posicionar-se criticamente enquanto discente que compartilha a sala de aula com um profissional surdo na condição de docente e refletir sobre o respeito e valorização dispensada a este profissional às pessoas surdas em geral; 2) Refletir criticamente sobre a pessoa surda como sujeito da enunciação; 3) Refletir sobre a importância e o valor linguístico, histórico, social e cultural da LIBRAS; 4) Refletir criticamente sobre o respeito e valorização dos hábitos, costumes e tradições culturais das pessoas surdas; 5) Reconhecer-se como sujeito que está a desenvolver enunciados em uma modalidade de língua gestual-visual, portanto diferente da modalidade oral que é utilizada predominantemente na sociedade\end{enumerate}}\\
\multicolumn{4}{|p{\dimexpr 16cm + 6\tabcolsep\relax}|}{}\\
\hline
\multicolumn{4}{|p{16cm}|}{Bibliografia Básica:}\\*
\multicolumn{4}{|p{\dimexpr 16cm + 6\tabcolsep\relax}|}{\begin{itemize}\item SOUZA, Tanya Amara Felipe. Libras em Contexto: livro do estudante/cursista. 2001.\end{itemize}}\\
\multicolumn{4}{|p{\dimexpr 16cm + 6\tabcolsep\relax}|}{\begin{itemize}\item BRITO, Lucinda Ferreira. Por uma gramática das línguas de sinais. Rio de Janeiro: Tempo Brasileiro, 1995. 273p.\end{itemize}}\\
\multicolumn{4}{|p{\dimexpr 16cm + 6\tabcolsep\relax}|}{\begin{itemize}\item Ronice Müller de Quadros, Lodenir Becker Karnopp. Língua de sinais brasileira: estudos linguísticos. Porto Alegre, RS: Artmed, 2004.\end{itemize}}\\
\multicolumn{4}{|p{\dimexpr 16cm + 6\tabcolsep\relax}|}{\begin{itemize}\item Audrei Gesser. Libras? Que língua é essa? Crenças e preconceitos em torno da língua de sinais e da realidade surda. São Paulo: Parábola, 2009.\end{itemize}}\\
\multicolumn{4}{|p{16cm}|}{}\\
\hline
\multicolumn{4}{|p{16cm}|}{Bibliografia Complementar:}\\*
\multicolumn{4}{|p{\dimexpr 16cm + 6\tabcolsep\relax}|}{\begin{itemize}\item SACKS, Oliver W. Vendo vozes: uma jornada pelo mundo dos surdos. São Paulo: Companhia das Letras, 1998. 196p.\end{itemize}}\\
\multicolumn{4}{|p{\dimexpr 16cm + 6\tabcolsep\relax}|}{\begin{itemize}\item Carlos Skliar (org.). A surdez: um olhar sobre as diferenças. 3. ed. Porto Alegre, RS: Mediação, 2005.\end{itemize}}\\
\multicolumn{4}{|p{\dimexpr 16cm + 6\tabcolsep\relax}|}{\begin{itemize}\item Strobel, Karin. As imagens do outro sobre a cultura surda. 3. ed. rev. Florianópolis, SC: Ed. UFSC, 2013.\end{itemize}}\\
\multicolumn{4}{|p{\dimexpr 16cm + 6\tabcolsep\relax}|}{\begin{itemize}\item Fernando César Capovilla, Walkíria Duarte Raphael, Aline Cristina Mauricio Novo. Deit-libras: dicionário enciclopédico ilustrado trilíngue da língua de sinais brasileira: baseado em linguística e neurociências cognitivas. 3. ed. rev. e ampl. São Paulo: EDUSP, 2013.\end{itemize}}\\
\hline
\end{longtable}
\end{center}

\clearpage
\begin{center}
\begin{tabular}{lccr}
 \multirow{3}{*}{\includegraphics[height=2.7cm]{brasao.png}} &
 \multicolumn{2}{c}{\bfseries UNIVERSIDADE FEDERAL DE OURO PRETO} &
 \ \ \ \ \multirow{3}{*}{\includegraphics[height=2.7cm]{ufop.png}} \\
 & \multicolumn{2}{c}{\bfseries PRÓ-REITORIA DE GRADUAÇÃO} & \\
 & \multicolumn{2}{c}{} & \\
 & \multicolumn{2}{c}{\Large\bfseries PROGRAMA DE DISCIPLINA} & \\
\end{tabular}
\end{center}

\begin{center}
\begin{longtable}{|p{4cm}|p{4cm}|p{4cm}|p{4cm}|}
\hline
\multicolumn{3}{|p{12cm}|}{Nome do Componente Curricular em Português:} &
\multicolumn{1}{p{4cm}|}{Código:} \\ 
\multicolumn{3}{|p{12cm}|}{\textbf{Ética}} &
\textbf{FIL652}\\ 
\multicolumn{3}{|p{12cm}|}{Nome do Componente Curricular em Inglês:} & \\ 
\multicolumn{3}{|p{12cm}|}{\textbf{Ethics}} & \\ 
\hline
\multicolumn{3}{|p{12cm}|}{Nome e Sigla do Departamento} & Unidade Acadêmica: \\ 
\multicolumn{3}{|p{12cm}|}{Departamento de Filosofia (DEFIL)} & {IFAC} \\ 
\hline
\multicolumn{4}{|p{16cm}|}{Modalidade de Oferta:
[X] presencial \hspace{0.6cm}
[ ] semipresencial \hspace{0.6cm}
[ ] à distância}\\
\hline
\multicolumn{2}{|p{8cm}|}{Carga horária semestral} &
\multicolumn{2}{p{8cm}|}{Carga horária semanal}\\
\hline
\multicolumn{1}{|p{4cm}|}{Total} &
\multicolumn{1}{p{4cm}|}{Extensionista} &
\multicolumn{1}{p{4cm}|}{Teórica} &
\multicolumn{1}{p{4cm}|}{Prática} \\ 
\multicolumn{1}{|p{4cm}|}{60\,horas} &
\multicolumn{1}{p{4cm}|}{0\;horas} &
\multicolumn{1}{p{4cm}|}{3\;horas/aula} &
\multicolumn{1}{p{4cm}|}{1\;horas/aula} \\ 
\hline
\multicolumn{4}{|p{16cm}|}{Ementa:}\\
\multicolumn{4}{|p{16cm}|}{}\\
\multicolumn{4}{|p{\dimexpr 16cm + 6\tabcolsep\relax}|}{Noções fundamentais no estudo da Ética. Ética Normativa: Éticas Teleológicas ou "Conseqüencialistas". A Ética de Virtude e o Utilitarismo. Éticas Deontológicas: ética kantiana, ética contratualista. Estudo nos textos dos filósofos.}\\
\multicolumn{4}{|p{16cm}|}{}\\
\hline
\multicolumn{4}{|p{16cm}|}{Conteúdo programático:}\\*
\multicolumn{4}{|p{\dimexpr 16cm + 6\tabcolsep\relax}|}{\begin{enumerate}[start=1]\item Noções fundamentais de ética: valor, norma, lei\end{enumerate}}\\
\multicolumn{4}{|p{\dimexpr 16cm + 6\tabcolsep\relax}|}{\begin{enumerate}[start=2]\item Éticas Normativas, Metaética e Bioética\end{enumerate}}\\
\multicolumn{4}{|p{\dimexpr 16cm + 6\tabcolsep\relax}|}{\begin{enumerate}[start=3]\item Éticas Normativas: Éticas Teleológicas – estudo de excertos dos textos dos filósofos; Éticas Deontológicas – estudo de excertos dos textos dos filósofos\end{enumerate}}\\
\multicolumn{4}{|p{\dimexpr 16cm + 6\tabcolsep\relax}|}{}\\
\hline
\multicolumn{4}{|p{16cm}|}{Bibliografia Básica:}\\*
\multicolumn{4}{|p{\dimexpr 16cm + 6\tabcolsep\relax}|}{\begin{itemize}\item Aristóteles. Ética a Nicômaco: tradução de Antonio de Castro Caieiro. São Paulo, SP: Atlas, 2009.\end{itemize}}\\
\multicolumn{4}{|p{\dimexpr 16cm + 6\tabcolsep\relax}|}{\begin{itemize}\item Kant, Immanuel. Crítica da razão prática: tradução de Valério Rohden. São Paulo, SP: Martins Fontes, 2005.\end{itemize}}\\
\multicolumn{4}{|p{\dimexpr 16cm + 6\tabcolsep\relax}|}{\begin{itemize}\item Macintyre, A. Depois da virtude: tradução de Jussara Simões. São Paulo, SP: EDUSC, 2001.\end{itemize}}\\
\multicolumn{4}{|p{\dimexpr 16cm + 6\tabcolsep\relax}|}{\begin{itemize}\item Platão. A república. 9. ed. Lisboa: Fundação Calouste Gulbenkian, 2001.\end{itemize}}\\
\multicolumn{4}{|p{\dimexpr 16cm + 6\tabcolsep\relax}|}{\begin{itemize}\item Platão. A república. 7. ed. Lisboa: Fundação Calouste Gulbenkian, 1993.\end{itemize}}\\
\multicolumn{4}{|p{\dimexpr 16cm + 6\tabcolsep\relax}|}{\begin{itemize}\item Taylor, Charles. A ética da autenticidade. Lisboa, Portugal: Edições 70, 2009.\end{itemize}}\\
\multicolumn{4}{|p{16cm}|}{}\\
\hline
\multicolumn{4}{|p{16cm}|}{Bibliografia Complementar:}\\*
\multicolumn{4}{|p{\dimexpr 16cm + 6\tabcolsep\relax}|}{\begin{itemize}\item Mill, John Stuart. Uma introdução aos princípios da moral e da legislação. São Paulo: Abril Cultural, 1974.\end{itemize}}\\
\multicolumn{4}{|p{\dimexpr 16cm + 6\tabcolsep\relax}|}{\begin{itemize}\item Tugendhat, Ernst. Lições sobre ética. 8. ed. Petrópolis, RJ: Vozes, 2010.\end{itemize}}\\
\multicolumn{4}{|p{\dimexpr 16cm + 6\tabcolsep\relax}|}{\begin{itemize}\item Tugendhat, Ernst. Lições sobre ética. 9. ed. Petrópolis: Vozes, 2012.\end{itemize}}\\
\multicolumn{4}{|p{\dimexpr 16cm + 6\tabcolsep\relax}|}{\begin{itemize}\item Dicionário de ética e filosofia moral. São Leopoldo, RS: Ed. UNISINOS, 2003.\end{itemize}}\\
\hline
\end{longtable}
\end{center}

\clearpage
\begin{center}
\begin{tabular}{lccr}
 \multirow{3}{*}{\includegraphics[height=2.7cm]{brasao.png}} &
 \multicolumn{2}{c}{\bfseries UNIVERSIDADE FEDERAL DE OURO PRETO} &
 \ \ \ \ \multirow{3}{*}{\includegraphics[height=2.7cm]{ufop.png}} \\
 & \multicolumn{2}{c}{\bfseries PRÓ-REITORIA DE GRADUAÇÃO} & \\
 & \multicolumn{2}{c}{} & \\
 & \multicolumn{2}{c}{\Large\bfseries PROGRAMA DE DISCIPLINA} & \\
\end{tabular}
\end{center}

\begin{center}
\begin{longtable}{|p{4cm}|p{4cm}|p{4cm}|p{4cm}|}
\hline
\multicolumn{3}{|p{12cm}|}{Nome do Componente Curricular em Português:} &
\multicolumn{1}{p{4cm}|}{Código:} \\ 
\multicolumn{3}{|p{12cm}|}{\textbf{Filosofia das Ciências: Século XX}} &
\textbf{FIL643}\\ 
\multicolumn{3}{|p{12cm}|}{Nome do Componente Curricular em Inglês:} & \\ 
\multicolumn{3}{|p{12cm}|}{\textbf{Philosophy of Science: 20th Century}} & \\ 
\hline
\multicolumn{3}{|p{12cm}|}{Nome e Sigla do Departamento} & Unidade Acadêmica: \\ 
\multicolumn{3}{|p{12cm}|}{Departamento de Filosofia (DEFIL)} & {IFAC} \\ 
\hline
\multicolumn{4}{|p{16cm}|}{Modalidade de Oferta:
[X] presencial \hspace{0.6cm}
[ ] semipresencial \hspace{0.6cm}
[ ] à distância}\\
\hline
\multicolumn{2}{|p{8cm}|}{Carga horária semestral} &
\multicolumn{2}{p{8cm}|}{Carga horária semanal}\\
\hline
\multicolumn{1}{|p{4cm}|}{Total} &
\multicolumn{1}{p{4cm}|}{Extensionista} &
\multicolumn{1}{p{4cm}|}{Teórica} &
\multicolumn{1}{p{4cm}|}{Prática} \\ 
\multicolumn{1}{|p{4cm}|}{60\,horas} &
\multicolumn{1}{p{4cm}|}{0\;horas} &
\multicolumn{1}{p{4cm}|}{3\;horas/aula} &
\multicolumn{1}{p{4cm}|}{1\;horas/aula} \\ 
\hline
\multicolumn{4}{|p{16cm}|}{Ementa:}\\
\multicolumn{4}{|p{16cm}|}{}\\
\multicolumn{4}{|p{\dimexpr 16cm + 6\tabcolsep\relax}|}{O curso consiste numa apresentação da Ciência Moderna a partir de pontos de vista ou perspectivas de acesso provenientes da Filosofia. Seu eixo teórico é o debate sobre o que é, afinal, o conhecimento científico. Em vista disso, dois objetivos são contemplados. Primeiro: buscar uma definição para a atividade científica que seja filosoficamente suficiente—e, caso não a obtenhamos, esclarecer porque isso não foi possível. Segundo: estimular o exercício de uma atitude crítica em relação à hegemonia cultural da ciência.}\\
\multicolumn{4}{|p{16cm}|}{}\\
\hline
\multicolumn{4}{|p{16cm}|}{Conteúdo programático:}\\*
\multicolumn{4}{|p{\dimexpr 16cm + 6\tabcolsep\relax}|}{\begin{enumerate}[start=1]\item I. Preâmbulo: Notícia histórica sobre o nascimento da Ciência Moderna\end{enumerate}}\\
\multicolumn{4}{|p{\dimexpr 16cm + 6\tabcolsep\relax}|}{\begin{enumerate}[start=2]\item II. O conhecimento científico como resultado da interação entre teoria e experimentação: 1. Em torno do problema da indução: um estatuto para a objetividade; 2. Duas concepções a propósito da cientificidade da ciência: a) Karl Popper e o racionalismo crítico; b) Thomas Kuhn, paradigmas e revoluções\end{enumerate}}\\
\multicolumn{4}{|p{\dimexpr 16cm + 6\tabcolsep\relax}|}{\begin{enumerate}[start=3]\item III. Conhecimento científico, valores e sociedade: 1. Anarquismo e humanitarismo: Feyerabend e a ciência; 2. Epistemologia e hermenêutica: Rorty e a ciência\end{enumerate}}\\
\multicolumn{4}{|p{\dimexpr 16cm + 6\tabcolsep\relax}|}{}\\
\hline
\multicolumn{4}{|p{16cm}|}{Bibliografia Básica:}\\*
\multicolumn{4}{|p{\dimexpr 16cm + 6\tabcolsep\relax}|}{\begin{itemize}\item Popper, Karl R. A lógica da pesquisa científica. São Paulo: Cultrix, 1972.\end{itemize}}\\
\multicolumn{4}{|p{\dimexpr 16cm + 6\tabcolsep\relax}|}{\begin{itemize}\item Kuhn, Thomas S. A estrutura das revoluções científicas. 9. ed. São Paulo: Perspectiva, 2005.\end{itemize}}\\
\multicolumn{4}{|p{\dimexpr 16cm + 6\tabcolsep\relax}|}{\begin{itemize}\item Karl Raimund Popper; tradução de Milton Amado. Conhecimento objetivo: uma abordagem evolucionária. Belo Horizonte, MG: Itatiaia, 1999.\end{itemize}}\\
\multicolumn{4}{|p{\dimexpr 16cm + 6\tabcolsep\relax}|}{\begin{itemize}\item Karl R. Popper; tradução de Sérgio Bath. Conjecturas e refutações: o progresso do conhecimento científico. 5. ed. Brasília: Ed. da UnB, 2008.\end{itemize}}\\
\multicolumn{4}{|p{\dimexpr 16cm + 6\tabcolsep\relax}|}{\begin{itemize}\item Bachelard, Gaston. O novo espírito científico; A poética do espaço. 3. ed. São Paulo: Nova Cultural, 1988.\end{itemize}}\\
\multicolumn{4}{|p{\dimexpr 16cm + 6\tabcolsep\relax}|}{\begin{itemize}\item Feyerabend, Paul K. Contra o método. 2. ed. São Paulo: Ed. UNESP, 2011.\end{itemize}}\\
\multicolumn{4}{|p{\dimexpr 16cm + 6\tabcolsep\relax}|}{\begin{itemize}\item Kuhn, Thomas S. A estrutura das revoluções científicas. 2. ed. São Paulo: Perspectiva, 1978.\end{itemize}}\\
\multicolumn{4}{|p{\dimexpr 16cm + 6\tabcolsep\relax}|}{\begin{itemize}\item Kuhn, Thomas S. A estrutura das revoluções científicas. 3. ed. São Paulo: Perspectiva, 1994.\end{itemize}}\\
\multicolumn{4}{|p{16cm}|}{}\\
\hline
\multicolumn{4}{|p{16cm}|}{Bibliografia Complementar:}\\*
\multicolumn{4}{|p{\dimexpr 16cm + 6\tabcolsep\relax}|}{\begin{itemize}\item Koyré, Alexandre. Estudos de história do pensamento científico. 2. rev. ed. Rio de Janeiro: Forense Universitária, 1991.\end{itemize}}\\
\multicolumn{4}{|p{\dimexpr 16cm + 6\tabcolsep\relax}|}{\begin{itemize}\item Van Fraassen, Bas C. A imagem científica. São Paulo: UNESP, 2007.\end{itemize}}\\
\multicolumn{4}{|p{\dimexpr 16cm + 6\tabcolsep\relax}|}{\begin{itemize}\item Sasaki, Chikara. Introdução à teoria da ciência. São Paulo: Edusp, 2010.\end{itemize}}\\
\multicolumn{4}{|p{\dimexpr 16cm + 6\tabcolsep\relax}|}{\begin{itemize}\item Ian Hacking; tradução Pedro Rocha de Oliveira; revisão técnica Antonio Augusto Passos Videira. Representar e intervir: tópicos introdutórios de filosofia da ciência natural. Rio de Janeiro, RJ: EdUERJ, 2012.\end{itemize}}\\
\multicolumn{4}{|p{\dimexpr 16cm + 6\tabcolsep\relax}|}{\begin{itemize}\item Chalmers, A. F. O que é ciência, afinal? São Paulo: Brasiliense, 1995.\end{itemize}}\\
\multicolumn{4}{|p{\dimexpr 16cm + 6\tabcolsep\relax}|}{\begin{itemize}\item Koyré, Alexandre. Estudos de história do pensamento científico. Rio de Janeiro: Forense Universitária.\end{itemize}}\\
\hline
\end{longtable}
\end{center}

\clearpage
\begin{center}
\begin{tabular}{lccr}
 \multirow{3}{*}{\includegraphics[height=2.7cm]{brasao.png}} &
 \multicolumn{2}{c}{\bfseries UNIVERSIDADE FEDERAL DE OURO PRETO} &
 \ \ \ \ \multirow{3}{*}{\includegraphics[height=2.7cm]{ufop.png}} \\
 & \multicolumn{2}{c}{\bfseries PRÓ-REITORIA DE GRADUAÇÃO} & \\
 & \multicolumn{2}{c}{} & \\
 & \multicolumn{2}{c}{\Large\bfseries PROGRAMA DE DISCIPLINA} & \\
\end{tabular}
\end{center}

\begin{center}
\begin{longtable}{|p{4cm}|p{4cm}|p{4cm}|p{4cm}|}
\hline
\multicolumn{3}{|p{12cm}|}{Nome do Componente Curricular em Português:} &
\multicolumn{1}{p{4cm}|}{Código:} \\ 
\multicolumn{3}{|p{12cm}|}{\textbf{Introdução à História da Filosofia}} &
\textbf{FIL612}\\ 
\multicolumn{3}{|p{12cm}|}{Nome do Componente Curricular em Inglês:} & \\ 
\multicolumn{3}{|p{12cm}|}{\textbf{Introduction to History of Philosophy}} & \\ 
\hline
\multicolumn{3}{|p{12cm}|}{Nome e Sigla do Departamento} & Unidade Acadêmica: \\ 
\multicolumn{3}{|p{12cm}|}{Departamento de Filosofia (DEFIL)} & {IFAC} \\ 
\hline
\multicolumn{4}{|p{16cm}|}{Modalidade de Oferta:
[X] presencial \hspace{0.6cm}
[ ] semipresencial \hspace{0.6cm}
[ ] à distância}\\
\hline
\multicolumn{2}{|p{8cm}|}{Carga horária semestral} &
\multicolumn{2}{p{8cm}|}{Carga horária semanal}\\
\hline
\multicolumn{1}{|p{4cm}|}{Total} &
\multicolumn{1}{p{4cm}|}{Extensionista} &
\multicolumn{1}{p{4cm}|}{Teórica} &
\multicolumn{1}{p{4cm}|}{Prática} \\ 
\multicolumn{1}{|p{4cm}|}{60\,horas} &
\multicolumn{1}{p{4cm}|}{0\;horas} &
\multicolumn{1}{p{4cm}|}{3\;horas/aula} &
\multicolumn{1}{p{4cm}|}{1\;horas/aula} \\ 
\hline
\multicolumn{4}{|p{16cm}|}{Ementa:}\\
\multicolumn{4}{|p{16cm}|}{}\\
\multicolumn{4}{|p{\dimexpr 16cm + 6\tabcolsep\relax}|}{Trata-se da apresentação de um panorama do pensamento filosófico ocidental através da discussão de questões e textos clássicos. O curso visa a oferecer condições para a formação de um entendimento básico do que é a filosofia, criando com isso a possibilidade de sua incorporação pelo estudante por meio de experiências próprias.}\\
\multicolumn{4}{|p{16cm}|}{}\\
\hline
\multicolumn{4}{|p{16cm}|}{Conteúdo programático:}\\*
\multicolumn{4}{|p{\dimexpr 16cm + 6\tabcolsep\relax}|}{\begin{enumerate}[start=1]\item I. Introdução: 1. Disposições conceituais preliminares; 2. O nascimento da Filosofia\end{enumerate}}\\
\multicolumn{4}{|p{\dimexpr 16cm + 6\tabcolsep\relax}|}{\begin{enumerate}[start=2]\item II. Sobre Platão: 1. Influxos formativos; 2. O idealismo; 3. A articulação sistemática da obra\end{enumerate}}\\
\multicolumn{4}{|p{\dimexpr 16cm + 6\tabcolsep\relax}|}{\begin{enumerate}[start=3]\item III. O pensamento moderno em duas versões: 1. Descartes e o grande racionalismo; 2. Kant e o programa do esclarecimento\end{enumerate}}\\
\multicolumn{4}{|p{\dimexpr 16cm + 6\tabcolsep\relax}|}{\begin{enumerate}[start=4]\item IV. A filosofia em crise: 1. Sobre Nietzsche: um novo significado para o filosofar\end{enumerate}}\\
\multicolumn{4}{|p{\dimexpr 16cm + 6\tabcolsep\relax}|}{}\\
\hline
\multicolumn{4}{|p{16cm}|}{Bibliografia Básica:}\\*
\multicolumn{4}{|p{\dimexpr 16cm + 6\tabcolsep\relax}|}{\begin{itemize}\item Descartes, René. Discurso do método. São Paulo, SP: Martins Fontes, 2006.\end{itemize}}\\
\multicolumn{4}{|p{\dimexpr 16cm + 6\tabcolsep\relax}|}{\begin{itemize}\item Nietzsche, Friedrich Wilhelm. Crepúsculo dos ídolos. Rio de Janeiro, RJ: Relume Dumará, 2000.\end{itemize}}\\
\multicolumn{4}{|p{\dimexpr 16cm + 6\tabcolsep\relax}|}{\begin{itemize}\item Kant, Immanuel. Textos seletos. 5. ed. Petrópolis: Vozes, 2009.\end{itemize}}\\
\multicolumn{4}{|p{\dimexpr 16cm + 6\tabcolsep\relax}|}{\begin{itemize}\item Vernant, Jean Pierre. As origens do pensamento grego. 10. ed. São Paulo: DIFEL, 1998.\end{itemize}}\\
\multicolumn{4}{|p{\dimexpr 16cm + 6\tabcolsep\relax}|}{\begin{itemize}\item Vernant, Jean Pierre. As origens do pensamento grego. 5. ed. São Paulo: DIFEL, 1986.\end{itemize}}\\
\multicolumn{4}{|p{\dimexpr 16cm + 6\tabcolsep\relax}|}{\begin{itemize}\item Platão. A república. 9. ed. Lisboa: Fundação Calouste Gulbenkian, 2001.\end{itemize}}\\
\multicolumn{4}{|p{16cm}|}{}\\
\hline
\multicolumn{4}{|p{16cm}|}{Bibliografia Complementar:}\\*
\multicolumn{4}{|p{\dimexpr 16cm + 6\tabcolsep\relax}|}{\begin{itemize}\item Ática, 1989.\end{itemize}}\\
\multicolumn{4}{|p{\dimexpr 16cm + 6\tabcolsep\relax}|}{\begin{itemize}\item Machado, Roberto. Nietzsche e a verdade. 2. ed. Rio de Janeiro: Rocco, 1985.\end{itemize}}\\
\multicolumn{4}{|p{\dimexpr 16cm + 6\tabcolsep\relax}|}{\begin{itemize}\item Koyre, Alexandre. Considerações sobre Descartes. 4. ed. Lisboa: Editorial Presença, 1992.\end{itemize}}\\
\multicolumn{4}{|p{\dimexpr 16cm + 6\tabcolsep\relax}|}{\begin{itemize}\item Williams, B. Platão. São Paulo, SP: UNESP, 2000.\end{itemize}}\\
\multicolumn{4}{|p{\dimexpr 16cm + 6\tabcolsep\relax}|}{\begin{itemize}\item Reale, Giovanni; Antiseri, D. História da filosofia: vol. 1; vol. 2; vol. 3. São Paulo, SP: Paulus, 2003.\end{itemize}}\\
\multicolumn{4}{|p{\dimexpr 16cm + 6\tabcolsep\relax}|}{\begin{itemize}\item Piettre, Bernard. Platão: república. São Paulo, SP.\end{itemize}}\\
\hline
\end{longtable}
\end{center}

\clearpage
\begin{center}
\begin{tabular}{lccr}
 \multirow{3}{*}{\includegraphics[height=2.7cm]{brasao.png}} &
 \multicolumn{2}{c}{\bfseries UNIVERSIDADE FEDERAL DE OURO PRETO} &
 \ \ \ \ \multirow{3}{*}{\includegraphics[height=2.7cm]{ufop.png}} \\
 & \multicolumn{2}{c}{\bfseries PRÓ-REITORIA DE GRADUAÇÃO} & \\
 & \multicolumn{2}{c}{} & \\
 & \multicolumn{2}{c}{\Large\bfseries PROGRAMA DE DISCIPLINA} & \\
\end{tabular}
\end{center}

\begin{center}
\begin{longtable}{|p{4cm}|p{4cm}|p{4cm}|p{4cm}|}
\hline
\multicolumn{3}{|p{12cm}|}{Nome do Componente Curricular em Português:} &
\multicolumn{1}{p{4cm}|}{Código:} \\ 
\multicolumn{3}{|p{12cm}|}{\textbf{Lógica I}} &
\textbf{FIL613}\\ 
\multicolumn{3}{|p{12cm}|}{Nome do Componente Curricular em Inglês:} & \\ 
\multicolumn{3}{|p{12cm}|}{\textbf{Logic I}} & \\ 
\hline
\multicolumn{3}{|p{12cm}|}{Nome e Sigla do Departamento} & Unidade Acadêmica: \\ 
\multicolumn{3}{|p{12cm}|}{Departamento de Filosofia (DEFIL)} & {IFAC} \\ 
\hline
\multicolumn{4}{|p{16cm}|}{Modalidade de Oferta:
[X] presencial \hspace{0.6cm}
[ ] semipresencial \hspace{0.6cm}
[ ] à distância}\\
\hline
\multicolumn{2}{|p{8cm}|}{Carga horária semestral} &
\multicolumn{2}{p{8cm}|}{Carga horária semanal}\\
\hline
\multicolumn{1}{|p{4cm}|}{Total} &
\multicolumn{1}{p{4cm}|}{Extensionista} &
\multicolumn{1}{p{4cm}|}{Teórica} &
\multicolumn{1}{p{4cm}|}{Prática} \\ 
\multicolumn{1}{|p{4cm}|}{60\,horas} &
\multicolumn{1}{p{4cm}|}{0\;horas} &
\multicolumn{1}{p{4cm}|}{3\;horas/aula} &
\multicolumn{1}{p{4cm}|}{1\;horas/aula} \\ 
\hline
\multicolumn{4}{|p{16cm}|}{Ementa:}\\
\multicolumn{4}{|p{16cm}|}{}\\
\multicolumn{4}{|p{\dimexpr 16cm + 6\tabcolsep\relax}|}{Estuda-se nesta disciplina elementos básicos de lógica proposicional e de predicados clássica. A abordagem é simplificada e sem perder de vista a linguagem natural e a aplicação da lógica à análise da argumentação e da linguagem filosófica. Inclui-se também alguns elementos de lógica informal.}\\
\multicolumn{4}{|p{16cm}|}{}\\
\hline
\multicolumn{4}{|p{16cm}|}{Conteúdo programático:}\\*
\multicolumn{4}{|p{\dimexpr 16cm + 6\tabcolsep\relax}|}{\begin{enumerate}[start=1]\item 1. Apresentação do curso: programa, bibliografia e avaliação. Lógica e filosofia: a importância do raciocínio; o papel da lógica na filosofia. Conceitos fundamentais de lógica: raciocínio e argumento, validade, solidez e cogência. Indicadores de premissa e de conclusão. Raciocínio dedutivo e não-dedutivo\end{enumerate}}\\
\multicolumn{4}{|p{\dimexpr 16cm + 6\tabcolsep\relax}|}{\begin{enumerate}[start=2]\item 2. Lógica proposicional clássica. Operadores verofuncionais. Forma lógica. Tabelas de verdade e inspetores de circunstâncias. Formalização da linguagem natural\end{enumerate}}\\
\multicolumn{4}{|p{\dimexpr 16cm + 6\tabcolsep\relax}|}{\begin{enumerate}[start=3]\item 3. Dedução natural proposicional clássica. O conceito de derivação ou demonstração. Regras, axiomas e teoremas. Regras primitivas e derivadas\end{enumerate}}\\
\multicolumn{4}{|p{\dimexpr 16cm + 6\tabcolsep\relax}|}{\begin{enumerate}[start=4]\item 4. Lógica quantificada clássica: a linguagem. Quantificadores. Nomes e predicados. Relações. Descrições definidas. Comparativos. Domínios de quantificação\end{enumerate}}\\
\multicolumn{4}{|p{\dimexpr 16cm + 6\tabcolsep\relax}|}{\begin{enumerate}[start=5]\item 5. Dedução natural quantificada clássica. Regras primitivas e derivadas\end{enumerate}}\\
\multicolumn{4}{|p{\dimexpr 16cm + 6\tabcolsep\relax}|}{}\\
\hline
\multicolumn{4}{|p{16cm}|}{Bibliografia Básica:}\\*
\multicolumn{4}{|p{\dimexpr 16cm + 6\tabcolsep\relax}|}{\begin{itemize}\item Warburton, Nigel. Pensar de A a Z: tradução V. Guerreiro. Lisboa, Portugal: Bizâncio, 2012.\end{itemize}}\\
\multicolumn{4}{|p{\dimexpr 16cm + 6\tabcolsep\relax}|}{\begin{itemize}\item Murcho, Desidério. O lugar da lógica na filosofia. Lisboa: Plátano, 2003.\end{itemize}}\\
\multicolumn{4}{|p{\dimexpr 16cm + 6\tabcolsep\relax}|}{\begin{itemize}\item Enciclopédia de termos lógico-filosóficos. São Paulo: Martins Fontes, 2006.\end{itemize}}\\
\multicolumn{4}{|p{\dimexpr 16cm + 6\tabcolsep\relax}|}{\begin{itemize}\item Weston, Anthony. A arte de argumentar. 2. ed. Lisboa: Gradiva, 1996.\end{itemize}}\\
\multicolumn{4}{|p{\dimexpr 16cm + 6\tabcolsep\relax}|}{\begin{itemize}\item W. H. Newton-Smith; tradução e notas Desidério Murcho. Lógica: um curso introdutório. 2. ed. Lisboa: Gradiva, 2005.\end{itemize}}\\
\multicolumn{4}{|p{16cm}|}{}\\
\hline
\multicolumn{4}{|p{16cm}|}{Bibliografia Complementar:}\\*
\multicolumn{4}{|p{\dimexpr 16cm + 6\tabcolsep\relax}|}{\begin{itemize}\item Epstein, R. L. e Carnielli, W. A. Pensamento crítico: o poder da lógica e da argumentação. São Paulo, SP: Rideel.\end{itemize}}\\
\multicolumn{4}{|p{\dimexpr 16cm + 6\tabcolsep\relax}|}{\begin{itemize}\item Pinto, Paulo Roberto Margutti. Introdução à lógica simbólica. Belo Horizonte, MG: Ed. UFMG, 2001.\end{itemize}}\\
\multicolumn{4}{|p{\dimexpr 16cm + 6\tabcolsep\relax}|}{\begin{itemize}\item Guido Imaguire e Cícero A. C. Barroso. Lógica: os jogos da razão. Ceará: Ed. UFC, 2006.\end{itemize}}\\
\multicolumn{4}{|p{\dimexpr 16cm + 6\tabcolsep\relax}|}{\begin{itemize}\item Graham Priest; tradução de Célia Teixeira, revisão científica Desidério Murcho. Lógica: para começar. Lisboa: Temas e Debates, 2002.\end{itemize}}\\
\multicolumn{4}{|p{\dimexpr 16cm + 6\tabcolsep\relax}|}{\begin{itemize}\item Mortari, Cezar Augusto. Introdução à lógica. 2. ed. São Paulo: UNESP.\end{itemize}}\\
\hline
\end{longtable}
\end{center}

\clearpage
\begin{center}
\begin{tabular}{lccr}
 \multirow{3}{*}{\includegraphics[height=2.7cm]{brasao.png}} &
 \multicolumn{2}{c}{\bfseries UNIVERSIDADE FEDERAL DE OURO PRETO} &
 \ \ \ \ \multirow{3}{*}{\includegraphics[height=2.7cm]{ufop.png}} \\
 & \multicolumn{2}{c}{\bfseries PRÓ-REITORIA DE GRADUAÇÃO} & \\
 & \multicolumn{2}{c}{} & \\
 & \multicolumn{2}{c}{\Large\bfseries PROGRAMA DE DISCIPLINA} & \\
\end{tabular}
\end{center}

\begin{center}
\begin{longtable}{|p{4cm}|p{4cm}|p{4cm}|p{4cm}|}
\hline
\multicolumn{3}{|p{12cm}|}{Nome do Componente Curricular em Português:} &
\multicolumn{1}{p{4cm}|}{Código:} \\ 
\multicolumn{3}{|p{12cm}|}{\textbf{Teoria do Conhecimento}} &
\textbf{FIL622}\\ 
\multicolumn{3}{|p{12cm}|}{Nome do Componente Curricular em Inglês:} & \\ 
\multicolumn{3}{|p{12cm}|}{\textbf{Theory of Knowledge}} & \\ 
\hline
\multicolumn{3}{|p{12cm}|}{Nome e Sigla do Departamento} & Unidade Acadêmica: \\ 
\multicolumn{3}{|p{12cm}|}{Departamento de Filosofia (DEFIL)} & {IFAC} \\ 
\hline
\multicolumn{4}{|p{16cm}|}{Modalidade de Oferta:
[X] presencial \hspace{0.6cm}
[ ] semipresencial \hspace{0.6cm}
[ ] à distância}\\
\hline
\multicolumn{2}{|p{8cm}|}{Carga horária semestral} &
\multicolumn{2}{p{8cm}|}{Carga horária semanal}\\
\hline
\multicolumn{1}{|p{4cm}|}{Total} &
\multicolumn{1}{p{4cm}|}{Extensionista} &
\multicolumn{1}{p{4cm}|}{Teórica} &
\multicolumn{1}{p{4cm}|}{Prática} \\ 
\multicolumn{1}{|p{4cm}|}{60\,horas} &
\multicolumn{1}{p{4cm}|}{0\;horas} &
\multicolumn{1}{p{4cm}|}{3\;horas/aula} &
\multicolumn{1}{p{4cm}|}{1\;horas/aula} \\ 
\hline
\multicolumn{4}{|p{16cm}|}{Ementa:}\\
\multicolumn{4}{|p{16cm}|}{}\\
\multicolumn{4}{|p{\dimexpr 16cm + 6\tabcolsep\relax}|}{O estudo da natureza, do método e dos limites do conhecimento humano, segundo as principais concepções de filosofia, com ênfase especial nas implicações deste estudo para a fundamentação das ciências naturais e humanas.}\\
\multicolumn{4}{|p{16cm}|}{}\\
\hline
\multicolumn{4}{|p{16cm}|}{Conteúdo programático:}\\*
\multicolumn{4}{|p{\dimexpr 16cm + 6\tabcolsep\relax}|}{\begin{enumerate}[start=1]\item Realismo e idealismo\end{enumerate}}\\
\multicolumn{4}{|p{\dimexpr 16cm + 6\tabcolsep\relax}|}{\begin{enumerate}[start=2]\item A ideia de conhecimento objetivo, gênese e estrutura do conhecimento, etc.\end{enumerate}}\\
\multicolumn{4}{|p{\dimexpr 16cm + 6\tabcolsep\relax}|}{\begin{enumerate}[start=3]\item Introdução ao problema do conhecimento\end{enumerate}}\\
\multicolumn{4}{|p{\dimexpr 16cm + 6\tabcolsep\relax}|}{\begin{enumerate}[start=4]\item Estudo da natureza, das fontes e da estrutura do conhecimento e da justificação\end{enumerate}}\\
\multicolumn{4}{|p{\dimexpr 16cm + 6\tabcolsep\relax}|}{\begin{enumerate}[start=5]\item Compreensão da discussão entre empiristas e racionalistas e entre fundacionistas e coerentistas\end{enumerate}}\\
\multicolumn{4}{|p{\dimexpr 16cm + 6\tabcolsep\relax}|}{}\\
\hline
\multicolumn{4}{|p{16cm}|}{Bibliografia Básica:}\\*
\multicolumn{4}{|p{\dimexpr 16cm + 6\tabcolsep\relax}|}{\begin{itemize}\item Russell, Bertrand. Lógica e conhecimento. São Paulo: Abril Cultural, 1974.\end{itemize}}\\
\multicolumn{4}{|p{\dimexpr 16cm + 6\tabcolsep\relax}|}{\begin{itemize}\item Wittgenstein, Ludwig. Investigações filosóficas. 2. ed. São Paulo: Abril Cultural, 1979.\end{itemize}}\\
\multicolumn{4}{|p{16cm}|}{}\\
\hline
\multicolumn{4}{|p{16cm}|}{Bibliografia Complementar:}\\*
\multicolumn{4}{|p{\dimexpr 16cm + 6\tabcolsep\relax}|}{\begin{itemize}\item Seleção de Oswaldo Porchat de Assis Pereira da Silva; traduções de Balthazar Barbosa Filho... [et al.]. Ensaios: Gilbert Ryle, John Langshaw Austin, Willard Van Orman Quine, Peter Frederick Strawson. 2. ed. São Paulo: Abril Cultural, 1980.\end{itemize}}\\
\multicolumn{4}{|p{\dimexpr 16cm + 6\tabcolsep\relax}|}{\begin{itemize}\item Schlick, Moritz. Coletânea de textos. São Paulo: Abril Cultural, 1980.\end{itemize}}\\
\hline
\end{longtable}
\end{center}

\clearpage
\begin{center}
\begin{tabular}{lccr}
 \multirow{3}{*}{\includegraphics[height=2.7cm]{brasao.png}} &
 \multicolumn{2}{c}{\bfseries UNIVERSIDADE FEDERAL DE OURO PRETO} &
 \ \ \ \ \multirow{3}{*}{\includegraphics[height=2.7cm]{ufop.png}} \\
 & \multicolumn{2}{c}{\bfseries PRÓ-REITORIA DE GRADUAÇÃO} & \\
 & \multicolumn{2}{c}{} & \\
 & \multicolumn{2}{c}{\Large\bfseries PROGRAMA DE DISCIPLINA} & \\
\end{tabular}
\end{center}

\begin{center}
\begin{longtable}{|p{4cm}|p{4cm}|p{4cm}|p{4cm}|}
\hline
\multicolumn{3}{|p{12cm}|}{Nome do Componente Curricular em Português:} &
\multicolumn{1}{p{4cm}|}{Código:} \\ 
\multicolumn{3}{|p{12cm}|}{\textbf{Filosofia Geral}} &
\textbf{FIL110}\\ 
\multicolumn{3}{|p{12cm}|}{Nome do Componente Curricular em Inglês:} & \\ 
\multicolumn{3}{|p{12cm}|}{\textbf{General Philosophy}} & \\ 
\hline
\multicolumn{3}{|p{12cm}|}{Nome e Sigla do Departamento} & Unidade Acadêmica: \\ 
\multicolumn{3}{|p{12cm}|}{Departamento de Filosofia (DEFIL)} & {IFAC} \\ 
\hline
\multicolumn{4}{|p{16cm}|}{Modalidade de Oferta:
[X] presencial \hspace{0.6cm}
[ ] semipresencial \hspace{0.6cm}
[ ] à distância}\\
\hline
\multicolumn{2}{|p{8cm}|}{Carga horária semestral} &
\multicolumn{2}{p{8cm}|}{Carga horária semanal}\\
\hline
\multicolumn{1}{|p{4cm}|}{Total} &
\multicolumn{1}{p{4cm}|}{Extensionista} &
\multicolumn{1}{p{4cm}|}{Teórica} &
\multicolumn{1}{p{4cm}|}{Prática} \\ 
\multicolumn{1}{|p{4cm}|}{60\,horas} &
\multicolumn{1}{p{4cm}|}{0\;horas} &
\multicolumn{1}{p{4cm}|}{4\;horas/aula} &
\multicolumn{1}{p{4cm}|}{0\;horas/aula} \\ 
\hline
\multicolumn{4}{|p{16cm}|}{Ementa:}\\
\multicolumn{4}{|p{16cm}|}{}\\
\multicolumn{4}{|p{\dimexpr 16cm + 6\tabcolsep\relax}|}{Panorama do pensamento filosófico ocidental desde seus primórdios até os dias de hoje, através de textos consagrados pela tradição e norteados por preocupações teórico-temáticas comuns. Eixos: ontologia, gnoseologia e ética. Noções fundamentais de Lógica.}\\
\multicolumn{4}{|p{16cm}|}{}\\
\hline
\multicolumn{4}{|p{16cm}|}{Conteúdo programático:}\\*
\multicolumn{4}{|p{\dimexpr 16cm + 6\tabcolsep\relax}|}{\begin{enumerate}[start=1]\item UNIDADE 1: Horizontes da Reflexão Filosófica Clássica. A alegoria da caverna como seu emblema. Desdobramentos da pesquisa em filosofia sobre diversos domínios da atividade humana\end{enumerate}}\\
\multicolumn{4}{|p{\dimexpr 16cm + 6\tabcolsep\relax}|}{\begin{enumerate}[start=2]\item UNIDADE 2: Filosofia Medieval a partir de um tópico exemplar: a querela dos universais\end{enumerate}}\\
\multicolumn{4}{|p{\dimexpr 16cm + 6\tabcolsep\relax}|}{\begin{enumerate}[start=3]\item UNIDADE 3: Linhas de Força do Pensamento Moderno. O século XVII: racionalismo cartesiano versus empirismo humano, a par da alternativa de Pascal. O século XVIII: Kant e o esclarecimento, a par da enciclopédia. Coramento e crise da filosofia no período: Hegel, hegelianismo e a crítica nietzcheana\end{enumerate}}\\
\multicolumn{4}{|p{\dimexpr 16cm + 6\tabcolsep\relax}|}{\begin{enumerate}[start=4]\item UNIDADE 4: Linhas de Força do Pensamento Contemporâneo. Empirismo lógico e Wittgenstein. Sartre e o existencialismo. Tentativa de retomada das sínteses sistemáticas: Habermas\end{enumerate}}\\
\multicolumn{4}{|p{\dimexpr 16cm + 6\tabcolsep\relax}|}{\begin{enumerate}[start=5]\item UNIDADE 5: Definição de Lógica. Argumentos. Dedução e indução. Verdade e validez\end{enumerate}}\\
\multicolumn{4}{|p{\dimexpr 16cm + 6\tabcolsep\relax}|}{\begin{enumerate}[start=6]\item UNIDADE 6: Silogismo. Formação do Silogismo. Propriedade do silogismo\end{enumerate}}\\
\multicolumn{4}{|p{\dimexpr 16cm + 6\tabcolsep\relax}|}{}\\
\hline
\multicolumn{4}{|p{16cm}|}{Bibliografia Básica:}\\*
\multicolumn{4}{|p{\dimexpr 16cm + 6\tabcolsep\relax}|}{\begin{itemize}\item LEBRUN, G. Passeios ao Léu. Brasiliense, 1983.\end{itemize}}\\
\multicolumn{4}{|p{\dimexpr 16cm + 6\tabcolsep\relax}|}{\begin{itemize}\item MORENTE, M. G. Fundamentos de Filosofia. Lições Preliminares. Mestre Jou, 1980.\end{itemize}}\\
\multicolumn{4}{|p{\dimexpr 16cm + 6\tabcolsep\relax}|}{\begin{itemize}\item Pascal, Blaise. Pensamentos. São Paulo: Cultrix, 1967.\end{itemize}}\\
\multicolumn{4}{|p{\dimexpr 16cm + 6\tabcolsep\relax}|}{\begin{itemize}\item Kant, Immanuel. Textos seletos. 2. ed. Petrópolis: Vozes, 1985.\end{itemize}}\\
\multicolumn{4}{|p{\dimexpr 16cm + 6\tabcolsep\relax}|}{\begin{itemize}\item Hegel, Georg Wilhelm Friedrich. Fenomenologia do espírito. 3. ed. Petrópolis: Vozes.\end{itemize}}\\
\multicolumn{4}{|p{\dimexpr 16cm + 6\tabcolsep\relax}|}{\begin{itemize}\item Tradução de Pietro Nasseti. A república: Platão. São Paulo: Martin Claret, 2002.\end{itemize}}\\
\multicolumn{4}{|p{\dimexpr 16cm + 6\tabcolsep\relax}|}{\begin{itemize}\item Aristóteles. Ética a Nicômaco. 2. ed. Bauru, SP: EDIPRO, 2007.\end{itemize}}\\
\multicolumn{4}{|p{\dimexpr 16cm + 6\tabcolsep\relax}|}{\begin{itemize}\item Descartes, René. Meditações metafísicas. 2. ed. São Paulo: Martins Fontes, 2005.\end{itemize}}\\
\multicolumn{4}{|p{\dimexpr 16cm + 6\tabcolsep\relax}|}{\begin{itemize}\item Jean-Paul Sartre. Conferências e escritos filosóficos / Martin Heidegger. O existencialismo é um humanismo: A imaginação; Questão de método. São Paulo: Abril Cultural, 1973.\end{itemize}}\\
\multicolumn{4}{|p{\dimexpr 16cm + 6\tabcolsep\relax}|}{\begin{itemize}\item Copi, Irving M. Introdução à lógica. 3. ed. São Paulo: Mestre Jou, 1981.\end{itemize}}\\
\multicolumn{4}{|p{\dimexpr 16cm + 6\tabcolsep\relax}|}{\begin{itemize}\item Russell, Bertrand. Os problemas da filosofia. 2. ed. Coimbra: Armênio Amado, 1959.\end{itemize}}\\
\multicolumn{4}{|p{\dimexpr 16cm + 6\tabcolsep\relax}|}{\begin{itemize}\item Gerd A. Bornheim. Introdução ao filosofar: o pensamento filosófico em bases existenciais. Porto Alegre, RS: Globo, 1970.\end{itemize}}\\
\multicolumn{4}{|p{\dimexpr 16cm + 6\tabcolsep\relax}|}{\begin{itemize}\item Nietzsche, Friedrich Wilhelm. O livro do filósofo. 6. ed. São Paulo: Centauro, 2005.\end{itemize}}\\
\multicolumn{4}{|p{\dimexpr 16cm + 6\tabcolsep\relax}|}{\begin{itemize}\item Hume, David. Investigações sobre o entendimento humano e sobre os princípios da moral. São Paulo: UNESP, 2004.\end{itemize}}\\
\multicolumn{4}{|p{16cm}|}{}\\
\hline
\multicolumn{4}{|p{16cm}|}{Bibliografia Complementar:}\\*
\multicolumn{4}{|p{\dimexpr 16cm + 6\tabcolsep\relax}|}{\begin{itemize}\item[] (nenhum item informado)\end{itemize}}\\
\hline
\end{longtable}
\end{center}

\clearpage
\begin{center}
\begin{tabular}{lccr}
 \multirow{3}{*}{\includegraphics[height=2.7cm]{brasao.png}} &
 \multicolumn{2}{c}{\bfseries UNIVERSIDADE FEDERAL DE OURO PRETO} &
 \ \ \ \ \multirow{3}{*}{\includegraphics[height=2.7cm]{ufop.png}} \\
 & \multicolumn{2}{c}{\bfseries PRÓ-REITORIA DE GRADUAÇÃO} & \\
 & \multicolumn{2}{c}{} & \\
 & \multicolumn{2}{c}{\Large\bfseries PROGRAMA DE DISCIPLINA} & \\
\end{tabular}
\end{center}

\begin{center}
\begin{longtable}{|p{4cm}|p{4cm}|p{4cm}|p{4cm}|}
\hline
\multicolumn{3}{|p{12cm}|}{Nome do Componente Curricular em Português:} &
\multicolumn{1}{p{4cm}|}{Código:} \\ 
\multicolumn{3}{|p{12cm}|}{\textbf{Ações Empreendedoras}} &
\textbf{PRO035}\\ 
\multicolumn{3}{|p{12cm}|}{Nome do Componente Curricular em Inglês:} & \\ 
\multicolumn{3}{|p{12cm}|}{\textbf{Entrepreneurship}} & \\ 
\hline
\multicolumn{3}{|p{12cm}|}{Nome e Sigla do Departamento} & Unidade Acadêmica: \\ 
\multicolumn{3}{|p{12cm}|}{Departamento de Engenharia de Produção (DEPRO)} & {EM} \\ 
\hline
\multicolumn{4}{|p{16cm}|}{Modalidade de Oferta:
[X] presencial \hspace{0.6cm}
[ ] semipresencial \hspace{0.6cm}
[ ] à distância}\\
\hline
\multicolumn{2}{|p{8cm}|}{Carga horária semestral} &
\multicolumn{2}{p{8cm}|}{Carga horária semanal}\\
\hline
\multicolumn{1}{|p{4cm}|}{Total} &
\multicolumn{1}{p{4cm}|}{Extensionista} &
\multicolumn{1}{p{4cm}|}{Teórica} &
\multicolumn{1}{p{4cm}|}{Prática} \\ 
\multicolumn{1}{|p{4cm}|}{60\,horas} &
\multicolumn{1}{p{4cm}|}{30\;horas} &
\multicolumn{1}{p{4cm}|}{2\;horas/aula} &
\multicolumn{1}{p{4cm}|}{2\;horas/aula} \\ 
\hline
\multicolumn{4}{|p{16cm}|}{Ementa:}\\
\multicolumn{4}{|p{16cm}|}{}\\
\multicolumn{4}{|p{\dimexpr 16cm + 6\tabcolsep\relax}|}{Carreira empreendedora, modelagem e implementação de negócios, instrumentos de apoio ao empreendedorismo, instituições de promoção de empreendedorismo, proteção intelectual, experiências de empreendedores. Práticas de extensão: gestão de material; participação em plataformas digitais abertas; eventos de empreendedorismo; concursos de empreendedorismo; assessorias em empresas.}\\
\multicolumn{4}{|p{16cm}|}{}\\
\hline
\multicolumn{4}{|p{16cm}|}{Conteúdo programático:}\\*
\multicolumn{4}{|p{\dimexpr 16cm + 6\tabcolsep\relax}|}{\begin{enumerate}[start=1]\item Carreira empreendedora: empreendedorismo como opção de carreira, perfil empreendedor, fatores a serem considerados na escolha\end{enumerate}}\\
\multicolumn{4}{|p{\dimexpr 16cm + 6\tabcolsep\relax}|}{\begin{enumerate}[start=2]\item Desenvolvimento e implementação de negócios: pré-requisitos necessários para um empreendimento, identificação de oportunidades, transformação da oportunidade em um conceito, análise financeira, avaliação do potencial de lucro e crescimento, dinâmica dos negócios, escolha da estratégia competitiva\end{enumerate}}\\
\multicolumn{4}{|p{\dimexpr 16cm + 6\tabcolsep\relax}|}{\begin{enumerate}[start=3]\item Instrumentos de apoio ao empreendedorismo: CANVAS, plano de negócio\end{enumerate}}\\
\multicolumn{4}{|p{\dimexpr 16cm + 6\tabcolsep\relax}|}{\begin{enumerate}[start=4]\item Instituições de promoção de empreendedorismo: incubadora, aceleradoras, parques tecnológicos\end{enumerate}}\\
\multicolumn{4}{|p{\dimexpr 16cm + 6\tabcolsep\relax}|}{\begin{enumerate}[start=5]\item Proteção intelectual: proteção à propriedade intelectual de programa de computador, patentes, Núcleo de Inovação Tecnológica\end{enumerate}}\\
\multicolumn{4}{|p{\dimexpr 16cm + 6\tabcolsep\relax}|}{\begin{enumerate}[start=6]\item Experiências de empreendedores: falas de empreendedores\end{enumerate}}\\
\multicolumn{4}{|p{\dimexpr 16cm + 6\tabcolsep\relax}|}{\begin{enumerate}[start=7]\item Conteúdo programático extensionista: Gestão de material: criação de material escrito e digital sobre empreendedorismo e disponibilização a comunidades; Participação em plataformas digitais abertas: criação e manutenção de canais de informação sobre empreendedorismo em plataformas digitais abertas; Eventos de empreendedorismo: criação e/ou participação em eventos sobre modelos de negócios; Concursos de empreendedorismo: participação em concursos sobre modelos de negócios; Assessorias em empresas: apoio a negócios que estejam em formação em colaboração com parceiros (incubadoras, aceleradoras e ecossistema de inovação)\end{enumerate}}\\
\multicolumn{4}{|p{\dimexpr 16cm + 6\tabcolsep\relax}|}{}\\
\hline
\multicolumn{4}{|p{16cm}|}{Bibliografia Básica:}\\*
\multicolumn{4}{|p{\dimexpr 16cm + 6\tabcolsep\relax}|}{\begin{itemize}\item FABRETE, Teresa Cristina Lopes. Empreendedorismo. 2. ed. São Paulo: Pearson, 2019.\end{itemize}}\\
\multicolumn{4}{|p{\dimexpr 16cm + 6\tabcolsep\relax}|}{\begin{itemize}\item DORNELAS, José. Empreendedorismo para visionários: desenvolvendo negócios inovadores para um mundo em transformação. 2. ed. São Paulo: Empreende, 2019.\end{itemize}}\\
\multicolumn{4}{|p{\dimexpr 16cm + 6\tabcolsep\relax}|}{\begin{itemize}\item MAÇÃES, Manuel Alberto Ramos. Empreendedorismo, inovação e mudança organizacional. Coimbra: Actual, 2017.\end{itemize}}\\
\multicolumn{4}{|p{16cm}|}{}\\
\hline
\multicolumn{4}{|p{16cm}|}{Bibliografia Complementar:}\\*
\multicolumn{4}{|p{\dimexpr 16cm + 6\tabcolsep\relax}|}{\begin{itemize}\item CANDIDO, Claudio Roberto. Empreendedorismo: uma perspectiva multidisciplinar. Rio de Janeiro, 2016.\end{itemize}}\\
\multicolumn{4}{|p{\dimexpr 16cm + 6\tabcolsep\relax}|}{\begin{itemize}\item SERTEK, Paula. Empreendedorismo. Curitiba: Intersaberes, 2013.\end{itemize}}\\
\multicolumn{4}{|p{\dimexpr 16cm + 6\tabcolsep\relax}|}{\begin{itemize}\item WILDAUER, Egon Walter. Plano de negócios. Intersaberes, 2012.\end{itemize}}\\
\multicolumn{4}{|p{\dimexpr 16cm + 6\tabcolsep\relax}|}{\begin{itemize}\item RAZZOLINI FILHO, Etelvino. Empreendedorismo: dicas e planos de negócios para o século XXI. Intersaberes, 2012.\end{itemize}}\\
\hline
\end{longtable}
\end{center}

\clearpage
\begin{center}
\begin{tabular}{lccr}
 \multirow{3}{*}{\includegraphics[height=2.7cm]{brasao.png}} &
 \multicolumn{2}{c}{\bfseries UNIVERSIDADE FEDERAL DE OURO PRETO} &
 \ \ \ \ \multirow{3}{*}{\includegraphics[height=2.7cm]{ufop.png}} \\
 & \multicolumn{2}{c}{\bfseries PRÓ-REITORIA DE GRADUAÇÃO} & \\
 & \multicolumn{2}{c}{} & \\
 & \multicolumn{2}{c}{\Large\bfseries PROGRAMA DE DISCIPLINA} & \\
\end{tabular}
\end{center}

\begin{center}
\begin{longtable}{|p{4cm}|p{4cm}|p{4cm}|p{4cm}|}
\hline
\multicolumn{3}{|p{12cm}|}{Nome do Componente Curricular em Português:} &
\multicolumn{1}{p{4cm}|}{Código:} \\ 
\multicolumn{3}{|p{12cm}|}{\textbf{Logística}} &
\textbf{PRO315}\\ 
\multicolumn{3}{|p{12cm}|}{Nome do Componente Curricular em Inglês:} & \\ 
\multicolumn{3}{|p{12cm}|}{\textbf{Logistics}} & \\ 
\hline
\multicolumn{3}{|p{12cm}|}{Nome e Sigla do Departamento} & Unidade Acadêmica: \\ 
\multicolumn{3}{|p{12cm}|}{Departamento de Engenharia de Produção (DEPRO)} & {EM} \\ 
\hline
\multicolumn{4}{|p{16cm}|}{Modalidade de Oferta:
[X] presencial \hspace{0.6cm}
[ ] semipresencial \hspace{0.6cm}
[ ] à distância}\\
\hline
\multicolumn{2}{|p{8cm}|}{Carga horária semestral} &
\multicolumn{2}{p{8cm}|}{Carga horária semanal}\\
\hline
\multicolumn{1}{|p{4cm}|}{Total} &
\multicolumn{1}{p{4cm}|}{Extensionista} &
\multicolumn{1}{p{4cm}|}{Teórica} &
\multicolumn{1}{p{4cm}|}{Prática} \\ 
\multicolumn{1}{|p{4cm}|}{60\,horas} &
\multicolumn{1}{p{4cm}|}{0\;horas} &
\multicolumn{1}{p{4cm}|}{4\;horas/aula} &
\multicolumn{1}{p{4cm}|}{0\;horas/aula} \\ 
\hline
\multicolumn{4}{|p{16cm}|}{Ementa:}\\
\multicolumn{4}{|p{16cm}|}{}\\
\multicolumn{4}{|p{\dimexpr 16cm + 6\tabcolsep\relax}|}{Conceitos de logística, cadeia de suprimentos, distribuição. Gerência de transportes (modais, realidade brasileira, custos, modelagem matemática). Serviço ao cliente (lead time, oferta de serviços, ECR). Gestão de estoques/armazéns (organização, processamento, embalagens). Distribuição Física (conceitos de canal, roteirização, localização industrial e de depósitos). Tecnologia e Sistema de Informação (EDI, GPS, GIS, E-commerce). Modelos de otimização aplicados à logística. Estudos de Caso. Objetivo: O aluno deverá ser capaz de conceituar logística, distribuição, cadeia de suprimentos, identificar funções e elementos da logística; analisar a importância da mesma na gerência de cadeia de suprimentos; analisar a inserção dos transportes, estoques, armazenagem, embalagem e localização nas atividades logísticas; e desenvolver modelos de otimização aplicados aos problemas inerentes à logística.}\\
\multicolumn{4}{|p{16cm}|}{}\\
\hline
\multicolumn{4}{|p{16cm}|}{Conteúdo programático:}\\*
\multicolumn{4}{|p{\dimexpr 16cm + 6\tabcolsep\relax}|}{\begin{enumerate}[start=1]\item Conceitos de logística. Histórico. Sub-áreas\end{enumerate}}\\
\multicolumn{4}{|p{\dimexpr 16cm + 6\tabcolsep\relax}|}{\begin{enumerate}[start=2]\item Atividades Primárias e de apoio. Logística e estratégia competitiva. Rede Logística\end{enumerate}}\\
\multicolumn{4}{|p{\dimexpr 16cm + 6\tabcolsep\relax}|}{\begin{enumerate}[start=3]\item Serviço ao cliente (conceito, significado, componentes; Redução do lead time, oferta de serviços)\end{enumerate}}\\
\multicolumn{4}{|p{\dimexpr 16cm + 6\tabcolsep\relax}|}{\begin{enumerate}[start=4]\item Localização (teorias, localização de fábricas, Armazéns). Localização no Brasil (teorias e realidade de Empresas). Estudo de caso\end{enumerate}}\\
\multicolumn{4}{|p{\dimexpr 16cm + 6\tabcolsep\relax}|}{\begin{enumerate}[start=5]\item Modelos de Localização\end{enumerate}}\\
\multicolumn{4}{|p{\dimexpr 16cm + 6\tabcolsep\relax}|}{\begin{enumerate}[start=6]\item Transportes (modais, multimodalidade, intermodalidade, externalidades, objetivos, planejamentos e custos)\end{enumerate}}\\
\multicolumn{4}{|p{\dimexpr 16cm + 6\tabcolsep\relax}|}{\begin{enumerate}[start=7]\item Frota própria ou de terceiros. Administração de transportes de terceiros (fretes, auditoria, monitoração). Frota própria (roteirização, rota, despacho de veículos). Estudo de Caso\end{enumerate}}\\
\multicolumn{4}{|p{\dimexpr 16cm + 6\tabcolsep\relax}|}{\begin{enumerate}[start=8]\item Operadores Logísticos\end{enumerate}}\\
\multicolumn{4}{|p{\dimexpr 16cm + 6\tabcolsep\relax}|}{\begin{enumerate}[start=9]\item Aplicações de modelos de Transporte\end{enumerate}}\\
\multicolumn{4}{|p{\dimexpr 16cm + 6\tabcolsep\relax}|}{\begin{enumerate}[start=10]\item Modelos de dimensionamento e operação de frota\end{enumerate}}\\
\multicolumn{4}{|p{\dimexpr 16cm + 6\tabcolsep\relax}|}{\begin{enumerate}[start=11]\item Modelos de roteirização\end{enumerate}}\\
\multicolumn{4}{|p{\dimexpr 16cm + 6\tabcolsep\relax}|}{\begin{enumerate}[start=12]\item Controle de estoques (demanda, custos, push e pull, Curva ABC, just-in-time, kanban)\end{enumerate}}\\
\multicolumn{4}{|p{\dimexpr 16cm + 6\tabcolsep\relax}|}{\begin{enumerate}[start=13]\item Armazenagem (localização, funções, dimensões, embalagem e manuseio de produtos). Estudo de Caso\end{enumerate}}\\
\multicolumn{4}{|p{\dimexpr 16cm + 6\tabcolsep\relax}|}{\begin{enumerate}[start=14]\item Modelos de controles\end{enumerate}}\\
\multicolumn{4}{|p{\dimexpr 16cm + 6\tabcolsep\relax}|}{\begin{enumerate}[start=15]\item Distribuição (conceitos, canal de distribuição, tipos, sistemas de distribuição, políticas de distribuição)\end{enumerate}}\\
\multicolumn{4}{|p{\dimexpr 16cm + 6\tabcolsep\relax}|}{\begin{enumerate}[start=16]\item Tecnologia e Sistema de Informação (conceitos, diferenças, GIS, GPS, EDI etc., uso na logística) e Competitividade. Estudo de Caso\end{enumerate}}\\
\multicolumn{4}{|p{\dimexpr 16cm + 6\tabcolsep\relax}|}{\begin{enumerate}[start=17]\item E-commerce e logística\end{enumerate}}\\
\multicolumn{4}{|p{\dimexpr 16cm + 6\tabcolsep\relax}|}{\begin{enumerate}[start=18]\item Visitas Técnicas\end{enumerate}}\\
\multicolumn{4}{|p{\dimexpr 16cm + 6\tabcolsep\relax}|}{\begin{enumerate}[start=19]\item Seminários\end{enumerate}}\\
\multicolumn{4}{|p{\dimexpr 16cm + 6\tabcolsep\relax}|}{}\\
\hline
\multicolumn{4}{|p{16cm}|}{Bibliografia Básica:}\\*
\multicolumn{4}{|p{\dimexpr 16cm + 6\tabcolsep\relax}|}{\begin{itemize}\item ALVARENGA, Antônio C.; NOVAES, Antônio G. N. Logística Aplicada – Suprimento e Distribuição Física. São Paulo: Edgard Blücher.\end{itemize}}\\
\multicolumn{4}{|p{\dimexpr 16cm + 6\tabcolsep\relax}|}{\begin{itemize}\item BALLOU, Ronald H. Logística Empresarial: Transportes, Administração de Materiais e Distribuição Física. São Paulo: Atlas, 1993.\end{itemize}}\\
\multicolumn{4}{|p{\dimexpr 16cm + 6\tabcolsep\relax}|}{\begin{itemize}\item BUFFA, Elwood S.; DYER, James S. Essentials of Management Science/Operations Research.\end{itemize}}\\
\multicolumn{4}{|p{\dimexpr 16cm + 6\tabcolsep\relax}|}{\begin{itemize}\item CHAPMAN, Paul T. In: ROBESON, J. F. et al. Logistics Networks Modeling. The Logistics Handbook. New York: The Free Press, 1994.\end{itemize}}\\
\multicolumn{4}{|p{\dimexpr 16cm + 6\tabcolsep\relax}|}{\begin{itemize}\item CHING, Hong Yuh. Gestão de estoques na cadeia de logística integrada – Supply Chain. São Paulo: Atlas, 1999.\end{itemize}}\\
\multicolumn{4}{|p{\dimexpr 16cm + 6\tabcolsep\relax}|}{\begin{itemize}\item CHRISTOPHER, Martin. Logística e gerenciamento da cadeia de suprimentos: estratégias para a redução de custos e melhoria dos serviços. São Paulo: Pioneira, 1997.\end{itemize}}\\
\multicolumn{4}{|p{\dimexpr 16cm + 6\tabcolsep\relax}|}{\begin{itemize}\item DIAS, Sérgio R. Estratégia e Canais de Distribuição. São Paulo: Atlas, 1993.\end{itemize}}\\
\multicolumn{4}{|p{\dimexpr 16cm + 6\tabcolsep\relax}|}{\begin{itemize}\item HARMON, Roy L. Reinventando a distribuição – Logística de Distribuição - Classe Mundial. Rio de Janeiro: Campus, 1994.\end{itemize}}\\
\multicolumn{4}{|p{\dimexpr 16cm + 6\tabcolsep\relax}|}{\begin{itemize}\item KOBAYASHI, Shun'ichi. Renovação da logística: como definir as estratégias de distribuição física. São Paulo: Atlas, 2000.\end{itemize}}\\
\multicolumn{4}{|p{\dimexpr 16cm + 6\tabcolsep\relax}|}{\begin{itemize}\item LIEBERMAN, Gerald J.; HILLIER, Frederick S. Introdução à Pesquisa Operacional. 3. ed. Rio de Janeiro: Campus.\end{itemize}}\\
\multicolumn{4}{|p{\dimexpr 16cm + 6\tabcolsep\relax}|}{\begin{itemize}\item WAGNER, Harvey M. Principles of Operations Research with Applications to Managerial Decisions. New Jersey: Prentice-Hall.\end{itemize}}\\
\multicolumn{4}{|p{\dimexpr 16cm + 6\tabcolsep\relax}|}{\begin{itemize}\item FLEURY, P.F.; WANKE, P.; FIGUEIREDO, K.F. (org.). Logística Empresarial – A Perspectiva Brasileira. São Paulo: Atlas/COPPEAD-UFRJ, 2000.\end{itemize}}\\
\multicolumn{4}{|p{16cm}|}{}\\
\hline
\multicolumn{4}{|p{16cm}|}{Bibliografia Complementar:}\\*
\multicolumn{4}{|p{\dimexpr 16cm + 6\tabcolsep\relax}|}{\begin{itemize}\item[] (nenhum item informado)\end{itemize}}\\
\hline
\end{longtable}
\end{center}

\clearpage
\begin{center}
\begin{tabular}{lccr}
 \multirow{3}{*}{\includegraphics[height=2.7cm]{brasao.png}} &
 \multicolumn{2}{c}{\bfseries UNIVERSIDADE FEDERAL DE OURO PRETO} &
 \ \ \ \ \multirow{3}{*}{\includegraphics[height=2.7cm]{ufop.png}} \\
 & \multicolumn{2}{c}{\bfseries PRÓ-REITORIA DE GRADUAÇÃO} & \\
 & \multicolumn{2}{c}{} & \\
 & \multicolumn{2}{c}{\Large\bfseries PROGRAMA DE DISCIPLINA} & \\
\end{tabular}
\end{center}

\begin{center}
\begin{longtable}{|p{4cm}|p{4cm}|p{4cm}|p{4cm}|}
\hline
\multicolumn{3}{|p{12cm}|}{Nome do Componente Curricular em Português:} &
\multicolumn{1}{p{4cm}|}{Código:} \\ 
\multicolumn{3}{|p{12cm}|}{\textbf{Internet das Coisas}} &
\textbf{BCC813}\\ 
\multicolumn{3}{|p{12cm}|}{Nome do Componente Curricular em Inglês:} & \\ 
\multicolumn{3}{|p{12cm}|}{\textbf{Internet of Things}} & \\ 
\hline
\multicolumn{3}{|p{12cm}|}{Nome e Sigla do Departamento} & Unidade Acadêmica: \\ 
\multicolumn{3}{|p{12cm}|}{Departamento de Computação (DECOM)} & {ICEB} \\ 
\hline
\multicolumn{4}{|p{16cm}|}{Modalidade de Oferta:
[X] presencial \hspace{0.6cm}
[ ] semipresencial \hspace{0.6cm}
[ ] à distância}\\
\hline
\multicolumn{2}{|p{8cm}|}{Carga horária semestral} &
\multicolumn{2}{p{8cm}|}{Carga horária semanal}\\
\hline
\multicolumn{1}{|p{4cm}|}{Total} &
\multicolumn{1}{p{4cm}|}{Extensionista} &
\multicolumn{1}{p{4cm}|}{Teórica} &
\multicolumn{1}{p{4cm}|}{Prática} \\ 
\multicolumn{1}{|p{4cm}|}{60\,horas} &
\multicolumn{1}{p{4cm}|}{0\;horas} &
\multicolumn{1}{p{4cm}|}{4\;horas/aula} &
\multicolumn{1}{p{4cm}|}{0\;horas/aula} \\ 
\hline
\multicolumn{4}{|p{16cm}|}{Ementa:}\\
\multicolumn{4}{|p{16cm}|}{}\\
\multicolumn{4}{|p{\dimexpr 16cm + 6\tabcolsep\relax}|}{Introdução aos conceitos, arquiteturas e tecnologias da Internet das Coisas (IoT). Protocolos de comunicação, sensores, atuadores, segurança e aplicações inteligentes em ambientes conectados.}\\
\multicolumn{4}{|p{16cm}|}{}\\
\hline
\multicolumn{4}{|p{16cm}|}{Conteúdo programático:}\\*
\multicolumn{4}{|p{\dimexpr 16cm + 6\tabcolsep\relax}|}{\begin{enumerate}[start=1]\item Fundamentos da Internet das Coisas: definições, evolução e panorama atual.\end{enumerate}}\\
\multicolumn{4}{|p{\dimexpr 16cm + 6\tabcolsep\relax}|}{\begin{enumerate}[start=2]\item Arquitetura IoT: camadas, dispositivos e sistemas embarcados.\end{enumerate}}\\
\multicolumn{4}{|p{\dimexpr 16cm + 6\tabcolsep\relax}|}{\begin{enumerate}[start=3]\item Sensores e atuadores: tipos, funcionamento e integração.\end{enumerate}}\\
\multicolumn{4}{|p{\dimexpr 16cm + 6\tabcolsep\relax}|}{\begin{enumerate}[start=4]\item Protocolos de comunicação IoT.\end{enumerate}}\\
\multicolumn{4}{|p{\dimexpr 16cm + 6\tabcolsep\relax}|}{\begin{enumerate}[start=5]\item Plataformas e ferramentas para desenvolvimento IoT.\end{enumerate}}\\
\multicolumn{4}{|p{\dimexpr 16cm + 6\tabcolsep\relax}|}{\begin{enumerate}[start=6]\item Segurança e privacidade em ambientes IoT.\end{enumerate}}\\
\multicolumn{4}{|p{\dimexpr 16cm + 6\tabcolsep\relax}|}{\begin{enumerate}[start=7]\item Coleta, transmissão e processamento de dados em IoT.\end{enumerate}}\\
\multicolumn{4}{|p{\dimexpr 16cm + 6\tabcolsep\relax}|}{\begin{enumerate}[start=8]\item Integração com inteligência artificial e computação em nuvem.\end{enumerate}}\\
\multicolumn{4}{|p{\dimexpr 16cm + 6\tabcolsep\relax}|}{\begin{enumerate}[start=9]\item Casos de uso: cidades inteligentes, agricultura de precisão, saúde e indústria 4.0.\end{enumerate}}\\
\multicolumn{4}{|p{\dimexpr 16cm + 6\tabcolsep\relax}|}{\begin{enumerate}[start=10]\item Tendências e desafios futuros da Internet das Coisas.\end{enumerate}}\\
\multicolumn{4}{|p{\dimexpr 16cm + 6\tabcolsep\relax}|}{}\\
\hline
\multicolumn{4}{|p{16cm}|}{Bibliografia Básica:}\\*
\multicolumn{4}{|p{\dimexpr 16cm + 6\tabcolsep\relax}|}{\begin{itemize}\item WILLETT, Richard; OLIVEIRA, Charalampos. Internet das Coisas: Conceitos, Tecnologias e Aplicações. Novatec, 2022.\end{itemize}}\\
\multicolumn{4}{|p{\dimexpr 16cm + 6\tabcolsep\relax}|}{\begin{itemize}\item VERGARA, Luiz Antônio. Internet das Coisas. 2. ed. Campus, 2020.\end{itemize}}\\
\multicolumn{4}{|p{\dimexpr 16cm + 6\tabcolsep\relax}|}{\begin{itemize}\item MINERVA, Roberto et al. Towards a definition of the Internet of Things (IoT). IEEE Internet Initiative, 2015.\end{itemize}}\\
\multicolumn{4}{|p{16cm}|}{}\\
\hline
\multicolumn{4}{|p{16cm}|}{Bibliografia Complementar:}\\*
\multicolumn{4}{|p{\dimexpr 16cm + 6\tabcolsep\relax}|}{\begin{itemize}\item GUBBI, Jayavardhana et al. Internet of Things (IoT): A Vision, Architectural Elements, and Future Directions. Future Generation Computer Systems, 2013.\end{itemize}}\\
\multicolumn{4}{|p{\dimexpr 16cm + 6\tabcolsep\relax}|}{\begin{itemize}\item KURUNDRAN, S.; BHASKAR, P. IoT Protocols and Standards: An Overview. Springer, 2021.\end{itemize}}\\
\multicolumn{4}{|p{\dimexpr 16cm + 6\tabcolsep\relax}|}{\begin{itemize}\item RAY, Partha Pratim. Internet of Things: Architecture and Applications. Wiley, 2018.\end{itemize}}\\
\multicolumn{4}{|p{\dimexpr 16cm + 6\tabcolsep\relax}|}{\begin{itemize}\item SINGH, D.; ZHANG, H.; KUMAR, N. Internet of Things: Principles and Paradigms. Wiley, 2018.\end{itemize}}\\
\multicolumn{4}{|p{\dimexpr 16cm + 6\tabcolsep\relax}|}{\begin{itemize}\item STANKOVIC, John A. Research Directions for the Internet of Things. IEEE Internet of Things Journal, 2014.\end{itemize}}\\
\hline
\end{longtable}
\end{center}

\clearpage
\end{document}